\input{preamble}

% OK, start here.
%
\begin{document}

\title{Sites et faisceaux}


\maketitle

\phantomsection
\label{section-phantom}

\tableofcontents

\section{Introduction}
\label{section-introduction}

\noindent
La notion de site a été introduite par Grothendieck afin de pouvoir étudier
les faisceaux pour la topologie \'etale des schémas. La référence fondamentale
pour cette notion est peut-être \cite{SGA4}. Notre notion de site diffère de
celle de \cite{SGA4} ; ce que nous appelons un site y est appelé une catégorie munie
d’une prétopologie dans \cite[Expos\'e II, D\'efinition 1.3]{SGA4}.
Nous faisons ce choix parce qu’en géométrie algébrique il est souvent commode de
travailler avec une classe donnée de recouvrements, par exemple lorsqu’on définit
ce que signifie, pour une propriété des schémas, être locale pour une topologie
donnée ; voir Descente, section \ref{descent-section-descending-properties}.
Notre exposé suivra de près \cite{ArtinTopologies}.
Nous n’utiliserons pas d’univers.












\section{Préfaisceaux}
\label{section-presheaves}

\noindent
Soit $\mathcal{C}$ une catégorie.
Un {\it préfaisceau d’ensembles} est un foncteur contravariant $\mathcal{F}$
de $\mathcal{C}$ vers $\textit{Ensembles}$ (voir Catégories, Remarque
\ref{categories-remark-functor-into-sets}).
Ainsi, pour tout objet $U$ de $\mathcal{C}$, nous avons un ensemble
$\mathcal{F}(U)$. Les éléments de cet ensemble sont appelés
les {\it sections} de $\mathcal{F}$ sur $U$. Pour tout morphisme
$f : V \to U$, l’application $\mathcal{F}(f) : \mathcal{F}(U) \to \mathcal{F}(V)$
est appelée l’{\it application de restriction} et est souvent notée
$f^\ast : \mathcal{F}(U) \to \mathcal{F}(V)$. Une autre façon
d’exprimer ceci consiste à dire que $f^*(s)$ est l’{\it image inverse}
de $s$ par $f$. La fonctorialité signifie que $g^* f^* (s) = (f \circ g)^*(s)$.
Nous employons parfois la notation $s|_V := f^\ast(s)$.
Cette notation est compatible avec la notion de restriction
des fonctions en topologie, car si $W \to V \to U$
sont des morphismes dans $\mathcal{C}$ et si $s$ est une section de
$\mathcal{F}$ sur $U$, alors $s|_W = (s|_V)|_W$ en vertu de la
nature fonctorielle de $\mathcal{F}$. Il faut bien entendu être
prudent, puisqu’il peut fort bien arriver
qu’il existe plus d’un morphisme $V \to U$, et il n’est
certainement pas vrai alors que les applications d’image inverse
correspondantes soient égales.

\begin{definition}
\label{definition-presheaves-sets}
Un {\it préfaisceau d’ensembles} sur $\mathcal{C}$ est un foncteur
contravariant de $\mathcal{C}$ vers $\textit{Ensembles}$. Les {\it morphismes
de préfaisceaux} sont les transformations naturelles de foncteurs. La catégorie
des préfaisceaux d’ensembles est notée $\textit{PSh}(\mathcal{C})$.
\end{definition}

\noindent
Remarquons que, pour tout objet $U$ de $\mathcal{C}$, le foncteur des
points $h_U$, voir Catégories, Exemple \ref{categories-example-hom-functor},
est un préfaisceau. Ceux-ci sont appelés les {\it préfaisceaux représentables}.
Ces préfaisceaux possèdent la propriété agréable suivante : pour tout
préfaisceau $\mathcal{F}$, nous avons
\begin{equation}
\label{equation-map-representable-into-presheaf}
\Mor_{\textit{PSh}(\mathcal{C})}(h_U, \mathcal{F})
=
\mathcal{F}(U).
\end{equation}
C’est le lemme de Yoneda (Catégories, Lemme \ref{categories-lemma-yoneda}).

\medskip\noindent
De même, nous pouvons définir la notion de préfaisceau de groupes abéliens,
d’anneaux, etc. Plus généralement, nous pouvons définir un préfaisceau à valeurs
dans une catégorie.

\begin{definition}
\label{definition-presheaf}
Soient $\mathcal{C}$, $\mathcal{A}$ des catégories.
Un {\it préfaisceau} $\mathcal{F}$ sur $\mathcal{C}$
à valeurs dans $\mathcal{A}$ est un foncteur
contravariant de $\mathcal{C}$ vers $\mathcal{A}$,
c’est-à-dire $\mathcal{F} : \mathcal{C}^{opp} \to \mathcal{A}$.
Un {\it morphisme} de préfaisceaux $\mathcal{F} \to \mathcal{G}$
sur $\mathcal{C}$ à valeurs dans $\mathcal{A}$ est une transformation
naturelle de foncteurs de $\mathcal{F}$ vers $\mathcal{G}$.
\end{definition}

\noindent
Ce sont les objets et les morphismes de la catégorie des préfaisceaux
sur $\mathcal{C}$ à valeurs dans $\mathcal{A}$.

\begin{remark}
\label{remark-big-presheaves}
Comme nous l’avons déjà signalé, nous pouvons considérer la catégorie des
préfaisceaux à valeurs dans l’une quelconque des « grandes » catégories
énumérées dans Catégories, Remarque \ref{categories-remark-big-categories}.
Ce seront elles aussi de « grandes » catégories, et elles seront
ajoutées à la remarque précitée au fur et à mesure.
\end{remark}
















\section{Morphismes injectifs et surjectifs de préfaisceaux}
\label{section-injective-surjective}

\begin{definition}
\label{definition-presheaves-injective-surjective}
Soit $\mathcal{C}$ une catégorie, et soit $\varphi : \mathcal{F}
\to \mathcal{G}$ un morphisme de préfaisceaux d’ensembles.
\begin{enumerate}
\item Nous disons que $\varphi$ est {\it injectif} si, pour tout objet
$U$ de $\mathcal{C}$, l’application $\varphi_U : \mathcal{F}(U)
\to \mathcal{G}(U)$ est injective.
\item Nous disons que $\varphi$ est {\it surjectif} si, pour tout objet
$U$ de $\mathcal{C}$, l’application $\varphi_U : \mathcal{F}(U)
\to \mathcal{G}(U)$ est surjective.
\end{enumerate}
\end{definition}

\begin{lemma}
\label{lemma-mono-epi}
Les morphismes injectifs (resp.\ surjectifs) définis ci-dessus
sont exactement les monomorphismes (resp.\ épimorphismes) de
$\textit{PSh}(\mathcal{C})$. Un morphisme est un isomorphisme
si et seulement s’il est à la fois injectif et surjectif.
\end{lemma}

\begin{proof}
Montrons que $\varphi : \mathcal{F} \to
\mathcal{G}$ est injectif si et seulement si c’est un monomorphisme
de $\textit{PSh}(\mathcal{C})$. Le sens direct
est immédiat ; montrons donc le sens réciproque.
Supposons que $\varphi$ soit un monomorphisme. Soit
$U \in \Ob(\mathcal{C})$ ; nous devons montrer que $\varphi_U$ est
injective. Soient donc $a, b \in \mathcal{F}(U)$ tels que
$\varphi_U (a) = \varphi_U (b)$ ; nous devons vérifier que $a = b$.
Sous l’isomorphisme
(\ref{equation-map-representable-into-presheaf}), les éléments
$a$ et $b$ de $\mathcal{F}(U)$ correspondent à deux transformations
naturelles
$a', b' \in \Mor_{\textit{PSh}(\mathcal{C})}(h_U, \mathcal{F})$.
De même, sous l’isomorphisme analogue
$\Mor_{\textit{PSh}(\mathcal{C})}(h_U, \mathcal{G})
= \mathcal{G}(U)$,
les deux éléments égaux $\varphi_U (a)$ et $\varphi_U (b)$ de
$\mathcal{G}(U)$ correspondent aux deux transformations naturelles
$\varphi \circ a', \varphi \circ b'
\in \Mor_{\textit{PSh}(\mathcal{C})}(h_U, \mathcal{G})$,
qui doivent donc elles aussi être égales. Ainsi
$\varphi \circ a' = \varphi \circ b'$, et par conséquent $a' = b'$
(puisque $\varphi$ est un monomorphisme), d’où $a = b$. Ceci achève (1).

\medskip\noindent
Montrons que $\varphi : \mathcal{F} \to
\mathcal{G}$ est surjectif si et seulement si c’est un épimorphisme
de $\textit{PSh}(\mathcal{C})$. Le sens direct
est immédiat ; montrons donc le sens réciproque.
Supposons que $\varphi$ soit un épimorphisme.

\medskip\noindent
Pour deux morphismes quelconques $f : A \to B$ et $g : A \to C$ de la
catégorie $\textit{Ensembles}$, notons $\text{inl}_{f,g}$ et
$\text{inr}_{f,g}$ les deux applications canoniques de
$B$ et $C$ vers $B \coprod_A C$. (Ici, la somme amalgamée est
calculée dans $\textit{Ensembles}$.)

\medskip\noindent
Définissons maintenant un préfaisceau d’ensembles $\mathcal{H}$ sur $\mathcal{C}$
en posant $\mathcal{H}(U)
= \mathcal{G}(U) \coprod_{\mathcal{F}(U)} \mathcal{G}(U)$ (où
la somme amalgamée est calculée dans $\textit{Ensembles}$ et induite par
l’application $\varphi_U : \mathcal{F}(U) \to \mathcal{G}(U)$) pour
tout $U \in \Ob(\mathcal{C})$ ; son action sur les morphismes est
définie de manière évidente (par la fonctorialité de la somme amalgamée).
Il existe alors deux transformations
naturelles $i_1 : \mathcal{G} \to \mathcal{H}$ et
$i_2 : \mathcal{G} \to \mathcal{H}$ dont les composantes en un objet
$U \in \Ob(\mathcal{C})$ sont données respectivement par les applications
$\text{inl}_{\varphi_U, \varphi_U}$ et
$\text{inr}_{\varphi_U, \varphi_U}$. La
définition d’une somme amalgamée montre que $i_1 \circ \varphi
= i_2 \circ \varphi$, d’où $i_1 = i_2$ (puisque $\varphi$ est un
épimorphisme). Ainsi, pour tout $U \in \Ob(\mathcal{C})$, nous avons
$\text{inl}_{\varphi_U, \varphi_U}
= \text{inr}_{\varphi_U, \varphi_U}$. Par conséquent, $\varphi_U$
doit être surjective (car un argument combinatoire élémentaire montre
que si $f : A \to B$ est un morphisme dans $\textit{Ensembles}$, alors
$\text{inl}_{f,f} = \text{inr}_{f,f}$ si et
seulement si $f$ est surjective). Autrement dit, $\varphi$ est
surjectif, ce qui prouve (2).

\medskip\noindent
Montrons que $\varphi : \mathcal{F} \to
\mathcal{G}$ est à la fois injectif et surjectif si et seulement si c’est
un isomorphisme de $\textit{PSh}(\mathcal{C})$. Cette fois,
le sens réciproque est immédiat. Pour prouver le sens direct,
il suffit d’observer que, si $\varphi$ est à la fois
injectif et surjectif, alors $\varphi_U$ est une application inversible
pour tout $U \in \Ob(\mathcal{C})$, et que les inverses de ces
applications, pour tous les $U$, se combinent en une transformation naturelle
$\mathcal{G} \to \mathcal{F}$ qui est inverse de $\varphi$.
\end{proof}

\begin{definition}
\label{definition-sub-presheaf}
Nous disons que $\mathcal{F}$ est un {\it sous-préfaisceau} de $\mathcal{G}$
si, pour tout objet $U \in \Ob(\mathcal{C})$, l’ensemble
$\mathcal{F}(U)$ est un sous-ensemble de $\mathcal{G}(U)$, de façon compatible
avec les applications de restriction.
\end{definition}

\noindent
Autrement dit, les applications
d’inclusion $\mathcal{F}(U) \to \mathcal{G}(U)$
se recollent pour donner un morphisme (injectif) de
préfaisceaux $\mathcal{F} \to \mathcal{G}$.

\begin{lemma}
\label{lemma-image}
Soit $\mathcal{C}$ une catégorie.
Supposons que $\varphi : \mathcal{F} \to \mathcal{G}$ soit un
morphisme de préfaisceaux d’ensembles sur $\mathcal{C}$.
Il existe un unique sous-préfaisceau $\mathcal{G}' \subset \mathcal{G}$
tel que $\varphi$ se factorise en
$\mathcal{F} \to \mathcal{G}' \to \mathcal{G}$
et que le premier morphisme soit surjectif.
\end{lemma}

\begin{proof}
Pour établir l’existence, il suffit de poser
$\mathcal{G}'(U) = \varphi_U \left(\mathcal{F}(U)\right)$
pour tout $U \in \Ob(C)$ (et d’hériter de l’action sur les morphismes
de $\mathcal{G}$), puis de vérifier que cela définit un
sous-préfaisceau de $\mathcal{G}$ et que $\varphi$ se factorise en
$\mathcal{F} \to \mathcal{G}' \to \mathcal{G}$, le
premier morphisme étant surjectif. L’unicité est immédiate.
\end{proof}

\begin{definition}
\label{definition-image}
Avec les notations du Lemme \ref{lemma-image}, nous
disons que $\mathcal{G}'$ est l’{\it image de $\varphi$}.
\end{definition}
















\section{Limites projectives et limites inductives de préfaisceaux}
\label{section-limits-colimits-PSh}

\noindent
Soit $\mathcal{C}$ une catégorie.
Les limites projectives et les limites inductives existent dans la catégorie
$\textit{PSh}(\mathcal{C})$. De plus, pour tout
$U \in \Ob(\mathcal{C})$, le foncteur
$$
\textit{PSh}(\mathcal{C})
\longrightarrow
\textit{Ensembles}, \quad
\mathcal{F}
\longmapsto
\mathcal{F}(U)
$$
commute aux limites projectives et aux limites inductives. La manière la plus simple
de prouver ces assertions est peut-être la suivante. Étant donné un diagramme
$
\mathcal{F} :
\mathcal{I}
\to
\textit{PSh}(\mathcal{C})
$
définissons les préfaisceaux
$$
\mathcal{F}_{\lim} :
U
\longmapsto
\lim_{i \in \mathcal{I}} \mathcal{F}_i(U)
\text{  et  }
\mathcal{F}_{\colim} :
U
\longmapsto
\colim_{i \in \mathcal{I}} \mathcal{F}_i(U)
$$
Il existe manifestement des morphismes de projection $\mathcal{F}_{\lim} \to \mathcal{F}_i$
et des morphismes canoniques $\mathcal{F}_i \to \mathcal{F}_{\colim}$. Ces
morphismes satisfont les conditions imposées aux morphismes d’une limite projective
(resp.\ d’une limite inductive) de Catégories, Définition \ref{categories-definition-limit}
(resp.\ Catégories, Définition \ref{categories-definition-colimit}).
En effet, ils forment manifestement un cône, resp. un cocône, sur $\mathcal{F}$.
En outre, si $(\mathcal{G}, q_i : \mathcal{G} \to \mathcal{F}_i)$
est un autre
système (comme dans la définition d’une limite projective), alors, pour tout
$U$, nous obtenons un système de morphismes $\mathcal{G}(U) \to \mathcal{F}_i(U)$
satisfaisant les conditions de fonctorialité appropriées. Il existe donc un unique
morphisme $\mathcal{G}(U) \to \mathcal{F}_{\lim}(U)$. On vérifie facilement
que ces morphismes sont compatibles lorsque $U$ varie et proviennent du
morphisme recherché $\mathcal{G} \to \mathcal{F}_{\lim}$.
Un argument analogue vaut dans le cas de la limite inductive.





















\section{Fonctorialité des catégories de préfaisceaux}
\label{section-functoriality-PSh}

\noindent
Soit $u : \mathcal{C} \to \mathcal{D}$ un foncteur entre catégories.
Dans ce cas, nous notons
$$
u^p :
\textit{PSh}(\mathcal{D})
\longrightarrow
\textit{PSh}(\mathcal{C})
$$
le foncteur qui associe à $\mathcal{G}$ sur $\mathcal{D}$ le préfaisceau
$u^p\mathcal{G} = \mathcal{G} \circ u$. Remarquons, d’après la section précédente,
que ce foncteur commute à toutes les limites projectives et inductives.

\medskip\noindent
Pour $V \in \Ob(\mathcal{D})$, notons $\mathcal{I}^u_V$
la catégorie définie par
\begin{equation}
\label{equation-colim-category}
\begin{matrix}
\Ob(\mathcal{I}^u_V)
&
=
&
\{
(U, \phi)
\mid
U \in \Ob(\mathcal{C}),
\phi : V \to u(U)
\}
\\
\Mor_{\mathcal{I}^u_V}((U, \phi), (U', \phi'))
&
=
&
\{
f : U \to U' \text{ dans }\mathcal{C}
\mid
u(f) \circ \phi = \phi'
\}
\end{matrix}
\end{equation}
Nous omettons parfois l’exposant ${}^u$ de la notation et écrivons simplement
$\mathcal{I}_V$.
Nous utiliserons ces catégories pour définir un adjoint à gauche du foncteur $u^p$.
Avant cela, démontrons quelques lemmes techniques.

\begin{lemma}
\label{lemma-almost-directed}
Soit $u : \mathcal{C} \to \mathcal{D}$ un foncteur entre catégories.
Supposons que $\mathcal{C}$ admette les produits fibrés et les égalisateurs, et que
$u$ commute à ceux-ci. Alors les catégories $(\mathcal{I}_V)^{opp}$
satisfont les hypothèses du
lemme de Catégories \ref{categories-lemma-split-into-directed}.
\end{lemma}

\begin{proof}
Il y a deux conditions à vérifier.

\medskip\noindent
Premièrement, supposons donnés trois objets
$\phi : V \to u(U)$, $\phi' : V \to u(U')$ et $\phi'' : V \to u(U'')$,
ainsi que des morphismes $a : U' \to U$, $b : U'' \to U$ tels que
$u(a) \circ \phi' = \phi$ et $u(b) \circ \phi'' = \phi$.
Nous devons montrer qu’il existe un autre objet $\phi''' : V \to u(U''')$
et des morphismes $c : U''' \to U'$ et $d : U''' \to U''$ tels que
$u(c) \circ \phi''' = \phi'$, $u(d) \circ \phi''' = \phi''$ et
$a \circ c = b \circ d$. Prenons $U''' = U' \times_U U''$,
$c$ et $d$ étant les morphismes de projection. Cela convient puisque $u$ commute
aux produits fibrés ; nous omettons la vérification.

\medskip\noindent
Deuxièmement, supposons donnés deux objets
$\phi : V \to u(U)$ et $\phi' : V \to u(U')$,
ainsi que des morphismes $a, b : (U, \phi) \to (U', \phi')$.
Nous devons trouver un morphisme $c : (U'', \phi'') \to (U, \phi)$
qui égalise $a$ et $b$. Soit $c : U'' \to U$ l’égalisateur de
$a$ et $b$ dans la catégorie $\mathcal{C}$. Puisque $u$ commute
aux égalisateurs et que $u(a) \circ \phi = u(b) \circ \phi = \phi'$,
nous obtenons un morphisme $\phi'' : V \to u(U'')$.
\end{proof}

\begin{lemma}
\label{lemma-directed}
Soit $u : \mathcal{C} \to \mathcal{D}$ un foncteur entre catégories.
Supposons que
\begin{enumerate}
\item la catégorie $\mathcal{C}$ possède un objet final $X$ et que
$u(X)$ soit un objet final de $\mathcal{D}$, et
\item la catégorie $\mathcal{C}$ admette les produits fibrés et que
$u$ commute à ceux-ci.
\end{enumerate}
Alors les catégories d’indices $(\mathcal{I}^u_V)^{opp}$ sont filtrantes (voir
Catégories, Définition \ref{categories-definition-directed}).
\end{lemma}

\begin{proof}
Les hypothèses impliquent que celles du
Lemme \ref{lemma-almost-directed}
sont satisfaites (voir la discussion de
Catégories, section \ref{categories-section-finite-limits}).
D’après le
lemme de Catégories \ref{categories-lemma-split-into-directed},
$\mathcal{I}_V$ est une réunion disjointe (éventuellement vide) de
catégories filtrantes.
Il suffit donc de montrer que $\mathcal{I}_V$ est connexe.

\medskip\noindent
Montrons d’abord que $\mathcal{I}_V$ est non vide.
Soit en effet $X$ l’objet final de $\mathcal{C}$,
qui existe par hypothèse.
Soit $V \to u(X)$ l’unique morphisme donné par le fait
que $u(X)$ est final dans $\mathcal{D}$ par hypothèse.
Cela fournit un objet de $\mathcal{I}_V$.

\medskip\noindent
Montrons ensuite que $\mathcal{I}_V$ est connexe.
Soient $\phi_1 : V \to u(U_1)$ et $\phi_2 : V \to u(U_2)$
des objets de $\Ob(\mathcal{I}_V)$. Par hypothèse, $U_1\times U_2$
existe et $u(U_1\times U_2) = u(U_1)\times u(U_2)$.
Considérons le morphisme $\phi : V \to u(U_1\times U_2)$
correspondant à $(\phi_1, \phi_2)$ par la propriété universelle
des produits. L’objet $\phi : V \to u(U_1\times U_2)$
admet manifestement des morphismes vers $\phi_1 : V \to u(U_1)$ et $\phi_2 : V \to u(U_2)$.
\end{proof}

\noindent
Étant donné $g : V' \to V$ dans $\mathcal{D}$, nous obtenons un foncteur
$\overline{g} : \mathcal{I}_V \to \mathcal{I}_{V'}$
en posant $\overline{g}(U, \phi) = (U, \phi \circ g)$
sur les objets. Étant donné un préfaisceau $\mathcal{F}$ sur $\mathcal{C}$,
nous obtenons un foncteur
$$
\mathcal{F}_V :
\mathcal{I}_V^{opp}
\longrightarrow
\textit{Ensembles}, \quad
(U, \phi)
\longmapsto
\mathcal{F}(U).
$$
Autrement dit, $\mathcal{F}_V$ est un préfaisceau d’ensembles sur $\mathcal{I}_V$.
Remarquons que $\mathcal{F}_{V'} \circ \overline{g} = \mathcal{F}_V$.
Nous définissons
$$
u_p\mathcal{F}(V) =
\colim_{\mathcal{I}_V^{opp}} \mathcal{F}_V
$$
La limite inductive fournit, pour chaque $(U, \phi) \in \Ob(\mathcal{I}_V)$,
une application canonique $\mathcal{F}(U)\xrightarrow{c(\phi)}u_p\mathcal{F}(V)$.
Pour $g : V' \to V$ comme ci-dessus, il existe une
application canonique de restriction
$g^* : u_p\mathcal{F}(V) \to u_p\mathcal{F}(V')$
compatible avec
$\mathcal{F}_{V'} \circ \overline{g} = \mathcal{F}_V$
d’après Catégories, Lemme \ref{categories-lemma-functorial-colimit}.
C’est l’unique application telle que, pour tout $(U, \phi) \in \Ob(\mathcal{I}_V)$,
le diagramme
$$
\xymatrix{
\mathcal{F}(U) \ar[r]^{c(\phi)} \ar[d]_{\text{id}}
&
u_p\mathcal{F}(V) \ar[d]^{g^*}
\\
\mathcal{F}(U) \ar[r]^{c(\phi \circ g)}
&
u_p\mathcal{F}(V')
}
$$
soit commutatif. L’unicité de ces applications entraîne que nous obtenons un
préfaisceau. Ce préfaisceau sera noté $u_p\mathcal{F}$.

\begin{lemma}
\label{lemma-recover}
Il existe une application canonique
$\mathcal{F}(U) \to u_p\mathcal{F}(u(U))$,
compatible avec les applications de restriction
(de $\mathcal{F}$ et de $u_p\mathcal{F}$).
\end{lemma}

\begin{proof}
C’est simplement l’application $c(\text{id}_{u(U)})$ introduite ci-dessus.
\end{proof}

\noindent
Remarquons que tout morphisme de préfaisceaux $\mathcal{F} \to \mathcal{F}'$
donne naissance à des systèmes compatibles de morphismes entre foncteurs
$\mathcal{F}_V \to \mathcal{F}'_V$, et donc à un morphisme
de préfaisceaux $u_p\mathcal{F} \to u_p\mathcal{F}'$. Autrement
dit, nous avons défini un foncteur
$$
u_p :
\textit{PSh}(\mathcal{C})
\longrightarrow
\textit{PSh}(\mathcal{D})
$$

\begin{lemma}
\label{lemma-adjoints-u}
Le foncteur $u_p$ est adjoint à gauche du foncteur $u^p$.
Autrement dit, la formule
$$
\Mor_{\textit{PSh}(\mathcal{C})}(\mathcal{F}, u^p\mathcal{G})
=
\Mor_{\textit{PSh}(\mathcal{D})}(u_p\mathcal{F}, \mathcal{G})
$$
est bifonctorielle en $\mathcal{F}$ et $\mathcal{G}$.
\end{lemma}

\begin{proof}
Soit $\mathcal{G}$ un préfaisceau sur $\mathcal{D}$ et soit
$\mathcal{F}$ un préfaisceau sur $\mathcal{C}$.
Nous allons démontrer la formule affichée
en construisant des applications dans les deux sens. Nous laissons
au lecteur le soin de vérifier qu’elles sont inverses l’une de l’autre.

\medskip\noindent
Étant donné un morphisme $\alpha : u_p \mathcal{F} \to \mathcal{G}$,
nous obtenons $u^p\alpha : u^p u_p \mathcal{F} \to u^p \mathcal{G}$.
Le Lemme \ref{lemma-recover} affirme qu’il existe un
morphisme $\mathcal{F} \to u^p u_p \mathcal{F}$. La composée
des deux donne le morphisme recherché. (L’intérêt de cette construction
est qu’elle est manifestement fonctorielle en toutes les données en présence.)

\medskip\noindent
Réciproquement, étant donné un morphisme $\beta : \mathcal{F} \to u^p\mathcal{G}$,
nous obtenons un morphisme $u_p\beta : u_p\mathcal{F} \to u_p u^p\mathcal{G}$.
Soit $V \in \Ob(\mathcal{D})$.
Nous affirmons que le foncteur $u^p\mathcal{G}_V$ sur $\mathcal{I}_V$
admet un morphisme canonique vers le foncteur constant de valeur $\mathcal{G}(V)$.
En effet, pour tout objet $(U, \phi)$ de $\mathcal{I}_V$,
la valeur de $u^p\mathcal{G}_V$ en cet objet est $\mathcal{G}(u(U))$,
qui s’envoie dans $\mathcal{G}(V)$ par $\mathcal{G}(\phi) = \phi^*$.
C’est une transformation naturelle de foncteurs parce que $\mathcal{G}$ est lui-même
un foncteur. Cela fournit un morphisme $u_p u^p \mathcal{G}(V) \to \mathcal{G}(V)$.
Une autre vérification immédiate montre que celui-ci est fonctoriel en $V$,
ce qui donne un morphisme de préfaisceaux $u_p u^p \mathcal{G} \to \mathcal{G}$.
La composée $u_p\mathcal{F} \to u_p u^p\mathcal{G} \to
\mathcal{G}$ est le morphisme recherché.
\end{proof}

\begin{remark}
\label{remark-functoriality-presheaves-values}
Supposons que $\mathcal{A}$ soit une catégorie telle que
tout diagramme $\mathcal{I}_Y \to \mathcal{A}$ admette une
limite inductive dans $\mathcal{A}$. Il est alors clair
qu’il existe des foncteurs $u^p$ et $u_p$, définis
exactement comme ci-dessus, entre les catégories
de préfaisceaux à valeurs dans $\mathcal{A}$.
De plus, l’adjonction du couple
$u^p$ et $u_p$ subsiste dans ce cadre.
\end{remark}

\begin{lemma}
\label{lemma-pullback-representable-presheaf}
Soit $u : \mathcal{C} \to \mathcal{D}$ un foncteur entre catégories.
Pour tout objet $U$ de $\mathcal{C}$, nous avons $u_ph_U = h_{u(U)}$.
\end{lemma}

\begin{proof}
Par l’adjonction de $u_p$ et $u^p$, nous avons
$$
\Mor_{\textit{PSh}(\mathcal{D})}(u_ph_U, \mathcal{G})
=
\Mor_{\textit{PSh}(\mathcal{C})}(h_U, u^p\mathcal{G})
=
u^p\mathcal{G}(U) =
\mathcal{G}(u(U))
$$
et le lemme de Yoneda montre donc que $u_ph_U = h_{u(U)}$ en tant que
préfaisceaux.
\end{proof}
   
    
 
 
 
  











\section{Sites}
\label{section-sites-definitions}

\noindent
Notre notion de site fait appel au type de structures suivant.

\begin{definition}
\label{definition-family-morphisms-fixed-target}
Soit $\mathcal{C}$ une catégorie ; voir Conventions,
section \ref{conventions-section-categories}.
Une {\it famille de morphismes à cible fixée} dans $\mathcal{C}$ est donnée
par un objet $U \in \Ob(\mathcal{C})$, un ensemble $I$ et, pour chaque
$i\in I$, un morphisme $U_i \to U$ de $\mathcal{C}$ de cible $U$.
Nous la noterons $\{U_i \to U\}_{i\in I}$.
\end{definition}

\noindent
Il se peut que l’ensemble $I$ soit vide ! Cette notation vise à évoquer
un recouvrement ouvert, comme en topologie.

\begin{definition}
\label{definition-site}
Un {\it site}\footnote{Cette notation diffère de celle de \cite{SGA4},
comme il est expliqué dans l’introduction.} est la donnée d’une catégorie
$\mathcal{C}$ et d’un ensemble $\text{Cov}(\mathcal{C})$ de familles de
morphismes à cible fixée $\{U_i \to U\}_{i \in I}$, appelées
{\it recouvrements de $\mathcal{C}$}, qui satisfont aux axiomes suivants :
\begin{enumerate}
\item Si $V \to U$ est un isomorphisme, alors $\{V \to U\} \in
\text{Cov}(\mathcal{C})$.
\item Si $\{U_i \to U\}_{i\in I} \in \text{Cov}(\mathcal{C})$ et si, pour
chaque $i$, on a $\{V_{ij} \to U_i\}_{j\in J_i} \in
\text{Cov}(\mathcal{C})$, alors
$\{V_{ij} \to U\}_{i \in I, j\in J_i} \in \text{Cov}(\mathcal{C})$.
\item Si $\{U_i \to U\}_{i\in I}\in \text{Cov}(\mathcal{C})$
et si $V \to U$ est un morphisme de $\mathcal{C}$, alors
$U_i \times_U V$ existe pour tout $i$ et
$\{U_i \times_U V \to V \}_{i\in I} \in \text{Cov}(\mathcal{C})$.
\end{enumerate}
\end{definition}

\noindent
Précisions. Dans l’axiome (1), on exige qu’il existe un recouvrement
$\{U_i \to U\}_{i \in I}$ de $\mathcal{C}$ tel que $I = \{i\}$ soit un
singleton et que le morphisme $U_i \to U$ soit égal au morphisme
$V \to U$ donné en (1). Dans la suite, nous désignerons souvent par
$\{V \to U\}$ une famille de morphismes à cible fixée dont l’ensemble
d’indices est un singleton. Dans l’axiome (3), on exige l’existence du
recouvrement pour un certain choix des produits fibrés
$U_i \times_U V$, $i \in I$.

\begin{remark}
\label{remark-no-big-sites}
(À propos des questions ensemblistes ; à omettre en première lecture.)
La raison principale pour laquelle nous introduisons les sites est l’étude
de la catégorie des faisceaux sur un site, car elle généralise la catégorie
des faisceaux sur un espace topologique, qui joue un rôle si important en
géométrie algébrique. Afin de ne pas avoir à considérer des objets tels que
des « classes de classes », et ainsi de suite, nous n’autoriserons pas les
sites à être des catégories « grandes », contrairement à ce que nous faisons
pour les catégories et les $2$-catégories.

\medskip\noindent
Supposons que $\mathcal{C}$ soit une catégorie et que
$\text{Cov}(\mathcal{C})$ soit une classe propre de recouvrements
satisfaisant (1), (2) et (3) ci-dessus. Nous ne l’autoriserons pas non plus
comme site, principalement parce que nous prendrons des limites sur les
recouvrements. Il existe toutefois plusieurs façons naturelles de remplacer
$\text{Cov}(\mathcal{C})$ par un ensemble de recouvrements, ou par une
structure légèrement différente, qui donnent lieu à la même catégorie de
faisceaux. Par exemple :
\begin{enumerate}
\item Dans Ensembles, section \ref{sets-section-coverings-site}, nous montrons
comment choisir un ensemble approprié de recouvrements qui donne la même
catégorie de faisceaux.
\item On peut aussi prendre la topologie associée (voir la
définition \ref{definition-topology-associated-site}). La topologie ainsi
obtenue sur $\mathcal{C}$ possède la même catégorie de faisceaux. Deux
topologies ont les mêmes catégories de faisceaux si et seulement si elles
sont égales ; voir le théorème \ref{theorem-topology-and-topos}. Une topologie
sur une catégorie est donnée par un choix de cribles sur les objets. La
collection de tous les cribles possibles, et même celle de toutes les
topologies possibles sur $\mathcal{C}$, est un ensemble.
\item On pourrait également modifier légèrement la notion de site ; voir la
remarque \ref{remark-shrink-coverings} ci-dessous. On obtiendrait alors un
ensemble canonique de recouvrements.
\end{enumerate}
Chacune de ces solutions présente un inconvénient mineur. Pour la première,
il faut vérifier que les constructions ultérieures ne dépendent pas du choix
de l’ensemble de recouvrements. Pour la deuxième, il faut se familiariser avec
les topologies et reprendre de nombreux arguments concernant les sites. Pour
la troisième, voir la dernière phrase de la
remarque \ref{remark-shrink-coverings}.

\medskip\noindent
Notre approche consistera à travailler avec les sites de la
définition \ref{definition-site} ci-dessus. Étant donnée une catégorie
$\mathcal{C}$ munie d’une classe propre de recouvrements comme ci-dessus,
nous remplacerons celle-ci par un ensemble de recouvrements donnant un site,
au moyen d’Ensembles, lemme \ref{sets-lemma-coverings-site}. Le
lemme \ref{lemma-choice-set-coverings-immaterial} ci-dessous montre que la
catégorie de faisceaux obtenue (le topos) est indépendante de ce choix. Nous
laissons au lecteur le soin d’employer, s’il le souhaite, l’une des deux
autres stratégies pour traiter ces questions.
\end{remark}

\begin{example}
\label{example-site-topological}
Soit $X$ un espace topologique. Soit $X_{Zar}$ la catégorie dont les objets
sont tous les ouverts $U$ de $X$ et dont les morphismes sont simplement les
applications d’inclusion. Autrement dit, il existe au plus un morphisme entre
deux objets quelconques de $X_{Zar}$. Posons maintenant
$\{U_i \to U\}_{i \in I}\in \text{Cov}(X_{Zar})$ si et seulement si
$\bigcup U_i = U$. Les conditions (1) et (2) ci-dessus sont claires, et (3)
l’est également dès que l’on remarque que, dans $X_{Zar}$, on a
$U \times V = U \cap V$. Notons en particulier que l’ensemble vide doit être
un élément de $X_{Zar}$, faute de quoi cela ne fonctionnerait pas en général.
De plus, comme nous le verrons plus tard, il est tout aussi important
d’autoriser {\it le recouvrement vide de l’ensemble vide comme recouvrement} !
Nous faisons de $X_{Zar}$ un site en choisissant un ensemble approprié de
recouvrements $\text{Cov}(X_{Zar})_{\kappa, \alpha}$ comme dans Ensembles,
lemme \ref{sets-lemma-coverings-site}. Les préfaisceaux et les faisceaux
(au sens défini ci-dessous) sur le site $X_{Zar}$ coïncident exactement avec
les notions usuelles de préfaisceau et de faisceau sur un espace topologique,
définies dans Faisceaux, section \ref{sheaves-section-introduction}.
\end{example}

\begin{example}
\label{example-site-on-group}
Soit $G$ un groupe. Considérons la catégorie $G\textit{-Ensembles}$ dont les objets
sont les ensembles $X$ munis d’une action à gauche de $G$ et dont les
morphismes sont les applications $G$-équivariantes. Un exemple important est
${}_GG$, le $G$-ensemble dont l’ensemble sous-jacent est $G$ et dont l’action
est donnée par la multiplication à gauche. Cette catégorie possède des
produits fibrés ; voir Catégories,
section \ref{categories-section-example-fibre-products}.
Nous déclarons que $\{\varphi_i : U_i \to U\}_{i\in I}$ est un recouvrement si
$\bigcup_{i\in I} \varphi_i(U_i) = U$. On obtient ainsi une classe de
recouvrements sur $G\textit{-Ensembles}$ dont on vérifie aisément qu’elle satisfait
les conditions (1), (2) et (3) de la définition \ref{definition-site}. Le
résultat n’est pas un site, car tant la collection des objets de la catégorie
sous-jacente que la collection des recouvrements forment une classe propre.
Nous remplaçons d’abord $G\textit{-Ensembles}$ par une sous-catégorie pleine
$G\textit{-Ensembles}_\alpha$ comme dans Ensembles,
lemme \ref{sets-lemma-sets-with-group-action}. Le site
$(G\textit{-Ensembles}_\alpha,
\text{Cov}_{\kappa, \alpha'}(G\textit{-Ensembles}_\alpha))$
obtenu ensuite en restreignant convenablement la collection des recouvrements,
comme dans Ensembles, lemme \ref{sets-lemma-coverings-site}, sera noté
$\mathcal{T}_G$.

\medskip\noindent
Dans le cas particulier où le groupe $G$ est dénombrable, on peut prendre
pour $\mathcal{T}_G$ la catégorie des $G$-ensembles dénombrables et pour
recouvrements les familles conjointement surjectives de morphismes
$\{\varphi_i : U_i \to U\}_{i \in I}$ telles que $I$ soit dénombrable.
\end{example}

\begin{example}
\label{example-indiscrete}
Soit $\mathcal{C}$ une catégorie. Il existe une manière canonique d’en faire
un site dans lequel les $\{f : V \to U \mid f\text{ est un isomorphisme}\}$
sont les recouvrements de $U$. Les faisceaux sur ce site sont les préfaisceaux
sur $\mathcal{C}$. La topologie correspondante est appelée {\it topologie
chaotique} ou {\it topologie grossière}.
\end{example}













\section{Faisceaux}
\label{section-sheaves}

\noindent
Soit $\mathcal{C}$ un site. Avant d’introduire la notion de faisceau à valeurs
dans une catégorie, expliquons ce que signifie, pour un préfaisceau
d’ensembles, être un faisceau. Soit $\mathcal{F}$ un préfaisceau d’ensembles
sur $\mathcal{C}$ et soit $\{U_i \to U\}_{i\in I}$ un élément de
$\text{Cov}(\mathcal{C})$. Par hypothèse, tous les produits fibrés
$U_i \times_U U_j$ existent dans $\mathcal{C}$. Il existe deux applications
naturelles
$$
\xymatrix{
\prod\nolimits_{i\in I}
\mathcal{F}(U_i)
\ar@<1ex>[r]^-{\text{pr}_0^*} \ar@<-1ex>[r]_-{\text{pr}_1^*}
&
\prod\nolimits_{(i_0, i_1) \in I \times I}
\mathcal{F}(U_{i_0} \times_U U_{i_1})
}
$$
que nous noterons $\text{pr}^*_i$, $i = 0, 1$, comme indiqué dans la formule
ci-dessus. En effet, un élément du membre de gauche correspond à une famille
$(s_i)_{i\in I}$, où chaque $s_i$ est une section de $\mathcal{F}$ sur $U_i$.
Pour chaque couple $(i_0, i_1) \in I \times I$, on dispose des morphismes de
projection
$$
\text{pr}^{(i_0, i_1)}_{i_0} :
U_{i_0} \times_U U_{i_1}
\longrightarrow
U_{i_0}
\text{ et }
\text{pr}^{(i_0, i_1)}_{i_1} :
U_{i_0} \times_U U_{i_1}
\longrightarrow
U_{i_1}.
$$
On peut donc prendre soit l’image inverse de la section $s_{i_0}$ par la
première de ces applications, soit celle de la section $s_{i_1}$ par la
seconde. Explicitement, les applications mentionnées ci-dessus sont
$$
\text{pr}_0^* :
(s_i)_{i\in I}
\longmapsto
\Big(
\text{pr}^{(i_0, i_1), *}_{i_0}(s_{i_0})
\Big)_{(i_0, i_1) \in I \times I}
$$
et
$$
\text{pr}_1^* :
(s_i)_{i\in I}
\longmapsto
\Big(
\text{pr}^{(i_0, i_1), *}_{i_1}(s_{i_1})
\Big)_{(i_0, i_1) \in I \times I}.
$$
Considérons enfin l’application naturelle
$$
\mathcal{F}(U)
\longrightarrow
\prod\nolimits_{i\in I}
\mathcal{F}(U_i), \quad
s
\longmapsto
(s|_{U_i})_{i \in I}
$$
où la notation $s|_{U_i}$ désigne l’image inverse de $s$ par l’application
$U_i \to U$. La fonctorialité de $\mathcal{F}$ et la commutativité des
diagrammes de produits fibrés montrent clairement que
$\text{pr}_0^*( (s|_{U_i})_{i \in I} ) =
\text{pr}_1^*( (s|_{U_i})_{i \in I} )$.

\begin{definition}
\label{definition-sheaf-sets}
Soit $\mathcal{C}$ un site et soit $\mathcal{F}$ un préfaisceau d’ensembles
sur $\mathcal{C}$. On dit que $\mathcal{F}$ est un {\it faisceau} si, pour
tout recouvrement $\{U_i \to U\}_{i \in I} \in
\text{Cov}(\mathcal{C})$, le diagramme
\begin{equation}
\label{equation-sheaf-condition}
\xymatrix{
\mathcal{F}(U) \ar[r]
&
\prod\nolimits_{i\in I}
\mathcal{F}(U_i)
\ar@<1ex>[r]^-{\text{pr}_0^*} \ar@<-1ex>[r]_-{\text{pr}_1^*}
&
\prod\nolimits_{(i_0, i_1) \in I \times I}
\mathcal{F}(U_{i_0} \times_U U_{i_1})
}
\end{equation}
présente la première flèche comme l’égalisateur de $\text{pr}_0^*$ et
$\text{pr}_1^*$.
\end{definition}

\noindent
Autrement dit, si l’on se donne des sections
$s_i \in \mathcal{F}(U_i)$ telles que
$$
s_i|_{U_i \times_U U_j} = s_j|_{U_i \times_U U_j}
$$
dans $\mathcal{F}(U_i \times_U U_j)$ pour tout couple
$(i, j) \in I \times I$, alors il existe un unique
$s \in \mathcal{F}(U)$ tel que $s_i = s|_{U_i}$.

\begin{remark}
\label{remark-sheaf-condition-empty-covering}
Si le recouvrement $\{U_i \to U\}_{i \in I}$ est la famille vide
(c’est-à-dire si $I = \emptyset$), la condition de faisceau signifie que
$\mathcal{F}(U) = \{*\}$ est un singleton. En effet, dans
(\ref{equation-sheaf-condition}), les deuxième et troisième ensembles sont
des produits vides dans la catégorie des ensembles ; ce sont des objets
finaux de cette catégorie, donc des singletons.
\end{remark}

\begin{example}
\label{example-sheaves-topological}
Soit $X$ un espace topologique. Soit $X_{Zar}$ le site construit dans
l’exemple \ref{example-site-topological}. La notion de faisceau sur $X_{Zar}$
coïncide avec celle de faisceau sur $X$ introduite dans Faisceaux,
définition \ref{sheaves-definition-sheaf}.
\end{example}

\begin{example}
\label{example-topological-wrong}
Soit $X$ un espace topologique. Considérons le site $X'_{Zar}$, identique au
site $X_{Zar}$ de l’exemple \ref{example-site-topological}, à ceci près que
nous n’autorisons pas le recouvrement vide de l’ensemble vide. Autrement dit,
nous autorisons le recouvrement $\{\emptyset \to \emptyset\}$, mais non le
recouvrement dont l’ensemble d’indices est vide. On vérifie aisément que cela
définit encore un site. Nous affirmons toutefois que les faisceaux sur
$X'_{Zar}$ diffèrent de ceux sur $X_{Zar}$. Comme cas extrême, considérons par
exemple la situation où $X = \{p\}$ est un singleton. Les objets de
$X'_{Zar}$ sont alors $\emptyset, X$, tout recouvrement de $\emptyset$ admet
$\{\emptyset \to \emptyset\}$ comme raffinement et tout recouvrement de $X$
admet $\{X \to X\}$ comme raffinement. Un faisceau sur ce site est clairement
donné par un choix arbitraire d’un ensemble $\mathcal{F}(\emptyset)$ et d’un
ensemble $\mathcal{F}(X)$, ainsi que d’une application de restriction
$\mathcal{F}(X) \to \mathcal{F}(\emptyset)$. Ainsi, les faisceaux sur
$X'_{Zar}$ sont les faisceaux usuels sur l’espace à deux points
$\{\eta, p\}$ dont les ouverts sont
$\{\emptyset, \{\eta\}, \{p, \eta\}\}$. En général, les faisceaux sur
$X'_{Zar}$ sont les faisceaux sur l’espace $X \amalg \{\eta\}$ dont les
ouverts sont l’ensemble vide et tous les ensembles de la forme
$U \cup \{\eta\}$, où $U \subset X$ est ouvert.
\end{example}


\begin{definition}
\label{definition-category-sheaves-sets}
La catégorie {\it $\Sh(\mathcal{C})$} des faisceaux d’ensembles est la
sous-catégorie pleine de la catégorie $\textit{PSh}(\mathcal{C})$ dont les
objets sont les faisceaux d’ensembles.
\end{definition}

\noindent
Soit $\mathcal{A}$ une catégorie. Si les produits indexés par $I$ et par
$I \times I$ existent dans $\mathcal{A}$ pour tout $I$ qui intervient comme
ensemble d’indices d’une famille couvrante, alors la
définition \ref{definition-sheaf-sets} ci-dessus a un sens et définit une
notion de faisceau sur $\mathcal{C}$ à valeurs dans $\mathcal{A}$. Remarquons
que le diagramme dans $\mathcal{A}$
$$
\xymatrix{
\mathcal{F}(U) \ar[r]
&
\prod\nolimits_{i\in I}
\mathcal{F}(U_i)
\ar@<1ex>[r]^-{\text{pr}_0^*} \ar@<-1ex>[r]_-{\text{pr}_1^*}
&
\prod\nolimits_{(i_0, i_1) \in I \times I}
\mathcal{F}(U_{i_0} \times_U U_{i_1})
}
$$
est un diagramme égalisateur si et seulement si, pour tout objet $X$ de
$\mathcal{A}$, le diagramme d’ensembles
$$
\xymatrix{
\Mor_\mathcal{A}(X, \mathcal{F}(U)) \ar[r]
&
\prod
\Mor_\mathcal{A}(X, \mathcal{F}(U_i))
\ar@<1ex>[r]^-{\text{pr}_0^*} \ar@<-1ex>[r]_-{\text{pr}_1^*}
&
\prod
\Mor_\mathcal{A}(X, \mathcal{F}(U_{i_0} \times_U U_{i_1}))
}
$$
est un diagramme égalisateur.

\medskip\noindent
Supposons $\mathcal{A}$ arbitraire. Soit $\mathcal{F}$ un préfaisceau à
valeurs dans $\mathcal{A}$. Choisissons un objet quelconque
$X\in \Ob(\mathcal{A})$. On obtient alors un préfaisceau d’ensembles
$\mathcal{F}_X$ défini par la règle
$$
\mathcal{F}_X(U) = \Mor_\mathcal{A}(X, \mathcal{F}(U)).
$$
Il résulte de ce qui précède qu’une bonne définition consiste à exiger que
tous les préfaisceaux $\mathcal{F}_X$ soient des faisceaux d’ensembles.

\begin{definition}
\label{definition-sheaf}
Soient $\mathcal{C}$ un site, $\mathcal{A}$ une catégorie et $\mathcal{F}$ un
préfaisceau sur $\mathcal{C}$ à valeurs dans $\mathcal{A}$. On dit que
$\mathcal{F}$ est un {\it faisceau} si, pour tout objet $X$ de $\mathcal{A}$,
le préfaisceau d’ensembles $\mathcal{F}_X$ (défini ci-dessus) est un faisceau.
\end{definition}










\section{Familles de morphismes à cible fixée}
\label{section-refinements}

\noindent
Cette section introduit quelques notions relatives aux familles de morphismes
ayant la même cible.

\begin{definition}
\label{definition-morphism-coverings}
Soit $\mathcal{C}$ une catégorie. Soit
$\mathcal{U} = \{U_i \to U\}_{i\in I}$ une famille de morphismes de
$\mathcal{C}$ à cible fixée. Soit
$\mathcal{V} = \{V_j \to V\}_{j\in J}$ une autre telle famille.
\begin{enumerate}
\item
Un {\it morphisme de familles de morphismes à cible fixée de $\mathcal{C}$ de
$\mathcal{U}$ vers $\mathcal{V}$}, ou simplement un {\it morphisme de
$\mathcal{U}$ vers $\mathcal{V}$}, est donné par un morphisme $U \to V$, une
application d’ensembles $\alpha : I \to J$ et, pour chaque $i\in I$, un
morphisme $U_i \to V_{\alpha(i)}$ tel que le diagramme
$$
\xymatrix{
U_i \ar[r] \ar[d]
&
V_{\alpha(i)} \ar[d]
\\
U \ar[r]
&
V
}
$$
soit commutatif.
\item Dans le cas particulier où $U = V$ et où $U \to V$ est l’identité,
nous disons que $\mathcal{U}$ est un {\it raffinement} de la famille
$\mathcal{V}$.
\end{enumerate}
\end{definition}

\noindent
Une remarque triviale mais importante est que, si
$\mathcal{V} = \{V_j \to V\}_{j \in J}$ est la {\it famille vide de
morphismes}, c’est-à-dire si $J = \emptyset$, aucune famille
$\mathcal{U} = \{U_i \to V\}_{i \in I}$ avec $I \not = \emptyset$ ne peut
raffiner $\mathcal{V}$ !

\begin{definition}
\label{definition-combinatorial-tautological}
Soit $\mathcal{C}$ une catégorie. Soient
$\mathcal{U} = \{\varphi_i : U_i \to U\}_{i\in I}$ et
$\mathcal{V} = \{\psi_j : V_j \to U\}_{j\in J}$ deux familles de morphismes
à cible fixée.
\begin{enumerate}
\item Nous disons que $\mathcal{U}$ et $\mathcal{V}$ sont
{\it combinatoirement équivalentes} s’il existe des applications
$\alpha : I \to J$ et $\beta : J\to I$ telles que
$\varphi_i = \psi_{\alpha(i)}$ et $\psi_j = \varphi_{\beta(j)}$.
\item Nous disons que $\mathcal{U}$ et $\mathcal{V}$ sont
{\it tautologiquement équivalentes} s’il existe des applications
$\alpha : I \to J$ et $\beta : J\to I$ ainsi que, pour tout $i\in I$ et tout
$j \in J$, des diagrammes commutatifs
$$
\xymatrix{
U_i \ar[rd] \ar[rr] & &
V_{\alpha(i)} \ar[ld] & &
V_j \ar[rd] \ar[rr] & &
U_{\beta(j)} \ar[ld] \\
&
U & & & &
U &
}
$$
dont les flèches horizontales sont des isomorphismes.
\end{enumerate}
\end{definition}

\begin{lemma}
\label{lemma-tautological-combinatorial}
Soit $\mathcal{C}$ une catégorie. Soient
$\mathcal{U} = \{\varphi_i : U_i \to U\}_{i\in I}$ et
$\mathcal{V} = \{\psi_j : V_j \to U\}_{j\in J}$ deux familles de morphismes
ayant la même cible fixée.
\begin{enumerate}
\item Si $\mathcal{U}$ et $\mathcal{V}$ sont combinatoirement équivalentes,
alors elles sont tautologiquement équivalentes.
\item Si $\mathcal{U}$ et $\mathcal{V}$ sont tautologiquement équivalentes,
alors $\mathcal{U}$ est un raffinement de $\mathcal{V}$ et $\mathcal{V}$ est
un raffinement de $\mathcal{U}$.
\item La relation « être combinatoirement équivalentes » est une relation
d’équivalence sur l’ensemble des familles de morphismes à cible fixée.
\item La relation « être tautologiquement équivalentes » est une relation
d’équivalence sur l’ensemble des familles de morphismes à cible fixée.
\item La relation « $\mathcal{U}$ raffine $\mathcal{V}$ et $\mathcal{V}$
raffine $\mathcal{U}$ » est une relation d’équivalence sur l’ensemble des
familles de morphismes à cible fixée.
\end{enumerate}
\end{lemma}

\begin{proof}
Omis.
\end{proof}

\noindent
Dans le lemme suivant, étant donnés une catégorie $\mathcal{C}$, un
préfaisceau $\mathcal{F}$ sur $\mathcal{C}$ et une famille
$\mathcal{U} = \{U_i \to U\}_{i\in I}$ telle que tous les produits fibrés
$U_i \times_U U_{i'}$ existent, nous dirons que {\it la condition de faisceau
pour $\mathcal{F}$ relativement à $\mathcal{U}$} est satisfaite si le
diagramme (\ref{equation-sheaf-condition}) est un diagramme égalisateur.

\begin{lemma}
\label{lemma-tautological-same-sheaf}
Soit $\mathcal{C}$ une catégorie. Soient
$\mathcal{U} = \{\varphi_i : U_i \to U\}_{i\in I}$ et
$\mathcal{V} = \{\psi_j : V_j \to U\}_{j\in J}$ deux familles de morphismes
ayant la même cible fixée. Supposons que les produits fibrés
$U_i \times_U U_{i'}$ et $V_j \times_U V_{j'}$ existent. Si $\mathcal{U}$ et
$\mathcal{V}$ sont tautologiquement équivalentes, alors, pour tout
préfaisceau $\mathcal{F}$ sur $\mathcal{C}$, la condition de faisceau pour
$\mathcal{F}$ relativement à $\mathcal{U}$ équivaut à la condition de
faisceau pour $\mathcal{F}$ relativement à $\mathcal{V}$.
\end{lemma}

\begin{proof}
Remarquons d’abord que, si $\varphi : A \to B$ est un isomorphisme dans la
catégorie $\mathcal{C}$, alors
$\varphi^* : \mathcal{F}(B) \to \mathcal{F}(A)$ est un isomorphisme. Soit
$\beta : J \to I$ une application et soient
$\chi_j : V_j \to U_{\beta(j)}$ des isomorphismes au-dessus de $U$, dont
l’existence résulte de l’hypothèse. Montrons que la condition de faisceau pour
$\mathcal{V}$ implique celle pour $\mathcal{U}$. Supposons données des
sections $s_i \in \mathcal{F}(U_i)$ telles que
$$
s_i|_{U_i \times_U U_{i'}} = s_{i'}|_{U_i \times_U U_{i'}}
$$
dans $\mathcal{F}(U_i \times_U U_{i'})$ pour tout couple
$(i, i') \in I \times I$. On peut alors poser
$s_j = \chi_j^*s_{\beta(j)}$. Pour tout couple
$(j, j') \in J \times J$, le morphisme
$\chi_j \times_{\text{id}_U} \chi_{j'} : V_j \times_U V_{j'} \to
U_{\beta(j)} \times_U U_{\beta(j')}$ est lui aussi un isomorphisme. Par
transport de structure, on voit donc que
$$
s_j|_{V_j \times_U V_{j'}} = s_{j'}|_{V_j \times_U V_{j'}}
$$
également. La condition de faisceau relativement à $\mathcal{V}$ implique
qu’il existe un unique $s$ tel que $s|_{V_j} = s_j$ pour tout $j \in J$.
D’après la première observation de la démonstration, cela implique aussi que
$s|_{U_i} = s_i$ pour tout $i \in \Im(\beta)$. Supposons que $i \in I$ et
$i \not \in \Im(\beta)$. Pour un tel $i$, on dispose d’isomorphismes
$U_i \to V_{\alpha(i)} \to U_{\beta(\alpha(i))}$ au-dessus de $U$. Cela donne
un morphisme $U_i \to U_i \times_U U_{\beta(\alpha(i))}$ qui est une section
de la projection. Puisque $s_i$ et $s_{\beta(\alpha(i))}$ ont la même
restriction au produit fibré, on en conclut que
$s_{\beta(\alpha(i))}$ a pour image inverse $s_i$ par le morphisme
$U_i \to U_{\beta(\alpha(i))}$. Ainsi,
on a également $s_i = s|_{U_i}$, comme voulu.
\end{proof}

\begin{lemma}
\label{lemma-compare-separated-presheaf-condition}
Soit $\mathcal{C}$ une catégorie. Soit
$\mathcal{V} = \{V_j \to U\}_{j \in J} \to
\mathcal{U} = \{U_i \to U\}_{i \in I}$ un morphisme de familles de
morphismes à cible fixée de $\mathcal{C}$ donné par
$\text{id} : U \to U$, $\alpha : J \to I$ et
$f_j : V_j \to U_{\alpha(j)}$. Soit $\mathcal{F}$ un préfaisceau sur
$\mathcal{C}$. Si
$\mathcal{F}(U) \to \prod_{j \in J} \mathcal{F}(V_j)$ est injective, alors
$\mathcal{F}(U) \to \prod_{i \in I} \mathcal{F}(U_i)$ est injective.
\end{lemma}

\begin{proof}
Omis.
\end{proof}

\begin{lemma}
\label{lemma-compare-sheaf-condition}
Soit $\mathcal{C}$ une catégorie. Soit
$\mathcal{V} = \{V_j \to U\}_{j \in J} \to
\mathcal{U} = \{U_i \to U\}_{i \in I}$ un morphisme de familles de
morphismes à cible fixée de $\mathcal{C}$ donné par
$\text{id} : U \to U$, $\alpha : J \to I$ et
$f_j : V_j \to U_{\alpha(j)}$. Soit $\mathcal{F}$ un préfaisceau sur
$\mathcal{C}$. Si
\begin{enumerate}
\item les produits fibrés $U_i \times_U U_{i'}$, $U_i \times_U V_j$ et
$V_j \times_U V_{j'}$ existent,
\item $\mathcal{F}$ satisfait à la condition de faisceau relativement à
$\mathcal{V}$, et
\item pour tout $i \in I$, l’application
$\mathcal{F}(U_i) \to \prod_{j \in J} \mathcal{F}(V_j \times_U U_i)$ est
injective,
\end{enumerate}
alors $\mathcal{F}$ satisfait à la condition de faisceau relativement à
$\mathcal{U}$.
\end{lemma}

\begin{proof}
D’après le lemme \ref{lemma-compare-separated-presheaf-condition},
l’application $\mathcal{F}(U) \to \prod \mathcal{F}(U_i)$ est injective.
Supposons donnés des $s_i \in \mathcal{F}(U_i)$ tels que
$s_i|_{U_i \times_U U_{i'}} = s_{i'}|_{U_i \times_U U_{i'}}$ pour tous
$i, i' \in I$. Posons
$s_j = f_j^*(s_{\alpha(j)}) \in \mathcal{F}(V_j)$. Puisque les morphismes
$f_j$ sont des morphismes au-dessus de $U$, on obtient des morphismes induits
$f_{jj'} : V_j \times_U V_{j'} \to
U_{\alpha(i)} \times_U U_{\alpha(i')}$ compatibles avec $f_j, f_{j'}$ par
les morphismes de projection. Il s’ensuit que
$$
s_j|_{V_j \times_U V_{j'}}
= f_{jj'}^*(s_{\alpha(j)}|_{U_{\alpha(j)} \times_U U_{\alpha(j')}})
= f_{jj'}^*(s_{\alpha(j')}|_{U_{\alpha(j)} \times_U U_{\alpha(j')}})
= s_{j'}|_{V_j \times_U V_{j'}}
$$
pour tous $j, j' \in J$. La condition de faisceau pour $\mathcal{F}$
relativement à $\mathcal{V}$ fournit donc une section
$s \in \mathcal{F}(U)$ dont la restriction est $s_j$ sur chaque $V_j$. Il
suffit de montrer que $s$ se restreint en $s_i$ sur $U_i$, quel que soit
$i \in I$. Comme $\mathcal{F}$ vérifie (3), il suffit de montrer que $s$ et
$s_i$ ont la même restriction à $U_i \times_U V_j$ pour tout $j \in J$.
Pour cela, on utilise
$$
s|_{U_i \times_U V_j} = s_j|_{U_i \times_U V_j} =
(\text{id} \times f_j)^*s_{\alpha(j)}|_{U_i \times_U U_{\alpha(j)}} =
(\text{id} \times f_j)^*s_i|_{U_i \times_U U_{\alpha(j)}} =
s_i|_{U_i \times_U V_j}
$$
comme voulu.
\end{proof}

\begin{lemma}
\label{lemma-refine-same-topology}
Soit $\mathcal{C}$ une catégorie. Soient $\text{Cov}_i$, $i = 1, 2$, deux
ensembles de familles de morphismes à cible fixée qui définissent chacun une
structure de site sur $\mathcal{C}$.
\begin{enumerate}
\item Si tout $\mathcal{U} \in \text{Cov}_1$ est tautologiquement équivalent
à un certain $\mathcal{V} \in \text{Cov}_2$, alors
$\Sh(\mathcal{C}, \text{Cov}_2) \subset
\Sh(\mathcal{C}, \text{Cov}_1)$.
Si, de plus, tout $\mathcal{U} \in \text{Cov}_2$ est tautologiquement
équivalent à un certain $\mathcal{V} \in \text{Cov}_1$, alors les catégories
de faisceaux sont égales.
\item Supposons que, pour tout $\mathcal{U} \in \text{Cov}_1$, il existe un
$\mathcal{V} \in \text{Cov}_2$ tel que $\mathcal{V}$ raffine $\mathcal{U}$.
Dans ce cas,
$\Sh(\mathcal{C}, \text{Cov}_2) \subset
\Sh(\mathcal{C}, \text{Cov}_1)$.
Si, de plus, pour tout $\mathcal{U} \in \text{Cov}_2$, il existe un
$\mathcal{V} \in \text{Cov}_1$ tel que $\mathcal{V}$ raffine $\mathcal{U}$,
alors les catégories de faisceaux sont égales.
\end{enumerate}
\end{lemma}

\begin{proof}
La partie (1) résulte directement du
lemme \ref{lemma-tautological-same-sheaf} et des définitions.

\medskip\noindent
Démonstration de (2). Soit $\mathcal{F}$ un faisceau d’ensembles pour le site
$(\mathcal{C}, \text{Cov}_2)$. Soit $\mathcal{U} \in \text{Cov}_1$, disons
$\mathcal{U} = \{U_i \to U\}_{i \in I}$. Par hypothèse, on peut choisir un
raffinement $\mathcal{V} \in \text{Cov}_2$ de $\mathcal{U}$, disons
$\mathcal{V} = \{V_j \to U\}_{j \in J}$, le raffinement étant donné par
$\alpha : J \to I$ et $f_j : V_j \to U_{\alpha(j)}$. Observons que
$\mathcal{F}$ satisfait à la condition de faisceau pour $\mathcal{V}$ et pour
les recouvrements $\{V_j \times_U U_i \to U_i\}_{j \in J}$, puisque ceux-ci
appartiennent à $\text{Cov}_2$. Le lemme
\ref{lemma-compare-sheaf-condition} montre donc que $\mathcal{F}$ satisfait à
la condition de faisceau pour $\mathcal{U}$.
\end{proof}

\begin{lemma}
\label{lemma-choice-set-coverings-immaterial}
Soit $\mathcal{C}$ une catégorie. Soit $\text{Cov}(\mathcal{C})$ une classe
propre de recouvrements satisfaisant les conditions (1), (2) et (3) de la
définition \ref{definition-site}. Soient
$\text{Cov}_1, \text{Cov}_2 \subset \text{Cov}(\mathcal{C})$ deux parties de
$\text{Cov}(\mathcal{C})$ qui munissent $\mathcal{C}$ d’une structure de site.
Si tout recouvrement $\mathcal{U} \in \text{Cov}(\mathcal{C})$ est
combinatoirement équivalent à un recouvrement de $\text{Cov}_1$ et
combinatoirement équivalent à un recouvrement de $\text{Cov}_2$, alors
$\Sh(\mathcal{C}, \text{Cov}_1) =
\Sh(\mathcal{C}, \text{Cov}_2)$.
\end{lemma}

\begin{proof}
C’est clair d’après les lemmes \ref{lemma-refine-same-topology} et
\ref{lemma-tautological-combinatorial} ci-dessus, puisque l’hypothèse implique
que tout recouvrement
$\mathcal{U} \in \text{Cov}_1 \subset \text{Cov}(\mathcal{C})$ est
combinatoirement équivalent à un élément de $\text{Cov}_2$, et de même après
échange des rôles de $\text{Cov}_1$ et de $\text{Cov}_2$.
\end{proof}












\section{L’exemple des G-ensembles}
\label{section-example-sheaf-G-sets}

\noindent
Comme exemple, considérons le site $\mathcal{T}_G$ de
l’Exemple \ref{example-site-on-group}. Nous allons décrire la
catégorie des faisceaux sur $\mathcal{T}_G$. La réponse se révélera
indépendante des choix faits dans la définition de $\mathcal{T}_G$.
En fait, au cours de la démonstration, nous n’aurons besoin que des
propriétés suivantes du site $\mathcal{T}_G$ :
\begin{enumerate}
\item[(a)] $\mathcal{T}_G$ est une sous-catégorie pleine de $G\textit{-Ensembles}$,
\item[(b)] $\mathcal{T}_G$ contient le $G$-ensemble ${}_GG$,
\item[(c)] $\mathcal{T}_G$ admet des produits fibrés et ceux-ci sont les mêmes que
dans $G\textit{-Ensembles}$,
\item[(d)] étant donné $U \in \Ob(\mathcal{T}_G)$ et un sous-ensemble
$G$-invariant $O \subset U$, il existe un objet de $\mathcal{T}_G$ isomorphe
à $O$, et
\item[(e)] toute famille surjective d’applications $\{U_i \to U\}_{i \in I}$, avec
$U, U_i \in \Ob(\mathcal{T}_G)$, est combinatoirement équivalente à un
recouvrement de $\mathcal{T}_G$.
\end{enumerate}
Ces propriétés découlent de Ensembles, Lemmes \ref{sets-lemma-what-is-in-it-G-sets}
et \ref{sets-lemma-coverings-site}.

\medskip\noindent
Remarquons que l’application
$$
\Hom_G({}_GG, {}_GG)
\longrightarrow
G^{opp},
\varphi
\longmapsto
\varphi(1)
$$
est un isomorphisme de groupes. L’application réciproque envoie $g \in G$
sur l’application $R_g : s \mapsto sg$ (c.-à-d.\ la multiplication à droite).
Notons que $R_{g_1g_2} = R_{g_2} \circ R_{g_1}$, de sorte que le passage au groupe
opposé est nécessaire.

\medskip\noindent
Il en résulte que, pour tout préfaisceau $\mathcal{F}$ sur $\mathcal{T}_G$,
l’ensemble $\mathcal{F}({}_GG)$ hérite d’une structure de $G$-ensemble
comme suit : $g \cdot s$, pour $g \in G$ et $s \in \mathcal{F}({}_GG)$,
est défini par $\mathcal{F}(R_g)(s)$. Il s’agit d’une action à gauche
car
$$
(g_1g_2) \cdot s  = \mathcal{F}(R_{g_1g_2})(s) =
\mathcal{F}(R_{g_2}\circ R_{g_1})(s) =
\mathcal{F}(R_{g_1})( \mathcal{F}(R_{g_2})(s)) =
g_1 \cdot (g_2 \cdot s).
$$
Nous avons utilisé ici le fait que $\mathcal{F}$
est contravariant. Notons que, si $\mathcal{F} \to \mathcal{G}$
est un morphisme de préfaisceaux d’ensembles sur $\mathcal{T}_G$,
nous obtenons alors une application $\mathcal{F}({}_GG) \to \mathcal{G}({}_GG)$
compatible avec les actions de $G$ que nous venons de définir.
En définitive, nous avons construit un foncteur
$$
\textit{PSh}(\mathcal{T}_G)
\longrightarrow
G\textit{-Ensembles}, \quad
\mathcal{F}
\longmapsto
\mathcal{F}({}_GG).
$$
Nous laissons au lecteur le soin de vérifier que cette construction
possède l’agréable propriété d’envoyer le préfaisceau représentable
$h_U$ sur un objet canoniquement isomorphe à $U$.
En formule, $h_U({}_GG) = \Hom_G({}_GG, U) \cong U$.

\medskip\noindent
Supposons que $S$ soit un $G$-ensemble. Nous définissons un préfaisceau
$\mathcal{F}_S$ par la formule\footnote{Il peut sembler qu’il s’agisse du
préfaisceau représentable défini par $S$. Ce n’est pas nécessairement le cas,
car $S$ peut ne pas être un objet de $\mathcal{T}_G$, lequel a été choisi
comme un ensemble suffisamment grand de $G$-ensembles.}
$$
\mathcal{F}_S(U)
=
\Mor_{G\textit{-Ensembles}}(U, S).
$$
C’est manifestement un préfaisceau. D’autre part, supposons que
$\{U_i \to U\}_{i\in I}$ soit un recouvrement de $\mathcal{T}_G$.
Il en résulte que $\coprod_i U_i \to U$ est surjective. Il est donc
clair que l’application
$$
\mathcal{F}_S(U)
=
\Mor_{G\textit{-Ensembles}}(U, S)
\longrightarrow
\prod \mathcal{F}_S(U_i)
=
\prod \Mor_{G\textit{-Ensembles}}(U_i, S)
$$
est injective. Et, étant donnée une famille d’applications
$G$-équivariantes $s_i : U_i \to S$, telle que tous les diagrammes
$$
\xymatrix{
U_i \times_U U_j \ar[d] \ar[r]
&
U_j \ar[d]^{s_j}
\\
U_i \ar[r]^{s_i}
&
S
}
$$
commutent, il existe une unique application $G$-équivariante
$s : U \to S$ telle que $s_i$ soit la composée
$U_i \to U \to S$. En effet, il suffit de définir $s(u) = s_i(u_i)$,
où $i\in I$ est un indice quelconque tel qu’il existe un
$u_i \in U_i$ s’envoyant sur $u$ par l’application $U_i \to U$.
La commutativité des diagrammes ci-dessus signifie exactement
que cette construction est bien définie. En définitive, nous avons
construit un foncteur
$$
G\textit{-Ensembles}
\longrightarrow
\Sh(\mathcal{T}_G), \quad
S
\longmapsto
\mathcal{F}_S
.
$$

\medskip\noindent
Nous avons maintenant le diagramme suivant de catégories et de foncteurs
$$
\xymatrix{
\textit{PSh}(\mathcal{T}_G) \ar[rr]^{\mathcal{F} \mapsto \mathcal{F}({}_GG)}
&
&
G\textit{-Ensembles} \ar[ld]_{S \mapsto \mathcal{F}_S}
\\
&
\Sh(\mathcal{T}_G) \ar[lu]
&
}
$$
Il découle immédiatement des définitions que $\mathcal{F}_S({}_GG)
= \Mor_G({}_GG, S) = S$, la dernière égalité résultant de l’évaluation en $1$.
Ceci démontre presque la proposition suivante.

\begin{proposition}
\label{proposition-sheaves-on-group}
Les foncteurs $\mathcal{F} \mapsto \mathcal{F}({}_GG)$
et $S \mapsto \mathcal{F}_S$ définissent des équivalences
quasi-inverses entre $\Sh(\mathcal{T}_G)$
et $G\textit{-Ensembles}$.
\end{proposition}

\begin{proof}
Nous avons déjà vu que la composée des foncteurs dans un sens
est isomorphe au foncteur identité.
Dans l’autre sens, pour tout faisceau $\mathcal{H}$, il existe un
morphisme naturel de faisceaux
$$
can :
\mathcal{H}
\longrightarrow
\mathcal{F}_{\mathcal{H}({}_GG)}.
$$
En effet, pour tout objet $U$ de $\mathcal{T}_G$, nous définissons $can_U$
comme l’application
$$
\begin{matrix}
\mathcal{H}(U)
&
\longrightarrow
&
\mathcal{F}_{\mathcal{H}({}_GG)}(U)
=
\Mor_G(U, \mathcal{H}({}_GG))
\\
s
&
\longmapsto
&
(u \mapsto \alpha_u^*s).
\end{matrix}
$$
Ici, $\alpha_u : {}_GG \to U$ est l’application
$\alpha_u(g) = gu$ et $\alpha_u^* : \mathcal{H}(U)
\to \mathcal{H}({}_GG)$ est l’application de restriction. Une vérification
élémentaire mais déroutante montre qu’il s’agit bien d’un morphisme
de préfaisceaux. Nous devons montrer que $can$ est un isomorphisme.
Pour cela, nous montrons que $can_U$ est un isomorphisme pour tout
$U \in \Ob(\mathcal{T}_G)$. Nous laissons au lecteur le cas (important mais
facile) où $U = {}_GG$.
Un objet général $U$ de $\mathcal{T}_G$ est une réunion disjointe de
$G$-orbites : $U = \coprod_{i\in I} O_i$. La famille d’applications
$\{O_i \to U\}_{i \in I}$ est tautologiquement équivalente
à un recouvrement de $\mathcal{T}_G$ (par les propriétés de $\mathcal{T}_G$
énumérées au début de cette section). Par conséquent, d’après le Lemme
\ref{lemma-tautological-same-sheaf}, le faisceau $\mathcal{H}$
satisfait la condition de faisceau relativement à
$\{O_i \to U\}_{i \in I}$. La condition de faisceau pour ce recouvrement
implique $\mathcal{H}(U) = \prod_i \mathcal{H}(O_i)$.
Il suffit donc de montrer que $can_U$ est un
isomorphisme lorsque $U$ consiste en une seule $G$-orbite. Soit $u \in U$
et soit $H \subset G$ son stabilisateur. Manifestement,
$\Mor_G(U, \mathcal{H}({}_GG)) = \mathcal{H}({}_GG)^H$
est le sous-ensemble des éléments invariants par $H$. Considérons d’autre part
le recouvrement $\{{}_GG \to U\}$ donné par $g \mapsto
gu$ (une fois encore, il est seulement combinatoirement équivalent à un recouvrement
de $\mathcal{T}_G$, ce qui, là encore, n’a aucune importance).
Notons que le produit fibré $({}_GG)\times_U ({}_GG)$
est égal à $\{(g, gh), g\in G, h\in H\} \cong \coprod_{h\in H} {}_GG$.
La condition de faisceau pour ce recouvrement s’écrit donc
$$
\xymatrix{
\mathcal{H}(U) \ar[r]
&
\mathcal{H}({}_GG)
\ar@<1ex>[r]^-{\text{pr}_0^*} \ar@<-1ex>[r]_-{\text{pr}_1^*}
&
\prod_{h \in H}
\mathcal{H}({}_GG).
}
$$
Les deux applications $\text{pr}_i^*$ à valeurs dans le facteur
$\mathcal{H}({}_GG)$ diffèrent par la multiplication par $h$.
Le résultat découle maintenant de ce fait et de ce que $can$
est un isomorphisme pour $U = {}_GG$.
\end{proof}





















\section{Passage au faisceau associé}
\label{section-sheafification}

\noindent
Afin de définir le passage au faisceau associé, nous étudions le groupe de
cohomologie de degré zéro de {\v C}ech d’un recouvrement et ses propriétés
de fonctorialité.

\medskip\noindent
Soit $\mathcal{F}$ un préfaisceau d’ensembles sur $\mathcal{C}$, et soit
$\mathcal{U} = \{U_i \to U\}_{i \in I}$ un recouvrement de $\mathcal{C}$.
Nous employons la notation $\mathcal{F}(\mathcal{U})$ pour désigner l’égalisateur
$$
H^0(\mathcal{U}, \mathcal{F})
=
\{
(s_i)_{i\in I} \in \prod\nolimits_i \mathcal{F}(U_i)
\mid
s_i|_{U_i \times_U U_j} = s_j|_{U_i \times_U U_j}
\ \forall i, j \in I
\}.
$$
Comme nous le verrons plus loin, il s’agit de la cohomologie de degré zéro de
{\v C}ech de $\mathcal{F}$ sur $U$ relativement au recouvrement $\mathcal{U}$.
Une petite remarque est que nous pouvons définir $H^0(\mathcal{U}, \mathcal{F})$
dès que tous les morphismes $U_i \to U$ sont représentables, c.-à-d.\ que
$\mathcal{U}$ n’a pas besoin d’être un recouvrement du site.
Il existe une application canonique $\mathcal{F}(U) \to H^0(\mathcal{U}, \mathcal{F})$.
Il est clair qu’un morphisme de recouvrements $\mathcal{U} \to \mathcal{V}$
induit des diagrammes commutatifs
$$
\xymatrix{
& U_i \ar[rr] & & V_{\alpha(i)} \\
U_i \times_U U_j \ar[rr] \ar[ur] \ar[dr] & &
V_{\alpha(i)} \times_V V_{\alpha(j)} \ar[ur] \ar[dr] & \\
& U_j \ar[rr] & & V_{\alpha(j)}
}.
$$
Ceux-ci produisent à leur tour une application $H^0(\mathcal{V}, \mathcal{F}) \to
H^0(\mathcal{U}, \mathcal{F})$, compatible à l’application $\mathcal{F}(V)
\to \mathcal{F}(U)$.

\medskip\noindent
Par construction, un préfaisceau $\mathcal{F}$ est un faisceau si et seulement si,
pour tout recouvrement $\mathcal{U}$ de $\mathcal{C}$, l’application naturelle
$\mathcal{F}(U) \to H^0(\mathcal{U}, \mathcal{F})$ est bijective.
Nous emploierons cette notion pour démontrer le lemme simple suivant
sur les limites de faisceaux.

\begin{lemma}
\label{lemma-limit-sheaf}
\begin{slogan}
La limite d’un diagramme de faisceaux existe et coïncide avec la limite dans les préfaisceaux
\end{slogan}
Soit $\mathcal{F} : \mathcal{I} \to \Sh(\mathcal{C})$
un diagramme. Alors $\lim_\mathcal{I} \mathcal{F}$ existe
et est égale à la limite dans la catégorie des préfaisceaux.
\end{lemma}

\begin{proof}
Soit $\lim_i \mathcal{F}_i$ la limite calculée comme préfaisceau.
Nous allons montrer qu’il s’agit d’un faisceau ; il en résultera alors trivialement
que c’est une limite dans la catégorie des faisceaux. Pour vérifier la condition
de faisceau, soit $\mathcal{V} = \{V_j \to V\}_{j\in J}$ un recouvrement.
Soit $(s_j)_{j\in J}$ un élément de $H^0(\mathcal{V}, \lim_i \mathcal{F}_i)$.
En utilisant les morphismes de projection, nous obtenons des éléments $(s_{j, i})_{j\in J}$
de $H^0(\mathcal{V}, \mathcal{F}_i)$. Par la condition de faisceau pour
$\mathcal{F}_i$, nous voyons qu’il existe un unique $s_i \in \mathcal{F}_i(V)$
tel que $s_{j, i} = s_i|_{V_j}$. Soit $\phi : i \to i'$ un morphisme
de la catégorie d’indices. Nous voulons montrer que
$\mathcal{F}(\phi) : \mathcal{F}_i \to \mathcal{F}_{i'}$
envoie $s_i$ sur $s_{i'}$. Nous savons que c’est vrai pour les sections
$s_{i, j}$ et $s_{i', j}$ pour tout $j$, et donc la condition de faisceau
pour $\mathcal{F}_{i'}$ montre que c’est vrai. Nous avons alors un
élément $s = (s_i)_{i \in \Ob(\mathcal{I})}$ de
$(\lim_i \mathcal{F}_i)(V)$. Nous laissons au lecteur le soin de vérifier
que cet élément possède la propriété requise $s_j = s|_{V_j}$.
\end{proof}

\begin{example}
\label{example-singleton-sheaf}
Un exemple particulier est la limite du diagramme vide.
Elle donne l’objet final de la catégorie des (pré)faisceaux.
Il s’agit du préfaisceau qui associe
à tout objet $U$ de $\mathcal{C}$ un ensemble réduit à un élément, avec d’uniques
applications de restriction ; de plus, ce préfaisceau est un faisceau.
Nous notons souvent ce faisceau par $*$.
\end{example}

\noindent
Soit $\mathcal{J}_U$ la catégorie de tous les recouvrements de $U$.
Autrement dit, les objets de $\mathcal{J}_U$ sont les recouvrements
de $U$ dans $\mathcal{C}$, et les morphismes sont les raffinements.
Compte tenu de nos conventions sur les sites, il s’agit bien d’une catégorie, c.-à-d.\ que
la collection des objets et des morphismes forme un ensemble.
Notons que $\Ob(\mathcal{J}_U)$ n’est pas vide puisque
$\{\text{id}_U\}$ en est un objet. D’après les remarques
ci-dessus, la construction $\mathcal{U} \mapsto H^0(\mathcal{U}, \mathcal{F})$
est un foncteur contravariant sur $\mathcal{J}_U$.
Nous définissons
$$
\mathcal{F}^{+}(U)
=
\colim_{\mathcal{J}_U^{opp}}
H^0(\mathcal{U}, \mathcal{F})
$$
Voir Catégories, section \ref{categories-section-limits}, pour
une discussion des limites projectives et inductives. Signalons que nous
verrons plus loin que $\mathcal{F}^{+}(U)$ est la cohomologie de degré zéro de
{\v C}ech de $\mathcal{F}$ sur $U$.

\medskip\noindent
Avant d’en dire davantage sur la structure de la limite inductive, faisons
de la collection d’ensembles
$\mathcal{F}^{+}(U)$, $U \in \Ob(\mathcal{C})$,
un préfaisceau. Soit donc $V \to U$ un morphisme de $\mathcal{C}$.
Les axiomes d’un site fournissent un foncteur\footnote{Cette construction
fait en réalité intervenir un choix des produits fibrés $U_i \times_U V$
et donc l’axiome du choix. L’application obtenue ne dépend pas
des choix effectués, voir ci-dessous.}
$$
\mathcal{J}_U
\longrightarrow
\mathcal{J}_V, \quad
\{U_i \to U\}
\longmapsto
\{U_i \times_U V \to V\}.
$$
Notons que les projections fournissent un
morphisme fonctoriel de recouvrements $\{U_i \times_U V \to V\} \to \{U_i \to U\}$
et donc, par la construction ci-dessus, une application fonctorielle d’ensembles
$H^0(\{U_i \to U\}, \mathcal{F}) \to
H^0(\{U_i \times_U V \to V\}, \mathcal{F})$.
Autrement dit, il existe une transformation de foncteurs
de $H^0(-, \mathcal{F}) : \mathcal{J}_U^{opp} \to \textit{Ensembles}$
vers la composée
$\mathcal{J}_U^{opp} \to \mathcal{J}_V^{opp}
\xrightarrow{H^0(-, \mathcal{F})} \textit{Ensembles}$. Par les
généralités sur les limites inductives, nous obtenons donc une application canonique
$\mathcal{F}^+(U) \to \mathcal{F}^+(V)$. En termes de la description
ci-dessus de l’ensemble $\mathcal{F}^+(U)$, elle envoie simplement l’élément
associé à $s = (s_i) \in H^0(\{U_i \to U\}, \mathcal{F})$ sur l’élément
associé à $(s_i|_{V \times_U U_i})
\in H^0(\{U_i \times_U V \to V\}, \mathcal{F})$.

\begin{lemma}
\label{lemma-plus-presheaf}
Les constructions ci-dessus définissent un préfaisceau
$\mathcal{F}^+$ muni d’un morphisme canonique
de préfaisceaux $\mathcal{F} \to \mathcal{F}^+$.
\end{lemma}

\begin{proof}
Il suffit de montrer qu’étant donnés des morphismes
$W \to V \to U$, la composée $\mathcal{F}^+(U)
\to \mathcal{F}^+(V) \to \mathcal{F}^+(W)$
est égale à l’application $\mathcal{F}^+(U) \to \mathcal{F}^+(W)$.
On peut le montrer directement en vérifiant que, étant donnés
un recouvrement $\{U_i \to U\}$ et
$s = (s_i) \in H^0(\{U_i \to U\}, \mathcal{F})$,
nous avons canoniquement
$W \times_U U_i \cong W \times_V (V \times_U U_i)$,
et que
$s_i|_{W \times_U U_i}$
correspond à
$(s_i|_{V \times_U U_i})|_{W \times_V (V \times_U U_i)}$
par cet isomorphisme.
\end{proof}

\noindent
Plus indirectement, le résultat du
Lemme \ref{lemma-independent-refinement} montre que
nous pouvons tirer en arrière un élément $s$ comme ci-dessus par tout morphisme
d’un recouvrement quelconque de $W$ vers $\{U_i \to U\}$,
et que nous aboutirons toujours au même élément de
$\mathcal{F}^+(W)$.

\begin{lemma}
\label{lemma-plus-functorial}
L’association $\mathcal{F} \mapsto
(\mathcal{F} \to \mathcal{F}^+)$
est un foncteur.
\end{lemma}

\begin{proof}
Au lieu de le démontrer, énonçons exactement ce qu’il faut prouver.
Soit $\mathcal{F} \to \mathcal{G}$ un morphisme de préfaisceaux. Il faut prouver
la commutativité de :
$$
\xymatrix{
\mathcal{F} \ar[r] \ar[d]
&
\mathcal{F}^{+} \ar[d]
\\
\mathcal{G} \ar[r]
&
\mathcal{G}^{+}
}
$$
\end{proof}

\noindent
Les deux lemmes suivants impliquent que les limites inductives ci-dessus sont indexées
par un ensemble ordonné filtrant.

\begin{lemma}
\label{lemma-common-refinement}
Étant donnés deux recouvrements $\{U_i \to U\}$
et $\{V_j \to U\}$ d’un objet donné $U$ du site
$\mathcal{C}$, il existe un recouvrement qui en est un
raffinement commun.
\end{lemma}

\begin{proof}
Puisque $\mathcal{C}$ est un site, pour tout $i$, la
famille $\{V_j \times_U U_i \to U_i\}_j$ est un recouvrement.
Un autre axiome implique alors que $\{V_j \times_U U_i \to U\}_{i, j}$
est un recouvrement de $U$. Manifestement, ce recouvrement raffine les deux
recouvrements donnés.
\end{proof}

\begin{lemma}
\label{lemma-independent-refinement}
Deux morphismes quelconques $f, g: \mathcal{U} \to \mathcal{V}$ de recouvrements
induisant le même morphisme $U \to V$ induisent la même
application $H^0(\mathcal{V}, \mathcal{F}) \to  H^0(\mathcal{U}, \mathcal{F})$.
\end{lemma}

\begin{proof}
Soient $\mathcal{U} = \{U_i \to U\}_{i\in I}$ et
$\mathcal{V} = \{V_j \to V\}_{j\in J}$.
Le morphisme $f$ consiste en une application $U\to V$, une application $\alpha : I \to J$ et
des morphismes $f_i : U_i \to V_{\alpha(i)}$.
De même, $g$~détermine une application $\beta : I \to J$ et des morphismes
$g_i : U_i \to V_{\beta(i)}$.
Comme $f$ et $g$ induisent la même application $U\to V$, le diagramme
$$
\xymatrix{
&
V_{\alpha(i)} \ar[dr]
\\
U_i \ar[ur]^{f_i} \ar[dr]_{g_i}
&
&
V
\\
&
V_{\beta(i)} \ar[ur]
}
$$
est commutatif pour tout $i\in I$. Par conséquent, $f$ et $g$ se factorisent par
le produit fibré
$$
\xymatrix{
&
V_{\alpha(i)}
\\
U_i \ar[r]^-\varphi \ar[ur]^{f_i} \ar[dr]_{g_i}
&
V_{\alpha(i)} \times_V V_{\beta(i)} \ar[u]_{\text{pr}_1} \ar[d]^{\text{pr}_2}
\\
&
V_{\beta(i)}.
}
$$
Soit maintenant $s = (s_j)_j \in H^0(\mathcal{V}, \mathcal{F})$.
Alors, pour tout $i\in I$ :
$$
(f^*s)_i
=
f_i^*(s_{\alpha(i)})
=
\varphi^*\text{pr}_1^*(s_{\alpha(i)})
=
\varphi^*\text{pr}_2^*(s_{\beta(i)})
=
g_i^*(s_{\beta(i)})
=
(g^*s)_i,
$$
où l’égalité du milieu résulte de la définition
de $H^0(\mathcal{V}, \mathcal{F})$.
Ceci montre que les applications
$H^0(\mathcal{V}, \mathcal{F}) \to H^0(\mathcal{U}, \mathcal{F})$
induites par $f$ et $g$ sont égales.
\end{proof}

\begin{remark}
\label{remark-both-refine-same-H0}
En particulier, ce lemme montre que, si $\mathcal{U}$ est
un raffinement de $\mathcal{V}$ et si $\mathcal{V}$ est un
raffinement de $\mathcal{U}$, alors il existe une identification
canonique $H^0(\mathcal{U}, \mathcal{F}) =
H^0(\mathcal{V}, \mathcal{F})$.
\end{remark}

\noindent
Ces deux lemmes et le fait que $\mathcal{J}_U$ est non vide
impliquent que le diagramme $H^0(-, \mathcal{F}) : \mathcal{J}_U^{opp}
\to \textit{Ensembles}$ est filtrant, voir
Catégories, Définition \ref{categories-definition-directed}.
Par conséquent, d’après Catégories,
section \ref{categories-section-directed-colimits},
la limite inductive $\mathcal{F}^{+}(U)$ peut être décrite
de la manière élémentaire suivante. Tout élément de l’ensemble
$\mathcal{F}^{+}(U)$ provient d’un élément
$s \in H^0(\mathcal{U}, \mathcal{F})$ pour un certain recouvrement
$\mathcal{U}$ de $U$. Étant donné un second élément $s' \in
H^0(\mathcal{U}', \mathcal{F})$, les éléments $s$ et $s'$ déterminent
le même élément de la limite inductive si et seulement s’il existe un recouvrement
$\mathcal{V}$ de $U$ et des raffinements $f : \mathcal{V} \to \mathcal{U}$ et
$f' : \mathcal{V} \to \mathcal{U}'$ tels que $f^*s = (f')^*s'$
dans $H^0(\mathcal{V}, \mathcal{F})$. Comme le recouvrement trivial
$\{\text{id}_U\}$ est un objet de $\mathcal{J}_U$, nous obtenons
une application canonique $\mathcal{F}(U) \to \mathcal{F}^+(U)$.

\begin{lemma}
\label{lemma-plus-surjective}
Le morphisme $\theta : \mathcal{F} \to \mathcal{F}^+$ possède la propriété
suivante : pour tout objet $U$ de $\mathcal{C}$ et toute section
$s \in \mathcal{F}^+(U)$, il existe un recouvrement $\{U_i \to U\}$
tel que $s|_{U_i}$ appartienne à l’image de $\theta : \mathcal{F}(U_i)
\to \mathcal{F}^{+}(U_i)$.
\end{lemma}

\begin{proof}
Soit en effet $\{U_i \to U\}$ un recouvrement tel que $s$ provienne
de l’élément $(s_i) \in H^0(\{U_i \to U\}, \mathcal{F})$.
D’après le Lemme \ref{lemma-independent-refinement}, nous pouvons
considérer le recouvrement $\{U_i \to U_i\}$ et le morphisme (évident)
de recouvrements $\{U_i \to U_i\} \to \{U_i \to U\}$ pour calculer la
restriction de $s$ dans $\mathcal{F}^+(U_i)$. Et en effet,
ce recouvrement donne exactement $\theta(s_i)$ comme restriction
de $s$ à $U_i$.
\end{proof}

\begin{definition}
\label{definition-separated}
Nous disons qu’un préfaisceau d’ensembles $\mathcal{F}$ sur un site
$\mathcal{C}$ est {\it séparé} si, pour tout recouvrement
$\{U_i \rightarrow U\}$, l’application
$\mathcal{F}(U) \to \prod \mathcal{F}(U_i)$ est injective.
\end{definition}

\begin{theorem}
\label{theorem-plus}
Avec $\mathcal{F}$ comme ci-dessus,
\begin{enumerate}
\item
\label{item-sep}
Le préfaisceau $\mathcal{F}^+$ est séparé.
\item
\label{item-sheaf}
Si $\mathcal{F}$ est séparé, alors $\mathcal{F}^+$ est un faisceau
et le morphisme de préfaisceaux $\mathcal{F} \to \mathcal{F}^+$ est injectif.
\item
\label{item-plus-iso}
Si $\mathcal{F}$ est un faisceau, alors $\mathcal{F} \to \mathcal{F}^+$
est un isomorphisme.
\item
\label{item-plusplus}
Le préfaisceau $\mathcal{F}^{++}$ est toujours un faisceau.
\end{enumerate}
\end{theorem}

\begin{proof}
Démonstration de (\ref{item-sep}).
Supposons que $s, s' \in \mathcal{F}^+(U)$ et qu’il
existe un recouvrement $\{U_i \to U\}$ tel que
$s|_{U_i} = s'|_{U_i}$ pour tout $i$. Nous avons alors trois recouvrements
de $U$ : le recouvrement $\{U_i \to U\}$ ci-dessus, un recouvrement $\mathcal{U}$
pour $s$ comme dans le Lemme \ref{lemma-plus-surjective},
et un recouvrement analogue $\mathcal{U}'$ pour $s'$. D’après le Lemme
\ref{lemma-common-refinement}, nous pouvons trouver un raffinement commun,
disons $\{W_j \to U\}$. Cela signifie qu’il existe $s_j, s'_j \in \mathcal{F}(W_j)$
tels que $s|_{W_j} = \theta(s_j)$, et de même pour $s'|_{W_j}$, et
tels que $\theta(s_j) = \theta(s'_j)$. Cette dernière égalité signifie
qu’il existe un recouvrement $\{W_{jk} \to W_j\}$ tel que
$s_j|_{W_{jk}} = s'_j|_{W_{jk}}$. Puisque $\{W_{jk} \to U\}$
est un recouvrement, nous voyons alors que $s, s'$ ont la même image dans
$H^0(\{W_{jk} \to U\}, \mathcal{F})$, comme voulu.

\medskip\noindent
Démonstration de (\ref{item-sheaf}). Il est clair que $\mathcal{F} \to
\mathcal{F}^+$ est injectif, car toutes les applications
$\mathcal{F}(U) \to H^0(\mathcal{U}, \mathcal{F})$
sont injectives. Il est également clair que, si $\mathcal{U} \to
\mathcal{U}'$ est un raffinement, alors $H^0(\mathcal{U}', \mathcal{F})
\to H^0(\mathcal{U}, \mathcal{F})$ est injective. Maintenant,
supposons que $\{U_i \to U\}$ soit un recouvrement, et soit
$(s_i)$ une famille d’éléments de $\mathcal{F}^+(U_i)$
satisfaisant la condition de faisceau
$s_i|_{U_i \times_U U_{i'}} = s_{i'}|_{U_i \times_U U_{i'}}$
pour tous $i, i' \in I$. Choisissons des recouvrements (comme dans le
Lemme \ref{lemma-plus-surjective}) $\{U_{ij} \to U_i\}$
tels que $s_i|_{U_{ij}}$ soit l’image de l’unique
élément $s_{ij} \in \mathcal{F}(U_{ij})$. La condition de faisceau
implique que $s_{ij}$ et $s_{i'j'}$ coïncident sur
$U_{ij} \times_U U_{i'j'}$, car ce produit s’envoie dans
$U_i \times_U U_{i'}$ et que nous y avons l’égalité.
Ainsi, $(s_{ij}) \in H^0(\{U_{ij} \to U\}, \mathcal{F})$
donne naissance à un élément $s \in \mathcal{F}^+(U)$. Nous laissons
au lecteur le soin de vérifier que $s|_{U_i} = s_i$.

\medskip\noindent
Démonstration de (\ref{item-plus-iso}). Cela résulte immédiatement des définitions,
car la condition de faisceau dit exactement que toute application
$\mathcal{F} \to H^0(\mathcal{U}, \mathcal{F})$ est bijective
(pour tout recouvrement $\mathcal{U}$ de $U$).

\medskip\noindent
L’assertion (\ref{item-plusplus}) est maintenant évidente.
\end{proof}

\begin{definition}
\label{definition-associated-sheaf}
Soit $\mathcal{C}$ un site et soit $\mathcal{F}$ un préfaisceau
d’ensembles sur $\mathcal{C}$. Le faisceau $\mathcal{F}^\# := \mathcal{F}^{++}$,
muni du morphisme canonique $\mathcal{F} \to \mathcal{F}^\#$,
est appelé le {\it faisceau associé à $\mathcal{F}$}.
\end{definition}

\begin{proposition}
\label{proposition-sheafification-adjoint}
Le morphisme canonique $\mathcal{F} \to \mathcal{F}^\#$ possède la
propriété universelle suivante : pour tout morphisme $\mathcal{F} \to \mathcal{G}$,
où $\mathcal{G}$ est un faisceau d’ensembles, il existe un unique morphisme
$\mathcal{F}^\# \to \mathcal{G}$ tel que $\mathcal{F} \to \mathcal{F}^\#
\to \mathcal{G}$ soit égal au morphisme donné.
\end{proposition}

\begin{proof}
Le Lemme \ref{lemma-plus-functorial} donne un diagramme commutatif
$$
\xymatrix{
\mathcal{F} \ar[r] \ar[d]
&
\mathcal{F}^{+} \ar[r] \ar[d]
&
\mathcal{F}^{++} \ar[d]
\\
\mathcal{G} \ar[r]
&
\mathcal{G}^{+} \ar[r]
&
\mathcal{G}^{++}
}
$$
et, d’après le Théorème \ref{theorem-plus}, les morphismes horizontaux inférieurs
sont des isomorphismes. L’unicité découle du Lemme
\ref{lemma-plus-surjective}, qui affirme que toute section de
$\mathcal{F}^\#$ provient localement de sections de $\mathcal{F}$.
\end{proof}

\noindent
Il résulte clairement de cette proposition que le foncteur $\mathcal{F}
\mapsto (\mathcal{F} \to \mathcal{F}^\#)$ est unique
à un unique isomorphisme de foncteurs près. En fait, notons temporairement
$i : \Sh(\mathcal{C}) \to \textit{PSh}(\mathcal{C})$
le foncteur d’inclusion. Le résultat ci-dessus affirme précisément que
$$
\Mor_{\textit{PSh}(\mathcal{C})}(\mathcal{F}, i(\mathcal{G}))
=
\Mor_{\Sh(\mathcal{C})}(\mathcal{F}^\#, \mathcal{G}).
$$
Autrement dit, le foncteur de passage au faisceau associé est l’adjoint à gauche
du foncteur d’inclusion $i$.

\begin{example}
\label{example-constant-sheaf}
Soit $\mathcal{C}$ un site. Soit $E$ un ensemble.
Le {\it faisceau constant $\underline{E}$ de valeur $E$}
est le faisceau associé au préfaisceau donné par le foncteur constant
de valeur $E$. Il est caractérisé par la règle
$\Mor_{\Sh(\mathcal{C})}(\underline{E}, \mathcal{F}) =
\Mor_{\textit{Ensembles}}(E, \Gamma(\mathcal{C}, \mathcal{F}))$
pour tout faisceau $\mathcal{F}$ sur $\mathcal{C}$, où la
notation est celle de la Définition \ref{definition-global-sections}.
\end{example}

\noindent
Nous terminons cette section par quelques lemmes.

\begin{lemma}
\label{lemma-colimit-sheaves}
\begin{slogan}
La limite inductive dans les faisceaux est le faisceau associé à la limite inductive dans les
préfaisceaux.
\end{slogan}
Soit $\mathcal{F} : \mathcal{I} \to \Sh(\mathcal{C})$
un diagramme. Alors $\colim_\mathcal{I} \mathcal{F}$ existe
et est le faisceau associé à la limite inductive dans la catégorie des préfaisceaux.
\end{lemma}

\begin{proof}
Puisque le foncteur de passage au faisceau associé est un adjoint à gauche, il commute
à la formation de toutes les limites inductives, voir Catégories,
Lemme \ref{categories-lemma-adjoint-exact}.
Ainsi, puisque $\textit{PSh}(\mathcal{C})$ admet des limites inductives, nous en déduisons
que $\Sh(\mathcal{C})$ admet des limites inductives (qui sont les
faisceaux associés aux limites inductives de préfaisceaux).
\end{proof}

\begin{lemma}
\label{lemma-sheafification-exact}
Le foncteur $\textit{PSh}(\mathcal{C}) \to \Sh(\mathcal{C})$,
$\mathcal{F} \mapsto \mathcal{F}^\#$, est exact.
\end{lemma}

\begin{proof}
Puisqu’il est adjoint à gauche, il est exact à droite, voir
Catégories, Lemme \ref{categories-lemma-exact-adjoint}.
D’autre part, d’après les Lemmes \ref{lemma-common-refinement}
et \ref{lemma-independent-refinement}, les limites inductives
intervenant dans la construction de $\mathcal{F}^+$ sont en réalité indexées par
l’ensemble ordonné filtrant $\Ob(\mathcal{J}_U)$, où
$\mathcal{U} \geq \mathcal{U}'$ si et seulement si
$\mathcal{U}$ est un raffinement de $\mathcal{U}'$. Il résulte alors de
Catégories, Lemme \ref{categories-lemma-directed-commutes},
que $\mathcal{F} \to \mathcal{F}^+$ commute
aux limites finies (comme foncteur des préfaisceaux vers les
préfaisceaux). Nous concluons alors en utilisant le
Lemme \ref{lemma-limit-sheaf}.
\end{proof}

\begin{lemma}
\label{lemma-sections-sheafification}
Soit $\mathcal{C}$ un site.
Soit $\mathcal{F}$ un préfaisceau d’ensembles sur $\mathcal{C}$.
Notons $\theta^2 : \mathcal{F} \to \mathcal{F}^\#$ le morphisme canonique
de $\mathcal{F}$ vers son faisceau associé.
Soit $U$ un objet de $\mathcal{C}$.
Soit $s \in \mathcal{F}^\#(U)$. Il existe
un recouvrement $\{U_i \to U\}$ et des sections
$s_i \in \mathcal{F}(U_i)$ tels que
\begin{enumerate}
\item $s|_{U_i} = \theta^2(s_i)$, et
\item pour tous $i, j$, il existe un recouvrement
$\{U_{ijk} \to U_i \times_U U_j\}$ de $\mathcal{C}$ tel que
les restrictions de $s_i$ et $s_j$ à chaque $U_{ijk}$ coïncident.
\end{enumerate}
Réciproquement, étant donnés un recouvrement $\{U_i \to U\}$ et des éléments
$s_i \in \mathcal{F}(U_i)$ tels que (2) soit satisfaite, il
existe une unique section $s \in \mathcal{F}^\#(U)$ telle
que (1) soit satisfaite.
\end{lemma}

\begin{proof}
Omis.
\end{proof}

\begin{lemma}
\label{lemma-isomorphism-sheafifications}
Soit $\mathcal{C}$ un site. Soit $\mathcal{F} \to \mathcal{G}$ un morphisme
de préfaisceaux d’ensembles sur $\mathcal{C}$. Notons $\mathcal{B}$ l’ensemble des
$U \in \Ob(\mathcal{C})$ tels que $\mathcal{F}(U) \to \mathcal{G}(U)$
soit bijective. Si tout objet de $\mathcal{C}$ admet un recouvrement par des éléments
de $\mathcal{B}$, alors $\mathcal{F}^\# \to \mathcal{G}^\#$ est un isomorphisme.
\end{lemma}

\begin{proof}
Soit $U \in \Ob(\mathcal{C})$.
Démontrons que $\mathcal{F}^\#(U) \to \mathcal{G}^\#(U)$ est
surjective. Pour ce faire, nous utiliserons le Lemme \ref{lemma-sections-sheafification}
sans autre mention. Pour tout $s \in \mathcal{G}^\#(U)$, il existe
un recouvrement $\{U_i \to U\}$ et des sections $s_i \in \mathcal{G}(U_i)$ tels que
\begin{enumerate}
\item $s|_{U_i}$ soit l’image de $s_i$ par $\mathcal{G} \to \mathcal{G}^\#$,
et
\item pour tous $i, j$, il existe un recouvrement
$\{U_{ijk} \to U_i \times_U U_j\}$ de $\mathcal{D}$ tel que
les restrictions de $s_i$ et $s_j$ à chaque $U_{ijk}$ coïncident.
\end{enumerate}
Par hypothèse, pour tout $i$, nous pouvons choisir un recouvrement $\{U_{i, a} \to U_i\}$
avec $U_{i, a} \in \mathcal{B}$. Alors $\{U_{i, a} \to U\}$ est un recouvrement.
Notant $s_{i, a}$ l’image de $s_i$ dans $\mathcal{G}(U_{i, a})$,
nous voyons que les restrictions de $s_{i, a}$ et $s_{j, b}$ aux membres
du recouvrement
$$
\{(U_{i, a} \times_U U_{j, b}) \times_{U_i \times_U U_j} U_{ijk}
\to U_{i, a} \times_U U_{j, b}\}
$$
coïncident. Nous pouvons donc supposer que $U_i \in \mathcal{B}$. En répétant
l’argument, nous pouvons aussi supposer que $U_{ijk} \in \mathcal{B}$ pour
tous $i, j, k$ (détails omis).
Alors, puisque $\mathcal{F}(U_i) \to \mathcal{G}(U_i)$
est bijective par notre définition de $\mathcal{B}$ dans l’énoncé
du lemme, nous obtenons un unique $t_i \in \mathcal{F}(U_i)$
qui s’envoie sur $s_i$. Les restrictions de $t_i$ et $t_j$
à chaque $U_{ijk}$ coïncident dans $\mathcal{F}(U_{ijk})$, car
$U_{ijk} \in \mathcal{B}$ et parce que nous avons l’égalité
pour $t_i$ et $t_j$.
Le lemme nous dit alors qu’il existe un unique
$t \in \mathcal{F}^\#(U)$ tel que $t|_{U_i}$ soit l’image de $t_i$.
Il en résulte que $t$ s’envoie sur $s$.
Nous omettons la démonstration de l’injectivité de
$\mathcal{F}^\#(U) \to \mathcal{G}^\#(U)$.
\end{proof}








\section{Applications injectives et surjectives de faisceaux}
\label{section-sheaves-injective}

\begin{definition}
\label{definition-sheaves-injective-surjective}
Soit $\mathcal{C}$ un site, et soit $\varphi : \mathcal{F}
\to \mathcal{G}$ un morphisme de faisceaux d’ensembles.
\begin{enumerate}
\item Nous disons que $\varphi$ est {\it injectif} si, pour tout objet
$U$ de $\mathcal{C}$, l’application $\varphi : \mathcal{F}(U)
\to \mathcal{G}(U)$ est injective.
\item Nous disons que $\varphi$ est {\it surjectif} si, pour tout objet
$U$ de $\mathcal{C}$ et toute section $s\in \mathcal{G}(U)$,
il existe un recouvrement $\{U_i \to U\}$ tel que, pour tout
$i$, la restriction $s|_{U_i}$ appartienne à l’image de
$\varphi : \mathcal{F}(U_i) \to \mathcal{G}(U_i)$.
\end{enumerate}
\end{definition}

\begin{lemma}
\label{lemma-mono-epi-sheaves}
Les morphismes injectifs (resp.\ surjectifs) définis ci-dessus
sont exactement les monomorphismes (resp.\ épimorphismes) de
la catégorie $\Sh(\mathcal{C})$. Un morphisme de faisceaux
est un isomorphisme si et seulement s’il est à la fois injectif
et surjectif.
\end{lemma}

\begin{proof}
Omis.
\end{proof}

\begin{lemma}
\label{lemma-coequalizer-surjection}
Soit $\mathcal{C}$ un site. Soit $\mathcal{F} \to \mathcal{G}$
un morphisme surjectif de faisceaux d’ensembles. Alors le diagramme
$$
\xymatrix{
\mathcal{F} \times_\mathcal{G} \mathcal{F}
\ar@<1ex>[r] \ar@<-1ex>[r]
&
\mathcal{F} \ar[r]
&
\mathcal{G}}
$$
présente $\mathcal{G}$ comme un coégalisateur.
\end{lemma}

\begin{proof}
Soit $\mathcal{H}$ un faisceau d’ensembles et soit
$\varphi : \mathcal{F} \to \mathcal{H}$ un morphisme de faisceaux qui égalise
les deux morphismes $\mathcal{F} \times_\mathcal{G} \mathcal{F} \to \mathcal{F}$.
Soit $\mathcal{G}' \subset \mathcal{G}$ l’image préfaisceautique du
morphisme $\mathcal{F} \to \mathcal{G}$. Comme le produit
$\mathcal{F} \times_\mathcal{G} \mathcal{F}$ peut être calculé dans la
catégorie des préfaisceaux, il est égal au produit préfaisceautique
$\mathcal{F} \times_{\mathcal{G}'} \mathcal{F}$. Ainsi, $\varphi$
induit un unique morphisme de préfaisceaux $\psi' : \mathcal{G}' \to \mathcal{H}$.
Puisque $\mathcal{G}$ est le faisceau associé à $\mathcal{G}'$ d’après
le Lemme \ref{lemma-mono-epi-sheaves},
nous en déduisons que $\psi'$ se prolonge de manière unique en un morphisme de faisceaux
$\psi : \mathcal{G} \to \mathcal{H}$. Nous omettons de vérifier que
$\varphi$ est égal à la composée de $\psi$ et du morphisme donné.
\end{proof}














\section{Faisceaux représentables}
\label{section-representable-sheaves}

\noindent
Soit $\mathcal{C}$ une catégorie. La topologie canonique est
la topologie la plus fine pour laquelle tous les préfaisceaux représentables
sont des faisceaux (elle est formellement définie dans la
Définition \ref{definition-canonical-topology}, mais nous n’en aurons pas
besoin).
Cette topologie n’est pas toujours celle qui est associée à la
structure de site sur $\mathcal{C}$.
Nous allons donner une famille de recouvrements qui engendre cette topologie
lorsque $\mathcal{C}$ possède des produits fibrés. Donnons d’abord
la définition générale suivante.

\begin{definition}
\label{definition-universal-effective-epimorphisms}
Soit $\mathcal{C}$ une catégorie. Nous disons qu’une famille $\{U_i \to U\}_{i \in I}$
est une {\it famille épimorphique effective} si tous les morphismes $U_i \to U$ sont
représentables (voir
Catégories, Définition \ref{categories-definition-representable-morphism}),
et si, pour tout $X\in \Ob(\mathcal{C})$, la suite
$$
\xymatrix{
\Mor_\mathcal{C}(U, X) \ar[r]
&
\prod\nolimits_{i \in I} \Mor_\mathcal{C}(U_i, X)
\ar@<1ex>[r] \ar@<-1ex>[r]
&
\prod\nolimits_{(i, j) \in I^2} \Mor_\mathcal{C}(U_i \times_U U_j, X)
}
$$
est un diagramme égalisateur. Nous disons qu’une famille $\{U_i \to U\}$ est une
{\it famille épimorphique effective universelle} si, pour tout morphisme $V \to U$,
le changement de base $\{U_i \times_U V \to V\}$ est une famille épimorphique effective.
\end{definition}

\noindent
La classe des familles épimorphiques effectives universelles
satisfait les axiomes de la Définition \ref{definition-site}.
Si $\mathcal{C}$ possède des produits fibrés, la topologie associée est
la topologie canonique. (Dans ce cas, pour obtenir un site, on raisonne comme dans Ensembles,
Lemme \ref{sets-lemma-coverings-site}.)

\medskip\noindent
Réciproquement, supposons que $\mathcal{C}$ soit un site tel que
tous les préfaisceaux représentables soient des faisceaux. Alors, manifestement, tous
les recouvrements sont des familles épimorphiques effectives universelles.
La définition suivante est donc la définition « correcte » dans le
cadre des sites.

\begin{definition}
\label{definition-weaker-than-canonical}
Nous disons que la topologie d’un site $\mathcal{C}$ est
{\it moins fine que la topologie canonique}, ou qu’elle est
{\it sous-canonique}, si tous les recouvrements
de $\mathcal{C}$ sont des familles épimorphiques effectives universelles.
\end{definition}

\noindent
Un faisceau représentable est un préfaisceau représentable qui est aussi un
faisceau. Comme il vaut peut-être mieux éviter cette terminologie lorsque la
topologie n’est pas sous-canonique, nous ne la définissons formellement que dans ce cas.

\begin{definition}
\label{definition-representable-sheaf}
Soit $\mathcal{C}$ un site dont la topologie est sous-canonique.
Le plongement de Yoneda $h$ (voir
Catégories, section \ref{categories-section-opposite})
présente $\mathcal{C}$ comme une sous-catégorie pleine de la
catégorie des faisceaux sur $\mathcal{C}$. Dans ce cas,
nous appelons les faisceaux de la forme $h_U$, avec $U \in \Ob(\mathcal{C})$,
des {\it faisceaux représentables} sur $\mathcal{C}$.
Notation : le faisceau représentable $h_U$ associé à $U$ est parfois
noté {\it $\underline{U}$}.
\end{definition}

\noindent
Remarquons que, dans la situation de la définition,
$$
\Mor_{\Sh(\mathcal{C})}(h_U, \mathcal{F}) = \mathcal{F}(U)
$$
pour tout faisceau $\mathcal{F}$, puisque cette égalité vaut pour les préfaisceaux ; voir
(\ref{equation-map-representable-into-presheaf}). En général, les
préfaisceaux $h_U$ ne sont pas des faisceaux et il faut leur associer un
faisceau. Dans ce cas, nous avons encore
\begin{equation}
\label{equation-map-representable-into-sheaf}
\Mor_{\Sh(\mathcal{C})}(h_U^\#, \mathcal{F}) =
\Mor_{\textit{PSh}(\mathcal{C})}(h_U, \mathcal{F}) =
\mathcal{F}(U)
\end{equation}
pour tout faisceau $\mathcal{F}$. En effet, la première égalité résulte
de la propriété d’adjonction de $\#$ et la seconde est
(\ref{equation-map-representable-into-presheaf}).

\begin{lemma}
\label{lemma-covering-surjective-after-sheafification}
\begin{slogan}
Les recouvrements deviennent surjectifs après passage au faisceau associé.
\end{slogan}
Soit $\mathcal{C}$ un site. Si
$\{U_i \to U\}_{i \in I}$ est un recouvrement du site
$\mathcal{C}$, alors le morphisme de préfaisceaux d’ensembles
$$
\coprod\nolimits_{i \in I} h_{U_i} \to h_U
$$
devient surjectif après passage aux faisceaux associés.
\end{lemma}

\begin{proof}
D’après le Lemme \ref{lemma-mono-epi-sheaves} ci-dessus, il suffit de montrer que
$\coprod\nolimits_{i \in I} h_{U_i}^\# \to h_U^\#$
est un épimorphisme. Soit $\mathcal{F}$ un faisceau d’ensembles.
Un morphisme $h_U^\# \to \mathcal{F}$
correspond à une section $s \in \mathcal{F}(U)$.
L’injectivité de $\Mor(h_U^\#, \mathcal{F})
\to \prod_i \Mor(h_{U_i}^\#, \mathcal{F})$ découle donc
directement de la propriété de faisceau de $\mathcal{F}$.
\end{proof}

\noindent
Le lemme suivant affirme, lorsque la topologie est moins fine que la
topologie canonique, que tout faisceau est en quelque sorte construit à partir de
faisceaux représentables.

\begin{lemma}
\label{lemma-sheaf-coequalizer-representable}
Soit $\mathcal{C}$ un site. Soit $E \subset \Ob(\mathcal{C})$ un
sous-ensemble tel que tout objet de $\mathcal{C}$ admette un recouvrement par
des éléments de $E$. Soit $\mathcal{F}$ un faisceau d’ensembles. Il existe un
diagramme de faisceaux d’ensembles
$$
\xymatrix{
\mathcal{F}_1 \ar@<1ex>[r] \ar@<-1ex>[r] &
\mathcal{F}_0 \ar[r] &
\mathcal{F}
}
$$
qui présente $\mathcal{F}$ comme un coégalisateur,
tel que les $\mathcal{F}_i$, $i = 0, 1$, soient des sommes disjointes
de faisceaux de la forme $h_U^\#$ avec $U \in E$.
\end{lemma}

\begin{proof}
Montrons d’abord qu’il existe un épimorphisme $\mathcal{F}_0 \to \mathcal{F}$
du type voulu. Il suffit de prendre
$$
\mathcal{F}_0 =
\coprod\nolimits_{U \in E, s \in \mathcal{F}(U)}
(h_U)^\# \longrightarrow \mathcal{F}
$$
Ici, la restriction de la flèche à la composante correspondant à $(U, s)$ envoie
l’élément $\text{id}_U \in h_U^\#(U)$ sur la section $s \in \mathcal{F}(U)$.
Il s’agit d’un épimorphisme d’après le Lemme \ref{lemma-mono-epi-sheaves} et
notre condition sur $E$. Pour construire $\mathcal{F}_1$, posons d’abord
$\mathcal{G} = \mathcal{F}_0 \times_\mathcal{F} \mathcal{F}_0$, puis
construisons comme ci-dessus un épimorphisme $\mathcal{F}_1 \to \mathcal{G}$.
Voir le Lemme \ref{lemma-coequalizer-surjection}.
\end{proof}

\begin{lemma}
\label{lemma-colimit-representable}
Soit $\mathcal{C}$ un site. Soit $\mathcal{F}$ un faisceau d’ensembles sur
$\mathcal{C}$. Il existe alors un diagramme $\mathcal{I} \to \mathcal{C}$,
$i \mapsto U_i$, tel que
$$
\mathcal{F} = \colim_{i \in \mathcal{I}} h_{U_i}^\#
$$
De plus, si $E \subset \Ob(\mathcal{C})$ est un sous-ensemble tel que tout
objet de $\mathcal{C}$ admette un recouvrement par des éléments de $E$, nous pouvons
supposer que $U_i$ est un élément de $E$ pour tout $i \in \Ob(\mathcal{I})$.
\end{lemma}

\begin{proof}
Soit $\mathcal{I}$ la catégorie dont les objets sont les couples
$(U, s)$, avec $U \in \Ob(\mathcal{C})$ et $s \in \mathcal{F}(U)$,
et dont les morphismes $(U, s) \to (U', s')$ sont les morphismes $f : U \to U'$
de $\mathcal{C}$ tels que $f^*s' = s$. Pour tout objet $(U, s)$ de $\mathcal{I}$,
l’élément $s$ définit, par le lemme de Yoneda, un morphisme $s : h_U \to \mathcal{F}$
de préfaisceaux, et donc, par la propriété universelle du faisceau associé, un morphisme
$h_U^\# \to \mathcal{F}$. On voit immédiatement que ces morphismes sont compatibles
avec les morphismes de la catégorie $\mathcal{I}$. Nous obtenons ainsi
un morphisme $\colim_{(U, s)} h_U \to \mathcal{F}$ de préfaisceaux (où la
limite inductive est prise dans la catégorie des préfaisceaux) et un morphisme
$\colim_{(U, s)} (h_U)^\# \to \mathcal{F}$ de faisceaux (où la
limite inductive est prise dans la catégorie des faisceaux).
Puisque le passage au faisceau associé est adjoint à gauche du foncteur d’inclusion
$\Sh(\mathcal{C}) \to \textit{PSh}(\mathcal{C})$
(Proposition \ref{proposition-sheafification-adjoint}), nous avons
$\colim (h_U)^\# = (\colim h_U)^\#$ d’après
Catégories, Lemme \ref{categories-lemma-adjoint-exact}.
Il suffit donc de montrer que $\colim_{(U, s)} h_U \to \mathcal{F}$
est un isomorphisme de préfaisceaux. Pour cela, montrons que, pour tout objet
$X$ de $\mathcal{C}$, l’application $\colim_{(U, s)} h_U(X) \to \mathcal{F}(X)$
est bijective. Elle admet en effet pour application inverse celle qui envoie l’élément
$t \in \mathcal{F}(X)$ sur l’élément de l’ensemble $\colim_{(U, s)} h_U(X)$
correspondant à $(X, t)$ et à $\text{id}_X \in h_X(X)$.

\medskip\noindent
Nous omettons la démonstration de la dernière assertion.
\end{proof}




\section{Foncteurs continus}
\label{section-continuous-functors}

\begin{definition}
\label{definition-continuous}
Soient $\mathcal{C}$ et $\mathcal{D}$ des sites.
Un foncteur $u : \mathcal{C} \to \mathcal{D}$ est dit
{\it continu} si, pour tout
$\{V_i \to V\}_{i\in I} \in \text{Cov}(\mathcal{C})$,
les conditions suivantes sont satisfaites :
\begin{enumerate}
\item $\{u(V_i) \to u(V)\}_{i\in I}$ appartient à $\text{Cov}(\mathcal{D})$, et
\item pour tout morphisme $T \to V$ dans $\mathcal{C}$, le morphisme
$u(T \times_V V_i) \to u(T) \times_{u(V)} u(V_i)$ est un isomorphisme.
\end{enumerate}
\end{definition}

\noindent
Rappelons qu’étant donné un foncteur $u$ comme ci-dessus et un préfaisceau d’ensembles
$\mathcal{F}$ sur $\mathcal{D}$, nous avons défini
$u^p\mathcal{F}$ comme le préfaisceau
$\mathcal{F} \circ u$, autrement dit
$$
u^p\mathcal{F} (V) = \mathcal{F}(u(V))
$$
pour tout objet $V$ de $\mathcal{C}$.

\begin{lemma}
\label{lemma-pushforward-sheaf}
Soient $\mathcal{C}$ et $\mathcal{D}$ des sites.
Soit $u : \mathcal{C} \to \mathcal{D}$ un foncteur continu.
Si $\mathcal{F}$ est un faisceau sur $\mathcal{D}$, alors
$u^p\mathcal{F}$ est également un faisceau.
\end{lemma}

\begin{proof}
Soit $\{V_i \to V\}$ un recouvrement.
Par hypothèse, $\{u(V_i) \to u(V)\}$ est un recouvrement
dans $\mathcal{D}$ et $u(V_i \times_V V_j) =
u(V_i)\times_{u(V)}u(V_j)$. La condition de faisceau pour
$u^p\mathcal{F}$ et le recouvrement $\{V_i \to V\}$
est donc précisément la même que la condition de faisceau pour $\mathcal{F}$
et le recouvrement $\{u(V_i) \to u(V)\}$.
\end{proof}

\noindent
Pour éviter toute confusion, nous notons parfois
$$
u^s :
\Sh(\mathcal{D})
\longrightarrow
\Sh(\mathcal{C})
$$
la restriction du foncteur $u^p$ à la sous-catégorie des faisceaux d’ensembles.
Rappelons que $u^p$ admet un adjoint à gauche
$u_p : \textit{PSh}(\mathcal{C}) \to \textit{PSh}(\mathcal{D})$ ; voir
la section \ref{section-functoriality-PSh}.

\begin{lemma}
\label{lemma-adjoint-sheaves}
Dans la situation du Lemme \ref{lemma-pushforward-sheaf},
le foncteur $u_s : \mathcal{G} \mapsto (u_p \mathcal{G})^\#$
est adjoint à gauche de $u^s$.
\end{lemma}

\begin{proof}
Cela résulte directement du Lemme \ref{lemma-adjoints-u} et de la
Proposition \ref{proposition-sheafification-adjoint}.
\end{proof}

\noindent
Voici un lemme technique.

\begin{lemma}
\label{lemma-technical-up}
Dans la situation du Lemme \ref{lemma-pushforward-sheaf},
pour tout préfaisceau $\mathcal{G}$ sur $\mathcal{C}$,
nous avons $(u_p\mathcal{G})^\# = (u_p(\mathcal{G}^\#))^\#$.
\end{lemma}

\begin{proof}
Pour tout faisceau $\mathcal{F}$ sur $\mathcal{D}$, nous avons
\begin{eqnarray*}
\Mor_{\Sh(\mathcal{D})}(u_s(\mathcal{G}^\#), \mathcal{F})
& = &
\Mor_{\Sh(\mathcal{C})}(\mathcal{G}^\#, u^s\mathcal{F}) \\
& = &
\Mor_{\textit{PSh}(\mathcal{C})}(\mathcal{G}^\#, u^p\mathcal{F}) \\
& = &
\Mor_{\textit{PSh}(\mathcal{C})}(\mathcal{G}, u^p\mathcal{F}) \\
& = &
\Mor_{\textit{PSh}(\mathcal{D})}(u_p\mathcal{G}, \mathcal{F}) \\
& = &
\Mor_{\Sh(\mathcal{D})}((u_p\mathcal{G})^\#, \mathcal{F})
\end{eqnarray*}
et le résultat découle du lemme de Yoneda.
\end{proof}

\begin{lemma}
\label{lemma-pullback-representable-sheaf}
Soit $u : \mathcal{C} \to \mathcal{D}$ un foncteur continu
entre sites.
Pour tout objet $U$ de $\mathcal{C}$, nous avons $u_sh_U^\# = h_{u(U)}^\#$.
\end{lemma}

\begin{proof}
Cela résulte des
Lemmes \ref{lemma-pullback-representable-presheaf}
et \ref{lemma-technical-up}.
\end{proof}


\begin{remark}
\label{remark-quasi-continuous}
(À omettre en première lecture.)
Soient $\mathcal{C}$ et $\mathcal{D}$ des sites. Utilisons
la définition des familles de morphismes tautologiquement équivalentes,
voir la Définition \ref{definition-combinatorial-tautological},
pour affaiblir (légèrement) les conditions qui définissent la continuité.
Soit $u : \mathcal{C} \to \mathcal{D}$ un foncteur.
Disons que $u$ est {\it quasi-continu} si, pour tout
$\mathcal{V} = \{V_i \to V\}_{i\in I} \in \text{Cov}(\mathcal{C})$,
les conditions suivantes sont satisfaites :
\begin{enumerate}
\item[(1')] la famille de morphismes
$\{u(V_i) \to u(V)\}_{i\in I}$ est tautologiquement équivalente
à un élément de $\text{Cov}(\mathcal{D})$, et
\item[(2)] pour tout morphisme $T \to V$ dans $\mathcal{C}$, le morphisme
$u(T \times_V V_i) \to u(T) \times_{u(V)} u(V_i)$ est un isomorphisme.
\end{enumerate}
Nous allons voir que les Lemmes \ref{lemma-pushforward-sheaf}
et \ref{lemma-adjoint-sheaves} restent valables lorsque
$u$ est quasi-continu.

\medskip\noindent
Remarquons d’abord que les morphismes $u(V_i) \to u(V)$ sont représentables, car
ils sont isomorphes à des morphismes représentables (d’après la première condition).
En particulier, la famille $u(\mathcal{V}) = \{u(V_i) \to u(V)\}_{i\in I}$
donne naissance à un groupe de cohomologie de {\v C}ech de degré zéro
$H^0(u(\mathcal{V}), \mathcal{F})$ pour tout préfaisceau $\mathcal{F}$ sur
$\mathcal{D}$.
Soit $\mathcal{U} = \{U_j \to u(V)\}_{j \in J}$ un élément
de $\text{Cov}(\mathcal{D})$ tautologiquement
équivalent à $\{u(V_i) \to u(V)\}_{i \in I}$. Remarquons que $u(\mathcal{V})$
est un raffinement de $\mathcal{U}$ et réciproquement. La Remarque
\ref{remark-both-refine-same-H0} donne donc
$H^0(u(\mathcal{V}), \mathcal{F}) = H^0(\mathcal{U}, \mathcal{F})$.
En particulier, si $\mathcal{F}$ est un faisceau, alors
$\mathcal{F}(u(V)) = H^0(u(\mathcal{V}), \mathcal{F})$ en vertu de
la propriété de faisceau exprimée au moyen des groupes de cohomologie de
{\v C}ech de degré zéro. Nous en déduisons que $u^p\mathcal{F}$ est un faisceau
si $\mathcal{F}$ est un faisceau, puisque $H^0(\mathcal{V}, u^p\mathcal{F}) =
H^0(u(\mathcal{V}), \mathcal{F})$, lequel est, comme nous venons de l’observer,
égal à $\mathcal{F}(u(V)) = u^p\mathcal{F}(V)$. Le Lemme
\ref{lemma-pushforward-sheaf} est donc valable. Le Lemme
\ref{lemma-adjoint-sheaves} en découle immédiatement.
\end{remark}






\section{Morphismes de sites}
\label{section-morphism-sites}

\begin{definition}
\label{definition-morphism-sites}
Soient $\mathcal{C}$ et $\mathcal{D}$ des sites.
Un {\it morphisme de sites} $f : \mathcal{D} \to \mathcal{C}$
est la donnée d’un foncteur continu $u : \mathcal{C} \to \mathcal{D}$
tel que le foncteur $u_s$ soit exact.
\end{definition}

\noindent
Notons que le foncteur $u$ va dans le sens {\it opposé}
au morphisme $f$. Si $f \leftrightarrow u$ est un morphisme de sites,
nous employons les notations $f^{-1} = u_s$ et $f_* = u^s$.
Le foncteur $f^{-1}$ est appelé {\it foncteur image inverse} et
le foncteur $f_*$ {\it foncteur image directe}.
Comme en topologie, nous avons la propriété d’adjonction suivante :
$$
\Mor_{\Sh(\mathcal{D})}(f^{-1}\mathcal{G}, \mathcal{F})
=
\Mor_{\Sh(\mathcal{C})}(\mathcal{G}, f_*\mathcal{F})
$$
La motivation de cette définition vient de l’exemple suivant.

\begin{example}
\label{example-continuous-map}
Soit $f : X  \to Y$ une application continue d’espaces topologiques.
Rappelons que nous disposons des sites $X_{Zar}$ et $Y_{Zar}$ ;
voir l’Exemple \ref{example-site-topological}. Considérons le foncteur
$u : Y_{Zar} \to X_{Zar}$, $V \mapsto f^{-1}(V)$.
Ce foncteur est manifestement continu puisque les images inverses de
recouvrements ouverts sont des recouvrements ouverts. (En réalité, cela dépend de la
manière dont on a choisi les ensembles de recouvrements de $X_{Zar}$ et de $Y_{Zar}$.
Mais, dans tous les cas, le foncteur est quasi-continu ; voir la Remarque
\ref{remark-quasi-continuous}.)
Il est facile de vérifier que le foncteur $u^s$ est égal au foncteur image directe
usuel $f_*$ de la topologie. Ainsi, puisque $u_s$ est un adjoint et que
le foncteur image inverse topologique usuel $f^{-1}$ est lui aussi un adjoint,
nous obtenons un isomorphisme canonique $f^{-1} \cong u_s$. Comme $f^{-1}$
est exact, nous en déduisons que $u_s$ est exact. Le foncteur $u$ définit donc
un morphisme de sites $f : X_{Zar} \to Y_{Zar}$, que nous pouvons également noter
$f$, puisque nous avons déjà vu que les foncteurs $u_s, u^s$ coïncident
avec les notions usuelles correspondantes.
\end{example}

\begin{example}
\label{example-finer-topology}
Soit $\mathcal{C}$ une catégorie. Soient
$$
\text{Cov}(\mathcal{C}) \supset \text{Cov}'(\mathcal{C})
$$
deux ensembles de familles de morphismes de même but
qui munissent $\mathcal{C}$ d’une structure de site. Notons $\mathcal{C}_\tau$
le site correspondant à $\text{Cov}(\mathcal{C})$ et
$\mathcal{C}_{\tau'}$ le site correspondant à $\text{Cov}'(\mathcal{C})$.
Nous affirmons que le foncteur identité de $\mathcal{C}$ définit un morphisme de sites
$$
\epsilon : \mathcal{C}_\tau \longrightarrow \mathcal{C}_{\tau'}
$$
En effet, $\text{id} : \mathcal{C}_{\tau'} \to \mathcal{C}_\tau$
est continu puisque tout recouvrement pour $\tau'$ est un recouvrement pour $\tau$.
Le foncteur $\epsilon_* = \text{id}^s$ est donc le foncteur identité sur les
préfaisceaux sous-jacents. Par conséquent, l’adjoint à gauche $\epsilon^{-1}$
de $\epsilon_*$ envoie un $\tau'$-faisceau $\mathcal{F}$ sur le
$\tau$-faisceau associé à $\mathcal{F}$ (d’après la
propriété universelle du faisceau associé).
Les limites finies de $\tau'$-faisceaux coïncident avec les limites finies
de préfaisceaux (Lemme \ref{lemma-limit-sheaf}) et le passage au
$\tau$-faisceau associé commute aux limites finies
(Lemme \ref{lemma-sheafification-exact}).
Ainsi, $\epsilon^{-1}$ est exact à gauche. Puisque $\epsilon^{-1}$
est adjoint à gauche, il est aussi exact à droite
(Catégories, Lemme \ref{categories-lemma-exact-adjoint}).
Par conséquent, $\epsilon^{-1}$ est exact et toutes les
conditions de la Définition \ref{definition-morphism-sites} sont vérifiées.
\end{example}

\begin{lemma}
\label{lemma-composition-morphisms-sites}
\begin{slogan}
La composition des foncteurs de sites préserve leur continuité.
\end{slogan}
Soient $\mathcal{C}_i$, $i = 1, 2, 3$, des sites. Soient
$u : \mathcal{C}_2 \to \mathcal{C}_1$ et
$v : \mathcal{C}_3 \to \mathcal{C}_2$ des foncteurs continus
qui induisent des morphismes de sites. Alors le foncteur
$u \circ v : \mathcal{C}_3 \to \mathcal{C}_1$ est continu et
définit un morphisme de sites $\mathcal{C}_1 \to \mathcal{C}_3$.
\end{lemma}

\begin{proof}
Il résulte immédiatement des définitions que $u \circ v$ est un foncteur continu.
De plus, nous avons manifestement $(u \circ v)^p = v^p \circ u^p$, et donc
$(u \circ v)^s = v^s \circ u^s$. Les foncteurs $(u \circ v)_s$ et
$u_s \circ v_s$ sont donc tous deux adjoints à gauche de $(u \circ v)^s$. Par conséquent,
$(u \circ v)_s \cong u_s \circ v_s$, et nous concluons que $(u \circ v)_s$
est exact en tant que composé de foncteurs exacts.
\end{proof}

\begin{definition}
\label{definition-composition-morphisms-sites}
Soient $\mathcal{C}_i$, $i = 1, 2, 3$, des sites. Soient
$f : \mathcal{C}_1 \to \mathcal{C}_2$ et
$g : \mathcal{C}_2 \to \mathcal{C}_3$ des morphismes de sites
donnés par les foncteurs continus $u : \mathcal{C}_2 \to \mathcal{C}_1$
et $v : \mathcal{C}_3 \to \mathcal{C}_2$. Le {\it composé}
$g \circ f$ est le morphisme de sites correspondant au
foncteur $u \circ v$.
\end{definition}

\noindent
Dans cette situation, nous avons $(g \circ f)_* = g_* \circ f_*$ et
$(g \circ f)^{-1} = f^{-1} \circ g^{-1}$ (voir la démonstration du
Lemme \ref{lemma-composition-morphisms-sites}).

\begin{lemma}
\label{lemma-directed-morphism}
Soient $\mathcal{C}$ et $\mathcal{D}$ des sites. Soit
$u : \mathcal{C} \to \mathcal{D}$ un foncteur continu.
Supposons que toutes les catégories $(\mathcal{I}_V^u)^{opp}$ de la
section \ref{section-functoriality-PSh}
soient filtrantes. Alors $u$ définit un morphisme de sites $\mathcal{D} \to
\mathcal{C}$ ; autrement dit, $u_s$ est exact.
\end{lemma}

\begin{proof}
Comme $u_s$ est adjoint à gauche de $u^s$, nous voyons que $u_s$ est exact à droite ;
voir Catégories, Lemme \ref{categories-lemma-exact-adjoint}.
Il suffit donc de montrer que $u_s$ est exact à gauche. Autrement dit,
nous devons montrer que $u_s$ commute aux limites finies.
Puisque les catégories $\mathcal{I}_Y^{opp}$ sont filtrantes,
$u_p$ commute aux limites finies ; voir
Catégories, Lemme \ref{categories-lemma-directed-commutes}
(on utilise également ici la description des limites dans $\textit{PSh}$ ;
voir la section \ref{section-limits-colimits-PSh}).
Comme le passage au faisceau associé commute lui aussi aux limites finies
(Lemme \ref{lemma-sheafification-exact}), nous concluons en utilisant
$u_s = \# \circ u_p$.
\end{proof}

\begin{proposition}
\label{proposition-get-morphism}
Soient $\mathcal{C}$ et $\mathcal{D}$ des sites. Soit
$u : \mathcal{C} \to \mathcal{D}$ un foncteur continu.
Supposons en outre que :
\begin{enumerate}
\item la catégorie $\mathcal{C}$ possède un objet final $X$ et
$u(X)$ est un objet final de $\mathcal{D}$, et
\item la catégorie $\mathcal{C}$ possède des produits fibrés et
$u$ commute à ceux-ci.
\end{enumerate}
Alors $u$ définit un morphisme de sites $\mathcal{D} \to
\mathcal{C}$ ; autrement dit, $u_s$ est exact.
\end{proposition}

\begin{proof}
Cela résulte des Lemmes \ref{lemma-directed} et
\ref{lemma-directed-morphism}.
\end{proof}

\begin{remark}
\label{remark-explain-left-exact}
Les conditions de la Proposition \ref{proposition-get-morphism} ci-dessus
équivalent à dire que $u$ est exact à gauche, c’est-à-dire qu’il commute
aux limites finies. Voir
Catégories, Lemmes
\ref{categories-lemma-finite-limits-exist} et
\ref{categories-lemma-characterize-left-exact}.
Il paraît plus naturel de les formuler au moyen des objets finaux
et des produits fibrés, car cette formulation semble avoir un sens plus géométrique
dans les exemples.
\end{remark}

\noindent
Le Lemme \ref{lemma-continuous-with-continuous-left-adjoint}
fournira une autre manière de démontrer qu’un foncteur continu
donne naissance à un morphisme de sites.

\begin{remark}
\label{remark-quasi-continuous-morphism-sites}
(À omettre en première lecture.)
Soient $\mathcal{C}$ et $\mathcal{D}$ des sites. Par analogie avec la
Définition \ref{definition-morphism-sites}, nous disons qu’un
{\it quasi-morphisme de sites $f : \mathcal{D} \to \mathcal{C}$}
est donné par un foncteur quasi-continu $u : \mathcal{C} \to \mathcal{D}$
(voir la Remarque \ref{remark-quasi-continuous}) tel que $u_s$ soit exact.
L’analogue de la Proposition \ref{proposition-get-morphism} dans ce
cadre s’obtient en remplaçant le mot « continu »
par « quasi-continu » et le mot
« morphisme » par « quasi-morphisme ». La démonstration est littéralement
la même.
\end{remark}

\noindent
Dans la Définition \ref{definition-morphism-sites}, la condition d’exactitude de
$u_s$ ne peut être omise. Par exemple, la conclusion du lemme suivant
peut être fausse si l’on suppose seulement que $u$ est continu.

\begin{lemma}
\label{lemma-morphism-of-sites-covering}
Soit $f : \mathcal{D} \to \mathcal{C}$ un morphisme de sites donné par le
foncteur $u : \mathcal{C} \to \mathcal{D}$. Pour tout objet $V$ de
$\mathcal{D}$, il existe un recouvrement $\{V_j \to V\}$ tel que, pour tout
$j$, il existe un morphisme $V_j \to u(U_j)$ pour un certain objet $U_j$
de $\mathcal{C}$.
\end{lemma}

\begin{proof}
Puisque $f^{-1} = u_s$ est exact, nous avons $f^{-1}* = *$, où $*$ désigne
l’objet final de la catégorie des faisceaux
(Exemple \ref{example-singleton-sheaf}).
Comme $f^{-1}* = u_s*$ est le faisceau associé à $u_p*$, il existe
un recouvrement $\{V_j \to V\}$ tel que $(u_p*)(V_j)$ soit non vide.
Or $(u_p*)(V_j)$ est une limite inductive indexée par la catégorie
$\mathcal{I}^u_{V_j}$, dont les objets sont les morphismes $V_j \to u(U)$ ;
le lemme en découle.
\end{proof}

























\section{Topos}
\label{section-topoi}

\noindent
Voici une définition du topos adaptée à nos besoins.
Un topos est la catégorie des faisceaux sur un site. Pour spécifier
un topos, il suffit donc de spécifier le site. La véritable différence entre un topos
et un site réside dans la définition des morphismes. En effet, il existe
de nombreux morphismes de topos qui ne proviennent pas de morphismes
des sites sous-jacents.

\begin{definition}[Topos]
\label{definition-topos}
Un {\it topos} est la catégorie $\Sh(\mathcal{C})$ des faisceaux
sur un site $\mathcal{C}$.
\begin{enumerate}
\item Soient $\mathcal{C}$, $\mathcal{D}$ des sites.
Un {\it morphisme de topos} $f$ de $\Sh(\mathcal{D})$
vers $\Sh(\mathcal{C})$ est la donnée de deux foncteurs
$f_* : \Sh(\mathcal{D}) \to \Sh(\mathcal{C})$
et
$f^{-1} : \Sh(\mathcal{C}) \to \Sh(\mathcal{D})$
tels que :
\begin{enumerate}
\item nous ayons
$$
\Mor_{\Sh(\mathcal{D})}(f^{-1}\mathcal{G}, \mathcal{F})
=
\Mor_{\Sh(\mathcal{C})}(\mathcal{G}, f_*\mathcal{F})
$$
bifonctoriellement, et
\item le foncteur $f^{-1}$ commute aux limites finies, c’est-à-dire
qu’il est exact à gauche.
\end{enumerate}
\item Soient $\mathcal{C}$, $\mathcal{D}$, $\mathcal{E}$ des sites.
Étant donnés des morphismes de topos
$f :\Sh(\mathcal{D}) \to \Sh(\mathcal{C})$ et
$g :\Sh(\mathcal{E}) \to \Sh(\mathcal{D})$, le
{\it composé $f\circ g$} est le morphisme de topos défini
par les foncteurs
$(f \circ g)_* = f_* \circ g_*$ et
$(f \circ g)^{-1} = g^{-1} \circ f^{-1}$.
\end{enumerate}
\end{definition}

\noindent
Supposons que $\alpha : \mathcal{S}_1 \to \mathcal{S}_2$ soit une équivalence de
catégories (éventuellement « grandes »). Si $\mathcal{S}_1$, $\mathcal{S}_2$
sont des topos, poser $f_* = \alpha$ et prendre pour $f^{-1}$ un quasi-inverse
de $\alpha$ donne un morphisme de topos $f : \mathcal{S}_1 \to \mathcal{S}_2$.
Ce morphisme est en outre une équivalence dans la $2$-catégorie des topos (voir
la section \ref{section-2-category}).
Il est donc légitime de dire que « $\mathcal{S}$ est un topos » si $\mathcal{S}$
est équivalente à la catégorie des faisceaux sur un site (sans être nécessairement
égale à la catégorie des faisceaux sur un site). Nous emploierons parfois
cet abus de notation.

\medskip\noindent
Le {\it topos vide} est le topos
des faisceaux sur le site $\mathcal{C}$, où $\mathcal{C}$
est la catégorie vide. Nous noterons parfois ce site $\emptyset$.
Ce site possède un unique faisceau (puisque $\emptyset$ n’a aucun
objet). Ainsi, $\Sh(\emptyset)$ est équivalente à la catégorie
qui possède un seul objet et un seul morphisme.

\medskip\noindent
Le {\it topos ponctuel} est le topos des faisceaux sur le site
$\mathcal{C}$ qui possède un unique objet $pt$, un unique morphisme
$\text{id}_{pt}$, et dont le seul recouvrement est
$\{\text{id}_{pt}\}$. Nous noterons simplement ce site $pt$.
Il est clair que la catégorie des faisceaux $ = $ la catégorie des
préfaisceaux $ = $ la catégorie des ensembles.
En formule, $\Sh(pt) = \textit{Ensembles}$.

\medskip\noindent
Soient $\mathcal{C}$ et $\mathcal{D}$ des sites. Soit
$f : \Sh(\mathcal{D}) \to \Sh(\mathcal{C})$ un morphisme de topos.
Remarquons que $f_*$ commute à toutes les limites et que
$f^{-1}$ commute à toutes les limites inductives ; voir Catégories,
Lemme \ref{categories-lemma-adjoint-exact}.
En particulier, la condition imposée à $f^{-1}$ dans la définition ci-dessus
garantit que $f^{-1}$ est exact. Les morphismes de topos sont souvent construits
à l’aide soit du Lemme \ref{lemma-cocontinuous-morphism-topoi},
soit du lemme suivant.

\begin{lemma}
\label{lemma-morphism-sites-topoi}
Étant donné un morphisme de sites $f : \mathcal{D} \to \mathcal{C}$
correspondant au foncteur $u : \mathcal{C} \to \mathcal{D}$,
le couple de foncteurs $(f^{-1} = u_s, f_* = u^s)$ est un morphisme de topos
$\Sh(\mathcal{D}) \to \Sh(\mathcal{C})$.
\end{lemma}

\begin{proof}
Cela résulte immédiatement de la Définition \ref{definition-morphism-sites}.
\end{proof}

\begin{remark}
\label{remark-pt-topos}
De nombreux sites donnent naissance au topos $\Sh(pt)$.
Voici un exemple utile. Supposons que $S$ soit un ensemble (d’ensembles)
qui contienne au moins un élément non vide. Soit $\mathcal{S}$ la
catégorie dont les objets sont les éléments de $S$ et dont les morphismes sont
les applications ensemblistes arbitraires. Supposons que $\mathcal{S}$ possède des produits fibrés.
Ce sera par exemple le cas si $S = \mathcal{P}(\text{ensemble infini})$
est l’ensemble des parties d’un ensemble infini quelconque (exercice de théorie des ensembles).
Munissons $\mathcal{S}$ d’une structure de site en déclarant que
les familles surjectives d’applications sont les recouvrements (et choisissons
un ensemble suffisamment grand convenable de familles de recouvrement, comme dans
Ensembles, section \ref{sets-section-coverings-site}).
Nous affirmons que $\Sh(\mathcal{S})$ est équivalente à la catégorie des
ensembles.

\medskip\noindent
Démontrons-le d’abord lorsque $S$ contient un singleton $e \in S$.
Dans ce cas, il existe une équivalence de topos
$i : \Sh(pt) \to \Sh(\mathcal{S})$ donnée par
les foncteurs
\begin{equation}
\label{equation-sheaves-pt-sets}
i^{-1}\mathcal{F} = \mathcal{F}(e), \quad
i_*E = (U \mapsto \Mor_{\textit{Ensembles}}(U, E))
\end{equation}
En effet, supposons que $\mathcal{F}$ soit un faisceau sur $\mathcal{S}$.
Pour tout $U \in \Ob(\mathcal{S}) = S$, nous pouvons trouver un recouvrement
$\{\varphi_u : e \to U\}_{u \in U}$, où $\varphi_u$
envoie l’unique élément de $e$ sur $u \in U$. Dans ce cas, la condition de faisceau
implique que
$\mathcal{F}(U) = \prod_{u \in U} \mathcal{F}(e)$.
Autrement dit,
$\mathcal{F}(U) = \Mor_{\textit{Ensembles}}(U, \mathcal{F}(e))$.
Cette règle est en outre compatible avec les applications de restriction. Le foncteur
$$
i_* :
\textit{Ensembles} = \Sh(pt)
\longrightarrow
\Sh(\mathcal{S}), \quad
E \longmapsto (U \mapsto \Mor_{\textit{Ensembles}}(U, E))
$$
est donc une équivalence de catégories, et son inverse est le foncteur
$i^{-1}$ donné ci-dessus.

\medskip\noindent
Si $\mathcal{S}$ ne contient aucun singleton, le foncteur
$i_*$ défini ci-dessus a encore un sens. Pour montrer qu’il reste
une équivalence dans ce cas, choisissons un élément non vide $\tilde e \in S$
et une application $\varphi : \tilde e \to \tilde e$ dont l’image est un singleton.
Pour tout faisceau $\mathcal{F}$, posons
$$
\mathcal{F}(e) :=
\Im(
\mathcal{F}(\varphi) :
\mathcal{F}(\tilde e)
\longrightarrow
\mathcal{F}(\tilde e)
)
$$
et montrons qu’il s’agit d’un quasi-inverse de $i_*$. Les détails sont omis.
\end{remark}

\begin{remark}
\label{remark-morphism-topoi-big}
(Questions ensemblistes relatives aux morphismes de topos. À omettre
en première lecture.)
Un morphisme de topos tel que défini ci-dessus n’est pas un ensemble, mais une classe.
Autrement dit, il est donné par une formule mathématique plutôt que par un
objet mathématique. Bien que nous puissions considérer
la classe de tous les morphismes entre deux topos donnés,
il n’est pas judicieux de l’introduire comme objet mathématique.
D’autre part, supposons que $\mathcal{C}$ et $\mathcal{D}$ soient
des sites donnés. Considérons un foncteur
$\Phi : \mathcal{C} \to \Sh(\mathcal{D})$.
Une telle donnée est un ensemble, autrement dit un objet mathématique.
Nous pouvons poser successivement les questions suivantes à propos de $\Phi$.
\begin{enumerate}
\item Est-il vrai que, pour tout faisceau $\mathcal{F}$ sur $\mathcal{D}$,
la règle
$U \mapsto \Mor_{\Sh(\mathcal{D})}(\Phi(U), \mathcal{F})$
définit un faisceau sur $\mathcal{C}$ ? Si oui, elle définit un foncteur
$\Phi_* : \Sh(\mathcal{D}) \to \Sh(\mathcal{C})$.
\item Est-il vrai que $\Phi_*$ admette un adjoint à gauche ? Si oui,
notons $\Phi^{-1}$ cet adjoint à gauche.
\item Est-il vrai que $\Phi^{-1}$ soit exact ?
\end{enumerate}
Si la réponse à la dernière question est encore « oui », nous obtenons
un morphisme de topos $(\Phi_*, \Phi^{-1})$. En outre, étant donné un
morphisme de topos $(f_*, f^{-1})$, nous pouvons poser
$\Phi(U) = f^{-1}(h_U^\#)$ et obtenir un foncteur $\Phi$ comme ci-dessus,
avec $f_* \cong \Phi_*$ et $f^{-1} \cong \Phi^{-1}$ (de manière compatible
avec la propriété d’adjonction).
En définitive, en travaillant avec la classe des $\Phi$
plutôt qu’avec les morphismes de topos, nous avons (a) remplacé la notion de
morphisme de topos par un objet mathématique et (b)
la classe des $\Phi$ est une classe (et non une classe
de classes). On peut bien sûr aller plus loin, par exemple en précisant
la signification des conditions (2) et (3) ci-dessus ;
nous le faisons pour les points de topos dans la section \ref{section-points}.
\end{remark}

\begin{remark}
\label{remark-quasi-continuous-morphism-topoi}
(À omettre en première lecture.)
Soient $\mathcal{C}$ et $\mathcal{D}$ des sites.
Un quasi-morphisme de sites $f : \mathcal{D} \to \mathcal{C}$
(voir la Remarque \ref{remark-quasi-continuous-morphism-sites})
donne naissance, exactement comme dans le
Lemme \ref{lemma-morphism-sites-topoi}, à un morphisme de topos $f$ de
$\Sh(\mathcal{D})$ vers $\Sh(\mathcal{C})$.
\end{remark}










\section{G-ensembles et morphismes}
\label{section-G-sets-morphisms}

\noindent
Soit $\varphi : G \to H$ un homomorphisme de groupes.
Choisissons des sites (convenables) $\mathcal{T}_G$ et $\mathcal{T}_H$ comme dans
l’Exemple \ref{example-site-on-group} et la
section \ref{section-example-sheaf-G-sets}.
Soit $u : \mathcal{T}_H \to \mathcal{T}_G$ le foncteur qui associe
à un $H$-ensemble $U$ le $G$-ensemble $U_\varphi$, de même ensemble sous-jacent,
mais où l’action de $G$ est définie par $g \cdot \xi = \varphi(g)\xi$ pour
$g \in G$ et $\xi \in U$.
Il est clair que $u$ commute aux limites finies et qu’il est
continu\footnote{Remarque ensembliste : choisissons d’abord $\mathcal{T}_H$.
Choisissons ensuite $\mathcal{T}_G$ contenant $u(\mathcal{T}_H)$ et de sorte que
tout recouvrement de $\mathcal{T}_H$ corresponde à un recouvrement de
$\mathcal{T}_G$. Cela est possible d’après
Ensembles, Lemmes
\ref{sets-lemma-sets-with-group-action},
\ref{sets-lemma-what-is-in-it-G-sets} et
\ref{sets-lemma-coverings-site}.}.
En appliquant
la Proposition \ref{proposition-get-morphism}
et le
Lemme \ref{lemma-morphism-sites-topoi},
nous obtenons un morphisme de topos
$$
f : \Sh(\mathcal{T}_G) \longrightarrow \Sh(\mathcal{T}_H)
$$
associé à $\varphi$. La
Proposition \ref{proposition-sheaves-on-group}
montre que nous obtenons un couple de foncteurs adjoints
$$
f_* : G\textit{-Ensembles} \longrightarrow H\textit{-Ensembles}, \quad
f^{-1} : H\textit{-Ensembles} \longrightarrow G\textit{-Ensembles}.
$$
Déterminons explicitement ces foncteurs dans ce cas.

\medskip\noindent
Commençons par établir une formule pour $f_*$.
Rappelons qu’au $G$-ensemble $S$ correspond le faisceau
$\mathcal{F}_S$ sur $\mathcal{T}_G$ donné par la règle
$\mathcal{F}_S(U) = \Mor_G(U, S)$. D’autre part, au
faisceau $\mathcal{G}$ sur $\mathcal{T}_H$ correspond le $H$-ensemble
donné par la règle $\mathcal{G}({}_HH)$. Nous obtenons donc
$$
f_*S = \Mor_{G\textit{-Ensembles}}(({}_HH)_\varphi, S).
$$
En développant un peu cette formule, nous obtenons
$$
f_*S = \{ a : H \to S \mid a(gh) = ga(h) \}
$$
avec l’action à gauche de $H$ donnée par
$(h \cdot a)(h') = a(h'h)$
pour tout élément $a \in f_*S$.

\medskip\noindent
Calculons ensuite explicitement $f^{-1}$. Comme la topologie
de $\mathcal{T}_G$ et celle de $\mathcal{T}_H$ sont sous-canoniques, tous les
préfaisceaux représentables sont des faisceaux. De plus, pour tout objet $V$
de $\mathcal{T}_H$, nous avons $f^{-1}h_V$ égal à $h_{u(V)}$ (voir le
Lemme \ref{lemma-pullback-representable-sheaf}). Ainsi,
$f^{-1}S = S_\varphi$ pour les faisceaux représentables. Comme tout faisceau sur
$\mathcal{T}_H$ est une somme disjointe de faisceaux représentables, nous concluons que
cette formule vaut en général. Par conséquent, pour tout $H$-ensemble $T$,
$$
f^{-1}T = T_\varphi.
$$
L’adjonction entre $f^{-1}$ et $f_*$ est exprimée par la formule
$$
\Mor_{G\textit{-Ensembles}}(T_\varphi, S) =
\Mor_{H\textit{-Ensembles}}(T, f_*S)
$$
où $f_*S$ est défini ci-dessus. On peut le démontrer directement. Il est alors clair,
en outre, que $(f^{-1}, f_*)$ forme un couple adjoint et que $f^{-1}$ est exact.
On peut donc, au lieu de procéder comme ci-dessus, définir directement de cette
manière le morphisme de topos
$f : G\textit{-Ensembles} \to H\textit{-Ensembles}$.












\section{Objets quasi-compacts et limites inductives}
\label{section-quasi-compact}

\noindent
Afin de pouvoir employer le même langage que pour les espaces topologiques,
nous introduisons la terminologie suivante.

\begin{definition}
\label{definition-quasi-compact}
Soit $\mathcal{C}$ un site. Un objet $U$ de $\mathcal{C}$ est
{\it quasi-compact} si, pour tout recouvrement $\mathcal{U} = \{U_i \to U\}_{i \in I}$
de $\mathcal{C}$, il existe un autre recouvrement
$\mathcal{V} = \{V_j \to U\}_{j \in J}$ et un morphisme
$\mathcal{V} \to \mathcal{U}$ de familles de morphismes de même but
donné par $\text{id} : U \to U$, $\alpha : J \to I$ et $V_j \to U_{\alpha(j)}$
(voir la Définition \ref{definition-morphism-coverings}),
tels que l’image de $\alpha$ soit une partie finie de $I$.
\end{definition}

\noindent
Bien entendu, la notion usuelle suffit pour conclure que
$U$ est quasi-compact.

\begin{lemma}
\label{lemma-conclude-quasi-compact}
Soit $\mathcal{C}$ un site. Soit $U$ un objet de $\mathcal{C}$.
Considérons les conditions suivantes :
\begin{enumerate}
\item $U$ est quasi-compact,
\item pour tout recouvrement $\{U_i \to U\}_{i \in I}$ de $\mathcal{C}$,
il existe un recouvrement fini $\{V_j \to U\}_{j = 1, \ldots, m}$
de $\mathcal{C}$ qui raffine $\mathcal{U}$, et
\item pour tout recouvrement $\{U_i \to U\}_{i \in I}$ de $\mathcal{C}$,
il existe une partie finie $I' \subset I$ telle que
$\{U_i \to U\}_{i \in I'}$ soit un recouvrement de $\mathcal{C}$.
\end{enumerate}
Nous avons toujours (3) $\Rightarrow$ (2) $\Rightarrow$ (1),
mais les implications réciproques sont fausses en général.
\end{lemma}

\begin{proof}
Les implications résultent immédiatement des définitions.
Soit $X = [0, 1] \subset \mathbf{R}$
muni de sa topologie usuelle (la topologie usuelle $\epsilon$-$\delta$).
Soit $\mathcal{C}$ la catégorie des sous-espaces ouverts de $X$, les
morphismes étant les inclusions, munie des recouvrements ouverts usuels (comparer à
l’Exemple \ref{example-site-topological}). Modifions cependant la notion
de recouvrement de $\mathcal{C}$ en excluant tous les recouvrements finis, sauf
ceux de la forme $\{U \to U\}$. On vérifie facilement que l’on obtient ainsi
un site au sens de la Définition \ref{definition-site}.
Dans ce site, l’objet $X = U = [0, 1]$ est quasi-compact au sens de la
Définition \ref{definition-quasi-compact}, mais $U$ ne satisfait pas (2).
Nous laissons au lecteur le soin de construire un exemple où (2) est vraie et (3) fausse.
\end{proof}

\noindent
Voici la signification, en théorie des topos, d’un objet quasi-compact.

\begin{lemma}
\label{lemma-quasi-compact}
Soit $\mathcal{C}$ un site. Soit $U$ un objet de $\mathcal{C}$.
Les conditions suivantes sont équivalentes :
\begin{enumerate}
\item $U$ est quasi-compact, et
\item pour tout morphisme surjectif de faisceaux
$\coprod_{i \in I} \mathcal{F}_i \to h_U^\#$,
il existe une partie finie $J \subset I$ telle que
$\coprod_{i \in J} \mathcal{F}_i \to h_U^\#$ soit surjectif.
\end{enumerate}
\end{lemma}

\begin{proof}
Supposons (1) et soit $\coprod_{i \in I} \mathcal{F}_i \to h_U^\#$
un morphisme surjectif. Alors $\text{id}_U$ est une section de
$h_U^\#$ sur $U$. Il existe donc un recouvrement
$\{U_a \to U\}_{a \in A}$ et, pour tout $a \in A$,
une section $s_a$ de $\coprod_{i \in I} \mathcal{F}_i$
sur $U_a$ qui s’envoie sur $\text{id}_U$. Par la construction des sommes disjointes
comme faisceaux associés aux sommes disjointes de préfaisceaux
(Lemme \ref{lemma-colimit-sheaves}), il existe, pour tout $a$,
un recouvrement $\{U_{ab} \to U_a\}_{b \in B_a}$ et,
pour tout $b \in B_a$, un indice $\iota(b) \in I$ et une section
$s_{b}$ de $\mathcal{F}_{\iota(b)}$ sur $U_{ab}$
qui s’envoie sur $\text{id}_U|_{U_{ab}}$. Ainsi, après avoir remplacé
le recouvrement $\{U_a \to U\}_{a \in A}$ par
$\{U_{ab} \to U\}_{a \in A, b \in B_a}$,
nous pouvons supposer que nous disposons d’une application $\iota : A \to I$
et, pour tout $a \in A$, d’une section $s_a$ de $\mathcal{F}_{\iota(a)}$
sur $U_a$ qui s’envoie sur $\text{id}_U$.
Puisque $U$ est quasi-compact, il existe un recouvrement
$\{V_c \to U\}_{c \in C}$, une application $\alpha : C \to A$
d’image finie et des morphismes $V_c \to U_{\alpha(c)}$ au-dessus de $U$.
Alors $J = \Im(\iota \circ \alpha) \subset I$ convient,
car $\coprod_{c \in C} h_{V_c}^\# \to h_U^\#$ est surjectif
(Lemme \ref{lemma-covering-surjective-after-sheafification})
et se factorise par $\coprod_{i \in J} \mathcal{F}_i \to h_U^\#$.
(Nous utilisons ici le fait que le composé
$h_{V_c}^\# \to h_{U_{\alpha(c)}}
\xrightarrow{s_{\alpha(c)}} \mathcal{F}_{\iota(\alpha(c))} \to h_U^\#$
est le morphisme $h_{V_c}^\# \to h_U^\#$ provenant du morphisme
$V_c \to U$, puisque $s_{\alpha(c)}$ s’envoie sur
$\text{id}_U|_{U_{\alpha(c)}}$.)

\medskip\noindent
Supposons (2). Soit $\{U_i \to U\}_{i \in I}$ un recouvrement.
D’après le Lemme \ref{lemma-covering-surjective-after-sheafification},
$\coprod_{i \in I} h_{U_i}^\# \to h_U^\#$ est surjectif.
Il existe donc une partie finie $J \subset I$ telle que
$\coprod_{j \in J} h_{U_j}^\# \to h_U^\#$ soit surjectif.
En raisonnant comme ci-dessus, nous trouvons un recouvrement
$\{V_c \to U\}_{c \in C}$ de $U$ dans $\mathcal{C}$
et une application $\iota : C \to J$
tels que $\text{id}_U$ se relève en une section
$s_c$ de $h_{U_{\iota(c)}}^\#$ sur $V_c$.
En raffinant encore le recouvrement, nous pouvons supposer que
$s_c \in h_{U_{\iota(c)}}(V_c)$ s’envoie sur $\text{id}_U$.
Alors $s_c : V_c \to U_{\iota(c)}$ est un morphisme au-dessus de $U$,
ce qui permet de conclure.
\end{proof}

\noindent
Le lemme précédent motive la définition suivante.

\begin{definition}
\label{definition-quasi-compact-topos}
Un objet $\mathcal{F}$ d’un topos $\Sh(\mathcal{C})$ est {\it quasi-compact}
si, pour tout morphisme surjectif $\coprod_{i \in I} \mathcal{F}_i \to \mathcal{F}$
de $\Sh(\mathcal{C})$, il existe une partie finie $J \subset I$ telle
que $\coprod_{i \in J} \mathcal{F}_i \to \mathcal{F}$ soit surjectif.
Un topos $\Sh(\mathcal{C})$ est dit {\it quasi-compact}
si son objet final $*$ est un objet quasi-compact.
\end{definition}

\noindent
D’après le Lemme \ref{lemma-quasi-compact},
si le site $\mathcal{C}$ possède un objet final $X$, alors
$\Sh(\mathcal{C})$ est quasi-compact si et seulement si $X$ est quasi-compact.

\begin{lemma}
\label{lemma-image-of-quasi-compact}
\begin{slogan}
Les morphismes surjectifs de faisceaux transmettent la quasi-compacité.
\end{slogan}
Soit $\mathcal{C}$ un site.
\begin{enumerate}
\item Si $U \to V$ est un morphisme de $\mathcal{C}$ tel que
$h_U^\# \to h_V^\#$ soit surjectif et si $U$ est quasi-compact, alors
$V$ est quasi-compact.
\item Si $\mathcal{F} \to \mathcal{G}$ est un morphisme surjectif de faisceaux
d’ensembles et si $\mathcal{F}$ est quasi-compact, alors $\mathcal{G}$
est quasi-compact.
\end{enumerate}
\end{lemma}

\begin{proof}
Omis.
\end{proof}

\begin{lemma}
\label{lemma-coproduct-quasi-compact}
\begin{slogan}
Les sommes disjointes finies de faisceaux préservent la quasi-compacité.
\end{slogan}
Soit $\mathcal{C}$ un site. Si $n \geq 1$ et si
$\mathcal{F}_1, \ldots, \mathcal{F}_n$ sont des faisceaux quasi-compacts
sur $\mathcal{C}$, alors $\coprod_{i = 1, \ldots, n} \mathcal{F}_i$
est quasi-compact.
\end{lemma}

\begin{proof}
Omis.
\end{proof}

\noindent
Les deux lemmes suivants sont les analogues, pour les sites, du
Lemme \ref{sheaves-lemma-directed-colimits-sections} de Faisceaux.

\begin{lemma}
\label{lemma-directed-colimits-sections}
Soit $\mathcal{C}$ un site. Soit
$\mathcal{I} \to \Sh(\mathcal{C})$, $i \mapsto \mathcal{F}_i$,
un diagramme filtrant de faisceaux d’ensembles.
Soit $U \in \Ob(\mathcal{C})$.
Considérons l’application canonique
$$
\Psi :
\colim_i \mathcal{F}_i(U)
\longrightarrow
\left(\colim_i \mathcal{F}_i\right)(U)
$$
Avec la terminologie introduite ci-dessus :
\begin{enumerate}
\item Si tous les morphismes de transition sont injectifs, alors
$\Psi$ est injective pour tout $U$.
\item Si $U$ est quasi-compact, alors $\Psi$ est injective.
\item Si $U$ est quasi-compact et si tous les morphismes de transition sont injectifs,
alors $\Psi$ est un isomorphisme.
\item Si $U$ admet un système cofinal de recouvrements
$\{U_j \to U\}_{j \in J}$ où
$J$ est fini et $U_j \times_U U_{j'}$ est quasi-compact
pour tous $j, j' \in J$, alors $\Psi$ est bijective.
\end{enumerate}
\end{lemma}

\begin{proof}
Supposons tous les morphismes de transition injectifs. Dans ce cas, le préfaisceau
$\mathcal{F}' : V \mapsto \colim_i \mathcal{F}_i(V)$ est
séparé (voir la Définition \ref{definition-separated}).
Par le Lemme \ref{lemma-colimit-sheaves}, nous avons
$(\mathcal{F}')^\# = \colim_i \mathcal{F}_i$.
Le Théorème \ref{theorem-plus}
montre que $\mathcal{F}' \to (\mathcal{F}')^\#$ est injectif.
Cela démontre (1).

\medskip\noindent
Supposons $U$ quasi-compact. Supposons que $s \in \mathcal{F}_i(U)$ et
$s' \in \mathcal{F}_{i'}(U)$ définissent des éléments du
membre de gauche ayant la même image par $\Psi$.
Cela signifie que nous pouvons choisir un recouvrement $\{U_a \to U\}_{a \in A}$
et, pour tout $a \in A$, un indice $i_a \in I$, $i_a \geq i$, $i_a \geq i'$,
tel que $\varphi_{ii_a}(s) = \varphi_{i'i_a}(s')$.
Comme $U$ est quasi-compact, nous pouvons choisir un recouvrement
$\{V_b \to U\}_{b \in B}$, une application $\alpha : B \to A$ d’image finie,
et des morphismes $V_b \to U_{\alpha(b)}$ au-dessus de $U$.
Choisissons $i''\in I$, supérieur ou égal ($\geq$) à tous les $i_{\alpha(b)}$,
ce qui est possible puisque l’image de $\alpha$ est finie.
Nous en déduisons que $\varphi_{ii''}(s)$ et $\varphi_{i'i''}(s)$
coïncident sur $V_b$ pour tout $b \in B$, et donc que
$\varphi_{ii''}(s) = \varphi_{i'i''}(s)$. Cela démontre (2).

\medskip\noindent
Supposons $U$ quasi-compact et tous les morphismes de transition injectifs.
Soit $s$ un élément du but de $\Psi$. Il existe un recouvrement
$\{U_a \to U\}_{a \in A}$ et, pour tout $a \in A$, un indice $i_a \in I$
et une section $s_a \in \mathcal{F}_{i_a}(U_a)$
tels que $s|_{U_a}$ provienne de $s_a$ pour tout $a \in A$.
Comme $U$ est quasi-compact, nous pouvons choisir un recouvrement
$\{V_b \to U\}_{b \in B}$, une application $\alpha : B \to A$ d’image finie,
et des morphismes $V_b \to U_{\alpha(b)}$ au-dessus de $U$.
Choisissons $i \in I$, supérieur ou égal ($\geq$) à tous les $i_{\alpha(b)}$,
ce qui est possible puisque l’image de $\alpha$ est finie.
D’après (1), les sections
$s_b = \varphi_{i_{\alpha(b)} i}(s_{\alpha(b)})|_{V_b}$
coïncident sur $V_b \times_U V_{b'}$.
Elles se recollent donc en une section
$s' \in \mathcal{F}_i(U)$ qui s’envoie sur $s$ par $\Psi$.
Cela démontre (3).

\medskip\noindent
Supposons vérifiée l’hypothèse de (4). D’après le Lemme
\ref{lemma-conclude-quasi-compact}, l’objet $U$ est quasi-compact ;
ainsi, $\Psi$ est injective par (2).
Pour démontrer la surjectivité, soit $s$ un élément du but de $\Psi$.
Par hypothèse, il existe un recouvrement fini
$\{U_j \to U\}_{j = 1, \ldots, m}$, tel que $U_j \times_U U_{j'}$
soit quasi-compact pour tous $1 \leq j, j' \leq m$, ainsi que,
pour tout $j$, un indice $i_j \in I$ et une section
$s_j \in \mathcal{F}_{i_j}(U_j)$ tels que $s|_{U_j}$ soit l’image de $s_j$
pour tout $j$.
Puisque $U_j \times_U U_{j'}$ est quasi-compact, nous pouvons appliquer (2) :
il existe $i_{jj'} \in I$, avec $i_{jj'} \geq i_j$ et
$i_{jj'} \geq i_{j'}$, tel que
$\varphi_{i_ji_{jj'}}(s_j)$ et $\varphi_{i_{j'}i_{jj'}}(s_{j'})$
coïncident sur $U_j \times_U U_{j'}$. Choisissons un indice $i \in I$
supérieur ou égal à tous les $i_{jj'}$. Les sections
$\varphi_{i_ji}(s_j)$ de $\mathcal{F}_i$ se recollent alors
en une section de $\mathcal{F}_i$ sur $U$. L’image de cette section
est bien l’élément $s$ voulu.
\end{proof}

\begin{lemma}
\label{lemma-directed-colimits-global-sections}
Soit $\mathcal{C}$ un site. Soit
$\mathcal{I} \to \Sh(\mathcal{C})$, $i \mapsto \mathcal{F}_i$,
un diagramme filtrant de faisceaux d’ensembles.
Considérons l’application canonique
$$
\Psi :
\colim_i \Gamma(\mathcal{C}, \mathcal{F}_i)
\longrightarrow
\Gamma(\mathcal{C}, \colim_i \mathcal{F}_i)
$$
On a les assertions suivantes :
\begin{enumerate}
\item Si tous les morphismes de transition sont injectifs, alors $\Psi$ est injective.
\item Si $\Sh(\mathcal{C})$ est quasi-compact, alors $\Psi$ est injective.
\item Si $\Sh(\mathcal{C})$ est quasi-compact et si tous les morphismes
de transition sont injectifs, alors $\Psi$ est un isomorphisme.
\item Supposons qu’il existe un ensemble $S \subset \Ob(\Sh(\mathcal{C}))$
ayant les propriétés suivantes :
\begin{enumerate}
\item pour toute surjection $\mathcal{F} \to *$, il existe
$\mathcal{K} \in S$ et un morphisme $\mathcal{K} \to \mathcal{F}$
tels que $\mathcal{K} \to *$ soit surjectif ;
\item pour $\mathcal{K} \in S$, le produit
$\mathcal{K} \times \mathcal{K}$ est quasi-compact.
\end{enumerate}
Alors $\Psi$ est bijective.
\end{enumerate}
\end{lemma}

\begin{proof}
Démonstration de (1). Supposons tous les morphismes de transition injectifs.
Dans ce cas, le préfaisceau
$\mathcal{F}' : V \mapsto \colim_i \mathcal{F}_i(V)$ est
séparé (voir la Définition \ref{definition-separated}).
D’après le Lemme \ref{lemma-colimit-sheaves}, on a
$(\mathcal{F}')^\# = \colim_i \mathcal{F}_i$.
Le Théorème \ref{theorem-plus} montre que
$\mathcal{F}' \to (\mathcal{F}')^\#$ est injectif.
Cela démontre (1).

\medskip\noindent
Démonstration de (2). Supposons $\Sh(\mathcal{C})$ quasi-compact.
Rappelons que $\Gamma(\mathcal{C}, \mathcal{F}) = \Mor(*, \mathcal{F})$
pour tout $\mathcal{F}$ dans $\Sh(\mathcal{C})$.
Soient $a_i, b_i : * \to \mathcal{F}_i$ et, pour
$i' \geq i$, notons $a_{i'}, b_{i'} : * \to \mathcal{F}_{i'}$
leurs composés avec les morphismes de transition du système.
Posons $a = \colim_{i' \geq i} a_{i'}$, et de même pour $b$.
Pour $i' \geq i$, notons
$$
E_{i'} = \text{Égalisateur}(a_{i'}, b_{i'}) \subset *
\quad\text{et}\quad
E = \text{Égalisateur}(a, b) \subset *
$$
D’après Catégories, Lemme \ref{categories-lemma-directed-commutes}, on a
$E = \colim_{i' \geq i} E_{i'}$. Il s’ensuit que
$\coprod_{i' \geq i} E_{i'} \to E$
est un morphisme surjectif de faisceaux. Par conséquent, si $E = *$,
c’est-à-dire si $a = b$, alors, puisque $*$ est quasi-compact,
on a $E_{i'} = *$ pour un certain $i' \geq i$, d’où
$a_{i'} = b_{i'}$ pour un certain $i' \geq i$.
Cela démontre (2).

\medskip\noindent
Démonstration de (3). Supposons $\Sh(\mathcal{C})$ quasi-compact et tous
les morphismes de transition injectifs. Soit
$a : * \to \colim \mathcal{F}_i$ un morphisme. Alors
$E_i = a^{-1}(\mathcal{F}_i) \subset *$ est un sous-faisceau et
$\colim E_i = *$ (d’après la référence ci-dessus). Il existe donc
un indice $i$ tel que $E_i = *$, et l’image de $a$ est alors
contenue dans $\mathcal{F}_i$, comme voulu.

\medskip\noindent
Démonstration de (4). Soit $S \subset \Ob(\Sh(\mathcal{C}))$ satisfaisant
(4)(a) et (b). En appliquant (4)(a) à $\text{id} : * \to *$, nous trouvons
$\mathcal{K} \in S$ tel que $\mathcal{K} \to *$ soit surjectif.
Les morphismes $\mathcal{K} \times \mathcal{K} \to \mathcal{K} \to *$
sont surjectifs.
D’après (4)(b) et le Lemme \ref{lemma-image-of-quasi-compact},
$\mathcal{K}$ et $\Sh(\mathcal{C})$ sont quasi-compacts.
Ainsi, $\Psi$ est injective par (2).
Posons $\mathcal{F} = \colim \mathcal{F}_i$.
Soit $s : * \to \mathcal{F}$ une section globale de la limite inductive.
Comme $\coprod \mathcal{F}_i \to \mathcal{F}$ est surjectif, la projection
$$
\coprod\nolimits_{i \in I} * \times_{s, \mathcal{F}} \mathcal{F}_i \to *
$$
est surjective. D’après (4)(a), il existe $\mathcal{K} \in S$ et un morphisme
$\mathcal{K} \to \coprod_{i \in I} * \times_{s, \mathcal{F}} \mathcal{F}_i$
tels que $\mathcal{K} \to *$ soit surjectif.
Puisque $\mathcal{K}$ est quasi-compact, ce morphisme se factorise par
$\mathcal{K} \to \coprod_{i' \in I'} * \times_{s, \mathcal{F}} \mathcal{F}_{i'}$
pour une partie finie $I' \subset I$. Soit $i \in I$ un majorant
de la partie finie $I'$. Les morphismes de transition définissent un morphisme
$\coprod_{i' \in I'} \mathcal{F}_{i'} \to \mathcal{F}_i$.
Celui-ci fournit à son tour un morphisme
$\mathcal{K} \to * \times_{s, \mathcal{F}} \mathcal{F}_i$.
Autrement dit, nous obtenons $\mathcal{K} \in S$, avec
$\mathcal{K} \to *$ surjectif, et un diagramme commutatif
$$
\xymatrix{
\mathcal{K} \times \mathcal{K} \ar@<1ex>[r] \ar@<-1ex>[r] &
\mathcal{K} \ar[d] \ar[r] & {*} \ar[d]^s \\
& \mathcal{F}_i \ar[r] & \mathcal{F} \ar@{=}[r] & \colim \mathcal{F}_i
}
$$
La ligne supérieure de ce diagramme est un coégalisateur.
Il suffit donc de montrer qu’après avoir augmenté $i$, les deux morphismes induits
$a_i, b_i : \mathcal{K} \times \mathcal{K} \to \mathcal{F}_i$
sont égaux. C’est ce que démontre le paragraphe suivant,
par exactement le même argument que dans la démonstration de (2) ;
nous invitons le lecteur à omettre le reste de la démonstration.

\medskip\noindent
Pour $i' \geq i$, notons
$a_{i'}, b_{i'} : \mathcal{K} \times \mathcal{K} \to \mathcal{F}_{i'}$
les composés de $a_i, b_i$ avec les morphismes de transition du système.
Posons $a = \colim_{i' \geq i} a_{i'} :
\mathcal{K} \times \mathcal{K} \to \mathcal{F}$, et de même pour $b$.
La commutativité du diagramme ci-dessus donne $a = b$.
Pour $i' \geq i$, notons
$$
E_{i'} = \text{Égalisateur}(a_{i'}, b_{i'}) \subset
\mathcal{K} \times \mathcal{K}
\quad\text{et}\quad
E = \text{Égalisateur}(a, b) \subset
\mathcal{K} \times \mathcal{K}
$$
D’après Catégories, Lemme \ref{categories-lemma-directed-commutes}, on a
$E = \colim_{i' \geq i} E_{i'}$. Il s’ensuit que
$\coprod_{i' \geq i} E_{i'} \to E$
est un morphisme surjectif de faisceaux. Comme $a = b$, on a
$E = \mathcal{K} \times \mathcal{K}$.
Puisque $\mathcal{K} \times \mathcal{K}$ est quasi-compact par (4)(b),
on a $E_{i'} = \mathcal{K} \times \mathcal{K}$ pour un certain
$i' \geq i$, et donc $a_{i'} = b_{i'}$ pour un certain $i' \geq i$.
\end{proof}

\begin{remark}
\label{remark-stronger-conditions}
Soit $\mathcal{C}$ un site. Plusieurs conditions permettent d’assurer
les hypothèses de la partie (4) du
Lemme \ref{lemma-directed-colimits-global-sections}.
En voici quelques-unes.
\begin{enumerate}
\item Supposons qu’il existe un ensemble $\mathcal{B} \subset \Ob(\mathcal{C})$
ayant les propriétés suivantes :
\begin{enumerate}
\item pour toute surjection $\mathcal{F} \to *$, il existe
$m \geq 0$ et $U_1, \ldots, U_m \in \mathcal{B}$ tels que
$\mathcal{F}(U_j)$ soit non vide et $\coprod h_{U_j}^\# \to *$ surjectif ;
\item pour $U, U' \in \mathcal{B}$, le faisceau
$h_U^\# \times h_{U'}^\#$ est quasi-compact.
\end{enumerate}
\item Supposons qu’il existe un ensemble $\mathcal{B} \subset \Ob(\mathcal{C})$
ayant les propriétés suivantes :
\begin{enumerate}
\item il existe $m \geq 0$ et $U_1, \ldots, U_m \in \mathcal{B}$ tels que
$\coprod h_{U_j}^\# \to *$ soit surjectif ;
\item pour $U \in \mathcal{B}$, tout recouvrement de $U$ peut être
raffiné par un recouvrement fini $\{U_j \to U\}_{j = 1, \ldots, m}$
avec $U_j \in \mathcal{B}$ ;
\item pour $U, U' \in \mathcal{B}$, il existe
$m \geq 0$, des objets $U_1, \ldots, U_m \in \mathcal{B}$ et
des morphismes $U_j \to U$ et $U_j \to U'$ tels que
$\coprod h_{U_j}^\# \to h_U^\# \times h_{U'}^\#$ soit surjectif.
\end{enumerate}
\item Supposons que :
\begin{enumerate}
\item $\Sh(\mathcal{C})$ soit quasi-compact ;
\item tout objet de $\mathcal{C}$ admette un recouvrement
dont les membres sont des objets quasi-compacts ;
\item si $U$ et $U'$ sont quasi-compacts, alors le faisceau
$h_U^\# \times h_{U'}^\#$ soit quasi-compact.
\end{enumerate}
\end{enumerate}
Dans les cas (1) et (2), nous prenons pour $S \subset \Ob(\Sh(\mathcal{C}))$
l’ensemble des coproduits finis des faisceaux $h_U^\#$, pour
$U \in \mathcal{B}$.
Dans le cas (3), nous prenons pour $S \subset \Ob(\Sh(\mathcal{C}))$
l’ensemble des coproduits finis des faisceaux $h_U^\#$, pour
$U \in \Ob(\mathcal{C})$ quasi-compact.
\end{remark}

\noindent
Nous aurons plus loin besoin d’une borne sur le comportement possible
des limites inductives, sous la forme suivante.

\begin{lemma}
\label{lemma-colimit-over-ordinal-sections}
Soit $\mathcal{C}$ un site. Soit $\beta$ un ordinal.
Soit $\beta \to \Sh(\mathcal{C})$, $\alpha \mapsto \mathcal{F}_\alpha$,
un système de faisceaux indexé par $\beta$. Pour $U \in \Ob(\mathcal{C})$,
considérons l’application canonique
$$
\colim_{\alpha < \beta} \mathcal{F}_\alpha(U)
\longrightarrow
\left(\colim_{\alpha < \beta} \mathcal{F}_\alpha \right)(U)
$$
Si la cofinalité de $\beta$ est assez grande, cette application
est bijective pour tout $U$.
\end{lemma}

\begin{proof}
Le membre de gauche est la valeur en $U$ de la limite inductive
$\mathcal{F}_{\colim}$ prise dans la catégorie des préfaisceaux ; voir
la section \ref{section-limits-colimits-PSh}.
Rappelons que $\colim_{\alpha < \beta} \mathcal{F}_\alpha$ est
le faisceau associé $\mathcal{F}_{\colim}^\#$ à $\mathcal{F}_{\colim}$ ;
voir le Lemme \ref{lemma-colimit-sheaves}.
Soit $\mathcal{U} = \{U_i \to U\}_{i \in I}$ un élément de
l’ensemble $\text{Cov}(\mathcal{C})$ des recouvrements de $\mathcal{C}$.
Si la cofinalité de $\beta$ est strictement supérieure au cardinal de $I$,
nous affirmons que
$$
H^0(\mathcal{U}, \mathcal{F}_{\colim}) =
\colim H^0(\mathcal{U}, \mathcal{F}_\alpha) =
\colim \mathcal{F}_\alpha(U) = \mathcal{F}_{\colim}(U)
$$
Les deuxième et troisième égalités sont claires. Pour la première, écrivons
$s = (s_i) \in H^0(\mathcal{U}, \mathcal{F}_{\colim})$.
Pour tout $i$, l’élément $s_i$ provient alors d’un élément
$s_{i, \alpha_i} \in \mathcal{F}_{\alpha_i}(U_i)$ pour un certain
$\alpha_i < \beta$. Grâce à l’hypothèse sur la cofinalité, nous pouvons
choisir $\alpha_i = \alpha$ indépendamment de $i$. Les éléments $s_i$ et $s_j$
ont alors la même image dans
$\mathcal{F}_{\alpha_{i, j}}(U_i \times_U U_j)$
pour un certain $\alpha_{i, j} < \beta$. Comme le cardinal de
$I \times I$ est lui aussi strictement inférieur à la cofinalité de $\beta$,
nous pouvons, quitte à augmenter $\alpha$, supposer
$\alpha_{i, j} = \alpha$ pour tous $i, j$. Cela démontre la surjectivité
de l’application naturelle
$\colim H^0(\mathcal{U}, \mathcal{F}_\alpha)
\to H^0(\mathcal{U}, \mathcal{F}_{\colim})$.
Un argument tout à fait semblable montre qu’elle est injective.
En particulier, $\mathcal{F}_{\colim}$ satisfait la condition de faisceau
pour $\mathcal{U}$.
Ainsi, si la cofinalité de $\beta$ est strictement supérieure au supremum
des cardinaux des ensembles d’indices $I$ des recouvrements, la conclusion suit.
\end{proof}

\section{Limites inductives de sites}
\label{section-colimit-sites}

\noindent
Il nous faut un analogue du Lemme \ref{lemma-directed-colimits-sections}
lorsque le site est la limite d’un système projectif de sites.
Par souci de simplicité, nous n’expliquons la construction que dans le cas où
les ensembles d’indices des recouvrements sont finis.

\begin{situation}
\label{situation-inverse-limit-sites}
On se donne ici :
\begin{enumerate}
\item une catégorie d’indices cofiltrante $\mathcal{I}$ ;
\item pour $i \in \Ob(\mathcal{I})$, un site $\mathcal{C}_i$ tel que tout
recouvrement de $\mathcal{C}_i$ possède un ensemble d’indices fini ;
\item pour un morphisme $a : i \to j$ de $\mathcal{I}$, un morphisme de sites
$f_a : \mathcal{C}_i \to \mathcal{C}_j$ défini par un foncteur continu
$u_a : \mathcal{C}_j \to \mathcal{C}_i$,
\end{enumerate}
tels que $f_a \circ f_b = f_c$ dès que $c = a \circ b$ dans $\mathcal{I}$.
\end{situation}

\begin{lemma}
\label{lemma-colimit-sites}
Dans la Situation \ref{situation-inverse-limit-sites}, nous pouvons construire
un site $(\mathcal{C}, \text{Cov}(\mathcal{C}))$ comme suit :
\begin{enumerate}
\item comme catégorie, $\mathcal{C} = \colim \mathcal{C}_i$ ;
\item $\text{Cov}(\mathcal{C})$ est la réunion des images
de $\text{Cov}(\mathcal{C}_i)$ par $u_i : \mathcal{C}_i \to \mathcal{C}$.
\end{enumerate}
\end{lemma}

\begin{proof}
Notre définition de la composition des morphismes de sites implique que
$u_b \circ u_a = u_c$ dès que $c = a \circ b$ dans $\mathcal{I}$.
La formule $\mathcal{C} = \colim \mathcal{C}_i$ signifie que
$\Ob(\mathcal{C}) = \colim \Ob(\mathcal{C}_i)$ et
$\text{Flèches}(\mathcal{C}) = \colim \text{Flèches}(\mathcal{C}_i)$.
La source, le but et la composition sont alors hérités de la source,
du but et de la composition sur $\text{Flèches}(\mathcal{C}_i)$.
Nous obtenons ainsi une catégorie.
Notons $u_i : \mathcal{C}_i \to \mathcal{C}$ le foncteur évident.
Remarquons que, pour tout diagramme fini dans $\mathcal{C}$,
il existe un $i$ tel que ce diagramme soit l’image d’un diagramme
dans $\mathcal{C}_i$.

\medskip\noindent
Soit $\{U^t \to U\}$ un recouvrement de $\mathcal{C}$. Montrons d’abord que,
si $V \to U$ est un morphisme de $\mathcal{C}$, alors $U^t \times_U V$ existe.
D’après la remarque précédente et notre définition des recouvrements,
nous pouvons trouver un $i$, un recouvrement $\{U_i^t \to U_i\}$ de
$\mathcal{C}_i$ et un morphisme $V_i \to U_i$ dont l’image par $u_i$
est la donnée considérée. Nous affirmons que $U^t \times_U V$ est l’image de
$U^t_i \times_{U_i} V_i$ par $u_i$.
En effet, pour tout $a : j \to i$ dans $\mathcal{I}$, le foncteur $u_a$
est continu, et donc
$u_a(U^t_i \times_{U_i} V_i) = u_a(U^t_i) \times_{u_a(U_i)} u_a(V_i)$.
En particulier, nous pouvons remplacer $i$ par $j$ si nous le souhaitons.
Ainsi, si $W$ est un autre objet de $\mathcal{C}$, nous pouvons supposer
$W = u_i(W_i)$, et nous obtenons
\begin{align*}
& \Mor_\mathcal{C}(W, u_i(U^t_i \times_{U_i} V_i)) \\
& =
\colim_{a : j \to i}
\Mor_{\mathcal{C}_j}(u_a(W_i), u_a(U^t_i \times_{U_i} V_i)) \\
& =
\colim_{a : j \to i}
\Mor_{\mathcal{C}_j}(u_a(W_i), u_a(U^t_i))
\times_{\Mor_{\mathcal{C}_j}(u_a(W_i), u_a(U_i))}
\Mor_{\mathcal{C}_j}(u_a(W_i), u_a(V_i)) \\
& =
\Mor_\mathcal{C}(W, U^t)
\times_{\Mor_\mathcal{C}(W, U)}
\Mor_\mathcal{C}(W, V)
\end{align*}
puisque les limites inductives filtrantes commutent aux limites finies
(Catégories, Lemme \ref{categories-lemma-directed-commutes}).
Il en résulte aussi que
$\{U^t \times_U V \to V\}$ est un recouvrement dans $\mathcal{C}$.
Nous voyons ainsi que l’axiome (3) de la Définition \ref{definition-site}
est satisfait.

\medskip\noindent
Pour vérifier l’axiome (2) de la Définition \ref{definition-site},
soit $\{U^t \to U\}_{t \in T}$ un recouvrement de $\mathcal{C}$ et,
pour tout $t$, soit $\{U^{ts} \to U^t\}$ un recouvrement de
$\mathcal{C}$. Nous pouvons alors trouver un $i$ et un recouvrement
$\{U^t_i \to U_i\}_{t \in T}$ de $\mathcal{C}_i$ dont l’image par $u_i$ est
$\{U^t \to U\}$. Comme $T$ est {\bf fini}, nous pouvons choisir
$a : j \to i$ dans $\mathcal{I}$ et des recouvrements
$\{U^{ts}_j \to u_a(U^t_i)\}$ de $\mathcal{C}_j$ dont l’image par $u_j$
donne $\{U^{ts} \to U^t\}$.
Nous en déduisons que $\{U^{ts} \to U\}$ est un recouvrement de $\mathcal{C}$
en appliquant l’axiome (2) au site $\mathcal{C}_j$.

\medskip\noindent
Nous omettons la démonstration de l’axiome (1) de la
Définition \ref{definition-site}.
\end{proof}

\begin{lemma}
\label{lemma-compute-pullback-to-limit}
Dans la Situation \ref{situation-inverse-limit-sites}, soit
$u_i : \mathcal{C}_i \to \mathcal{C}$ le foncteur construit au
Lemme \ref{lemma-colimit-sites}. Alors $u_i$ définit un morphisme
de sites $f_i : \mathcal{C} \to \mathcal{C}_i$. Pour
$U_i \in \Ob(\mathcal{C}_i)$ et un faisceau $\mathcal{F}$ sur $\mathcal{C}_i$,
on a
\begin{equation}
\label{equation-compute-pullback-to-limit}
f_i^{-1}\mathcal{F}(u_i(U_i)) =
\colim_{a : j \to i} f_a^{-1}\mathcal{F}(u_a(U_i))
\end{equation}
\end{lemma}

\begin{proof}
Les arguments de la démonstration du Lemme \ref{lemma-colimit-sites}
montrent immédiatement que les foncteurs $u_i$ sont continus.
Pour achever la démonstration, il faut montrer que
$f_i^{-1} := u_{i, s}$ est un foncteur exact
$\Sh(\mathcal{C}_i) \to \Sh(\mathcal{C})$.
Il suffit en fait de montrer que $f_i^{-1}$ est exact à gauche,
car il est exact à droite en tant qu’adjoint à gauche
(Catégories, Lemme \ref{categories-lemma-exact-adjoint}).
Nous démontrons d’abord (\ref{equation-compute-pullback-to-limit}),
puis nous en déduisons l’exactitude.

\medskip\noindent
Pour un objet arbitraire $V$ de $\mathcal{C}$, nous pouvons choisir
un $a : j \to i$ et un objet $V_j \in \Ob(\mathcal{C})$ tels que
$V = u_j(V_j)$. Nous pouvons alors poser
$$
\mathcal{G}(V) = \colim_{b : k \to j} f_{a \circ b}^{-1}\mathcal{F}(u_b(V_j))
$$
La valeur $\mathcal{G}(V)$ de la limite inductive ne dépend ni du choix
de $b : j \to i$, ni de l’objet $V_j$ tel que $u_j(V_j) = V$ ;
nous omettons la vérification. En outre, si
$\alpha : V \to V'$ est un morphisme de $\mathcal{C}$, nous pouvons choisir
$b : j \to i$ et un morphisme $\alpha_j : V_j \to V'_j$ tels que
$u_j(\alpha_j) = \alpha$. Celui-ci induit une application
$\mathcal{G}(V') \to \mathcal{G}(V)$ à l’aide des restrictions le long des
morphismes $u_b(\alpha_j) : u_b(V_j) \to u_b(V'_j)$. Une vérification montre
que $\mathcal{G}$ est un préfaisceau (nous l’omettons).
En fait, $\mathcal{G}$ satisfait la condition de faisceau. En effet,
tout recouvrement $\mathcal{U} = \{U^t \to U\}$ de $\mathcal{C}$
provient d’un niveau fini. Supposons que
$\mathcal{U}_j = \{U^t_j \to U_j\}$ soit envoyé sur $\mathcal{U}$ par $u_j$
pour un certain $a : j \to i$ de $\mathcal{I}$. Alors
$$
H^0(\mathcal{U}, \mathcal{G}) =
\colim_{b : k \to j} H^0(u_b(\mathcal{U}_j), f_{b \circ a}^{-1}\mathcal{F}) =
\colim_{b : k \to j} f_{b \circ a}^{-1}\mathcal{F}(u_b(U_j)) =
\mathcal{G}(U)
$$
comme voulu. La première égalité est vraie parce que les limites inductives
filtrantes commutent aux limites finies
(Catégories, Lemme \ref{categories-lemma-directed-commutes}).
Par construction, $\mathcal{G}(U)$ est donné par le membre de droite de
(\ref{equation-compute-pullback-to-limit}).
Ainsi, (\ref{equation-compute-pullback-to-limit}) est vraie si nous pouvons
montrer que $\mathcal{G}$ est égal à $f_i^{-1}\mathcal{F}$.

\medskip\noindent
Vérifions dans ce paragraphe que $\mathcal{G}$ est canoniquement isomorphe à
$f_i^{-1}\mathcal{F}$. Nous recommandons vivement au lecteur de sauter ce paragraphe.
Pour cette vérification, nous devons établir une bijection
$\Mor_{\Sh(\mathcal{C})}(\mathcal{G}, \mathcal{H}) =
\Mor_{\Sh(\mathcal{C}_i)}(\mathcal{F}, f_{i, *}\mathcal{H})$
fonctorielle en le faisceau $\mathcal{H}$ sur $\mathcal{C}$,
où $f_{i, *} = u_i^p$. Un morphisme
$\mathcal{G} \to \mathcal{H}$ équivaut à un système compatible de morphismes
$$
\varphi_{a, b, V_j} :
f_{a \circ b}^{-1}\mathcal{F}(u_b(V_j))
\longrightarrow
\mathcal{H}(u_j(V_j))
$$
pour tous $a : j \to i$, $b : k \to j$ et
$V_j \in \Ob(\mathcal{C}_j)$.
Les compatibilités imposent que les morphismes $\varphi_{a, b, V_j}$ soient
égaux à $\varphi_{a \circ b, \text{id}, u_b(V_j)}$. Pour $a : j \to i$ fixé,
la famille de morphismes $\varphi_{a, \text{id}, V_j}$ correspond à un morphisme
de faisceaux $\varphi_a : f_a^{-1}\mathcal{F} \to f_{j, *}\mathcal{H}$.
Les compatibilités entre les $\varphi_{a, \text{id}, u_a(V_i)}$ et les
$\varphi_{\text{id}, \text{id}, V_i}$ impliquent que $\varphi_a$
est l’adjoint du morphisme $\varphi_{id}$ par les égalités
$$
\Mor_{\Sh(\mathcal{C}_j)}(f_a^{-1}\mathcal{F}, f_{j, *}\mathcal{H}) =
\Mor_{\Sh(\mathcal{C}_i)}(\mathcal{F}, f_{a, *}f_{j, *}\mathcal{H}) =
\Mor_{\Sh(\mathcal{C}_i)}(\mathcal{F}, f_{i, *}\mathcal{H})
$$
Nous voyons finalement que tout le système de morphismes
$\varphi_{a, b, V_j}$ est déterminé par le morphisme
$\varphi_{id} : \mathcal{F} \to f_{i, *}\mathcal{H}$.
Réciproquement, étant donné un tel morphisme
$\psi : \mathcal{F} \to f_{i, *}\mathcal{H}$, nous pouvons remonter
l’argument précédent pour construire la famille de morphismes
$\varphi_{a, b, V_j}$. Cela achève la démonstration de l’égalité
$\mathcal{G} = f_i^{-1}\mathcal{F}$.

\medskip\noindent
Supposons (\ref{equation-compute-pullback-to-limit}) vérifiée. Alors le foncteur
$\mathcal{F} \mapsto f_i^{-1}\mathcal{F}(U)$ commute aux limites finies,
car les limites finies de faisceaux se calculent dans la catégorie des préfaisceaux
(Lemme \ref{lemma-limit-sheaf}), les foncteurs $f_a^{-1}$ commutent aux limites
finies, et les limites inductives filtrantes commutent aux limites finies.
Pour voir que $\mathcal{F} \mapsto f_i^{-1}\mathcal{F}(V)$ commute aux limites
finies pour un objet général $V$ de $\mathcal{C}$, nous pouvons employer le même
argument avec la formule donnée ci-dessus pour
$f_i^{-1}\mathcal{F}(V) = \mathcal{G}(V)$.
Ainsi, $f_i^{-1}$ est exact à gauche, ce qui achève la démonstration du lemme.
\end{proof}

\begin{lemma}
\label{lemma-colimit}
Dans la Situation \ref{situation-inverse-limit-sites}, supposons donnés :
\begin{enumerate}
\item un faisceau $\mathcal{F}_i$ sur $\mathcal{C}_i$ pour tout
$i \in \Ob(\mathcal{I})$ ;
\item pour $a : j \to i$, un morphisme
$\varphi_a : f_a^{-1}\mathcal{F}_i \to \mathcal{F}_j$
de faisceaux sur $\mathcal{C}_j$.
\end{enumerate}
Supposons $\varphi_c = \varphi_b \circ f_b^{-1}\varphi_a$
dès que $c = a \circ b$. Posons
$\mathcal{F} = \colim f_i^{-1}\mathcal{F}_i$
sur le site $\mathcal{C}$ du Lemme \ref{lemma-colimit-sites}.
Soient $i \in \Ob(\mathcal{I})$ et $X_i \in \text{Ob}(\mathcal{C}_i)$. Alors
$$
\colim_{a : j \to i} \mathcal{F}_j(u_a(X_i)) = \mathcal{F}(u_i(X_i))
$$
\end{lemma}

\begin{proof}
Un argument formel montre que
$$
\colim_{a : j \to i} \mathcal{F}_i(u_a(X_i)) =
\colim_{a : j \to i} \colim_{b : k \to j}
f_b^{-1}\mathcal{F}_j(u_{a \circ b}(X_i))
$$
D’après (\ref{equation-compute-pullback-to-limit}), la limite inductive
intérieure est égale à $f_j^{-1}\mathcal{F}_j(u_i(X_i))$ ;
on conclut donc par le Lemme \ref{lemma-directed-colimits-sections}.
\end{proof}

\begin{lemma}
\label{lemma-colimit-push-pull}
Dans la Situation \ref{situation-inverse-limit-sites}, supposons donné
un faisceau $\mathcal{F}$ sur $\mathcal{C}$. Alors
$$
\mathcal{F} = \colim f_i^{-1}f_{i, *}\mathcal{F}
$$
où les morphismes de transition sont les $f_j^{-1}\varphi_a$
pour $a : j \to i$, avec
$\varphi_a : f_a^{-1}f_{i, *}\mathcal{F} \to f_{j, *}\mathcal{F}$
un morphisme canonique satisfaisant une condition de cocycle comme au
Lemme \ref{lemma-colimit}.
\end{lemma}

\begin{proof}
Comme morphisme
$$
\varphi_a : f_a^{-1}f_{i, *}\mathcal{F} \to f_{j, *}\mathcal{F}
$$
nous choisissons l’adjoint du morphisme identité
$$
f_{i, *}\mathcal{F} \to f_{a, *}f_{j, *}\mathcal{F}
$$
Ainsi, $\varphi_a$ est la co-unité de l’adjonction
$(f_a^{-1}, f_{a, *})$. Nous devons montrer que, pour tous
$a : j \to i$ et $b : k \to i$ de composé $c = a \circ b$, on a
$\varphi_c = \varphi_b \circ f_b^{-1}\varphi_a$.
Cela résulte de Catégories, Lemme \ref{categories-lemma-compose-counits}.
Enfin, il faut montrer que le morphisme fourni par adjonction
$$
\colim_{i \in I} f_i^{-1}f_{i, *}\mathcal{F}
\longrightarrow
\mathcal{F}
$$
est un isomorphisme. Pour un objet $U$ de $\mathcal{C}$,
nous devons montrer que l’application
$$
(\colim_{i \in I} f_i^{-1}\mathcal{F}_i)(U) \to \mathcal{F}(U)
$$
est bijective. Choisissons un $i$ et un objet $U_i$ de $\mathcal{C}_i$
tels que $u_i(U_i) = U$. Le membre de gauche est alors égal à
$$
(\colim_{i \in I} f_i^{-1}\mathcal{F}_i)(U) =
\colim_{a : j \to i} f_{j, *}\mathcal{F}(u_a(U_i))
$$
d’après le Lemme \ref{lemma-colimit}. Comme $u_j(u_a(U_i)) = U$,
on a $f_{j, *}\mathcal{F}(u_a(U_i)) = \mathcal{F}(U)$
pour tout $a : j \to i$, par définition. La valeur de la limite inductive
est donc $\mathcal{F}(U)$, ce qui achève la démonstration.
\end{proof}











\section{Fonctorialité supplémentaire des préfaisceaux}
\label{section-more-functoriality-PSh}

\noindent
Dans cette section, nous revenons sur le contenu de la
section \ref{section-functoriality-PSh}.
Soit $u : \mathcal{C} \to \mathcal{D}$ un foncteur entre catégories.
Rappelons que
$$
u^p :
\textit{PSh}(\mathcal{D})
\longrightarrow
\textit{PSh}(\mathcal{C})
$$
est le foncteur qui associe à $\mathcal{G}$ sur $\mathcal{D}$ le préfaisceau
$u^p\mathcal{G} = \mathcal{G} \circ u$. Ce foncteur possède non seulement
un adjoint à gauche, à savoir $u_p$, mais aussi un adjoint à droite.

\medskip\noindent
Plus précisément, pour tout $V \in \Ob(\mathcal{D})$,
nous définissons une catégorie ${}_V\mathcal{I} = {}_V^u\mathcal{I}$.
Ses objets sont les couples $(U, \psi : u(U) \to V)$.
Remarquons que la flèche est orientée dans le sens opposé à celle utilisée
pour définir la catégorie $\mathcal{I}_V^u$ dans la
section \ref{section-functoriality-PSh}.
Un morphisme $(U, \psi) \to (U', \psi')$ est donné par un morphisme
$\alpha : U \to U'$ tel que $\psi = \psi' \circ u(\alpha)$.
De plus, pour tout préfaisceau d’ensembles $\mathcal{F}$ sur $\mathcal{C}$,
nous introduisons le foncteur
${}_V\mathcal{F} : {}_V\mathcal{I}^{opp} \to \textit{Ensembles}$,
défini par la règle ${}_V\mathcal{F}(U, \psi) = \mathcal{F}(U)$.
Nous définissons
$$
{}_pu(\mathcal{F})(V) := \lim_{{}_V\mathcal{I}^{opp}} {}_V\mathcal{F}
$$
En tant que limite, cet objet possède des morphismes de projection
$c(\psi) : {}_pu(\mathcal{F})(V) \to \mathcal{F}(U)$
pour tout objet $(U, \psi)$ de ${}_V\mathcal{I}$.
En fait,
$$
{}_pu(\mathcal{F})(V)
=
\left\{
\begin{matrix}
\text{familles }
s_{(U, \psi)} \in \mathcal{F}(U) \\
\forall \beta : (U_1, \psi_1) \to (U_2, \psi_2)
\text{ dans }{}_V\mathcal{I} \\
\text{ telles que } \beta^*s_{(U_2, \psi_2)} = s_{(U_1, \psi_1)}
\end{matrix}
\right\}
$$
où la correspondance est donnée par
$s \mapsto s_{(U, \psi)} = c(\psi)(s)$.
Nous laissons au lecteur le soin de définir les applications de restriction
${}_pu(\mathcal{F})(V) \to {}_pu(\mathcal{F})(V')$
associées à tout morphisme $V' \to V$ de $\mathcal{D}$.
Le préfaisceau ainsi obtenu sera noté ${}_pu\mathcal{F}$.

\begin{lemma}
\label{lemma-recover-pu}
Il existe une application canonique
${}_pu\mathcal{F}(u(U)) \to \mathcal{F}(U)$,
compatible aux applications de restriction.
\end{lemma}

\begin{proof}
C’est simplement le morphisme de projection $c(\text{id}_{u(U)})$ ci-dessus.
\end{proof}

\noindent
Remarquons que tout morphisme de préfaisceaux
$\mathcal{F} \to \mathcal{F}'$ donne naissance à des systèmes compatibles
de morphismes entre les foncteurs
${}_V\mathcal{F} \to {}_V\mathcal{F}'$, et donc à un morphisme de préfaisceaux
${}_pu\mathcal{F} \to {}_pu\mathcal{F}'$.
Autrement dit, nous avons défini un foncteur
$$
{}_pu :
\textit{PSh}(\mathcal{C})
\longrightarrow
\textit{PSh}(\mathcal{D})
$$

\begin{lemma}
\label{lemma-adjoints-pu}
Le foncteur ${}_pu$ est adjoint à droite du foncteur $u^p$.
Autrement dit, la formule
$$
\Mor_{\textit{PSh}(\mathcal{C})}(u^p\mathcal{G}, \mathcal{F})
=
\Mor_{\textit{PSh}(\mathcal{D})}(\mathcal{G}, {}_pu\mathcal{F})
$$
est bifonctorielle en $\mathcal{F}$ et $\mathcal{G}$.
\end{lemma}

\begin{proof}
La démonstration est exactement la même que celle du
Lemme \ref{lemma-adjoints-u}. Remarquons que le morphisme
$u^p{}_pu \mathcal{F} \to \mathcal{F}$ du
Lemme \ref{lemma-recover-pu} est celui qui permet de passer
du membre de droite au membre de gauche.

\medskip\noindent
On peut aussi considérer un préfaisceau d’ensembles $\mathcal{F}$ sur
$\mathcal{C}$ comme un préfaisceau $\mathcal{F}'$ sur
$\mathcal{C}^{opp}$ à valeurs dans $\textit{Ensembles}^{opp}$,
et procéder de même sur $\mathcal{D}$.
Vérifions que $({}_pu \mathcal{F})' = u_p(\mathcal{F}')$ et que
$(u^p\mathcal{G})' = u^p(\mathcal{G}')$.
D’après la Remarque \ref{remark-functoriality-presheaves-values},
$u_p$ et $u^p$ sont adjoints pour les préfaisceaux à valeurs dans
$\textit{Ensembles}^{opp}$. Le résultat s’en déduit formellement.
\end{proof}

\noindent
Ainsi, un foncteur $u : \mathcal{C} \to \mathcal{D}$ de catégories
fournit une suite de foncteurs
$$
u_p, u^p, {}_pu
$$
entre les catégories de préfaisceaux, dans laquelle, pour chaque couple
consécutif, le premier foncteur est adjoint à gauche du second.

\begin{lemma}
\label{lemma-adjoint-functors}
Soient $u : \mathcal{C} \to \mathcal{D}$ et
$v : \mathcal{D} \to \mathcal{C}$ des foncteurs de catégories.
Supposons que $v$ soit adjoint à droite de $u$. Alors :
\begin{enumerate}
\item $u^ph_V = h_{v(V)}$ pour tout $V$ dans $\mathcal{D}$ ;
\item la catégorie $\mathcal{I}^v_U$ possède un objet initial ;
\item la catégorie ${}_V^u\mathcal{I}$ possède un objet final ;
\item ${}_pu = v^p$ ;
\item $u^p = v_p$.
\end{enumerate}
\end{lemma}

\begin{proof}
Démonstration de (1). Soit $V$ un objet de $\mathcal{D}$. On a
$u^ph_V = h_{v(V)}$ parce que
$u^ph_V(U) = \Mor_\mathcal{D}(u(U), V) = \Mor_\mathcal{C}(U, v(V))$
par hypothèse.

\medskip\noindent
Démonstration de (2). Soit $U$ un objet de $\mathcal{C}$. Soit
$\eta : U \to v(u(U))$ le morphisme adjoint du morphisme
$\text{id} : u(U) \to u(U)$. Nous affirmons que $(u(U), \eta)$
est un objet initial de $\mathcal{I}_U^v$. En effet, étant donné
un objet $(V, \phi : U \to v(V))$ de $\mathcal{I}_U^v$,
le morphisme $\phi$ est adjoint à un morphisme $\psi : u(U) \to V$,
qui définit alors un morphisme $(u(U), \eta) \to (V, \phi)$.

\medskip\noindent
Démonstration de (3). Soit $V$ un objet de $\mathcal{D}$. Soit
$\xi : u(v(V)) \to V$ le morphisme adjoint du morphisme
$\text{id} : v(V) \to v(V)$. Nous affirmons que $(v(V), \xi)$
est un objet final de ${}_V^u\mathcal{I}$. En effet, étant donné
un objet $(U, \psi : u(U) \to V)$ de ${}_V^u\mathcal{I}$,
le morphisme $\psi$ est adjoint à un morphisme $\phi : U \to v(V)$,
qui définit alors un morphisme $(U, \psi) \to (v(V), \xi)$.

\medskip\noindent
Par conséquent, pour tout préfaisceau $\mathcal{F}$ sur $\mathcal{C}$,
on a
\begin{eqnarray*}
v^p\mathcal{F}(V)
& = &
\mathcal{F}(v(V)) \\
& = &
\Mor_{\textit{PSh}(\mathcal{C})}(h_{v(V)}, \mathcal{F}) \\
& = &
\Mor_{\textit{PSh}(\mathcal{C})}(u^ph_V, \mathcal{F}) \\
& = &
\Mor_{\textit{PSh}(\mathcal{D})}(h_V, {}_pu\mathcal{F}) \\
& = &
{}_pu\mathcal{F}(V)
\end{eqnarray*}
ce qui démontre la partie (4). La partie (5) résulte de l’unicité
des foncteurs adjoints.
\end{proof}





\section{Foncteurs cocontinus}
\label{section-cocontinuous-functors}

\noindent
Il existe une autre manière de construire des morphismes de topos.
Elle fait intervenir des foncteurs cocontinus entre sites, définis comme suit.

\begin{definition}
\label{definition-cocontinuous}
Soient $\mathcal{C}$ et $\mathcal{D}$ des sites.
Soit $u : \mathcal{C} \to \mathcal{D}$ un foncteur.
Le foncteur $u$ est dit {\it cocontinu} si, pour tout
$U \in \Ob(\mathcal{C})$ et tout recouvrement
$\{V_j \to u(U)\}_{j \in J}$ de $\mathcal{D}$, il existe un recouvrement
$\{U_i \to U\}_{i\in I}$ de $\mathcal{C}$ tel que la famille de morphismes
$\{u(U_i) \to u(U)\}_{i \in I}$ raffine le recouvrement
$\{V_j \to u(U)\}_{j \in J}$.
\end{definition}

\noindent
Remarquons que $\{u(U_i) \to u(U)\}_{i \in I}$ n’est en général {\it pas}
un recouvrement du site $\mathcal{D}$.

\begin{lemma}
\label{lemma-pu-sheaf}
Soient $\mathcal{C}$ et $\mathcal{D}$ des sites.
Soit $u : \mathcal{C} \to \mathcal{D}$ cocontinu.
Soit $\mathcal{F}$ un faisceau sur $\mathcal{C}$.
Alors ${}_pu\mathcal{F}$ est un faisceau sur $\mathcal{D}$,
que nous noterons ${}_su\mathcal{F}$.
\end{lemma}

\begin{proof}
Soit $\{V_j \to V\}_{j \in J}$ un recouvrement du site $\mathcal{D}$.
Nous devons montrer que
$$
\xymatrix{
&\ \phantom{}_pu\mathcal{F}(V) \ar[r] &
\prod {}_pu\mathcal{F}(V_j) \ar@<1ex>[r] \ar@<-1ex>[r] &
\prod {}_pu\mathcal{F}(V_j \times_V V_{j'}) &
}
$$
est un diagramme égalisateur. Comme ${}_pu$ est adjoint à droite de $u^p$,
nous avons
$$
{}_pu\mathcal{F}(V) =
\Mor_{\textit{PSh}(\mathcal{D})}(h_V, {}_pu\mathcal{F}) =
\Mor_{\textit{PSh}(\mathcal{C})}(u^ph_V, \mathcal{F}) =
\Mor_{\Sh(\mathcal{C})}((u^ph_V)^\#, \mathcal{F})
$$
Il suffit donc de montrer que
\begin{equation}
\label{equation-coequalizer}
\xymatrix{
\coprod u^p h_{V_j \times_V V_{j'}} \ar@<1ex>[r] \ar@<-1ex>[r] &
\coprod u^p h_{V_j} \ar[r] &
u^p h_V
}
\end{equation}
devient un diagramme coégalisateur après passage aux faisceaux associés.
(Rappelons qu’un coproduit dans la catégorie des faisceaux est le faisceau
associé au coproduit dans la catégorie des préfaisceaux ; voir le
Lemme \ref{lemma-colimit-sheaves}.)

\medskip\noindent
Montrons d’abord que la deuxième flèche de
(\ref{equation-coequalizer}) devient surjective après passage aux faisceaux
associés. Nous utilisons pour cela le Lemme \ref{lemma-mono-epi-sheaves}.
Il suffit donc de montrer qu’une section $s$ de $u^ph_V$ sur $U$ se relève
en une section de $\coprod u^p h_{V_j}$ sur les membres d’un recouvrement de $U$.
Remarquons que $s$ est un morphisme $s : u(U) \to V$. Alors
$\{V_j \times_{V, s} u(U) \to u(U)\}$ est un recouvrement de $\mathcal{D}$.
Comme $u$ est cocontinu, il existe donc un recouvrement $\{U_i \to U\}$
tel que $\{u(U_i) \to u(U)\}$ raffine
$\{V_j \times_{V, s} u(U) \to u(U)\}$.
Cela signifie que chaque restriction $s|_{U_i} : u(U_i) \to V$ se factorise
par un morphisme $s_i : u(U_i) \to V_j$ pour un certain $j$ ; autrement dit,
$s|_{U_i}$ appartient à l’image de
$u^ph_{V_j}(U_i) \to u^ph_V(U_i)$, comme voulu.

\medskip\noindent
Soient $s, s' \in (\coprod u^ph_{V_j})^\#(U)$ ayant la même image
dans $(u^ph_V)^\#(U)$. Pour achever la démonstration du lemme, montrons
qu’après avoir remplacé $U$ par les membres d’un recouvrement, $s, s'$ sont
l’image de la même section de $\coprod u^p h_{V_j \times_V V_{j'}}$
par les deux morphismes de (\ref{equation-coequalizer}). Nous pouvons d’abord
remplacer $U$ par les membres d’un recouvrement et supposer
$s \in u^ph_{V_j}(U)$ et $s' \in u^ph_{V_{j'}}(U)$.
Un second remplacement de ce type assure que $s$ et $s'$ ont la même image
dans $u^ph_V(U)$, et non seulement dans le faisceau associé.
Ainsi, $s : u(U) \to V_j$ et $s' : u(U) \to V_{j'}$ sont des morphismes
de $\mathcal{D}$ tels que
$$
\xymatrix{
u(U) \ar[r]_{s'} \ar[d]_s & V_{j'} \ar[d] \\
V_j \ar[r] & V
}
$$
soit commutatif. Nous obtenons donc
$t = (s, s') : u(U) \to V_j \times_V V_{j'}$,
c’est-à-dire une section $t \in u^ph_{V_j \times_V V_{j'}}(U)$
qui s’envoie sur $s, s'$, comme voulu.
\end{proof}

\begin{lemma}
\label{lemma-exact-cocontinuous}
Soient $\mathcal{C}$ et $\mathcal{D}$ des sites.
Soit $u : \mathcal{C} \to \mathcal{D}$ cocontinu.
Le foncteur
$\Sh(\mathcal{D}) \to \Sh(\mathcal{C})$,
$\mathcal{G} \mapsto (u^p\mathcal{G})^\#$
est adjoint à gauche du foncteur ${}_su$ introduit au
Lemme \ref{lemma-pu-sheaf}. De plus, il est exact.
\end{lemma}

\begin{proof}
Démontrons la propriété d’adjonction par les égalités suivantes :
\begin{eqnarray*}
\Mor_{\Sh(\mathcal{C})}
((u^p\mathcal{G})^\#, \mathcal{F})
& = &
\Mor_{\textit{PSh}(\mathcal{C})}
(u^p\mathcal{G}, \mathcal{F}) \\
& = &
\Mor_{\textit{PSh}(\mathcal{D})}
(\mathcal{G}, {}_pu\mathcal{F}) \\
& = &
\Mor_{\Sh(\mathcal{D})}
(\mathcal{G}, {}_su\mathcal{F}).
\end{eqnarray*}
Il est donc adjoint à gauche, et par conséquent exact à droite ; voir
Catégories, Lemme \ref{categories-lemma-exact-adjoint}.
Nous avons vu que le passage au faisceau associé est exact à gauche ; voir le
Lemme \ref{lemma-sheafification-exact}.
En outre, l’inclusion
$i : \Sh(\mathcal{D}) \to \textit{PSh}(\mathcal{D})$
est exacte à gauche d’après le Lemme \ref{lemma-limit-sheaf}.
Enfin, le foncteur $u^p$ est exact à gauche parce qu’il est adjoint à droite,
à savoir de $u_p$. Le foncteur considéré est donc le composé
${}^\# \circ u^p \circ i$ de foncteurs exacts à gauche ;
il est par conséquent exact à gauche.
\end{proof}

\noindent
Terminons cette section par un lemme technique.

\begin{lemma}
\label{lemma-technical-pu}
Dans la situation du Lemme \ref{lemma-exact-cocontinuous},
pour tout préfaisceau $\mathcal{G}$ sur $\mathcal{D}$, on a
$(u^p\mathcal{G})^\# = (u^p(\mathcal{G}^\#))^\#$.
\end{lemma}

\begin{proof}
Pour tout faisceau $\mathcal{F}$ sur $\mathcal{C}$, on a
\begin{eqnarray*}
\Mor_{\Sh(\mathcal{C})}((u^p(\mathcal{G}^\#))^\#, \mathcal{F})
& = &
\Mor_{\Sh(\mathcal{D})}(\mathcal{G}^\#, {}_su\mathcal{F}) \\
& = &
\Mor_{\Sh(\mathcal{D})}(\mathcal{G}^\#, {}_pu\mathcal{F}) \\
& = &
\Mor_{\textit{PSh}(\mathcal{D})}(\mathcal{G}, {}_pu\mathcal{F}) \\
& = &
\Mor_{\textit{PSh}(\mathcal{C})}(u^p\mathcal{G}, \mathcal{F}) \\
& = &
\Mor_{\Sh(\mathcal{C})}((u^p\mathcal{G})^\#, \mathcal{F})
\end{eqnarray*}
et le résultat résulte du lemme de Yoneda.
\end{proof}

\begin{remark}
\label{remark-cartesian-cocontinuous}
Soit $u : \mathcal{C} \to \mathcal{D}$ un foncteur entre catégories.
Étant donnés des morphismes $g : u(U) \to V$ et $f : W \to V$
de $\mathcal{D}$, nous pouvons considérer le foncteur
$$
\mathcal{C}^{opp} \longrightarrow \textit{Ensembles},\quad
T \longmapsto
\Mor_\mathcal{C}(T, U)
\times_{\Mor_\mathcal{D}(u(T), V)}
\Mor_\mathcal{D}(u(T), W)
$$
Si ce foncteur est représentable, notons
$U \times_{g, V, f} W$ l’objet correspondant de $\mathcal{C}$.
Supposons que $\mathcal{C}$ et $\mathcal{D}$ soient des sites.
Considérons la propriété $P$ suivante : pour tout recouvrement
$\{f_j : V_j \to V\}$ de $\mathcal{D}$ et tout morphisme
$g : u(U) \to V$, on a :
\begin{enumerate}
\item $U \times_{g, V, f_i} V_i$ existe pour tout $i$ ;
\item $\{U \times_{g, V, f_i} V_i \to U\}$ est un recouvrement de $\mathcal{C}$.
\end{enumerate}
Remarquons la similitude avec la définition des foncteurs continus.
Si $u$ possède la propriété $P$, alors $u$ est cocontinu (détails omis).
De nombreux foncteurs cocontinus que nous rencontrerons possèdent $P$.
\end{remark}

\section{Foncteurs cocontinus et morphismes de topos}
\label{section-cocontinuous-morphism-topoi}

\noindent
Il ressort clairement de ce qui précède qu’un foncteur cocontinu $u$
définit un morphisme de topos dans le même sens que $u$.
Ce sens est donc opposé à celui du morphisme de topos associé
(sous certaines conditions) à un foncteur continu $u$, comme dans la
Définition \ref{definition-morphism-sites}, la
Proposition \ref{proposition-get-morphism} et le
Lemme \ref{lemma-morphism-sites-topoi}.

\begin{lemma}
\label{lemma-cocontinuous-morphism-topoi}
Soient $\mathcal{C}$ et $\mathcal{D}$ des sites.
Soit $u : \mathcal{C} \to \mathcal{D}$ cocontinu.
Les foncteurs $g_* = {}_su$ et $g^{-1} = (u^p\ )^\#$
définissent un morphisme de topos $g$ de $\Sh(\mathcal{C})$ vers
$\Sh(\mathcal{D})$.
\end{lemma}

\begin{proof}
C’est exactement le contenu du Lemme \ref{lemma-exact-cocontinuous}.
\end{proof}

\begin{lemma}
\label{lemma-composition-cocontinuous}
\begin{slogan}
La composition des foncteurs de sites préserve leur cocontinuité.
\end{slogan}
Soient $u : \mathcal{C} \to \mathcal{D}$ et
$v : \mathcal{D} \to \mathcal{E}$ des foncteurs cocontinus. Alors
$v \circ u$ est cocontinu et l’on a $h = g \circ f$, où
$f : \Sh(\mathcal{C}) \to \Sh(\mathcal{D})$, respectivement
$g : \Sh(\mathcal{D}) \to \Sh(\mathcal{E})$, respectivement
$h : \Sh(\mathcal{C}) \to \Sh(\mathcal{E})$, est le morphisme de topos
associé à $u$, respectivement à $v$, respectivement à $v \circ u$.
\end{lemma}

\begin{proof}
Soit $U \in \Ob(\mathcal{C})$.
Soit $\{E_i \to v(u(U))\}$ un recouvrement de $U$ dans $\mathcal{E}$.
Par hypothèse, il existe un recouvrement $\{D_j \to u(U)\}$ dans
$\mathcal{D}$ tel que $\{v(D_j) \to v(u(U))\}$ raffine
$\{E_i \to v(u(U))\}$. Toujours par hypothèse, il existe un recouvrement
$\{C_l \to U\}$ dans $\mathcal{C}$ tel que
$\{u(C_l) \to u(U)\}$ raffine $\{D_j \to u(U)\}$. Alors
$\{v(u(C_l)) \to v(u(U))\}$ raffine bien le recouvrement
$\{E_i \to v(u(U))\}$. Cela montre que $v \circ u$ est cocontinu.
Pour démontrer la dernière assertion, il suffit de montrer que
${}_sv \circ {}_su = {}_s(v \circ u)$. Il suffit de démontrer que
${}_pv \circ {}_pu = {}_p(v \circ u)$ ; voir le
Lemme \ref{lemma-pu-sheaf}. Comme ${}_pu$, respectivement ${}_pv$,
respectivement ${}_p(v \circ u)$, est adjoint à droite de $u^p$,
respectivement de $v^p$, respectivement de $(v \circ u)^p$, il suffit de
montrer que $u^p \circ v^p = (v \circ u)^p$. Cela résulte directement des
définitions.
\end{proof}

\begin{example}
\label{example-open-immersion-cocontinuous}
Soit $X$ un espace topologique.
Soit $j : U  \to X$ l’inclusion d’un sous-espace ouvert.
Rappelons que nous disposons des sites $X_{Zar}$ et $U_{Zar}$ ; voir
l’Exemple \ref{example-site-topological}.
Rappelons également que le foncteur $u : X_{Zar} \to U_{Zar}$ associé à $j$
est continu et donne lieu à un morphisme de sites
$U_{Zar} \to X_{Zar}$ ; voir l’Exemple \ref{example-continuous-map}.
Cela donne aussi un morphisme de topos $(j_*, j^{-1})$.
Considérons ensuite le foncteur
$v : U_{Zar} \to X_{Zar}$, $V \mapsto v(V) = V$
(le même ouvert, mais considéré maintenant comme un objet de $X_{Zar}$).
Ce foncteur est cocontinu. En effet, si
$v(V) = \bigcup_{j \in J} W_j$ est un recouvrement ouvert dans $X$, alors
chaque $W_j$ est nécessairement contenu dans $U$, donc est de la forme
$v(V_j)$, et $V = \bigcup_{j \in J} V_j$ est trivialement un recouvrement
ouvert dans $U$. Le Lemme \ref{lemma-cocontinuous-morphism-topoi} ci-dessus
montre donc qu’il existe un morphisme de topos associé à $v$
$$
\Sh(U) \longrightarrow \Sh(X)
$$
donné par ${}_sv$ et $(v^p\ )^\#$. Nous affirmons qu’en fait
$(v^p\ )^\# = j^{-1}$ et que ${}_sv = j_*$ ; autrement dit, c’est le même
morphisme de topos que celui donné ci-dessus. La manière la plus simple de le
voir est peut-être de remarquer que, pour tout faisceau $\mathcal{G}$ sur $X$,
on a $v^p\mathcal{G}(V) = \mathcal{G}(V)$, ce qui, d’après Faisceaux,
Lemme \ref{sheaves-lemma-j-pullback}, décrit $j^{-1}\mathcal{G}$ (le passage au
faisceau associé est donc superflu dans ce cas). L’égalité de ${}_sv$ et de
$j_*$ résulte de l’unicité des foncteurs adjoints (mais on peut également la
calculer directement).
\end{example}

\begin{example}
\label{example-open-map-cocontinuous}
Cet exemple généralise légèrement
l’Exemple \ref{example-open-immersion-cocontinuous}.
Soit $f : X \to Y$ une application continue d’espaces topologiques.
Supposons que $f$ soit ouverte.
Rappelons que nous disposons des sites $X_{Zar}$ et $Y_{Zar}$ ; voir
l’Exemple \ref{example-site-topological}.
Rappelons également que le foncteur $u : Y_{Zar} \to X_{Zar}$ associé à $f$
est continu et donne lieu à un morphisme de sites
$X_{Zar} \to Y_{Zar}$ ; voir l’Exemple \ref{example-continuous-map}.
Cela donne aussi un morphisme de topos $(f_*, f^{-1})$.
Considérons ensuite le foncteur
$v : X_{Zar} \to Y_{Zar}$, $U \mapsto v(U) = f(U)$.
Ce foncteur est cocontinu. En effet, si
$f(U) = \bigcup_{j \in J} V_j$ est un recouvrement ouvert dans $Y$, alors, en
posant $U_j = f^{-1}(V_j) \cap U$, on obtient un recouvrement ouvert
$U = \bigcup U_j$ tel que $f(U) = \bigcup f(U_j)$ raffine
$f(U) = \bigcup V_j$.
Le Lemme \ref{lemma-cocontinuous-morphism-topoi} ci-dessus montre donc qu’il
existe un morphisme de topos associé à $v$
$$
\Sh(X) \longrightarrow \Sh(Y)
$$
donné par ${}_sv$ et $(v^p\ )^\#$. Nous affirmons qu’en fait
$(v^p\ )^\# = f^{-1}$ et que ${}_sv = f_*$ ; autrement dit, c’est le même
morphisme de topos que celui donné ci-dessus. Pour tout faisceau
$\mathcal{G}$ sur $Y$, on a
$v^p\mathcal{G}(U) = \mathcal{G}(f(U))$.
D’autre part, on peut calculer
$u_p\mathcal{G}(U) = \colim_{f(U) \subset V} \mathcal{G}(V)
= \mathcal{G}(f(U))$, car $(f(U), U \subset f^{-1}(f(U)))$ est manifestement
un objet initial de la catégorie $\mathcal{I}_U^u$ de la
section \ref{section-functoriality-PSh}.
Ainsi $u_p = v^p$ et l’on en conclut que
$f^{-1} = u_s = (v^p\ )^\#$.
L’égalité de ${}_sv$ et de $f_*$ résulte de l’unicité des foncteurs adjoints
(mais on peut également la calculer directement).
\end{example}

\noindent
Dans le premier Exemple \ref{example-open-immersion-cocontinuous}, le
foncteur $v$ est également continu. Mais, dans le second
Exemple \ref{example-open-map-cocontinuous}, il n’est généralement pas
continu, car la condition (2) de la Définition \ref{definition-continuous}
peut ne pas être satisfaite. Le lemme suivant s’applique donc au premier
exemple, mais non au second.

\begin{lemma}
\label{lemma-when-shriek}
\begin{slogan}
Le passage au faisceau associé est superflu dans les morphismes de topos
associés à des foncteurs de sites à la fois continus et cocontinus.
\end{slogan}
Soient $\mathcal{C}$ et $\mathcal{D}$ des sites.
Soit $u : \mathcal{C} \to \mathcal{D}$ un foncteur.
Supposons que
\begin{enumerate}
\item[(a)] $u$ soit cocontinu, et
\item[(b)] $u$ soit continu.
\end{enumerate}
Soit $g : \Sh(\mathcal{C}) \to \Sh(\mathcal{D})$
le morphisme de topos associé. Alors
\begin{enumerate}
\item le passage au faisceau associé dans la formule
$g^{-1} = (u^p\ )^\#$ est inutile ; autrement dit,
$g^{-1}(\mathcal{G})(U) = \mathcal{G}(u(U))$,
\item $g^{-1}$ possède un adjoint à gauche $g_{!} = (u_p\ )^\#$, et
\item $g^{-1}$ commute à toutes les limites projectives et inductives.
\end{enumerate}
\end{lemma}

\begin{proof}
D’après le Lemme \ref{lemma-pushforward-sheaf}, pour tout faisceau
$\mathcal{G}$ sur $\mathcal{D}$, le préfaisceau $u^p\mathcal{G}$ est un
faisceau sur $\mathcal{C}$ (et est noté $u^s\mathcal{G}$). Le foncteur
$\mathcal{F} \mapsto (u_p\mathcal{F})^\#$ est adjoint à gauche de $u^s$
d’après le Lemme \ref{lemma-adjoint-sheaves}. L’assertion relative aux limites
projectives et inductives résulte de la discussion de Catégories,
section \ref{categories-section-adjoint}.
\end{proof}

\noindent
Dans la situation du Lemme \ref{lemma-when-shriek} ci-dessus, on voit que
l’on dispose d’une suite de foncteurs adjoints
$$
g_{!}, \ g^{-1}, \ g_*.
$$
Le foncteur $g_!$ n’est {\it pas} exact en général, car il ne transforme pas,
en général, un objet final de $\Sh(\mathcal{C})$ en un objet final de
$\Sh(\mathcal{D})$.
Voir Faisceaux, Remarque \ref{sheaves-remark-j-shriek-not-exact}.
En revanche, dans le cadre topologique de
l’Exemple \ref{example-open-immersion-cocontinuous}, le foncteur $j_!$ est
exact sur les faisceaux abéliens ; voir Modules,
Lemme \ref{modules-lemma-j-shriek-exact}. Le lemme suivant donne la
généralisation au cas des sites.

\begin{lemma}
\label{lemma-preserve-equalizers}
Soient $\mathcal{C}$ et $\mathcal{D}$ des sites.
Soit $u : \mathcal{C} \to \mathcal{D}$ un foncteur.
Supposons que
\begin{enumerate}
\item[(a)] $u$ soit cocontinu,
\item[(b)] $u$ soit continu, et
\item[(c)] les produits fibrés et les égalisateurs existent dans
$\mathcal{C}$ et que $u$ commute à ceux-ci.
\end{enumerate}
Dans ce cas, le foncteur $g_!$ ci-dessus commute aux produits fibrés et aux
égalisateurs (et, plus généralement, aux limites projectives finies connexes).
\end{lemma}

\begin{proof}
Supposons (a), (b) et (c).
On a $g_! = (u_p\ )^\#$. Rappelons (Lemme \ref{lemma-limit-sheaf}) que les
limites projectives de faisceaux sont égales aux limites projectives
correspondantes calculées comme préfaisceaux. De plus, le passage au faisceau
associé commute aux limites projectives finies
(Lemme \ref{lemma-sheafification-exact}). Il suffit donc de montrer que $u_p$
commute aux produits fibrés et aux égalisateurs. Pour cela, il suffit que les
limites inductives indexées par les catégories
$(\mathcal{I}_V^u)^{opp}$ de la
section \ref{section-functoriality-PSh}
commutent aux produits fibrés et aux égalisateurs. Cela résulte du
Lemme \ref{lemma-almost-directed} et de Catégories,
Lemme \ref{categories-lemma-almost-directed-commutes-equalizers}.
\end{proof}

\noindent
Le lemme suivant traite d’un cas qui ressemble encore davantage au morphisme
associé à une immersion ouverte d’espaces topologiques.

\begin{lemma}
\label{lemma-back-and-forth}
Soient $\mathcal{C}$ et $\mathcal{D}$ des sites.
Soit $u : \mathcal{C} \to \mathcal{D}$ un foncteur.
Supposons que
\begin{enumerate}
\item[(a)] $u$ soit cocontinu,
\item[(b)] $u$ soit continu, et
\item[(c)] $u$ soit pleinement fidèle.
\end{enumerate}
Pour $g_!, g^{-1}, g_*$ comme ci-dessus, les morphismes canoniques
$\mathcal{F} \to g^{-1}g_!\mathcal{F}$ et
$g^{-1}g_*\mathcal{F} \to \mathcal{F}$ sont des isomorphismes pour tout
faisceau $\mathcal{F}$ sur $\mathcal{C}$.
\end{lemma}

\begin{proof}
Soit $X$ un objet de $\mathcal{C}$.
Nous avons vu dans les Lemmes \ref{lemma-pu-sheaf} et
\ref{lemma-when-shriek} que le passage au faisceau associé n’est pas
nécessaire pour les foncteurs
$g^{-1} = (u^p\ )^\#$ et $g_{*} = ({}_pu\ )^\#$.
On peut calculer
$(g^{-1}g_{*}\mathcal{F})(X) = g_{*}\mathcal{F}(u(X))
= \lim \mathcal{F}(Y)$. Ici, la limite est prise sur la catégorie des couples
$(Y, u(Y) \to u(X))$, où l’on n’exige pas que les morphismes
$u(Y) \to u(X)$ soient de la forme $u(\alpha)$, avec $\alpha$ un morphisme de
$\mathcal{C}$. D’après l’hypothèse (c), ils proviennent automatiquement de
morphismes de $\mathcal{C}$, et l’on en déduit que la limite est la valeur en
$(X, u(\text{id}_X))$, c’est-à-dire $\mathcal{F}(X)$.
Cela montre que $g^{-1}g_{*}\mathcal{F} = \mathcal{F}$.

\medskip\noindent
D’autre part, $(g^{-1}g_{!}\mathcal{F})(X) =
g_{!}\mathcal{F}(u(X)) = (u_p\mathcal{F})^\#(u(X))$, et
$u_p\mathcal{F}(u(X)) = \colim \mathcal{F}(Y)$.
Ici, la limite inductive est prise sur la catégorie des couples
$(Y, u(X) \to u(Y))$, où l’on n’exige pas que les morphismes
$u(X) \to u(Y)$ soient de la forme $u(\alpha)$, avec $\alpha$ un morphisme de
$\mathcal{C}$. D’après l’hypothèse (c), ils proviennent automatiquement de
morphismes de $\mathcal{C}$, et l’on en déduit que la limite inductive est la
valeur en $(X, u(\text{id}_X))$, c’est-à-dire $\mathcal{F}(X)$. Ainsi, pour
tout $X \in \Ob(\mathcal{C})$, on a
$u_p\mathcal{F}(u(X)) = \mathcal{F}(X)$.
Comme $u$ est cocontinu et continu, tout recouvrement de $u(X)$ dans
$\mathcal{D}$ peut être raffiné par un recouvrement (!)
$\{u(X_i) \to u(X)\}$ de $\mathcal{D}$, où $\{X_i \to X\}$ est un
recouvrement dans $\mathcal{C}$. Cela implique également
$(u_p\mathcal{F})^+(u(X)) = \mathcal{F}(X)$, car, dans la limite inductive qui
définit la valeur de $(u_p\mathcal{F})^+$ en $u(X)$, on peut se restreindre au
système cofinal des recouvrements $\{u(X_i) \to u(X)\}$ ci-dessus. On voit
donc que $(u_p\mathcal{F})^+(u(X)) = \mathcal{F}(X)$ pour tout objet $X$ de
$\mathcal{C}$. En répétant encore une fois cet argument, on obtient l’égalité
$(u_p\mathcal{F})^\#(u(X)) = \mathcal{F}(X)$ pour tout objet $X$ de
$\mathcal{C}$. Cela donne l’égalité voulue
$g^{-1}g_!\mathcal{F} = \mathcal{F}$.
\end{proof}

\noindent
Enfin, voici un cas qui ne possède aucun exemple topologique correspondant.
Nous utiliserons ce lemme pour voir ce qui se produit lorsque l’on agrandit
un « univers partiel » de schémas tout en conservant la même topologie. Dans
la situation du lemme, le morphisme de topos
$g : \Sh(\mathcal{C}) \to \Sh(\mathcal{D})$ identifie
$\Sh(\mathcal{C})$ à un sous-topos de $\Sh(\mathcal{D})$
(section \ref{section-subtopoi}) et, de plus, le plongement donné admet une
rétraction.

\begin{lemma}
\label{lemma-bigger-site}
Soient $\mathcal{C}$ et $\mathcal{D}$ des sites.
Soit $u : \mathcal{C} \to \mathcal{D}$ un foncteur.
Supposons que
\begin{enumerate}
\item[(a)] $u$ soit cocontinu,
\item[(b)] $u$ soit continu,
\item[(c)] $u$ soit pleinement fidèle,
\item[(d)] les produits fibrés existent dans $\mathcal{C}$ et que $u$ commute
à ceux-ci, et
\item[(e)] il existe des objets finaux
$e_\mathcal{C} \in \Ob(\mathcal{C})$,
$e_\mathcal{D} \in \Ob(\mathcal{D})$ tels que
$u(e_\mathcal{C}) = e_\mathcal{D}$.
\end{enumerate}
Soient $g_!, g^{-1}, g_*$ comme ci-dessus. Alors $u$ définit un morphisme de
sites $f : \mathcal{D} \to \mathcal{C}$ avec $f_* = g^{-1}$ et
$f^{-1} = g_!$. Le composé
$$
\xymatrix{
\Sh(\mathcal{C}) \ar[r]^g &
\Sh(\mathcal{D}) \ar[r]^f &
\Sh(\mathcal{C})
}
$$
est isomorphe au morphisme identité du topos $\Sh(\mathcal{C})$. De plus, le
foncteur $f^{-1}$ est pleinement fidèle.
\end{lemma}

\begin{proof}
Par hypothèse, le foncteur $u$ satisfait aux hypothèses de la
Proposition \ref{proposition-get-morphism}. Ainsi, $u$ définit un morphisme de
sites, donc un morphisme de topos $f$ comme dans le
Lemme \ref{lemma-morphism-sites-topoi}. Les formules $f_* = g^{-1}$ et
$f^{-1} = g_!$ résultent clairement du lemme cité et du
Lemme \ref{lemma-when-shriek}.
On a
$f_* \circ g_* = g^{-1} \circ g_* \cong \text{id}$, et
$g^{-1} \circ f^{-1} = g^{-1} \circ g_! \cong \text{id}$
d’après le Lemme \ref{lemma-back-and-forth}.

\medskip\noindent
Il reste à montrer que $f^{-1}$ est pleinement fidèle.
Soient $\mathcal{F}, \mathcal{G} \in \Ob(\Sh(\mathcal{C}))$.
Nous devons montrer que l’application
$$
\Mor_{\Sh(\mathcal{C})}(\mathcal{F}, \mathcal{G})
\longrightarrow
\Mor_{\Sh(\mathcal{D})}(f^{-1}\mathcal{F}, f^{-1}\mathcal{G})
$$
est bijective. Mais le membre de droite est égal à
\begin{align*}
\Mor_{\Sh(\mathcal{D})}(f^{-1}\mathcal{F}, f^{-1}\mathcal{G})
& =
\Mor_{\Sh(\mathcal{C})}(\mathcal{F}, f_*f^{-1}\mathcal{G}) \\
& =
\Mor_{\Sh(\mathcal{C})}(\mathcal{F}, g^{-1}f^{-1}\mathcal{G}) \\
& =
\Mor_{\Sh(\mathcal{C})}(\mathcal{F}, \mathcal{G})
\end{align*}
(la première égalité résultant de l’adjonction), ce qui démontre l’assertion.
\end{proof}

\begin{example}
\label{example-closed-map-cocontinuous-false}
Soit $X$ un espace topologique.
Soit $i : Z  \to X$ l’inclusion d’une partie (munie de la topologie induite).
Considérons le foncteur $u : X_{Zar} \to Z_{Zar}$,
$U \mapsto u(U) = Z \cap U$.
À première vue, il pourrait sembler que ce foncteur soit également
cocontinu. Après tout, puisque $Z$ est muni de la topologie induite, tout
recouvrement de $U\cap Z$ ne devrait-il pas provenir d’un recouvrement de
$U$ dans $X$ ? Il n’en est rien ! En effet, que se passe-t-il si
$U \cap Z = \emptyset$ ? Dans ce cas, le recouvrement vide est un recouvrement
de $U \cap Z$, et le recouvrement vide ne peut être raffiné que par le
recouvrement vide. On en conclut donc que
$u$ cocontinu $\Rightarrow$ tout ouvert non vide $U$ de $X$ rencontre $Z$.
Mais cela ne suffit pas. Par exemple, si $X = \mathbf{R}$ est la droite réelle
munie de la topologie usuelle et si $Z = \mathbf{R} \setminus \{0\}$, alors
il existe un recouvrement ouvert de $Z$, à savoir
$Z = \{x < 0\} \cup \bigcup_n \{1/n < x\}$, qui ne peut être raffiné par la
restriction d’aucun recouvrement ouvert de $X$.
\end{example}






\section{Foncteurs cocontinus possédant un adjoint à droite}
\label{section-cocontinuous-adjoint}

\noindent
Soient $\mathcal{C}$ et $\mathcal{D}$ des sites. Soient
$u : \mathcal{C} \to \mathcal{D}$ et $v : \mathcal{D} \to \mathcal{C}$ des
foncteurs des catégories sous-jacentes tels que $v$ soit adjoint à droite de
$u$. Dans ce cas, si $v$ est continu, alors $u$ est cocontinu
(Lemme \ref{lemma-continuous-with-continuous-left-adjoint}).
Si $u$ est cocontinu, alors il arrive souvent (mais pas toujours ; voir
l’Exemple \ref{example-not-equivalent}) que $v$ soit continu ; dans ce cas,
$v$ définit un morphisme de sites dont le morphisme de topos associé est le
même que celui défini par $u$.

\begin{lemma}
\label{lemma-have-functor-other-way}
Soient $\mathcal{C}$ et $\mathcal{D}$ des sites. Soient
$u : \mathcal{C} \to \mathcal{D}$ et $v : \mathcal{D} \to \mathcal{C}$ des
foncteurs. Supposons que $u$ soit cocontinu et que $v$ soit adjoint à droite de
$u$. Soit $g : \Sh(\mathcal{C}) \to \Sh(\mathcal{D})$ le morphisme de topos
associé à $u$ ; voir le Lemme \ref{lemma-cocontinuous-morphism-topoi}.
Alors
\begin{enumerate}
\item pour un faisceau $\mathcal{F}$ sur $\mathcal{C}$, le faisceau
$g_*\mathcal{F}$ est égal au préfaisceau $v^p\mathcal{F}$ ; autrement dit,
$(g_*\mathcal{F})(V) = \mathcal{F}(v(V))$, et
\item pour un faisceau $\mathcal{G}$ sur $\mathcal{D}$, on a
$g^{-1}\mathcal{G} = (v_p\mathcal{G})^\#$.
\end{enumerate}
\end{lemma}

\begin{proof}
Pour $\mathcal{F}$ comme en (1), on a
$$
g_*\mathcal{F} =
{}_su\mathcal{F} =
{}_pu\mathcal{F} =
v^p\mathcal{F} =
\mathcal{F} \circ v
$$
La première égalité résulte du Lemme \ref{lemma-cocontinuous-morphism-topoi}.
La deuxième égalité résulte du Lemme \ref{lemma-pu-sheaf}.
La troisième égalité résulte du Lemme \ref{lemma-adjoint-functors}.
La dernière égalité est la définition de $v^p$ dans la
section \ref{section-functoriality-PSh}. Cela démontre (1).
Pour $\mathcal{G}$ comme en (2), on a
$$
g^{-1}\mathcal{G} = (u^p\mathcal{G})^\# = (v_p\mathcal{G})^\#
$$
La première égalité résulte du Lemme \ref{lemma-cocontinuous-morphism-topoi}.
La deuxième égalité résulte du Lemme \ref{lemma-adjoint-functors}.
\end{proof}

\begin{lemma}
\label{lemma-have-functor-other-way-morphism}
Les notations et les hypothèses sont celles du
Lemme \ref{lemma-have-functor-other-way}.
Si, de plus, $v$ est continu, alors $v$ définit un morphisme de sites
$f : \mathcal{C} \to \mathcal{D}$ dont le morphisme de topos associé est égal
à $g$.
\end{lemma}

\begin{proof}
Nous utiliserons sans plus les mentionner les résultats du
Lemme \ref{lemma-have-functor-other-way}. Pour montrer que $v$ définit un
morphisme de sites $f$ comme dans l’énoncé du lemme, nous devons montrer que
$v_s$ est un foncteur exact (voir la
Définition \ref{definition-morphism-sites}). Comme
$v_s\mathcal{G} = (v_p\mathcal{G})^\# = g^{-1}\mathcal{G}$, cela résulte du
fait que $g$ est un morphisme de topos. On voit alors que
$f^{-1} = v_s = g^{-1}$ et l’on obtient $f = g$ comme morphismes de topos.
\end{proof}

\begin{example}
\label{example-finer-topology-bis}
Cet exemple poursuit la discussion de l’Exemple \ref{example-finer-topology},
auquel nous empruntons les notations $\mathcal{C}, \tau, \tau', \epsilon$.
Observons que le foncteur identité
$v : \mathcal{C}_{\tau'} \to \mathcal{C}_\tau$ est continu et que le foncteur
identité $u : \mathcal{C}_\tau \to \mathcal{C}_{\tau'}$ est cocontinu. De
plus, $u$ est adjoint à gauche de $v$. Les résultats des
Lemmes \ref{lemma-have-functor-other-way} et
\ref{lemma-have-functor-other-way-morphism} s’appliquent donc, et l’on conclut
que $v$ définit un morphisme de sites, à savoir
$$
\epsilon : \mathcal{C}_\tau \longrightarrow \mathcal{C}_{\tau'}
$$
dont le morphisme de topos correspondant est le même que le morphisme de topos
associé au foncteur cocontinu $u$.
\end{example}

\begin{lemma}
\label{lemma-continuous-with-continuous-left-adjoint}
Soient $\mathcal{C}$ et $\mathcal{D}$ des sites et soit
$v : \mathcal{D} \to \mathcal{C}$ un foncteur continu.
Supposons que $v$ possède un adjoint à gauche
$u : \mathcal{C} \to \mathcal{D}$. Alors
\begin{enumerate}
\item $u$ est cocontinu,
\item les conclusions des Lemmes \ref{lemma-have-functor-other-way} et
\ref{lemma-have-functor-other-way-morphism} sont valables.
\end{enumerate}
En particulier, $v$ définit un morphisme de sites
$f : \mathcal{C} \to \mathcal{D}$.
\end{lemma}

\begin{proof}
Soit $U$ un objet de $\mathcal{C}$ et soit
$\{V_j \to u(U)\}_{j \in J}$ un recouvrement dans $\mathcal{D}$. Alors
$\{v(V_j) \to v(u(U))\}_{i \in I}$ est un recouvrement dans $\mathcal{C}$.
Par changement de base au moyen du morphisme d’adjonction
$U \to v(u(U))$, on obtient un recouvrement
$\{W_j \to U\}_{j \in J}$, où
$W_j = v(V_j) \times_{v(u(U))} U$. En notant
$p_j : W_j \to v(V_j)$ la première projection, on obtient des morphismes
$$
u(W_j) \xrightarrow{u(p_j)} u(v(V_j)) \longrightarrow V_j
$$
où la deuxième flèche est le morphisme d’adjonction. Cela détermine un
morphisme $\{u(W_j) \to u(U)\} \to \{V_j \to u(U)\}$ de familles de
morphismes à cible fixée, ce qui montre que $u$ est bien cocontinu. Les autres
assertions résultent immédiatement des
Lemmes \ref{lemma-have-functor-other-way} et
\ref{lemma-have-functor-other-way-morphism}.
\end{proof}

\begin{example}
\label{example-not-equivalent}
Soient $\mathcal{C}$ et $\mathcal{D}$ des sites. Soient
$u : \mathcal{C} \to \mathcal{D}$ et $v : \mathcal{D} \to \mathcal{C}$ des
foncteurs des catégories sous-jacentes tels que $v$ soit adjoint à droite de
$u$. Le Lemme \ref{lemma-continuous-with-continuous-left-adjoint} montre que,
si $v$ est continu, alors $u$ est cocontinu. Réciproquement, si $u$ est
cocontinu, on ne peut pas conclure que $v$ est continu. Nous donnerons un
exemple de ce phénomène à l’aide des grands sites \'etale et lisse d’un
schéma, mais il existe vraisemblablement aussi un exemple élémentaire.
Considérons en effet un schéma $S$ et les sites $(\Sch/S)_\etale$ et
$(\Sch/S)_{smooth}$. On peut supposer que ces sites ont la même catégorie
sous-jacente ; voir Topologies,
Remarque \ref{topologies-remark-choice-sites}.
Soit $u = v = \text{id}$. Alors $u$, considéré comme foncteur de
$(\Sch/S)_\etale$ vers $(\Sch/S)_{smooth}$, est cocontinu, car tout
recouvrement lisse d’un schéma peut être raffiné par un recouvrement \'etale ;
voir Compléments sur les morphismes,
Lemme \ref{more-morphisms-lemma-etale-dominates-smooth}.
Réciproquement, le foncteur $v$ de $(\Sch/S)_{smooth}$ vers
$(\Sch/S)_\etale$ n’est pas continu, car un recouvrement lisse n’est pas, en
général, un recouvrement \'etale.
\end{example}




\section{Foncteurs cocontinus possédant un adjoint à gauche}
\label{section-cocontinuous-left-adjoint}

\noindent
Il peut arriver qu’un foncteur cocontinu $u$ possède un adjoint à gauche $w$.

\begin{lemma}
\label{lemma-have-left-adjoint}
Soient $\mathcal{C}$ et $\mathcal{D}$ des sites. Soit
$g : \Sh(\mathcal{C}) \to \Sh(\mathcal{D})$ le morphisme de topos
associé à un foncteur continu et cocontinu
$u : \mathcal{C} \to \mathcal{D}$ ; voir les
Lemmes \ref{lemma-cocontinuous-morphism-topoi} et
\ref{lemma-when-shriek}.
\begin{enumerate}
\item Si $w : \mathcal{D} \to \mathcal{C}$ est adjoint à gauche de $u$, alors :
\begin{enumerate}
\item $g_!\mathcal{F}$ est le faisceau associé au préfaisceau
$w^p\mathcal{F}$ ;
\item $g_!$ est exact.
\end{enumerate}
\item Si $w$ est un adjoint à gauche continu, alors $g_!$
possède un adjoint à gauche.
\item Si $w$ est un adjoint à gauche cocontinu, alors $g_! = h^{-1}$ et
$g^{-1} = h_*$, où $h : \Sh(\mathcal{D}) \to \Sh(\mathcal{C})$ est
le morphisme de topos associé à $w$.
\end{enumerate}
\end{lemma}

\begin{proof}
Rappelons que $g_!\mathcal{F}$ est le faisceau associé à $u_p\mathcal{F}$.
Ainsi, (1)(a) résulte de l’égalité $u_p = w^p$ donnée par le
Lemme \ref{lemma-adjoint-functors}.

\medskip\noindent
Pour établir (1)(b), remarquons que $g_!$ commute à toutes les limites
inductives, puisque $g_!$ est adjoint à gauche (Catégories, Lemme
\ref{categories-lemma-adjoint-exact}).
Soit $i \mapsto \mathcal{F}_i$ un diagramme fini dans $\Sh(\mathcal{C})$.
Alors $\lim \mathcal{F}_i$ se calcule dans la catégorie des préfaisceaux
(Lemme \ref{lemma-limit-sheaf}). Comme $w^p$ est adjoint à droite
(Lemme \ref{lemma-adjoints-u}), on a
$w^p \lim \mathcal{F}_i = \lim w^p\mathcal{F}_i$.
Puisque le passage au faisceau associé est exact
(Lemme \ref{lemma-sheafification-exact}), on conclut par (1)(a).

\medskip\noindent
Supposons $w$ continu. Alors $g_! = (w^p\ )^\# = w^s$, mais le passage
au faisceau associé n’est pas nécessaire, et $w_s$ est un adjoint à gauche ;
voir les Lemmes \ref{lemma-pushforward-sheaf} et
\ref{lemma-adjoint-sheaves}.

\medskip\noindent
Supposons $w$ cocontinu. L’égalité $g_! = h^{-1}$ résulte de (1)(a)
et des définitions. L’égalité $g^{-1} = h_*$ résulte de
$g_! = h^{-1}$ et de l’unicité du foncteur adjoint.
On peut aussi la déduire du Lemme \ref{lemma-have-functor-other-way}.
\end{proof}









\section{Existence de l’image directe extraordinaire}
\label{section-lower-shriek}

\noindent
Dans cette section, nous étudions certains cas de morphismes de topos $f$
pour lesquels $f^{-1}$ possède un adjoint à gauche $f_!$.

\begin{lemma}
\label{lemma-existence-lower-shriek}
Soient $\mathcal{C}$ et $\mathcal{D}$ deux sites.
Soit $f : \Sh(\mathcal{D}) \to \Sh(\mathcal{C})$ un morphisme de topos.
Soit $E \subset \Ob(\mathcal{D})$ une partie telle que :
\begin{enumerate}
\item pour $V \in E$, il existe un faisceau $\mathcal{G}$ sur
$\mathcal{C}$ tel que $f^{-1}\mathcal{F}(V) =
\Mor_{\Sh(\mathcal{C})}(\mathcal{G}, \mathcal{F})$ fonctoriellement
en $\mathcal{F}$ dans $\Sh(\mathcal{C})$ ;
\item tout objet de $\mathcal{D}$ admet un recouvrement par des objets de $E$.
\end{enumerate}
Alors $f^{-1}$ possède un adjoint à gauche $f_!$.
\end{lemma}

\begin{proof}
D’après le lemme de Yoneda (Catégories, Lemme
\ref{categories-lemma-yoneda}), le faisceau $\mathcal{G}_V$
correspondant à $V \in E$ est défini à isomorphisme unique près par la formule
$f^{-1}\mathcal{F}(V) =
\Mor_{\Sh(\mathcal{C})}(\mathcal{G}_V, \mathcal{F})$.
Rappelons que
$f^{-1}\mathcal{F}(V) =
\Mor_{\Sh(\mathcal{D})}(h_V^\#, f^{-1}\mathcal{F})$.
Notons $i_V : h_V^\# \to f^{-1}\mathcal{G}_V$ le morphisme correspondant à
$\text{id}$ dans $\Mor(\mathcal{G}_V, \mathcal{G}_V)$.
La fonctorialité de (1) implique que la bijection est donnée par
$$
\Mor_{\Sh(\mathcal{C})}(\mathcal{G}_V, \mathcal{F}) \to
\Mor_{\Sh(\mathcal{D})}(h_V^\#, f^{-1}\mathcal{F}),\quad
\varphi \mapsto f^{-1}\varphi \circ i_V
$$
Pour tous $V_1, V_2 \in E$, il existe une application canonique
$$
\Mor_{\Sh(\mathcal{D})}(h^\#_{V_2}, h^\#_{V_1})
\to
\Hom_{\Sh(\mathcal{C})}(\mathcal{G}_{V_2}, \mathcal{G}_{V_1}),\quad
\varphi \mapsto f_!(\varphi)
$$
caractérisée par
$f^{-1}(f_!(\varphi)) \circ i_{V_2} = i_{V_1} \circ \varphi$.
Remarquons que $\varphi \mapsto f_!(\varphi)$ est compatible à la composition ;
cela se voit directement à partir de la caractérisation.
Ainsi, $h_V^\# \mapsto \mathcal{G}_V$ et
$\varphi \mapsto f_!\varphi$ définissent un foncteur sur la sous-catégorie
pleine de $\Sh(\mathcal{D})$ dont les objets sont les $h_V^\#$.

\medskip\noindent
Soit $J$ un ensemble et soit $J \to E$, $j \mapsto V_j$, une application.
Nous avons alors une bijection fonctorielle
$$
\Mor_{\Sh(\mathcal{C})}(\coprod \mathcal{G}_{V_j}, \mathcal{F})
\longrightarrow
\Mor_{\Sh(\mathcal{D})}(\coprod h_{V_j}^\#, f^{-1}\mathcal{F})
$$
obtenue comme produit des bijections précédentes. Nous pouvons donc prolonger
le foncteur $f_!$ à la sous-catégorie pleine de $\Sh(\mathcal{D})$
dont les objets sont les coproduits de $h_V^\#$ avec $V \in E$.

\medskip\noindent
Étant donné un faisceau arbitraire $\mathcal{H}$ sur $\mathcal{D}$,
choisissons un diagramme coégalisateur
$$
\xymatrix{
\mathcal{H}_1 \ar@<1ex>[r] \ar@<-1ex>[r] &
\mathcal{H}_0 \ar[r] &
\mathcal{H}
}
$$
où $\mathcal{H}_i = \coprod h_{V_{i, j}}^\#$
est un coproduit avec $V_{i, j} \in E$.
Cela est possible par l’hypothèse (2) ; voir le
Lemme \ref{lemma-sheaf-coequalizer-representable}
(pour qui s’inquiéterait des questions ensemblistes, remarquons que
la construction donnée dans le
Lemme \ref{lemma-sheaf-coequalizer-representable} est canonique).
Définissons $f_!(\mathcal{H})$ comme le faisceau sur $\mathcal{C}$ qui fait de
$$
\xymatrix{
f_!\mathcal{H}_1 \ar@<1ex>[r] \ar@<-1ex>[r] &
f_!\mathcal{H}_0 \ar[r] &
f_!\mathcal{H}
}
$$
un diagramme coégalisateur. Alors
\begin{align*}
\Mor(f_!\mathcal{H}, \mathcal{F})
& =
\text{Égalisateur}(
\xymatrix{
\Mor(f_!\mathcal{H}_0, \mathcal{F}) \ar@<1ex>[r] \ar@<-1ex>[r] &
\Mor(f_!\mathcal{H}_1, \mathcal{F})
}
) \\
& =
\text{Égalisateur}(
\xymatrix{
\Mor(\mathcal{H}_0, f^{-1}\mathcal{F}) \ar@<1ex>[r] \ar@<-1ex>[r] &
\Mor(\mathcal{H}_1, f^{-1}\mathcal{F})
}
) \\
& =
\Hom(\mathcal{H}, f^{-1}\mathcal{F})
\end{align*}
Nous voyons ainsi que $f_!$ se prolonge à toute la catégorie des faisceaux
sur $\mathcal{D}$.
\end{proof}








\section{Localisation}
\label{section-localize}

\noindent
Soit $\mathcal{C}$ un site.
Soit $U \in \Ob(\mathcal{C})$.
Voir Catégories, Exemple \ref{categories-example-category-over-X},
pour la définition de la catégorie $\mathcal{C}/U$ des objets au-dessus de $U$.
Nous munissons $\mathcal{C}/U$ d’une structure de site en déclarant qu’une
famille de morphismes $\{V_j \to V\}$ d’objets au-dessus de $U$ est un
recouvrement de $\mathcal{C}/U$ si et seulement si elle est un recouvrement
dans $\mathcal{C}$. Considérons le foncteur d’oubli
$$
j_U : \mathcal{C}/U \longrightarrow \mathcal{C}.
$$
Il est manifestement cocontinu et continu. Par le
Lemme \ref{lemma-when-shriek}, nous obtenons donc un morphisme de topos
$$
j_U : \Sh(\mathcal{C}/U) \longrightarrow \Sh(\mathcal{C})
$$
donné par $j_U^{-1}$ et $j_{U*}$, ainsi qu’un foncteur $j_{U!}$
adjoint à gauche de $j_U^{-1}$.

\begin{definition}
\label{definition-localize}
Soit $\mathcal{C}$ un site.
Soit $U \in \Ob(\mathcal{C})$.
\begin{enumerate}
\item Le site $\mathcal{C}/U$ est appelé la {\it localisation du site
$\mathcal{C}$ en l’objet $U$}.
\item Le morphisme de topos
$j_U : \Sh(\mathcal{C}/U) \to \Sh(\mathcal{C})$
est appelé le {\it morphisme de localisation}.
\item Le foncteur $j_{U*}$ est appelé le {\it foncteur image directe}.
\item Pour un faisceau $\mathcal{F}$ sur $\mathcal{C}$, le faisceau
$j_U^{-1}\mathcal{F}$ est appelé la {\it restriction de $\mathcal{F}$
à $\mathcal{C}/U$}.
\item Pour un faisceau $\mathcal{G}$ sur $\mathcal{C}/U$, le faisceau
$j_{U!}\mathcal{G}$ est appelé le
{\it prolongement de $\mathcal{G}$ par l’ensemble vide}.
\end{enumerate}
\end{definition}

\noindent
La restriction $j_U^{-1}\mathcal{F}$ est le faisceau défini par la règle
$j_U^{-1}\mathcal{F}(X/U) = \mathcal{F}(X)$, comme prévu.
Le prolongement par l’ensemble vide admet lui aussi une description très simple
dans ce cas ; la voici.

\begin{lemma}
\label{lemma-describe-j-shriek}
Soit $\mathcal{C}$ un site.
Soit $U \in \Ob(\mathcal{C})$.
Soit $\mathcal{G}$ un préfaisceau sur $\mathcal{C}/U$.
Alors $j_{U!}(\mathcal{G}^\#)$ est le faisceau associé au préfaisceau
$$
V
\longmapsto
\coprod\nolimits_{\varphi \in \Mor_\mathcal{C}(V, U)}
\mathcal{G}(V \xrightarrow{\varphi} U)
$$
muni des applications de restriction évidentes.
\end{lemma}

\begin{proof}
D’après le Lemme \ref{lemma-when-shriek}, on a
$j_{U!}(\mathcal{G}^\#) = ((j_U)_p\mathcal{G}^\#)^\#$.
D’après le Lemme \ref{lemma-technical-up}, ceci est égal à
$((j_U)_p\mathcal{G})^\#$.
Il suffit donc de démontrer que $(j_U)_p$ est donné par la formule ci-dessus
pour tout préfaisceau $\mathcal{G}$ sur $\mathcal{C}/U$.
D’après la définition de la
section \ref{section-functoriality-PSh}, on a
$$
(j_U)_p\mathcal{G}(V)
=
\colim_{(W/U, V \to W)} \mathcal{G}(W)
$$
Il est maintenant clair que la catégorie des couples $(W/U, V \to W)$
possède un objet $O_\varphi = (\varphi : V \to U, \text{id} : V \to V)$
pour tout $\varphi : V \to U$ et que, de plus, tout objet reçoit un unique
morphisme depuis l’un des $O_\varphi$. Le résultat s’ensuit.
\end{proof}

\begin{lemma}
\label{lemma-describe-j-shriek-representable}
Soit $\mathcal{C}$ un site.
Soit $U \in \Ob(\mathcal{C})$.
Soit $X/U$ un objet de $\mathcal{C}/U$.
Alors $j_{U!}(h_{X/U}^\#) = h_X^\#$.
\end{lemma}

\begin{proof}
Notons $p : X \to U$ le morphisme structural de $X$.
D’après le Lemme \ref{lemma-describe-j-shriek},
$j_{U!}(h_{X/U}^\#)$ est le faisceau associé au préfaisceau
$$
V
\longmapsto
\coprod\nolimits_{\varphi \in \Mor_\mathcal{C}(V, U)}
\{\psi : V \to X \mid p \circ \psi = \varphi\}
$$
Celui-ci n’est autre que $\Mor_\mathcal{C}(V, X)$.
Le lemme en résulte.
\end{proof}

\noindent
On a $j_{U!}(*) = h_U^\#$ par chacun des deux lemmes précédents.
Ainsi, pour tout faisceau $\mathcal{G}$ sur $\mathcal{C}/U$, il existe
un morphisme canonique de faisceaux $j_{U!}\mathcal{G} \to h_U^\#$.
Cela caractérise les faisceaux appartenant à l’image essentielle de $j_{U!}$.

\begin{lemma}
\label{lemma-essential-image-j-shriek}
Soit $\mathcal{C}$ un site.
Soit $U \in \Ob(\mathcal{C})$.
Le foncteur $j_{U!}$ définit une équivalence de catégories
$$
\Sh(\mathcal{C}/U)
\longrightarrow
\Sh(\mathcal{C})/h_U^\#
$$
\end{lemma}

\begin{proof}
Notons les objets de $\mathcal{C}/U$ sous la forme de couples $(X, a)$,
où $X$ est un objet de $\mathcal{C}$ et $a : X \to U$ un morphisme de
$\mathcal{C}$. De même, les objets de $\Sh(\mathcal{C})/h_U^\#$ sont
des couples $(\mathcal{F}, \varphi)$.
Le foncteur $\Sh(\mathcal{C}/U) \to \Sh(\mathcal{C})/h_U^\#$
envoie $\mathcal{G}$ sur le couple $(j_{U!}\mathcal{G}, \gamma)$,
où $\gamma$ est le composé de $j_{U!}\mathcal{G} \to j_{U!}*$
avec l’identification $j_{U!}* = h_U^\#$.

\medskip\noindent
Construisons un foncteur de $\Sh(\mathcal{C})/h_U^\#$ vers
$\Sh(\mathcal{C}/U)$.
Supposons donné $(\mathcal{F}, \varphi)$.
Pour un objet $(X, a)$ de $\mathcal{C}/U$, considérons l’ensemble
$\mathcal{F}_\varphi(X, a)$ des éléments $s \in \mathcal{F}(X)$ dont l’image
par $\varphi$ est l’image de
$a \in \Mor_\mathcal{C}(X, U) = h_U(X)$ dans $h_U^\#(X)$.
Il est immédiat que $(X, a) \mapsto \mathcal{F}_\varphi(X, a)$ est un faisceau
sur $\mathcal{C}/U$. Il est clair que la règle
$(\mathcal{F}, \varphi) \mapsto \mathcal{F}_\varphi$
définit un foncteur $\Sh(\mathcal{C})/h_U^\# \to \Sh(\mathcal{C}/U)$.

\medskip\noindent
Considérons aussi le foncteur
$\textit{PSh}(\mathcal{C})/h_U \to \textit{PSh}(\mathcal{C}/U)$,
$(\mathcal{F}, \varphi) \mapsto \mathcal{F}_\varphi$,
où $\mathcal{F}_\varphi(X, a)$ est défini comme l’ensemble des éléments
de $\mathcal{F}(X)$ qui s’envoient sur $a \in h_U(X)$.
Nous affirmons que le diagramme
$$
\xymatrix{
\textit{PSh}(\mathcal{C})/h_U \ar[r] \ar[d] &
\textit{PSh}(\mathcal{C}/U) \ar[d] \\
\Sh(\mathcal{C})/h_U^\# \ar[r] &
\Sh(\mathcal{C}/U)
}
$$
est commutatif, où les flèches verticales sont données par le passage aux
faisceaux associés.
Pour le voir\footnote{Une autre méthode consiste à décrire
$\mathcal{F}_\varphi$ par le diagramme cartésien
$$
\vcenter{
\xymatrix{
\mathcal{F}_\varphi \ar[r] \ar[d] & {*} \ar[d] \\
\mathcal{F}|_{\mathcal{C}/U} \ar[r] & h_U|_{\mathcal{C}/U}
}
}
\quad\text{pour les préfaisceaux et}\quad
\vcenter{
\xymatrix{
\mathcal{F}_\varphi \ar[r] \ar[d] & {*} \ar[d] \\
\mathcal{F}|_{\mathcal{C}/U} \ar[r] & h_U^\#|_{\mathcal{C}/U}
}
}
$$
pour les faisceaux, puis à utiliser le fait que la restriction à
$\mathcal{C}/U$ commute au passage au faisceau associé.}, il suffit de
montrer que la construction commute au foncteur
$\mathcal{F} \mapsto \mathcal{F}^+$ des
Lemmes \ref{lemma-plus-presheaf} et \ref{lemma-plus-functorial},
ainsi que du Théorème \ref{theorem-plus}.
La commutation avec $\mathcal{F} \mapsto \mathcal{F}^+$ résulte du fait que,
pour $(X, a)$ donné, la catégorie des recouvrements de $(X, a)$ dans
$\mathcal{C}/U$ s’identifie canoniquement à celle des recouvrements de $X$
dans $\mathcal{C}$.

\medskip\noindent
Ensuite, faisons envoyer par
$\textit{PSh}(\mathcal{C}/U) \to \textit{PSh}(\mathcal{C})/h_U$
le préfaisceau $\mathcal{G}$ sur le couple
$(j_{U!}^{PSh}\mathcal{G}, \gamma)$, où
$j_{U!}^{PSh}\mathcal{G}$ est le préfaisceau défini par la formule du
Lemme \ref{lemma-describe-j-shriek}, et où $\gamma$ est le composé de
$j_{U!}^{PSh}\mathcal{G} \to j_{U!}*$ avec l’identification
$j_{U!}^{PSh}* = h_U$ (évidente d’après la formule).
Il est alors immédiatement clair que le diagramme
$$
\xymatrix{
\textit{PSh}(\mathcal{C}/U) \ar[r] \ar[d] &
\textit{PSh}(\mathcal{C})/h_U \ar[d] \\
\Sh(\mathcal{C}/U) \ar[r] &
\Sh(\mathcal{C})/h_U^\#
}
$$
est commutatif, où les flèches verticales sont les passages aux faisceaux
associés. En réunissant ces résultats, il suffit de montrer qu’il existe
des isomorphismes fonctoriels
$(j_{U!}^{PSh}\mathcal{G})_\gamma = \mathcal{G}$
pour $\mathcal{G}$ dans $\textit{PSh}(\mathcal{C}/U)$, et
$j_{U!}^{PSh}\mathcal{F}_\varphi = \mathcal{F}$
pour $(\mathcal{F}, \varphi)$ dans $\textit{PSh}(\mathcal{C})/h_U$.
La valeur du préfaisceau $(j_{U!}^{PSh}\mathcal{G})_\gamma$
en $(X, a)$ est la fibre de l’application
$$
\coprod\nolimits_{a' : X \to U} \mathcal{G}(X, a') \to \Mor_\mathcal{C}(X, U)
$$
au-dessus de $a$, laquelle est $\mathcal{G}(X, a)$.
Cela démontre la première égalité.
La valeur du préfaisceau $j_{U!}^{PSh}\mathcal{F}_\varphi$ en $X$ est
$$
\coprod\nolimits_{a : X \to U} \mathcal{F}_\varphi(X, a) =
\mathcal{F}(X)
$$
car, pour toute application d’ensembles $S \to S'$, l’ensemble $S$ est
la réunion disjointe de ses fibres.
\end{proof}

\noindent
Le Lemme \ref{lemma-essential-image-j-shriek} affirme que le foncteur
$j_{U!}$ est le composé
$$
\Sh(\mathcal{C}/U) \rightarrow
\Sh(\mathcal{C})/h_U^\# \rightarrow
\Sh(\mathcal{C})
$$
où la première flèche est une équivalence.

\begin{lemma}
\label{lemma-j-shriek-commutes-equalizers-fibre-products}
Soit $\mathcal{C}$ un site. Soit $U \in \Ob(\mathcal{C})$.
Le foncteur $j_{U!}$ commute aux produits fibrés et aux égalisateurs
(et, plus généralement, aux limites finies connexes). En particulier, si
$\mathcal{F} \subset \mathcal{F}'$ dans $\Sh(\mathcal{C}/U)$, alors
$j_{U!}\mathcal{F} \subset j_{U!}\mathcal{F}'$.
\end{lemma}

\begin{proof}
Grâce au Lemme \ref{lemma-essential-image-j-shriek} et au fait qu’une
équivalence de catégories commute à toutes les limites, cela se ramène
au fait que le foncteur
$\Sh(\mathcal{C})/h_U^\# \rightarrow \Sh(\mathcal{C})$
commute aux produits fibrés et aux égalisateurs.
On peut également le démontrer directement à partir de la description de
$j_{U!}$ donnée au Lemme \ref{lemma-describe-j-shriek}, en utilisant
l’exactitude du passage au faisceau associé. (De même, lorsque $\mathcal{C}$
possède des produits fibrés et des égalisateurs, le résultat découle du
Lemme \ref{lemma-preserve-equalizers}.)
\end{proof}

\begin{lemma}
\label{lemma-j-shriek-reflects-injections-surjections}
Soit $\mathcal{C}$ un site. Soit $U \in \Ob(\mathcal{C})$.
Le foncteur $j_{U!}$ reflète les injections et les surjections.
\end{lemma}

\begin{proof}
Nous devons montrer que $j_{U!}$ reflète les monomorphismes et les
épimorphismes ; voir le Lemme \ref{lemma-mono-epi-sheaves}.
Grâce au Lemme \ref{lemma-essential-image-j-shriek}, cela se ramène au fait
que le foncteur $\Sh(\mathcal{C})/h_U^\# \to \Sh(\mathcal{C})$
reflète les monomorphismes et les épimorphismes.
\end{proof}

\begin{lemma}
\label{lemma-compute-j-shriek-restrict}
Soit $\mathcal{C}$ un site. Soit $U \in \Ob(\mathcal{C})$.
Pour tout faisceau $\mathcal{F}$ sur $\mathcal{C}$, on a
$j_{U!}j_U^{-1}\mathcal{F} = \mathcal{F} \times h_U^\#$.
\end{lemma}

\begin{proof}
Cela résulte immédiatement de la description de $j_{U!}$ donnée au
Lemme \ref{lemma-describe-j-shriek}.
\end{proof}

\begin{lemma}
\label{lemma-relocalize}
Soit $\mathcal{C}$ un site.
Soit $f : V \to U$ un morphisme de $\mathcal{C}$.
Il existe alors un diagramme commutatif
$$
\xymatrix{
\mathcal{C}/V \ar[rd]_{j_V} \ar[rr]_j & &
\mathcal{C}/U \ar[ld]^{j_U} \\
& \mathcal{C} &
}
$$
de foncteurs continus et cocontinus.
Le foncteur $j : \mathcal{C}/V \to \mathcal{C}/U$,
$(a : W \to V) \mapsto (f \circ a : W \to U)$,
s’identifie au foncteur
$j_{V/U} : (\mathcal{C}/U)/(V/U) \to \mathcal{C}/U$
par l’identification $(\mathcal{C}/U)/(V/U) = \mathcal{C}/V$.
En outre, on a $j_{V!} = j_{U!} \circ j_!$,
$j_V^{-1} = j^{-1} \circ j_U^{-1}$ et
$j_{V*} = j_{U*} \circ j_*$.
\end{lemma}

\begin{proof}
La commutativité du diagramme est immédiate.
L’identification de $j$ à $j_{V/U}$ résulte des définitions.
D’après le Lemme \ref{lemma-composition-cocontinuous}, le diagramme suivant
de morphismes de topos
\begin{equation}
\label{equation-relocalize}
\vcenter{
\xymatrix{
\Sh(\mathcal{C}/V) \ar[rd]_{j_V} \ar[rr]_j & &
\Sh(\mathcal{C}/U) \ar[ld]^{j_U} \\
& \Sh(\mathcal{C}) &
}
}
\end{equation}
est commutatif. Cela démontre les égalités
$j_V^{-1} = j^{-1} \circ j_U^{-1}$ et
$j_{V*} = j_{U*} \circ j_*$.
L’égalité $j_{V!} = j_{U!} \circ j_!$ résulte formellement des propriétés
d’adjonction.
\end{proof}

\begin{lemma}
\label{lemma-relocalize-explicit}
Notons $\mathcal{C}$, $f : V \to U$, $j_U$, $j_V$ et $j$ comme au
Lemme \ref{lemma-relocalize}. Par les identifications
$\Sh(\mathcal{C}/V) = \Sh(\mathcal{C})/h_V^\#$ et
$\Sh(\mathcal{C}/U) = \Sh(\mathcal{C})/h_U^\#$ du
Lemme \ref{lemma-essential-image-j-shriek}, on a :
\begin{enumerate}
\item le foncteur $j^{-1}$ admet la description suivante :
$$
j^{-1}(\mathcal{H} \xrightarrow{\varphi} h_U^\#)
=
(\mathcal{H} \times_{\varphi, h_U^\#, f} h_V^\# \to h_V^\#).
$$
\item le foncteur $j_!$ admet la description suivante :
$$
j_!(\mathcal{H} \xrightarrow{\varphi} h_V^\#) =
(\mathcal{H} \xrightarrow{h_f \circ \varphi} h_U^\#)
$$
\end{enumerate}
\end{lemma}

\begin{proof}
Démonstration de (2). Rappelons que l’identification
$\Sh(\mathcal{C}/V) \to \Sh(\mathcal{C})/h_V^\#$
envoie $\mathcal{G}$ sur
$j_{V!}\mathcal{G} \to j_{V!}(*) = h_V^\#$,
et qu’il en va de même pour
$\Sh(\mathcal{C}/U) \to \Sh(\mathcal{C})/h_U^\#$.
Ainsi, $j_!\mathcal{G}$ s’envoie sur
$j_{U!}(j_!\mathcal{G}) \to j_{U!}(*) = h_U^\#$,
et (2) résulte de l’égalité $j_{U!}j_! = j_{V!}$ donnée par le
Lemme \ref{lemma-relocalize}.

\medskip\noindent
Le lecteur peut maintenant démontrer (1) en utilisant le fait que $j^{-1}$
est adjoint à droite de $j_!$, et que la règle de (1) est adjointe à droite
de celle de (2). Voici une démonstration directe.
Supposons que $\varphi : \mathcal{H} \to h_U^\#$ soit un objet de
$\Sh(\mathcal{C})/h_U^\#$. D’après la démonstration du
Lemme \ref{lemma-essential-image-j-shriek}, il correspond au faisceau
$\mathcal{H}_\varphi$ sur $\mathcal{C}/U$ défini par la règle
$$
(a : W \to U)
\longmapsto
\{ s \in \mathcal{H}(W) \mid \varphi(s) = a\}
$$
sur $\mathcal{C}/U$. L’image inverse $j^{-1}\mathcal{H}_\varphi$ sur
$\mathcal{C}/V$ est donnée par la règle
$$
(a : W \to V)
\longmapsto
\{ s \in \mathcal{H}(W) \mid \varphi(s) = f \circ a\}
$$
d’après la description de $j^{-1} = j_{U/V}^{-1}$ comme restriction de
$\mathcal{H}_\varphi$ à $\mathcal{C}/V$.
D’autre part, en appliquant la règle à l’objet
$$
\xymatrix{
\mathcal{H}' = \mathcal{H} \times_{\varphi, h_U^\#, f} h_V^\#
\ar[rr]^-{\varphi'} & & h_V^\#
}
$$
de $\Sh(\mathcal{C})/h_V^\#$, nous obtenons
$\mathcal{H}'_{\varphi'}$ donné par
\begin{align*}
(a : W \to V)
\longmapsto
&  \{ s' \in \mathcal{H}'(W) \mid \varphi'(s') = a\} \\
= &
\{ (s, a') \in \mathcal{H}(W) \times h_V^\#(W) \mid
a' = a \text{ et } \varphi(s) = f \circ a'\}
\end{align*}
qui est exactement la règle décrivant $j^{-1}\mathcal{H}_\varphi$ ci-dessus.
\end{proof}

\begin{remark}
\label{remark-localize-presheaves}
Localisation et préfaisceaux. Soit $\mathcal{C}$ une catégorie.
Soit $U$ un objet de $\mathcal{C}$. À strictement parler, les foncteurs
$j_U^{-1}$, $j_{U*}$ et $j_{U!}$ n’ont pas été définis pour les préfaisceaux.
Nous pouvons cependant considérer un préfaisceau comme un faisceau pour la
topologie chaotique sur $\mathcal{C}$ (voir l’Exemple
\ref{example-indiscrete}). Nous obtenons donc aussi un foncteur
$$
j_U^{-1} :
\textit{PSh}(\mathcal{C})
\longrightarrow
\textit{PSh}(\mathcal{C}/U)
$$
et des foncteurs
$$
j_{U*}, j_{U!} :
\textit{PSh}(\mathcal{C}/U)
\longrightarrow
\textit{PSh}(\mathcal{C})
$$
qui sont respectivement adjoints à droite et à gauche de $j_U^{-1}$.
D’après le Lemme \ref{lemma-describe-j-shriek}, le préfaisceau
$j_{U!}\mathcal{G}$ est donné par
$$
V \longmapsto
\coprod\nolimits_{\varphi \in \Mor_\mathcal{C}(V, U)}
\mathcal{G}(V \xrightarrow{\varphi} U)
$$
De plus, le foncteur $j_{U!}$ commute aux produits fibrés et aux égalisateurs.
\end{remark}

\begin{remark}
\label{remark-localization-cartesian-cocontinuous}
Soit $\mathcal{C}$ un site. Soit $U \to V$ un morphisme de $\mathcal{C}$.
Les foncteurs cocontinus $\mathcal{C}/U \to \mathcal{C}$ et
$j : \mathcal{C}/U \to \mathcal{C}/V$ (Lemme \ref{lemma-relocalize})
satisfont la propriété $P$ de la Remarque
\ref{remark-cartesian-cocontinuous}.
Par exemple, si l’on se donne des objets $(X/U)$, $(W/V)$, un morphisme
$g : j(X/U) \to (W/V)$ et un recouvrement
$\{f_i  : (W_i/V) \to (W/V)\}$, alors
$(X \times_W W_i/U)$ est une incarnation de
$(X/U) \times_{g, (W/V), f_i} (W_i/V)$, et la famille
$\{(X \times_W W_i/U) \to (X/U)\}$ est un recouvrement de $\mathcal{C}/U$.
\end{remark}



\section{Recollement de faisceaux}
\label{section-glueing-sheaves}

\noindent
Cette section est l’analogue de
Faisceaux, section \ref{sheaves-section-glueing-sheaves}.

\begin{lemma}
\label{lemma-glue-maps}
\begin{slogan}
Les morphismes de faisceaux se recollent.
\end{slogan}
Soit $\mathcal{C}$ un site.
Soit $\{U_i \to U\}$ un recouvrement du site $\mathcal{C}$.
Soient $\mathcal{F}$, $\mathcal{G}$ des faisceaux sur $\mathcal{C}$.
Étant donnée une famille
$$
\varphi_i :
\mathcal{F}|_{\mathcal{C}/U_i}
\longrightarrow
\mathcal{G}|_{\mathcal{C}/U_i}
$$
de morphismes de faisceaux telle que, pour tous $i, j \in I$, les morphismes
$\varphi_i, \varphi_j$ induisent par restriction le même morphisme
$\varphi_{ij} : \mathcal{F}|_{\mathcal{C}/U_i \times_U U_j} \to
\mathcal{G}|_{\mathcal{C}/U_i \times_U U_j}$,
il existe un unique morphisme de faisceaux
$$
\varphi :
\mathcal{F}|_{\mathcal{C}/U}
\longrightarrow
\mathcal{G}|_{\mathcal{C}/U}
$$
dont la restriction à chaque $\mathcal{C}/U_i$ coïncide avec $\varphi_i$.
\end{lemma}

\begin{proof}
Les restrictions employées dans le lemme sont celles du
Lemme \ref{lemma-relocalize}.
Soit $V/U$ un objet de $\mathcal{C}/U$.
Posons $V_i = U_i \times_U V$ et notons $\mathcal{V} = \{V_i \to V\}$.
Remarquons que $(U_i \times_U U_j) \times_U V = V_i \times_V V_j$.
Nous avons alors
$\mathcal{F}|_{\mathcal{C}/U_i}(V_i/U_i) = \mathcal{F}(V_i)$
et
$\mathcal{F}|_{\mathcal{C}/U_i \times_U U_j}(V_i \times_V V_j/U_i \times_U U_j)
= \mathcal{F}(V_i \times_V V_j)$,
et de même pour $\mathcal{G}$.
Nous pouvons donc définir $\varphi$ sur les sections au-dessus de $V$
par les flèches pointillées du diagramme
$$
\xymatrix{
\mathcal{F}(V) \ar@{=}[r] &
H^0(\mathcal{V}, \mathcal{F}) \ar@{..>}[d] \ar[r] &
\prod \mathcal{F}(V_i)
\ar[d]_{\prod \varphi_i}
\ar@<1ex>[r] \ar@<-1ex>[r] &
\prod \mathcal{F}(V_i \times_V V_j) \ar[d]_{\prod \varphi_{ij}} \\
\mathcal{G}(V) \ar@{=}[r] &
H^0(\mathcal{V}, \mathcal{G}) \ar[r] &
\prod \mathcal{G}(V_i)
\ar@<1ex>[r] \ar@<-1ex>[r] &
\prod \mathcal{G}(V_i \times_V V_j)
}
$$
Les signes d’égalité proviennent de la condition de faisceau ; voir la
section \ref{section-sheafification} pour la notation
$H^0(\mathcal{V}, -)$.
Nous omettons la vérification que ces morphismes sont compatibles
avec les applications de restriction.
\end{proof}

\noindent
Le lemme précédent implique qu’étant donnés deux faisceaux $\mathcal{F}$,
$\mathcal{G}$ sur un site $\mathcal{C}$, la règle
$$
U \longmapsto
\Mor_{\Sh(\mathcal{C}/U)}(
\mathcal{F}|_{\mathcal{C}/U}, \mathcal{G}|_{\mathcal{C}/U})
$$
définit un faisceau $\SheafHom(\mathcal{F},\mathcal{G})$.
Il s’agit d’une sorte de {\it faisceau Hom interne}. On l’emploie
rarement dans le cadre des faisceaux d’ensembles, et plus souvent
dans celui des faisceaux de modules ; voir
Modules sur les sites, section \ref{sites-modules-section-internal-hom}.

\begin{lemma}
\label{lemma-internal-hom-sheaf}
\begin{slogan}
La catégorie des faisceaux sur un site est cartésienne fermée
\end{slogan}
Soit $\mathcal{C}$ un site. Soient $\mathcal{F}$, $\mathcal{G}$ et
$\mathcal{H}$ des faisceaux sur $\mathcal{C}$. Il existe une bijection canonique
$$
\Mor_{\Sh(\mathcal{C})}(\mathcal{F}\times\mathcal{G},\mathcal{H}) =
\Mor_{\Sh(\mathcal{C})}(\mathcal{F},\SheafHom(\mathcal{G},\mathcal{H}))
$$
qui est fonctorielle en ses trois arguments.
\end{lemma}

\begin{proof}
Le lemme affirme que les foncteurs $-\times\mathcal{G}$ et
$\SheafHom(\mathcal{G},-)$ sont adjoints l’un de l’autre.
Pour le montrer, nous utilisons les notions d’unité et de counité ; voir
Catégories, section \ref{categories-section-adjoint}.
L’unité
$$
\eta_\mathcal{F} :
\mathcal{F}
\longrightarrow
\SheafHom(\mathcal{G},\mathcal{F}\times\mathcal{G})
$$
envoie $s \in \mathcal{F}(U)$ sur le morphisme
$\mathcal{G}|_{\mathcal{C}/U} \to
\mathcal{F}|_{\mathcal{C}/U}\times\mathcal{G}|_{\mathcal{C}/U}$
qui, au-dessus de $V/U$, est donné par
$$
\mathcal{G}(V) \longrightarrow \mathcal{F}(V)\times \mathcal{G}(V), \quad
t \longmapsto (s|_{V},t).
$$
La counité
$$
\epsilon_{\mathcal{H}} :
\SheafHom(\mathcal{G}, \mathcal{H}) \times \mathcal{G}
\longrightarrow
\mathcal{H}
$$
est le morphisme d’évaluation. Il est donné par la règle
$$
\Mor_{\Sh(\mathcal{C}/U)}(
\mathcal{G}|_{\mathcal{C}/U}, \mathcal{H}|_{\mathcal{C}/U})
\times \mathcal{G}(U)
\longrightarrow
\mathcal{H}(U),\quad
(\varphi, s) \longmapsto \varphi(s).
$$
Alors, pour tout $\varphi : \mathcal{F} \times \mathcal{G} \to \mathcal{H}$,
le morphisme correspondant
$\mathcal{F} \to \SheafHom(\mathcal{G},\mathcal{H})$
est défini en envoyant chaque section $s \in \mathcal{F}(U)$
sur le morphisme de faisceaux sur $\mathcal{C}/U$ qui, sur les
sections au-dessus de $V/U$, est donné par
$$
\mathcal{G}(V) \longrightarrow \mathcal{H}(V),\quad
t \longmapsto \varphi(s|_V, t).
$$
Réciproquement, pour tout
$\psi : \mathcal{F} \to \SheafHom(\mathcal{G}, \mathcal{H})$,
le morphisme correspondant
$\mathcal{F} \times \mathcal{G} \to \mathcal{H}$ est donné par
$$
\mathcal{F}(U) \times \mathcal{G}(U) \longrightarrow \mathcal{H}(U),\quad
(s, t) \longmapsto \psi(s)(t)
$$
sur les sections au-dessus d’un objet $U$. Nous omettons les détails montrant
que ces constructions sont inverses l’une de l’autre.
\end{proof}

\begin{lemma}
\label{lemma-hom-sheaf-hU}
Soit $\mathcal{C}$ un site et soit $U \in \Ob(\mathcal{C})$.
Alors $\SheafHom(h_U^\#, \mathcal{F}) = j_*(\mathcal{F}|_{\mathcal{C}/U})$
pour $\mathcal{F}$ dans $\Sh(\mathcal{C})$.
\end{lemma}

\begin{proof}
On peut le montrer en construisant directement un isomorphisme
de faisceaux. Nous raisonnons plutôt comme suit.
Soit $\mathcal{G}$ un faisceau sur $\mathcal{C}$.
Alors
\begin{align*}
\Mor(\mathcal{G}, j_*(\mathcal{F}|_{\mathcal{C}/U}))
& =
\Mor(\mathcal{G}|_{\mathcal{C}/U}, \mathcal{F}|_{\mathcal{C}/U}) \\
& =
\Mor(j_!(\mathcal{G}|_{\mathcal{C}/U}), \mathcal{F}) \\
& =
\Mor(\mathcal{G} \times h_U^\#, \mathcal{F}) \\
& =
\Mor(\mathcal{G}, \SheafHom(h_U^\#, \mathcal{F}))
\end{align*}
et nous concluons par le lemme de Yoneda. Nous avons utilisé ici les
Lemmes \ref{lemma-internal-hom-sheaf} et
\ref{lemma-compute-j-shriek-restrict}.
\end{proof}

\noindent
Soit $\mathcal{C}$ un site.
Soit $\{U_i \to U\}_{i \in I}$ un recouvrement du site $\mathcal{C}$.
Pour tout $i \in I$, soit $\mathcal{F}_i$ un faisceau d’ensembles sur
$\mathcal{C}/U_i$.
Pour toute paire $i, j \in I$, soit
$$
\varphi_{ij} :
\mathcal{F}_i|_{\mathcal{C}/U_i \times_U U_j}
\longrightarrow
\mathcal{F}_j|_{\mathcal{C}/U_i \times_U U_j}
$$
un isomorphisme de faisceaux d’ensembles. Supposons en outre
que, pour tout triplet d’indices $i, j, k \in I$, le
diagramme suivant soit commutatif :
$$
\xymatrix{
\mathcal{F}_i|_{\mathcal{C}/U_i \times_U U_j \times_U U_k}
\ar[rr]_{\varphi_{ik}}
\ar[rd]_{\varphi_{ij}} & &
\mathcal{F}_k|_{\mathcal{C}/U_i \times_U U_j \times_U U_k} \\
&
\mathcal{F}_j|_{\mathcal{C}/U_i \times_U U_j \times_U U_k}
\ar[ru]_{\varphi_{jk}}
}
$$
Nous appellerons une telle famille de données
$(\mathcal{F}_i, \varphi_{ij})$
une {\it donnée de recollement de faisceaux d’ensembles relativement
au recouvrement $\{U_i \to U\}_{i \in I}$}.

\begin{lemma}
\label{lemma-glue-sheaves}
Soit $\mathcal{C}$ un site.
Soit $\{U_i \to U\}_{i \in I}$ un recouvrement du site $\mathcal{C}$.
Pour toute donnée de recollement $(\mathcal{F}_i, \varphi_{ij})$
de faisceaux d’ensembles relativement au recouvrement $\{U_i \to U\}_{i \in I}$,
il existe un faisceau d’ensembles $\mathcal{F}$ sur $\mathcal{C}/U$
muni d’isomorphismes
$$
\varphi_i : \mathcal{F}|_{\mathcal{C}/U_i} \to \mathcal{F}_i
$$
tels que les diagrammes
$$
\xymatrix{
\mathcal{F}|_{\mathcal{C}/U_i \times_U U_j}
\ar[d]_{\text{id}} \ar[r]_{\varphi_i} &
\mathcal{F}_i|_{\mathcal{C}/U_i \times_U U_j} \ar[d]^{\varphi_{ij}} \\
\mathcal{F}|_{\mathcal{C}/U_i \times_U U_j} \ar[r]^{\varphi_j} &
\mathcal{F}_j|_{\mathcal{C}/U_i \times_U U_j}
}
$$
soient commutatifs.
\end{lemma}

\begin{proof}
Décrivons la construction du faisceau $\mathcal{F}$ sur
$\mathcal{C}/U$. Soit $a : V \to U$ un objet de $\mathcal{C}/U$.
Alors
$$
\mathcal{F}(V/U) = \{
(s_i)_{i \in I} \in \prod_{i \in I} \mathcal{F}_i(U_i \times_U V/U_i)
\mid
\varphi_{ij}(s_i|_{U_i \times_U U_j \times_U V})
=
s_j|_{U_i \times_U U_j \times_U V}
\}
$$
Nous omettons la construction des applications de restriction.
Nous omettons la vérification qu’il s’agit d’un faisceau.
Nous omettons la construction des isomorphismes $\varphi_i$,
ainsi que la preuve de la commutativité des diagrammes
du lemme.
\end{proof}

\noindent
Soit $\mathcal{C}$ un site.
Soit $\{U_i \to U\}_{i \in I}$ un recouvrement du site $\mathcal{C}$.
Soit $\mathcal{F}$ un faisceau sur $\mathcal{C}/U$.
À $\mathcal{F}$ est associée sa
{\it donnée de recollement canonique}, constituée des restrictions
$\mathcal{F}|_{\mathcal{C}/U_i}$ et des isomorphismes canoniques
$$
\left(\mathcal{F}|_{\mathcal{C}/U_i}\right)|_{\mathcal{C}/U_i \times_U U_j}
=
\left(\mathcal{F}|_{\mathcal{C}/U_j}\right)|_{\mathcal{C}/U_i \times_U U_j}
$$
provenant du fait que les composés des foncteurs
$\mathcal{C}/U_i \times_U U_j \to \mathcal{C}/U_i \to \mathcal{C}/U$
et
$\mathcal{C}/U_i \times_U U_j \to \mathcal{C}/U_j \to \mathcal{C}/U$
sont égaux.

\begin{lemma}
\label{lemma-mapping-property-glue}
Soit $\mathcal{C}$ un site.
Soit $\{U_i \to U\}_{i \in I}$ un recouvrement du site $\mathcal{C}$.
La catégorie $\Sh(\mathcal{C}/U)$ est équivalente
à la catégorie des données de recollement par le foncteur qui associe
à $\mathcal{F}$ sur $\mathcal{C}/U$ sa donnée de recollement canonique.
\end{lemma}

\begin{proof}
Dans le
Lemme \ref{lemma-glue-maps},
nous avons vu que le foncteur est pleinement fidèle, et dans le
Lemme \ref{lemma-glue-sheaves},
nous avons démontré qu’il est essentiellement surjectif (en construisant
explicitement un foncteur quasi-inverse).
\end{proof}

\noindent
Soit $\mathcal{C}$ un site. Nous allons étudier
une version du recollement de faisceaux sur le site $\mathcal{C}$ tout entier.
Pour tout objet $U$ de $\mathcal{C}$, soit $\mathcal{F}_U$
un faisceau sur $\mathcal{C}/U$. Rappelons qu’il existe un foncteur
$j_f : \mathcal{C}/V \rightarrow \mathcal{C}/U$ associé à chaque
morphisme $f : V \rightarrow U$ de $\mathcal{C}$, défini par
$(a : W \to V) \mapsto (f \circ a : W \to U)$. Pour chaque tel $f$, soit
$$
c_f : j_f^{-1} \mathcal{F}_U \to \mathcal{F}_V
$$
un isomorphisme de faisceaux. Supposons que, pour deux flèches quelconques
$f : V \to U$ et $g : W \to V$ de $\mathcal{C}$, le composé
$c_g \circ j_{g}^{-1} c_f$ soit égal à $c_{f \circ g}$. Nous appellerons
une telle famille de données $(\mathcal{F}_{U}, c_f)$ une
{\it donnée de recollement absolue de faisceaux d’ensembles sur $\mathcal{C}$}.
Un morphisme de données de recollement absolues
$(\mathcal{F}_{U}, c_f) \to (\mathcal{G}_{U}, c'_f)$
est donné par une famille $(\varphi_U)$
de morphismes de faisceaux $\varphi_U : \mathcal{F}_U \to \mathcal{G}_U$,
telle que
$$
\xymatrix{
j_f^{-1} \mathcal{F}_U \ar[r]_{c_f} \ar[d]_{j_f^{-1} \varphi_U} &
\mathcal{F}_V \ar[d]^{\varphi_V} \\
j_f^{-1} \mathcal{G}_U \ar[r]^{c'_f} &
\mathcal{G}_V
}
$$
soit commutatif pour tout morphisme $f : V \to U$ de $\mathcal{C}$.

\medskip\noindent
À tout faisceau $\mathcal{F}$ sur $\mathcal{C}$ est associée sa
{\it donnée canonique de recollement absolue} $(\mathcal{F}|_{\mathcal{C}/U},c_f)$,
où les isomorphismes canoniques
$c_f : j_f^{-1} \mathcal{F}|_{\mathcal{C}/U} \to
\mathcal{F}|_{\mathcal{C}/V}$
pour $f : V \to U$ proviennent de la relation
$j_V = j_U \circ j_f$ comme dans le Lemme \ref{lemma-relocalize}.
Tout morphisme $\varphi : \mathcal{F} \to \mathcal{G}$
de faisceaux sur $\mathcal{C}$ induit un morphisme
$(\varphi|_{\mathcal{C}/U})$ de données canoniques de recollement absolues.

\begin{lemma}
\label{lemma-glue-sheaves-absolute}
Soit $\mathcal{C}$ un site. La catégorie
$\Sh(\mathcal{C})$ est équivalente à la catégorie
des données de recollement absolues par le foncteur qui
associe à $\mathcal{F}$ sur $\mathcal{C}$ sa
donnée canonique de recollement absolue.
\end{lemma}

\begin{proof}
Étant donnée une donnée de recollement absolue $(\mathcal{F}_U, c_f)$,
nous construisons un faisceau $\mathcal{F}$ sur $\mathcal{C}$ en
posant $\mathcal{F}(U) = \mathcal{F}_U(U)$, la restriction
le long de $f : V \to U$ étant donnée par le diagramme commutatif
$$
\xymatrix{
\mathcal{F}_U(U) \ar[r] \ar@{=}[d] &
\mathcal{F}_U(V) \ar[r]^{c_f} &
\mathcal{F}_V(V) \ar@{=}[d] \\
\mathcal{F}(U) \ar[rr] &  &
\mathcal{F}(V)
}
$$
La condition de compatibilité $c_g \circ j_{g}^{-1} c_f = c_{f \circ g}$
garantit que $\mathcal{F}$ est un préfaisceau, ainsi que le fait que les morphismes
$c_f : \mathcal{F}_U(V) \to \mathcal{F}(V)$ définissent un isomorphisme
$\mathcal{F}_U \to \mathcal{F}|_{\mathcal{C}/U}$ pour tout $U$.
Puisque chaque $\mathcal{F}_U$ est un faisceau, il en résulte que $\mathcal{F}$
est lui aussi un faisceau. Le foncteur $(\mathcal{F}_U, c_f) \mapsto \mathcal{F}$
que nous venons de construire est quasi-inverse du foncteur qui envoie un faisceau
sur $\mathcal{C}$ sur sa donnée de recollement canonique. Nous omettons les détails.
\end{proof}

\begin{remark}
\label{remark-variant-glue-sheaves-absolute}
Il existe une variante du Lemme \ref{lemma-glue-sheaves-absolute}
qui intervient en géométrie algébrique. Supposons en effet que $\mathcal{C}$
soit un site admettant tous les produits fibrés et que, pour tout $U \in \Ob(\mathcal{C})$,
on se donne une sous-catégorie pleine $U_\tau \subset \mathcal{C}/U$ possédant
les propriétés suivantes :
\begin{enumerate}
\item $U/U$ appartient à $U_\tau$,
\item si $V/U$ appartient à $U_\tau$ et si $\{V_j \to V\}$ est un recouvrement
du site $\mathcal{C}$, alors $V_j/U$ appartient à $U_\tau$, et
\item pour un morphisme $U' \to U$ de $\mathcal{C}$ et $V/U$ dans $U_\tau$,
le changement de base $V' = V \times_U U'$ appartient à $U'_\tau$.
\end{enumerate}
Dans ce cadre, $U_\tau$ est un site pour tout $U$ de $\mathcal{C}$,
et le foncteur de changement de base $U_\tau \to U'_\tau$ définit un morphisme
$f_\tau : U_\tau \to U'_\tau$ de sites pour tout morphisme $f : U' \to U$
de $\mathcal{C}$. L’énoncé de recollement obtenu est alors le suivant :
Un faisceau $\mathcal{F}$ sur $\mathcal{C}$ est donné par les données suivantes :
\begin{enumerate}
\item pour tout $U \in \Ob(\mathcal{C})$, un faisceau
$\mathcal{F}_U$ sur $U_\tau$,
\item pour tout $f : U' \to U$ de $\mathcal{C}$, un morphisme
$c_f : f_\tau^{-1}\mathcal{F}_U \to \mathcal{F}_{U'}$.
\end{enumerate}
Ces données sont soumises aux conditions suivantes :
\begin{enumerate}
\item[(a)] étant donnés $f : U' \to U$ et $g : U'' \to U'$ dans
$\mathcal{C}$, le composé
$c_g \circ g_\tau^{-1}c_f$ est égal à $c_{f \circ g}$, et
\item[(b)] si $f : U' \to U$ appartient à $U_\tau$, alors $c_f$ est un isomorphisme.
\end{enumerate}
Si nous en avons un jour besoin, nous en donnerons ici un énoncé précis et une preuve.
(Remarquons que ce résultat diffère légèrement des énoncés
ci-dessus, puisque nous ne demandons pas que tous les morphismes $c_f$ soient des isomorphismes !)
\end{remark}



\section{Compléments sur la localisation}
\label{section-more-localize}

\noindent
Dans cette section, nous démontrons quelques lemmes de localisation en
imposant des hypothèses supplémentaires sur le site ou sur l’objet
en lequel nous localisons.

\begin{lemma}
\label{lemma-describe-j-shriek-good-site}
Soit $\mathcal{C}$ un site.
Soit $U \in \Ob(\mathcal{C})$.
Si la topologie de $\mathcal{C}$ est sous-canonique, voir la
Définition \ref{definition-weaker-than-canonical},
et si $\mathcal{G}$ est un faisceau sur $\mathcal{C}/U$, alors
$$
j_{U!}(\mathcal{G})(V)
=
\coprod\nolimits_{\varphi \in \Mor_\mathcal{C}(V, U)}
\mathcal{G}(V \xrightarrow{\varphi} U),
$$
autrement dit, le passage au faisceau associé n’est pas nécessaire dans le
Lemme \ref{lemma-describe-j-shriek}.
\end{lemma}

\begin{proof}
Soit $\mathcal{V} = \{V_i \to V\}_{i \in I}$ un recouvrement de $V$
dans le site $\mathcal{C}$.
Nous allons vérifier directement la condition de faisceau pour le préfaisceau $\mathcal{H}$
du Lemme \ref{lemma-describe-j-shriek}.
Soit $(s_i, \varphi_i)_{i \in I} \in \prod_i \mathcal{H}(V_i)$.
Cela signifie que $\varphi_i : V_i \to U$ est un morphisme de $\mathcal{C}$ et que
$s_i \in \mathcal{G}(V_i \xrightarrow{\varphi_i} U)$.
La restriction du couple $(s_i, \varphi_i)$ à
$V_i \times_V V_j$ est le couple
$(s_i|_{V_i \times_V V_j/U}, \text{pr}_1 \circ \varphi_i)$, et
de même, la restriction du couple $(s_j, \varphi_j)$ à
$V_i \times_V V_j$ est le couple
$(s_j|_{V_i \times_V V_j/U}, \text{pr}_2 \circ \varphi_j)$.
Par conséquent, si la famille $(s_i, \varphi_i)$ appartient à
$\check{H}^0(\mathcal{V}, \mathcal{H})$, alors
$\text{pr}_1 \circ \varphi_i = \text{pr}_2 \circ \varphi_j$.
Le fait que la topologie de $\mathcal{C}$ soit moins fine que la topologie canonique
implique alors qu’il existe un unique morphisme
$\varphi : V \to U$ tel que $\varphi_i$ soit le composé
de $V_i \to V$ avec $\varphi$. La condition de faisceau pour
$\mathcal{G}$ garantit alors que les sections $s_i$ se recollent en une unique
section $s \in \mathcal{G}(V \xrightarrow{\varphi} U)$.
Ainsi $(s, \varphi) \in \mathcal{H}(V)$, comme voulu.
\end{proof}

\begin{lemma}
\label{lemma-localize-given-products}
Soit $\mathcal{C}$ un site.
Soit $U \in \Ob(\mathcal{C})$.
Supposons que $\mathcal{C}$ admette les produits de deux objets.
Alors
\begin{enumerate}
\item le foncteur $j_U$ admet un adjoint à droite continu,
à savoir le foncteur $v(X) = X \times U / U$,
\item le foncteur $v$ définit un morphisme de sites
$\mathcal{C}/U \to \mathcal{C}$ dont le morphisme de topos associé est égal à
$j_U : \Sh(\mathcal{C}/U) \to \Sh(\mathcal{C})$, et
\item nous avons $j_{U*}\mathcal{F}(X) = \mathcal{F}(X \times U/U)$.
\end{enumerate}
\end{lemma}

\begin{proof}
Dire que le foncteur $v$ est adjoint à droite de $j_U$ signifie que, pour $Y/U$ et $X$,
nous avons
$$
\Mor_\mathcal{C}(Y, X)
=
\Mor_{\mathcal{C}/U}(Y/U, X \times U/U)
$$
ce qui est clair. Pour vérifier que $v$ est continu, soit $\{X_i \to X\}$
un recouvrement du site $\mathcal{C}$. Par le troisième axiome d’un site
(Définition \ref{definition-site}),
nous voyons que
$$
\{X_i \times_X (X \times U) \to X \times_X (X \times U)\}
=
\{X_i \times U \to X \times U\}
$$
est encore un recouvrement du site $\mathcal{C}$. Ainsi $v$ est continu. Les autres
assertions du lemme résultent des Lemmes \ref{lemma-have-functor-other-way}
et \ref{lemma-have-functor-other-way-morphism}.
\end{proof}

\begin{lemma}
\label{lemma-relocalize-given-fibre-products}
Soit $\mathcal{C}$ un site. Soit $U \to V$ un morphisme de $\mathcal{C}$.
Supposons que $\mathcal{C}$ admette les produits fibrés. Soit $j$ comme dans le
Lemme \ref{lemma-relocalize}. Alors
\begin{enumerate}
\item le foncteur $j : \mathcal{C}/U \to \mathcal{C}/V$
admet un adjoint à droite continu, à savoir le foncteur
$v : (X/V) \mapsto (X \times_V U/U)$,
\item le foncteur $v$ définit un morphisme de sites
$\mathcal{C}/U \to \mathcal{C}/V$ dont le morphisme de topos associé est égal à
$j$, et
\item nous avons $j_*\mathcal{F}(X/V) = \mathcal{F}(X \times_V U/U)$.
\end{enumerate}
\end{lemma}

\begin{proof}
Cela résulte du Lemme \ref{lemma-localize-given-products}, puisque $j$ peut être vu
comme un foncteur de localisation par le Lemme \ref{lemma-relocalize}.
\end{proof}

\noindent
Une propriété fondamentale d’une immersion ouverte est
que la restriction de l’image directe et la restriction
du prolongement par l’ensemble vide redonnent le faisceau initial.
Ce n’est pas toujours vrai pour les foncteurs associés à $j_U$
ci-dessus. Cela est vrai lorsque $U$ est un « sous-objet de l’objet final ».

\begin{lemma}
\label{lemma-restrict-back}
Soit $\mathcal{C}$ un site.
Soit $U \in \Ob(\mathcal{C})$.
Supposons que tout $X$ de $\mathcal{C}$ admette au plus
un morphisme vers $U$. Soit $\mathcal{F}$ un faisceau sur $\mathcal{C}/U$.
Les morphismes canoniques $\mathcal{F} \to j_U^{-1}j_{U!}\mathcal{F}$
et $j_U^{-1}j_{U*}\mathcal{F} \to \mathcal{F}$ sont des
isomorphismes.
\end{lemma}

\begin{proof}
C’est un cas particulier du Lemme \ref{lemma-back-and-forth},
car l’hypothèse sur $U$ équivaut à la pleine fidélité
du foncteur de localisation $\mathcal{C}/U \to \mathcal{C}$.
\end{proof}

\begin{lemma}
\label{lemma-localize-cartesian-square}
Soit $\mathcal{C}$ un site. Soit
$$
\xymatrix{
U' \ar[d] \ar[r] & U \ar[d] \\
V' \ar[r] & V
}
$$
un diagramme commutatif de $\mathcal{C}$. Les
morphismes du Lemme \ref{lemma-relocalize}
fournissent les diagrammes commutatifs
$$
\vcenter{
\xymatrix{
\mathcal{C}/U' \ar[d]_{j_{U'/V'}} \ar[r]_{j_{U'/U}} &
\mathcal{C}/U \ar[d]^{j_{U/V}} \\
\mathcal{C}/V' \ar[r]^{j_{V'/V}} & \mathcal{C}/V
}
}
\quad\text{et}\quad
\vcenter{
\xymatrix{
\Sh(\mathcal{C}/U') \ar[d]_{j_{U'/V'}} \ar[r]_{j_{U'/U}} &
\Sh(\mathcal{C}/U) \ar[d]^{j_{U/V}} \\
\Sh(\mathcal{C}/V') \ar[r]^{j_{V'/V}} &
\Sh(\mathcal{C}/V)
}
}
$$
de foncteurs continus et cocontinus, et de topos.
De plus, si le diagramme initial de $\mathcal{C}$ est cartésien,
alors nous avons
$j_{V'/V}^{-1} \circ j_{U/V, *} = j_{U'/V', *} \circ j_{U'/U}^{-1}$.
\end{lemma}

\begin{proof}
La commutativité du carré de gauche dans la première assertion du lemme
résulte immédiatement des définitions. Elle implique la commutativité
du diagramme de topos par le Lemme \ref{lemma-composition-cocontinuous}.
Supposons le diagramme cartésien.
Par l’unicité des foncteurs adjoints, montrer
$j_{V'/V}^{-1} \circ j_{U/V, *} = j_{U'/V', *} \circ j_{U'/U}^{-1}$
équivaut à montrer
$j_{U/V}^{-1} \circ j_{V'/V!} = j_{U'/U!} \circ j_{U'/V'}^{-1}$.
Par les identifications du Lemme \ref{lemma-essential-image-j-shriek},
nous pouvons identifier notre diagramme de topos à
$$
\xymatrix{
\Sh(\mathcal{C})/h_{U'}^\# \ar[d] \ar[r] &
\Sh(\mathcal{C})/h_U^\# \ar[d] \\
\Sh(\mathcal{C})/h_{V'}^\# \ar[r] &
\Sh(\mathcal{C})/h_V^\#
}
$$
et nous savons interpréter les foncteurs $j^{-1}$ et $j_!$
grâce au Lemme \ref{lemma-relocalize-explicit}. Nous devons donc montrer,
étant donné $\mathcal{F} \to h_{V'}^\#$, que
$$
\mathcal{F} \times_{h_{V'}^\#} h_{U'}^\# =
\mathcal{F} \times_{h_V^\#} h_U^\#
$$
en tant que faisceaux munis d’un morphisme vers $h_U^\#$.
C’est vrai parce que $h_{U'} = h_{V'} \times_{h_V} h_U$
et donc aussi
$$
h_{U'}^\# = h_{V'}^\# \times_{h_V^\#} h_U^\#
$$
puisque le passage au faisceau associé est exact.
\end{proof}











\section{Localisation et morphismes}
\label{section-localize-morphisms}

\noindent
Le lemme suivant est important pour comprendre la relation
entre la localisation et les morphismes de sites et de topos.

\begin{lemma}
\label{lemma-localize-morphism}
Soit $f : \mathcal{C} \to \mathcal{D}$ un morphisme de sites
correspondant au foncteur continu $u : \mathcal{D} \to \mathcal{C}$.
Soit $V \in \Ob(\mathcal{D})$ et posons $U = u(V)$.
Alors le foncteur $u' : \mathcal{D}/V \to \mathcal{C}/U$,
$V'/V \mapsto u(V')/U$, détermine un morphisme de sites
$f' : \mathcal{C}/U \to \mathcal{D}/V$.
Le morphisme $f'$ s’insère dans un diagramme commutatif de topos
$$
\xymatrix{
\Sh(\mathcal{C}/U) \ar[r]_{j_U} \ar[d]_{f'} &
\Sh(\mathcal{C}) \ar[d]^f \\
\Sh(\mathcal{D}/V) \ar[r]^{j_V} &
\Sh(\mathcal{D}).
}
$$
En utilisant les identifications
$\Sh(\mathcal{C}/U) = \Sh(\mathcal{C})/h_U^\#$ et
$\Sh(\mathcal{D}/V) = \Sh(\mathcal{D})/h_V^\#$ du
Lemme \ref{lemma-essential-image-j-shriek},
le foncteur $(f')^{-1}$ est décrit par la règle
$$
(f')^{-1}(\mathcal{H} \xrightarrow{\varphi} h_V^\#)
=
(f^{-1}\mathcal{H} \xrightarrow{f^{-1}\varphi} h_U^\#).
$$
Enfin, nous avons $f'_*j_U^{-1} = j_V^{-1}f_*$.
\end{lemma}

\begin{proof}
Il est clair que $u'$ est continu, et nous obtenons donc les foncteurs
$f'_* = (u')^s = (u')^p$ (voir les
sections \ref{section-functoriality-PSh}
et \ref{section-continuous-functors})
et un adjoint $(f')^{-1} = (u')_s = ((u')_p\ )^\#$. L’assertion
$f'_*j_U^{-1} = j_V^{-1}f_*$ résulte de
$$
(j_V^{-1}f_*\mathcal{F})(V'/V)
= f_*\mathcal{F}(V') = \mathcal{F}(u(V'))
= (j_U^{-1}\mathcal{F})(u(V')/U)
= (f'_*j_U^{-1}\mathcal{F})(V'/V)
$$
ce qui vaut même pour les préfaisceaux. Ce qui n’est pas clair a priori, c’est
que $(f')^{-1}$ est exact, que le diagramme commute et que
la description de $(f')^{-1}$ est valable.

\medskip\noindent
Soit $\mathcal{H}$ un faisceau sur $\mathcal{D}/V$.
Calculons $j_{U!}(f')^{-1}\mathcal{H}$. Nous avons
\begin{eqnarray*}
j_{U!}(f')^{-1}\mathcal{H}
& =
((j_U)_p(u'_p\mathcal{H})^\#)^\# \\
& =
((j_U)_pu'_p\mathcal{H})^\# \\
& =
(u_p(j_V)_p\mathcal{H})^\# \\
& =
f^{-1}j_{V!}\mathcal{H}
\end{eqnarray*}
La première égalité résulte du développement des définitions.
La deuxième résulte du Lemme \ref{lemma-technical-up}.
La troisième vient de ce que $u \circ j_V = j_U \circ u'$.
La quatrième résulte à nouveau du Lemme \ref{lemma-technical-up}.
Toutes les égalités ci-dessus sont des isomorphismes de foncteurs,
et nous pouvons donc les interpréter en disant que le diagramme
suivant de catégories et de foncteurs est commutatif :
$$
\xymatrix{
\Sh(\mathcal{C}/U) \ar[r]_{j_{U!}} &
\Sh(\mathcal{C})/h_U^\# \ar[r] &
\Sh(\mathcal{C}) \\
\Sh(\mathcal{D}/V) \ar[r]^{j_{V!}} \ar[u]^{(f')^{-1}} &
\Sh(\mathcal{D})/h_V^\# \ar[r] \ar[u]^{f^{-1}} &
\Sh(\mathcal{D}) \ar[u]^{f^{-1}}
}
$$
La flèche médiane a un sens puisque $f^{-1}h_V^\# = (h_{u(V)})^\# = h_U^\#$, voir le
Lemme \ref{lemma-pullback-representable-sheaf}.
Ceci démontre en particulier la description de $(f')^{-1}$ donnée
dans l’énoncé du lemme.
Puisque, d’après le
Lemme \ref{lemma-essential-image-j-shriek},
les flèches horizontales de gauche sont des équivalences,
et puisque $f^{-1}$ est exact par hypothèse, nous concluons que
$(f')^{-1} = u'_s$ est exact. En effet, comme c’est un adjoint à gauche,
il est déjà exact à droite
(Catégories, Lemme \ref{categories-lemma-adjoint-exact}).
Il nous suffit donc de montrer qu’il
transforme un objet final en un objet final et commute
aux produits fibrés
(Catégories, Lemme \ref{categories-lemma-characterize-left-exact}).
Ces deux propriétés sont claires pour le foncteur induit
$f^{-1} : \Sh(\mathcal{D})/h_V^\# \to \Sh(\mathcal{C})/h_U^\#$.
Ceci démontre que $f'$ est un morphisme de sites.

\medskip\noindent
Il nous reste à vérifier que $(f')^{-1}j_V^{-1} = j_U^{-1}f^{-1}$.
Pour le voir, combinons la formule ci-dessus avec la description
du Lemme \ref{lemma-compute-j-shriek-restrict}. Nous obtenons ainsi,
pour tout faisceau $\mathcal{G}$ sur $\mathcal{D}$,
$$
j_{U!}(f')^{-1}j_V^{-1}\mathcal{G}
=
f^{-1}j_{V!}j_V^{-1}\mathcal{G}
=
f^{-1}(\mathcal{G} \times h_V^\#)
=
f^{-1}\mathcal{G} \times h_U^\#
=
j_{U!}j_U^{-1}f^{-1}\mathcal{G}.
$$
Puisque le foncteur $j_{U!}$ induit une équivalence
$\Sh(\mathcal{C}/U) \to \Sh(\mathcal{C})/h_U^\#$,
nous concluons.
\end{proof}

\noindent
Le lemme suivant est un cas particulier du
Lemme \ref{lemma-localize-morphism}
ci-dessus, plus général.

\begin{lemma}
\label{lemma-localize-morphism-strong}
Soient $\mathcal{C}$ et $\mathcal{D}$ des sites.
Soit $u : \mathcal{D} \to \mathcal{C}$ un foncteur.
Soit $V \in \Ob(\mathcal{D})$. Posons $U = u(V)$.
Supposons que
\begin{enumerate}
\item $\mathcal{C}$ et $\mathcal{D}$ admettent
toutes les limites finies,
\item $u$ soit continu, et
\item $u$ commute aux limites finies.
\end{enumerate}
Il existe un diagramme commutatif de morphismes de sites
$$
\xymatrix{
\mathcal{C}/U \ar[r]_{j_U} \ar[d]_{f'} & \mathcal{C} \ar[d]^f \\
\mathcal{D}/V \ar[r]^{j_V} & \mathcal{D}
}
$$
où la flèche verticale de droite correspond à $u$,
la flèche verticale de gauche correspond au
foncteur $u' : \mathcal{D}/V \to \mathcal{C}/U$, $V'/V \mapsto u(V')/u(V)$,
et les flèches horizontales correspondent aux foncteurs
$\mathcal{C} \to \mathcal{C}/U$, $X \mapsto X \times U$,
et $\mathcal{D} \to \mathcal{D}/V$, $Y \mapsto Y \times V$,
comme dans le Lemme \ref{lemma-localize-given-products}.
De plus, le diagramme associé de morphismes de topos est
égal au diagramme du
Lemme \ref{lemma-localize-morphism}.
En particulier, nous avons $f'_*j_U^{-1} = j_V^{-1}f_*$.
\end{lemma}

\begin{proof}
Notons que $u$ satisfait aux hypothèses de la
Proposition \ref{proposition-get-morphism} et induit donc
un morphisme de sites $f : \mathcal{C} \to \mathcal{D}$ par cette proposition.
Il est clair que $u$ induit un foncteur $u'$ comme indiqué.
Il est clair que ce foncteur satisfait également aux hypothèses de la
Proposition \ref{proposition-get-morphism}.
Nous obtenons donc un morphisme de sites $f' : \mathcal{C}/U \to \mathcal{D}/V$.
Le diagramme commute par notre définition de la composition des morphismes de
sites (voir la Définition \ref{definition-composition-morphisms-sites})
et par l’égalité
$$
u(Y \times V) = u(Y) \times u(V) = u(Y) \times U
$$
qui montre que le diagramme de catégories et de foncteurs opposé
à celui du lemme est commutatif.
\end{proof}

\noindent
À ce stade, nous pouvons localiser un site, nous savons le relocaliser,
et nous pouvons localiser un morphisme de sites en un objet du site d’arrivée.
En combinant ces opérations, nous obtenons le type de diagramme suivant.

\begin{lemma}
\label{lemma-relocalize-morphism}
Soit $f : \mathcal{C} \to \mathcal{D}$ un morphisme de sites
correspondant au foncteur continu $u : \mathcal{D} \to \mathcal{C}$.
Soient $V \in \Ob(\mathcal{D})$, $U \in \Ob(\mathcal{C})$
et $c : U \to u(V)$ un morphisme de $\mathcal{C}$.
Il existe un diagramme commutatif de topos
$$
\xymatrix{
\Sh(\mathcal{C}/U) \ar[r]_{j_U} \ar[d]_{f_c} &
\Sh(\mathcal{C}) \ar[d]^f \\
\Sh(\mathcal{D}/V) \ar[r]^{j_V} &
\Sh(\mathcal{D}).
}
$$
Nous avons $f_c = f' \circ j_{U/u(V)}$, où
$f' : \Sh(\mathcal{C}/u(V)) \to \Sh(\mathcal{D}/V)$
est celui du
Lemme \ref{lemma-localize-morphism},
et où
$j_{U/u(V)} : \Sh(\mathcal{C}/U) \to \Sh(\mathcal{C}/u(V))$
est celui du
Lemme \ref{lemma-relocalize}.
En utilisant les identifications
$\Sh(\mathcal{C}/U) = \Sh(\mathcal{C})/h_U^\#$ et
$\Sh(\mathcal{D}/V) = \Sh(\mathcal{D})/h_V^\#$ du
Lemme \ref{lemma-essential-image-j-shriek},
le foncteur $(f_c)^{-1}$ est décrit par la règle
$$
(f_c)^{-1}(\mathcal{H} \xrightarrow{\varphi} h_V^\#)
=
(f^{-1}\mathcal{H} \times_{f^{-1}\varphi, h_{u(V)}^\#, c} h_U^\#
\rightarrow h_U^\#).
$$
Enfin, étant donnés des morphismes $b : V' \to V$, $a : U' \to U$ et
$c' : U' \to u(V')$ tels que
$$
\xymatrix{
U' \ar[r]_-{c'} \ar[d]_a & u(V') \ar[d]^{u(b)} \\
U \ar[r]^-c & u(V)
}
$$
commute, le diagramme
$$
\xymatrix{
\Sh(\mathcal{C}/U') \ar[r]_{j_{U'/U}} \ar[d]_{f_{c'}} &
\Sh(\mathcal{C}/U) \ar[d]^{f_c} \\
\Sh(\mathcal{D}/V') \ar[r]^{j_{V'/V}} &
\Sh(\mathcal{D}/V).
}
$$
est commutatif.
\end{lemma}

\begin{proof}
Ce lemme se démontre de lui-même et constitue plutôt une synthèse des faits connus
à ce stade du développement de la théorie. Par exemple, la commutativité
du premier carré résulte de la commutativité du
Diagramme (\ref{equation-relocalize})
et de celle du diagramme du
Lemme \ref{lemma-localize-morphism}.
La description de $f_c^{-1}$ résulte de la combinaison du
Lemme \ref{lemma-relocalize-explicit}
avec le
Lemme \ref{lemma-localize-morphism}.
La commutativité du dernier carré résulte alors de
l’égalité
$$
f^{-1}\mathcal{H} \times_{h_{u(V)}^\#, c} h_U^\# \times_{h_U^\#} h_{U'}^\#
=
f^{-1}(\mathcal{H} \times_{h_V^\#} h_{V'}^\#)
\times_{h_{u(V'), c'}^\#} h_{U'}^\#
$$
qui est formelle puisque $f^{-1}h_V^\# = h_{u(V)}^\#$ et
$f^{-1}h_{V'}^\# = h_{u(V')}^\#$, voir le
Lemme \ref{lemma-pullback-representable-sheaf}.
\end{proof}

\noindent
Dans le lemme suivant, nous rencontrons une autre forme de fonctorialité de la
localisation, lorsque le morphisme de topos provient d’un
foncteur cocontinu. Ce type de diagramme diffère
du diagramme du
Lemme \ref{lemma-localize-morphism},
et, en particulier, l’égalité $f'_*j_U^{-1} = j_V^{-1}f_*$ vue dans le
Lemme \ref{lemma-localize-morphism}
n’est en général pas valable dans la situation du lemme suivant.

\begin{lemma}
\label{lemma-localize-cocontinuous}
Soient $\mathcal{C}$ et $\mathcal{D}$ des sites.
Soit $u : \mathcal{C} \to \mathcal{D}$ un foncteur cocontinu.
Soit $U$ un objet de $\mathcal{C}$, et posons $V = u(U)$.
Nous avons un diagramme commutatif
$$
\xymatrix{
\mathcal{C}/U \ar[r]_{j_U} \ar[d]_{u'} & \mathcal{C} \ar[d]^u \\
\mathcal{D}/V \ar[r]^-{j_V} & \mathcal{D}
}
$$
où la flèche verticale de gauche est
$u' : \mathcal{C}/U \to \mathcal{D}/V$, $U'/U \mapsto V'/V$.
Alors $u'$ est également cocontinu et nous obtenons un diagramme commutatif de topos
$$
\xymatrix{
\Sh(\mathcal{C}/U) \ar[r]_{j_U} \ar[d]_{f'} &
\Sh(\mathcal{C}) \ar[d]^f \\
\Sh(\mathcal{D}/V) \ar[r]^-{j_V} &
\Sh(\mathcal{D})
}
$$
où $f$ (resp.\ $f'$) correspond à $u$ (resp.\ $u'$).
\end{lemma}

\begin{proof}
La commutativité du premier diagramme est claire.
Elle implique celle du second diagramme pourvu que nous
montrions que $u'$ est cocontinu.

\medskip\noindent
Soit $U'/U$ un objet de $\mathcal{C}/U$.
Soit $\{V_j/V \to u(U')/V\}_{j \in J}$ un recouvrement de $u(U')/V$
dans $\mathcal{D}/V$. Puisque $u$ est cocontinu, il existe un
recouvrement $\{U_i' \to U'\}_{i \in I}$ tel que la famille
$\{u(U_i') \to u(U')\}$ raffine le recouvrement
$\{V_j \to u(U')\}$ dans $\mathcal{D}$. Autrement dit, il existe
une application $\alpha : I \to J$ et des morphismes
$\phi_i : u(U_i') \to V_{\alpha(i)}$ au-dessus de $U'$.
Considérons $U_i'$ comme un objet au-dessus de
$U$ par la composée $U'_i \to U' \to U$. Alors
$\{U'_i/U \to U'/U\}$ est un recouvrement de $\mathcal{C}/U$ tel que
$\{u(U_i')/V \to u(U')/V\}$ raffine $\{V_j/V \to u(U')/V\}$
(en utilisant les mêmes $\alpha$ et $\phi_i$). Ainsi,
$u' : \mathcal{C}/U \to \mathcal{D}/V$ est cocontinu.
\end{proof}

\begin{lemma}
\label{lemma-localize-cocontinuous-downstairs}
Soient $\mathcal{C}$ et $\mathcal{D}$ des sites.
Soit $u : \mathcal{C} \to \mathcal{D}$ un foncteur cocontinu.
Soit $V$ un objet de $\mathcal{D}$. Soit
${}^u_V\mathcal{I}$ la catégorie introduite dans la
section \ref{section-more-functoriality-PSh}.
Nous avons un diagramme commutatif
$$
\vcenter{
\xymatrix{
\,_V^u\mathcal{I} \ar[r]_j \ar[d]_{u'} &
\mathcal{C} \ar[d]^u \\
\mathcal{D}/V \ar[r]^-{j_V} &
\mathcal{D}
}
}
\quad\text{où}\quad
\begin{matrix}
j : (U, \psi) \mapsto U \\
u' : (U, \psi) \mapsto (\psi : u(U) \to V)
\end{matrix}
$$
Déclarons qu’une famille de morphismes $\{(U_i, \psi_i) \to (U, \psi)\}$
de ${}^u_V\mathcal{I}$ est un recouvrement si et seulement si
$\{U_i \to U\}$ est un recouvrement de $\mathcal{C}$.
Alors
\begin{enumerate}
\item ${}^u_V\mathcal{I}$ est un site,
\item $j$ est continu et cocontinu,
\item $u'$ est cocontinu,
\item nous obtenons un diagramme commutatif de topos
$$
\xymatrix{
\Sh({}^u_V\mathcal{I}) \ar[r]_j \ar[d]_{f'} &
\Sh(\mathcal{C}) \ar[d]^f \\
\Sh(\mathcal{D}/V) \ar[r]^-{j_V} &
\Sh(\mathcal{D})
}
$$
où $f$ (resp.\ $f'$) correspond à $u$ (resp.\ $u'$), et
\item nous avons $f'_*j^{-1} = j_V^{-1}f_*$.
\end{enumerate}
\end{lemma}

\begin{proof}
Les parties (1), (2), (3) et (4) sont des conséquences immédiates des
définitions et du fait que le foncteur $j$ commute aux produits fibrés.
Nous omettons les détails. Pour voir (5), rappelons que $f_*$ est donné par
${}_su = {}_pu$. Ainsi, la valeur de $j_V^{-1}f_*\mathcal{F}$ en
$V'/V$ est la valeur de ${}_pu\mathcal{F}$ en $V'$, laquelle est la
limite des valeurs de $\mathcal{F}$ sur la catégorie
${}^u_{V'}\mathcal{I}$. Il existe manifestement une équivalence
de catégories
$$
{}^u_{V'}\mathcal{I} \to {}^{u'}_{V'/V}\mathcal{I}
$$
Puisque la valeur de $f'_*j^{-1}\mathcal{F}$ en $V'/V$ est
donnée par la limite des valeurs de $j^{-1}\mathcal{F}$
sur la catégorie ${}^{u'}_{V'/V}\mathcal{I}$, et puisque
les valeurs de $j^{-1}\mathcal{F}$ sur les objets de
${}^u_V\mathcal{I}$ sont simplement les valeurs de $\mathcal{F}$
(d’après le Lemme \ref{lemma-when-shriek}, puisque $j$ est continu et
cocontinu),
nous voyons que (5) est vraie.
\end{proof}

\noindent
Les deux résultats suivants sont d’une nature légèrement différente.

\begin{lemma}
\label{lemma-special-square-cocontinuous}
Supposons donnés des sites $\mathcal{C}', \mathcal{C}, \mathcal{D}', \mathcal{D}$
et des foncteurs
$$
\xymatrix{
\mathcal{C}' \ar[r]_{v'} \ar[d]_{u'} &
\mathcal{C} \ar[d]^u \\
\mathcal{D}' \ar[r]^v &
\mathcal{D}
}
$$
Supposons que
\begin{enumerate}
\item $u$, $u'$, $v$ et $v'$ soient cocontinus et donnent naissance aux morphismes de
topos $f$, $f'$, $g$ et $g'$ par le Lemme \ref{lemma-cocontinuous-morphism-topoi},
\item $v \circ u' = u \circ v'$,
\item $v$ et $v'$ soient continus aussi bien que cocontinus, et
\item pour tout objet $V'$ de $\mathcal{D}'$, le foncteur
${}^{u'}_{V'}\mathcal{I} \to {}^{\ \ \ u}_{v(V')}\mathcal{I}$
donné par $v$ soit cofinal.
\end{enumerate}
Alors $f'_* \circ (g')^{-1} = g^{-1} \circ f_*$ et
$g'_! \circ (f')^{-1} = f^{-1} \circ g_!$.
\end{lemma}

\begin{proof}
Les catégories ${}^{u'}_{V'}\mathcal{I}$ et
${}^{\ \ \ u}_{v(V')}\mathcal{I}$ sont définies dans la
section \ref{section-more-functoriality-PSh}.
Le foncteur de la condition (4) envoie l’objet
$\psi : u'(U') \to V'$ de ${}^{u'}_{V'}\mathcal{I}$ sur l’objet
$v(\psi) : uv'(U') = vu'(U') \to v(V')$ de ${}^{\ \ \ u}_{v(V')}\mathcal{I}$.
Rappelons que $g^{-1}$ est donné par $v^p$ (Lemme \ref{lemma-when-shriek}) et
que $f_*$ est donné par ${}_su = {}_pu$. La valeur de
$g^{-1}f_*\mathcal{F}$ en $V'$ est donc la valeur de ${}_pu\mathcal{F}$
en $v(V')$, c’est-à-dire la limite
$$
\lim_{u(U) \to v(V') \in \Ob({}^{\ \ \ u}_{v(V')}\mathcal{I}^{opp})}
\mathcal{F}(U)
$$
Par le même raisonnement, la valeur de $f'_*(g')^{-1}\mathcal{F}$
en $V'$ est donnée par la limite
$$
\lim_{u'(U') \to V' \in \Ob({}^{u'}_{V'}\mathcal{I}^{opp})} \mathcal{F}(v'(U'))
$$
Ainsi, l’hypothèse (4) et Catégories, Lemme \ref{categories-lemma-initial},
montrent que ces limites coïncident, ce qui démontre la première égalité du
lemme. La seconde égalité découle de la première par unicité des
adjoints.
\end{proof}

\begin{lemma}
\label{lemma-special-square-continuous}
Supposons donnés des sites $\mathcal{C}', \mathcal{C}, \mathcal{D}', \mathcal{D}$
et des foncteurs
$$
\xymatrix{
\mathcal{C}' \ar[r]_{v'} &
\mathcal{C} \\
\mathcal{D}' \ar[r]^v \ar[u]^{u'} &
\mathcal{D} \ar[u]_u
}
$$
Avec les notations des sections \ref{section-morphism-sites}
et \ref{section-cocontinuous-morphism-topoi}, supposons que
\begin{enumerate}
\item $u$ et $u'$ soient continus et donnent naissance aux morphismes de
sites $f$ et $f'$,
\item $v$ et $v'$ soient cocontinus et donnent naissance aux morphismes
de topos $g$ et $g'$,
\item $u \circ v = v' \circ u'$, et
\item $v$ et $v'$ soient continus aussi bien que cocontinus.
\end{enumerate}
Alors\footnote{Dans cette généralité,
nous ne savons pas si $f \circ g'$ est égal à $g \circ f'$
comme morphismes de topos (il existe une $2$-flèche canonique du
premier vers le second, qui peut ne pas être un isomorphisme).}
$f'_* \circ (g')^{-1} = g^{-1} \circ f_*$ et
$g'_! \circ (f')^{-1} = f^{-1} \circ g_!$.
\end{lemma}

\begin{proof}
En effet, nous avons
$$
f'_*(g')^{-1}\mathcal{F} = (u')^p((v')^p\mathcal{F})^\# =
(u')^p(v')^p\mathcal{F}
$$
La première égalité résulte de la définition, et la
seconde du Lemme \ref{lemma-when-shriek}. Nous avons
$$
g^{-1}f_*\mathcal{F} = (v^pu^p\mathcal{F})^\# =
((u')^p(v')^p\mathcal{F})^\# = (u')^p(v')^p\mathcal{F}
$$
La première égalité résulte de la définition, la deuxième de ce que
$u \circ v = v' \circ u'$, et la troisième de ce que nous avons déjà vu que
$(u')^p(v')^p\mathcal{F}$ est un faisceau. Cela démontre
$f'_* \circ (g')^{-1} = g^{-1} \circ f_*$ et l’égalité
$g'_! \circ (f')^{-1} = f^{-1} \circ g_!$ résulte de
l’unicité des adjoints à gauche.
\end{proof}
















\section{Morphismes de topos}
\label{section-morphisms-topoi}

\noindent
Dans cette section, nous montrons que tout morphisme de topos est équivalent
à un morphisme de topos qui provient d’un morphisme de sites.
Comparer avec \cite[Expos\'e III, Proposition 4.9.4]{SGA4}.

\begin{lemma}
\label{lemma-equivalence}
Soient $\mathcal{C}$, $\mathcal{D}$ des sites.
Soit $u : \mathcal{C} \to \mathcal{D}$ un foncteur.
Supposons que
\begin{enumerate}
\item $u$ soit cocontinu,
\item $u$ soit continu,
\item étant donnés $a, b : U' \to U$ dans $\mathcal{C}$ tels que
$u(a) = u(b)$, il existe un recouvrement $\{f_i : U'_i \to U'\}$
dans $\mathcal{C}$ tel que $a \circ f_i = b \circ f_i$,
\item étant donnés $U', U \in \Ob(\mathcal{C})$ et
un morphisme $c : u(U') \to u(U)$ dans $\mathcal{D}$, il existe
un recouvrement $\{f_i : U_i' \to U'\}$ dans $\mathcal{C}$
et des morphismes $c_i : U_i' \to U$ tels que $u(c_i) = c \circ u(f_i)$, et
\item pour tout $V \in \Ob(\mathcal{D})$, il existe un recouvrement
de $V$ dans $\mathcal{D}$ de la forme $\{u(U_i) \to V\}_{i \in I}$.
\end{enumerate}
Alors le morphisme de topos
$$
g : \Sh(\mathcal{C}) \longrightarrow \Sh(\mathcal{D})
$$
associé au foncteur cocontinu $u$ par le
Lemme \ref{lemma-cocontinuous-morphism-topoi}
est une équivalence.
\end{lemma}

\begin{proof}
Supposons que $u$ satisfasse les propriétés (1) -- (5). Nous allons montrer que
les morphismes d’adjonction
$$
\mathcal{G} \longrightarrow g_*g^{-1}\mathcal{G}
\quad\text{et}\quad
g^{-1}g_*\mathcal{F} \longrightarrow \mathcal{F}
$$
sont des isomorphismes.

\medskip\noindent
Remarquons que le Lemme \ref{lemma-when-shriek} s’applique et que nous avons
$g^{-1}\mathcal{G}(U) = \mathcal{G}(u(U))$ pour tout faisceau $\mathcal{G}$
sur $\mathcal{D}$. Ensuite, soit $\mathcal{F}$ un faisceau sur $\mathcal{C}$,
et soit $V$ un objet de $\mathcal{D}$. Par définition, nous avons
$g_*\mathcal{F}(V) = \lim_{u(U) \to V} \mathcal{F}(U)$.
Par conséquent
$$
g^{-1}g_*\mathcal{F}(U) = \lim_{U', u(U') \to u(U)} \mathcal{F}(U')
$$
où les morphismes $\psi : u(U') \to u(U)$ ne sont pas nécessairement de la forme
$u(\alpha)$. La catégorie de ces couples $(U', \psi)$ possède un objet final,
à savoir $(U, \text{id})$, qui donne le morphisme de
la limite vers $\mathcal{F}(U)$. Soit $(s_{(U', \psi)})$ un élément
de la limite. Nous voulons montrer que $s_{(U', \psi)}$ est uniquement déterminé
par la valeur $s_{(U, \text{id})} \in \mathcal{F}(U)$. D’après la propriété (4), pour
tout $(U', \psi)$, il existe un recouvrement $\{U'_i \to U'\}$ tel que les
composés $u(U'_i) \to u(U') \to u(U)$ soient de la forme $u(c_i)$
pour certains $c_i : U'_i \to U$ dans $\mathcal{C}$. Par conséquent
$$
s_{(U', \psi)}|_{U'_i} = c_i^*(s_{(U, \text{id})}).
$$
Comme $\mathcal{F}$ est un faisceau, il s’ensuit que $s_{(U', \psi)}$
est bien déterminé par $s_{(U, \text{id})}$. Cela démontre l’unicité.
Pour l’existence, supposons donnés
$s \in \mathcal{F}(U)$, $\psi : u(U') \to u(U)$, $\{f_i : U_i' \to U'\}$
et $c_i : U_i' \to U$ tels que $\psi \circ u(f_i) = u(c_i)$ comme ci-dessus.
Nous affirmons qu’il existe un (unique) élément
$s_{(U', \psi)} \in \mathcal{F}(U')$ tel que
$$
s_{(U', \psi)}|_{U'_i} = c_i^*(s).
$$
En effet, a priori, il n’est pas clair que les éléments
$c_i^*(s)|_{U_i' \times_{U'} U_j'}$
et $c_j^*(s)|_{U_i' \times_{U'} U_j'}$ coïncident, car
le diagramme
$$
\xymatrix{
U_i' \times_{U'} U_j' \ar[r]_-{\text{pr}_2} \ar[d]_{\text{pr}_1} &
U_j' \ar[d]^{c_j} \\
U_i' \ar[r]^{c_i} & U}
$$
ne commute pas nécessairement. Mais la condition (3) du lemme garantit qu’il
existe des recouvrements
$\{f_{ijk} : U'_{ijk} \to U_i' \times_{U'} U_j'\}_{k \in K_{ij}}$ tels que
$c_i \circ \text{pr}_1 \circ f_{ijk} = c_j \circ \text{pr}_2 \circ f_{ijk}$.
Par conséquent
$$
f_{ijk}^* \left(c_i^*s|_{U_i' \times_{U'} U_j'}\right)
=
f_{ijk}^* \left(c_j^*s|_{U_i' \times_{U'} U_j'}\right)
$$
Par conséquent $c_i^*(s)|_{U_i' \times_{U'} U_j'} = c_j^*(s)|_{U_i' \times_{U'} U_j'}$
par la condition de faisceau pour $\mathcal{F}$, et l’existence de
$s_{(U', \psi)}$ en résulte également par la condition de faisceau pour $\mathcal{F}$. L’unicité
garantit que la famille $(s_{(U', \psi)})$ ainsi obtenue est un élément
de la limite avec $s_{(U, \psi)} = s$. Cela démontre
que $g^{-1}g_*\mathcal{F} \to \mathcal{F}$ est un isomorphisme.

\medskip\noindent
Soit $\mathcal{G}$ un faisceau sur $\mathcal{D}$. Soit $V$ un
objet de $\mathcal{D}$. Alors nous voyons que
$$
g_*g^{-1}\mathcal{G}(V) =
\lim_{U, \psi : u(U) \to V} \mathcal{G}(u(U))
$$
D’après le paragraphe précédent, nous voyons que la valeur du faisceau
$g_*g^{-1}\mathcal{G}$ sur un objet $V$ de la forme $V = u(U)$
est égale à $\mathcal{G}(u(U))$. (Formellement, cela découle de
$g^{-1}g_*g^{-1} \cong g^{-1}$ et de la description
de $g^{-1}$ donnée au début de la démonstration ; informellement, il suffit de
comparer les limites ici et ci-dessus.)
Ainsi, le morphisme d’adjonction $\mathcal{G} \to g_*g^{-1}\mathcal{G}$
a la propriété d’être une bijection sur les sections au-dessus de tout objet de
la forme $u(U)$. Comme, par l’axiome (5), il
existe un recouvrement de $V$ par des objets de la forme $u(U)$, nous voyons
facilement que le morphisme d’adjonction est un isomorphisme.
\end{proof}

\noindent
Il sera commode de donner un nom aux foncteurs cocontinus comme dans le
Lemme \ref{lemma-equivalence}.

\begin{definition}
\label{definition-special-cocontinuous-functor}
Soient $\mathcal{C}$, $\mathcal{D}$ des sites.
Un {\it foncteur cocontinu spécial $u$ de $\mathcal{C}$ vers $\mathcal{D}$}
est un foncteur cocontinu $u : \mathcal{C} \to \mathcal{D}$ satisfaisant
aux hypothèses et aux conclusions du Lemme \ref{lemma-equivalence}.
\end{definition}

\begin{lemma}
\label{lemma-localize-special-cocontinuous}
Soient $\mathcal{C}$, $\mathcal{D}$ des sites.
Soit $u : \mathcal{C} \to \mathcal{D}$ un foncteur cocontinu spécial.
Pour tout objet $U$ de $\mathcal{C}$, nous avons un diagramme commutatif
$$
\xymatrix{
\mathcal{C}/U \ar[r]_{j_U} \ar[d] & \mathcal{C} \ar[d]^u \\
\mathcal{D}/u(U) \ar[r]^-{j_{u(U)}} & \mathcal{D}
}
$$
comme dans le Lemme \ref{lemma-localize-cocontinuous}.
La flèche verticale de gauche est un foncteur cocontinu spécial.
Ainsi, dans le diagramme commutatif de topos
$$
\xymatrix{
\Sh(\mathcal{C}/U) \ar[r]_{j_U} \ar[d] &
\Sh(\mathcal{C}) \ar[d]^u \\
\Sh(\mathcal{D}/u(U)) \ar[r]^-{j_{u(U)}} &
\Sh(\mathcal{D})
}
$$
les flèches verticales sont des équivalences.
\end{lemma}

\begin{proof}
Nous avons établi l’existence et la commutativité des diagrammes dans le
Lemme \ref{lemma-localize-cocontinuous}. Il reste à vérifier
les hypothèses (1) -- (5) du Lemme \ref{lemma-equivalence} pour le
foncteur induit $u : \mathcal{C}/U \to \mathcal{D}/u(U)$.
Cette vérification est entièrement formelle.

\medskip\noindent
Propriété (1). C’est le Lemme \ref{lemma-localize-cocontinuous}.

\medskip\noindent
Propriété (2). Soit $\{U_i'/U \to U'/U\}_{i \in I}$ un recouvrement
de $U'/U$ dans $\mathcal{C}/U$. Comme $u$ est continu, nous voyons que
$\{u(U_i')/u(U) \to u(U')/u(U)\}_{i \in I}$ est un recouvrement
de $u(U')/u(U)$ dans $\mathcal{D}/u(U)$. Ainsi, (2) est vérifiée
pour $u : \mathcal{C}/U \to \mathcal{D}/u(U)$.

\medskip\noindent
Propriété (3). Soient $a, b : U''/U \to U'/U$ dans $\mathcal{C}/U$
des morphismes tels que $u(a) = u(b)$ dans $\mathcal{D}/u(U)$.
Comme $u$ satisfait (3), nous voyons qu’il existe un recouvrement
$\{f_i : U''_i \to U''\}$ dans $\mathcal{C}$ tel que
$a \circ f_i = b \circ f_i$. Cela donne un recouvrement
$\{f_i : U''_i/U \to U''/U\}$ dans $\mathcal{C}/U$ tel que
$a \circ f_i = b \circ f_i$. Ainsi, (3) est vérifiée
pour $u : \mathcal{C}/U \to \mathcal{D}/u(U)$.

\medskip\noindent
Propriété (4). Soient $U''/U, U'/U \in \Ob(\mathcal{C}/U)$ et
soit donné un morphisme $c : u(U'')/u(U) \to u(U')/u(U)$ dans $\mathcal{D}/u(U)$.
Comme $u$ satisfait la propriété (4), il existe
un recouvrement $\{f_i : U_i'' \to U''\}$ dans $\mathcal{C}$
et des morphismes $c_i : U_i'' \to U'$ tels que $u(c_i) = c \circ u(f_i)$.
Nous considérons $U_i''$ comme un objet au-dessus de $U$ par la composée
$U_i'' \to U'' \to U$.
Il n’est pas nécessairement vrai que $c_i$ soit un morphisme au-dessus de $U$ !
Mais puisque $u(c_i)$ est un morphisme au-dessus de $u(U)$, nous pouvons appliquer
la propriété (3) à $u$ et trouver des recouvrements $\{f_{ik} : U''_{ik} \to U''_i\}$
tels que $c_{ik} = c_i \circ f_{ik} : U''_{ik} \to U'$ soient des morphismes au-dessus de $U$.
Ainsi, $\{f_i \circ f_{ik} : U''_{ik}/U \to U''/U\}$ est un recouvrement
dans $\mathcal{C}/U$ tel que $u(c_{ik}) = c \circ u(f_{ik})$.
Ainsi, (4) est vérifiée
pour $u : \mathcal{C}/U \to \mathcal{D}/u(U)$.

\medskip\noindent
Propriété (5). Soit $h : V \to u(U)$ un objet de $\mathcal{D}/u(U)$.
Comme $u$ satisfait la propriété (5), il existe un recouvrement
$\{c_i : u(U_i) \to V\}$ dans $\mathcal{D}$. Par la propriété (4), nous pouvons trouver
des recouvrements $\{f_{ij} : U_{ij} \to U_i\}$ et des morphismes
$c_{ij} : U_{ij} \to U$ tels que $u(c_{ij}) = h \circ c_i \circ u(f_{ij})$.
Ainsi, $\{u(U_{ij})/u(U) \to V/u(U)\}$ est un recouvrement dans
$\mathcal{D}/u(U)$ de la forme voulue, et nous concluons que
(5) est vérifiée pour $u : \mathcal{C}/U \to \mathcal{D}/u(U)$.
\end{proof}

\begin{lemma}
\label{lemma-special-equivalence}
Soit $\mathcal{C}$ un site. Soit
$\mathcal{C}' \subset \Sh(\mathcal{C})$
une sous-catégorie pleine (possédant un ensemble d’objets) telle que
\begin{enumerate}
\item $h_U^\# \in \Ob(\mathcal{C}')$ pour tout
$U \in \Ob(\mathcal{C})$, et
\item $\mathcal{C}'$ soit stable par produits fibrés dans
$\Sh(\mathcal{C})$.
\end{enumerate}
Déclarons qu’un recouvrement de $\mathcal{C}'$ est toute famille
$\{\mathcal{F}_i \to \mathcal{F}\}_{i \in I}$ de morphismes telle que
$\coprod_{i \in I} \mathcal{F}_i \to \mathcal{F}$ soit un morphisme surjectif
de faisceaux. Alors
\begin{enumerate}
\item $\mathcal{C}'$ est un site (après avoir
choisi un ensemble de recouvrements, voir Ensembles, Lemme \ref{sets-lemma-coverings-site}),
\item les préfaisceaux représentables sur $\mathcal{C}'$ sont des faisceaux
(c’est-à-dire que la topologie sur $\mathcal{C}'$ est sous-canonique, voir
Définition \ref{definition-weaker-than-canonical}),
\item le foncteur $v : \mathcal{C} \to \mathcal{C}'$,
$U \mapsto h_U^\#$ est un foncteur cocontinu spécial, et induit donc une
équivalence $g : \Sh(\mathcal{C}) \to \Sh(\mathcal{C}')$,
\item pour tout $\mathcal{F} \in \Ob(\mathcal{C}')$, nous avons
$g^{-1}h_\mathcal{F} = \mathcal{F}$, et
\item pour tout $U \in \Ob(\mathcal{C})$, nous avons
$g_*h_U^\# = h_{v(U)} = h_{h_U^\#}$.
\end{enumerate}
\end{lemma}

\begin{proof}
Avertissement : certains des énoncés ci-dessus peuvent sembler un peu déroutants au premier abord ;
cela vient de ce que les objets de $\mathcal{C}'$ peuvent aussi être considérés comme des faisceaux sur
$\mathcal{C}$ ! Nous omettons de vérifier que les recouvrements de $\mathcal{C}'$
décrits dans le lemme satisfont les conditions de la
Définition \ref{definition-site}.

\medskip\noindent
Supposons que $\{\mathcal{F}_i \to \mathcal{F}\}$ soit une famille surjective
de morphismes de faisceaux. Soit $\mathcal{G}$ un autre faisceau.
La partie (2) du lemme affirme que l’égalisateur de
$$
\xymatrix{
\Mor_{\Sh(\mathcal{C})}(
\coprod_{i \in I} \mathcal{F}_i, \mathcal{G})
\ar@<1ex>[r] \ar@<-1ex>[r]
&
\Mor_{\Sh(\mathcal{C})}(
\coprod_{(i_0, i_1) \in I \times I}
\mathcal{F}_{i_0} \times_\mathcal{F} \mathcal{F}_{i_1}, \mathcal{G})
}
$$
est $\Mor_{\Sh(\mathcal{C})}(\mathcal{F}, \mathcal{G}).$
C’est clair (utiliser par exemple le Lemme \ref{lemma-coequalizer-surjection}).

\medskip\noindent
Pour démontrer (3), nous
devons vérifier les conditions (1) -- (5) du Lemme \ref{lemma-equivalence}.
Le fait que $v$ soit cocontinu est
équivalent à la description des morphismes surjectifs de faisceaux dans le
Lemme \ref{lemma-mono-epi-sheaves}.
Le foncteur $v$ est continu parce que
$U \mapsto h_U^\#$ commute aux produits fibrés,
et transforme les recouvrements en recouvrements (voir le
Lemme \ref{lemma-sheafification-exact}, et le
Lemme \ref{lemma-covering-surjective-after-sheafification}).
Les propriétés (3), (4) du Lemme \ref{lemma-equivalence}
sont des énoncés portant sur les morphismes $f : h_{U'}^\# \to h_U^\#$.
Un tel morphisme revient à un élément de $h_U^\#(U')$.
Ainsi, (3) et (4) découlent immédiatement de la construction du faisceau associé.
La propriété (5) du Lemme \ref{lemma-equivalence} est le
Lemme \ref{lemma-sheaf-coequalizer-representable}.
Soit $g : \Sh(\mathcal{C}) \to \Sh(\mathcal{C}')$
l’équivalence de topos associée à $v$ par le Lemme \ref{lemma-equivalence}.

\medskip\noindent
Soit $\mathcal{F}$ comme dans la partie (4) du lemme.
Pour tout $U \in \Ob(\mathcal{C})$, nous avons
$$
g^{-1}h_\mathcal{F}(U) = h_\mathcal{F}(v(U))
= \Mor_{\Sh(\mathcal{C})}(h_U^\#, \mathcal{F})
= \mathcal{F}(U)
$$
La première égalité résulte du
Lemme \ref{lemma-when-shriek}.
La partie (4) est donc démontrée.

\medskip\noindent
Soit $\mathcal{F} \in \Ob(\mathcal{C}')$.
Soit $U \in \Ob(\mathcal{C})$.
Alors
\begin{align*}
g_*h_U^\#(\mathcal{F})
& =
\Mor_{\Sh(\mathcal{C}')}(h_\mathcal{F}, g_*h_U^\#) \\
& =
\Mor_{\Sh(\mathcal{C})}(g^{-1}h_\mathcal{F}, h_U^\#) \\
& =
\Mor_{\Sh(\mathcal{C})}(\mathcal{F}, h_U^\#) \\
& =
\Mor_{\mathcal{C}'}(\mathcal{F}, h_U^\#)
\end{align*}
comme voulu (où la troisième égalité a été démontrée ci-dessus).
\end{proof}

\noindent
Cela nous permet de réaliser tout topos sur un site possédant
toutes les limites finies.

\begin{lemma}
\label{lemma-topos-good-site}
Soit $\Sh(\mathcal{C})$ un topos. Soit $\{\mathcal{F}_i\}_{i \in I}$
un ensemble de faisceaux sur $\mathcal{C}$. Il existe une équivalence de topos
$g : \Sh(\mathcal{C}) \to \Sh(\mathcal{C}')$ induite par un foncteur
cocontinu spécial $u : \mathcal{C} \to \mathcal{C}'$ de sites
telle que
\begin{enumerate}
\item $\mathcal{C}'$ possède une topologie sous-canonique,
\item une famille $\{V_j \to V\}$ de morphismes de $\mathcal{C}'$
est (combinatoirement équivalente à) un recouvrement de $\mathcal{C}'$
si et seulement si $\coprod h_{V_j} \to h_V$ est surjectif,
\item $\mathcal{C}'$ possède des produits fibrés et un objet final
(c’est-à-dire que $\mathcal{C}'$ possède toutes les limites finies),
\item tout sous-faisceau d’un faisceau représentable sur $\mathcal{C}'$
est représentable, et
\item chaque $g_*\mathcal{F}_i$ est un faisceau représentable.
\end{enumerate}
\end{lemma}

\begin{proof}
Considérons la sous-catégorie pleine
$\mathcal{C}_1 \subset \Sh(\mathcal{C})$ formée de tous les
$h_U^\#$ pour tout $U \in \Ob(\mathcal{C})$, des faisceaux donnés
$\mathcal{F}_i$ et du faisceau final $*$ (voir
Exemple \ref{example-singleton-sheaf}). Nous allons définir par récurrence
des sous-catégories pleines
$$
\mathcal{C}_1 \subset \mathcal{C}_2 \subset \mathcal{C}_2 \subset \ldots
\subset \Sh(\mathcal{C})
$$
Plus précisément, $\mathcal{C}_n$ étant donnée, soit $\mathcal{C}_{n + 1}$ la sous-catégorie pleine
formée de tous les produits fibrés et sous-faisceaux d’objets
de $\mathcal{C}_n$. (Remarquons que $\mathcal{C}_{n + 1}$ possède un ensemble
d’objets.) Posons
$\mathcal{C}' = \bigcup_{n \geq 1} \mathcal{C}_n$.
Un recouvrement de $\mathcal{C}'$ est toute famille
$\{\mathcal{G}_j \to \mathcal{G}\}_{j \in J}$ de morphismes d’objets
de $\mathcal{C}'$ telle que
$\coprod \mathcal{G}_j \to \mathcal{G}$ soit surjectif
comme morphisme de faisceaux sur $\mathcal{C}$.
Le foncteur $v : \mathcal{C} \to \mathcal{C'}$ est donné par
$U \mapsto h_U^\#$. Appliquons le
Lemme \ref{lemma-special-equivalence}.
\end{proof}

\noindent
Voici le but de la présente section.

\begin{lemma}
\label{lemma-morphism-topoi-comes-from-morphism-sites}
\begin{reference}
Cet énoncé est étroitement lié à
\cite[Proposition 4.9.4. Expos\'e IV]{SGA4}.
Pour obtenir le résultat complet, il faut aussi utiliser
\cite[Remarque 4.7.4, Expos\'e IV]{SGA4}.
\end{reference}
Soient $\mathcal{C}$, $\mathcal{D}$ des sites.
Soit $f : \Sh(\mathcal{C}) \to \Sh(\mathcal{D})$ un
morphisme de topos.
Alors il existe un site $\mathcal{C}'$ et un diagramme de foncteurs
$$
\xymatrix{
\mathcal{C} \ar[r]_v & \mathcal{C}' & \mathcal{D} \ar[l]^u
}
$$
tels que
\begin{enumerate}
\item le foncteur $v$ soit un foncteur cocontinu spécial,
\item le foncteur $u$ commute aux produits fibrés, soit
continu et définisse un morphisme de sites
$\mathcal{C}' \to \mathcal{D}$, et
\item le morphisme de topos $f$ coïncide avec le composé
des morphismes de topos
$$
\Sh(\mathcal{C}) \longrightarrow
\Sh(\mathcal{C}') \longrightarrow
\Sh(\mathcal{D})
$$
où la première flèche provient de $v$ par le Lemme \ref{lemma-equivalence}
et la seconde de $u$ par le Lemme \ref{lemma-morphism-sites-topoi}.
\end{enumerate}
\end{lemma}

\begin{proof}
Considérons la sous-catégorie pleine
$\mathcal{C}_1 \subset \Sh(\mathcal{C})$ formée de tous les
$h_U^\#$ et de tous les $f^{-1}h_V^\#$ pour tout
$U \in \Ob(\mathcal{C})$ et tout $V \in \Ob(\mathcal{D})$.
Soit $\mathcal{C}_{n + 1}$ la sous-catégorie pleine formée de tous les
produits fibrés d’objets de $\mathcal{C}_n$. Posons
$\mathcal{C}' = \bigcup_{n \geq 1} \mathcal{C}_n$.
Un recouvrement de $\mathcal{C}'$ est toute famille
$\{\mathcal{F}_i \to \mathcal{F}\}_{i \in I}$ telle que
$\coprod_{i \in I} \mathcal{F}_i \to \mathcal{F}$ soit surjectif
comme morphisme de faisceaux sur $\mathcal{C}$.
Le foncteur $v : \mathcal{C} \to \mathcal{C'}$ est donné par
$U \mapsto h_U^\#$.
Le foncteur $u : \mathcal{D} \to \mathcal{C'}$ est donné par
$V \mapsto f^{-1}h_V^\#$.

\medskip\noindent
La partie (1) résulte du Lemme \ref{lemma-special-equivalence}.

\medskip\noindent
Démonstration de (2) et (3) du lemme. Le foncteur $u$ commute aux produits
fibrés, puisque $V \mapsto h_V^\#$ et $f^{-1}$ le font tous deux. En outre,
comme $f^{-1}$ est exact et commute aux limites inductives arbitraires, nous voyons
qu’il transforme un recouvrement en une famille surjective de morphismes de
faisceaux. Ainsi, $u$ est continu. Pour voir qu’il définit un morphisme de
sites, il reste à vérifier que $u_s$ est exact. À cette fin,
nous allons montrer que $g^{-1} \circ u_s = f^{-1}$. En effet, puisque $g^{-1}$
est une équivalence et que $f^{-1}$ est exact, nous pourrons conclure.
Comme $g^{-1}$ est adjoint à $g_*$, que $u_s$ est adjoint à
$u^s$, et que $f^{-1}$ est adjoint à $f_*$, il suffit aussi de démontrer que
$u^s \circ g_* = f_*$. Soit $\mathcal{F}$ un faisceau sur $\mathcal{C}$
et soit $V$ un objet de $\mathcal{D}$. Alors
\begin{align*}
(u^s g_{\ast}\mathcal{F})(V)
& =
(g_{\ast}\mathcal{F})(f^{-1}h_V^\#) \\
& =
\Mor_{Sh(\mathcal{C}')}(h_{f^{-1}h_V^\#},g_{\ast}\mathcal{F}) \\
& =
\Mor_{Sh(\mathcal{C})}(g^{-1}h_{f^{-1}h_V^\#},\mathcal{F}) \\
& =
\Mor_{Sh(\mathcal{C})}(f^{-1}h_V^\#,\mathcal{F}) \\
& =
\Mor_{Sh(\mathcal{D})}(h_V^\#,f_{\ast}\mathcal{F}) \\
& =
f_{\ast}\mathcal{F}(V)
\end{align*}
La première égalité vient de ce que $u^s = u^p$.
La deuxième égalité est le lemme de Yoneda.
La troisième égalité résulte de l’adjonction entre $g^{-1}$ et $g_*$.
La quatrième égalité résulte du Lemme \ref{lemma-special-equivalence} (4).
La cinquième égalité résulte de l’adjonction entre $f^{-1}$ et $f_*$.
La sixième égalité est le lemme de Yoneda.
Ainsi, $u^s g_*\mathcal{F} = f_*\mathcal{F}$ et
cela achève la démonstration du lemme.
\end{proof}

\begin{remark}
\label{remark-morphism-topoi-comes-from-morphism-sites}
Notations et hypothèses
comme dans le Lemme \ref{lemma-morphism-topoi-comes-from-morphism-sites}.
Si le site $\mathcal{D}$ possède un objet final et des produits fibrés,
alors le foncteur $u : \mathcal{D} \to \mathcal{C}'$ satisfait
toutes les hypothèses de la Proposition \ref{proposition-get-morphism}.
En effet, outre les propriétés mentionnées dans le lemme, $u$
transforme aussi l’objet final de $\mathcal{D}$ en l’objet final
de $\mathcal{C}'$. Cela ressort clairement de la construction de $u$.
Ainsi, si nous appliquons d’abord le
Lemme \ref{lemma-topos-good-site}
à $\mathcal{D}$
puis le
Lemme \ref{lemma-morphism-topoi-comes-from-morphism-sites}
au morphisme de topos qui en résulte
$\Sh(\mathcal{C}) \to \Sh(\mathcal{D}')$,
nous obtenons l’énoncé suivant :
Tout morphisme de topos
$f : \Sh(\mathcal{C}) \to \Sh(\mathcal{D})$
s’insère dans un diagramme commutatif
$$
\xymatrix{
\Sh(\mathcal{C}) \ar[d]_g \ar[r]_f &
\Sh(\mathcal{D}) \ar[d]^e \\
\Sh(\mathcal{C}') \ar[r]^{f'} &
\Sh(\mathcal{D}')
}
$$
où les propriétés suivantes sont satisfaites :
\begin{enumerate}
\item les morphismes $e$ et $g$ sont des équivalences données par des
foncteurs cocontinus spéciaux $\mathcal{C} \to \mathcal{C}'$ et
$\mathcal{D} \to \mathcal{D}'$,
\item les sites $\mathcal{C}'$ et $\mathcal{D}'$ possèdent des produits fibrés et un objet
final et sont munis de topologies sous-canoniques,
\item le morphisme $f' : \mathcal{C}' \to \mathcal{D}'$ provient d’un
morphisme de sites correspondant à un foncteur
$u : \mathcal{D}' \to \mathcal{C}'$ auquel la
Proposition \ref{proposition-get-morphism}
s’applique, et
\item étant donné un ensemble quelconque de faisceaux $\mathcal{F}_i$ (resp.\ $\mathcal{G}_j$)
sur $\mathcal{C}$ (resp.\ $\mathcal{D}$), nous pouvons supposer que chacun d’eux est
un faisceau représentable sur $\mathcal{C}'$ (resp.\ $\mathcal{D}'$).
\end{enumerate}
Il est souvent utile de remplacer $\mathcal{C}$ et $\mathcal{D}$ par
$\mathcal{C}'$ et $\mathcal{D}'$.
\end{remark}

\begin{remark}
\label{remark-equivalence-topoi-comes-from-morphism-sites}
Notations et hypothèses
comme dans le Lemme \ref{lemma-morphism-topoi-comes-from-morphism-sites}.
Supposons en outre que le morphisme de topos d’origine
$\Sh(\mathcal{C}) \to \Sh(\mathcal{D})$ soit une équivalence.
Alors la construction de la démonstration du
Lemme \ref{lemma-morphism-topoi-comes-from-morphism-sites}
donne deux foncteurs
$$
\mathcal{C} \rightarrow \mathcal{C}' \leftarrow \mathcal{D}
$$
qui sont tous deux des foncteurs cocontinus spéciaux.
Ainsi, dans ce cas, nous pouvons effectivement
factoriser le morphisme de topos comme un composé
$$
\Sh(\mathcal{C}) \rightarrow
\Sh(\mathcal{C}') =
\Sh(\mathcal{D}') \leftarrow
\Sh(\mathcal{D})
$$
comme dans la Remarque \ref{remark-morphism-topoi-comes-from-morphism-sites}, mais
avec le morphisme du milieu égal à l’identité.
\end{remark}








\section{Localisation des topos}
\label{section-localize-topoi}

\noindent
Nous reprenons certains résultats sur la localisation dans le cadre, en
apparence plus général, des topos. En réalité, ce cadre n’est pas plus général,
car nous pouvons toujours agrandir les sites sous-jacents pour supposer que
nous localisons en des objets du site.

\begin{lemma}
\label{lemma-localize-topos}
Soit $\mathcal{C}$ un site.
Soit $\mathcal{F}$ un faisceau sur $\mathcal{C}$.
Alors la catégorie $\Sh(\mathcal{C})/\mathcal{F}$ est un topos. Il existe un
morphisme canonique de topos
$$
j_\mathcal{F} :
\Sh(\mathcal{C})/\mathcal{F}
\longrightarrow
\Sh(\mathcal{C})
$$
qui est une localisation au sens de la section \ref{section-localize} et tel
que
\begin{enumerate}
\item le foncteur $j_\mathcal{F}^{-1}$ soit le foncteur
$\mathcal{H} \mapsto \mathcal{H} \times \mathcal{F}/\mathcal{F}$, et
\item le foncteur $j_{\mathcal{F}!}$ soit le foncteur d’oubli
$\mathcal{G}/\mathcal{F} \mapsto \mathcal{G}$.
\end{enumerate}
\end{lemma}

\begin{proof}
Appliquons le Lemme \ref{lemma-topos-good-site}.
Cela signifie que nous pouvons supposer que $\mathcal{C}$ est un site muni
d’une topologie sous-canonique et que $\mathcal{F} = h_U = h_U^\#$ pour un
certain $U \in \Ob(\mathcal{C})$.
Les résultats de la section \ref{section-localize} s’appliquent donc. En
particulier, il existe une équivalence
$\Sh(\mathcal{C}/U) = \Sh(\mathcal{C})/h_U^\#$ telle que le composé
$$
\Sh(\mathcal{C}/U)
\to
\Sh(\mathcal{C})/h_U^\#
\to \Sh(\mathcal{C})
$$
soit égal à $j_{U!}$ ; voir le
Lemme \ref{lemma-essential-image-j-shriek}.
Notons $a : \Sh(\mathcal{C})/h_U^\# \to \Sh(\mathcal{C}/U)$ le foncteur
inverse, de sorte que $j_{\mathcal{F}!} = j_{U!} \circ a$,
$j_\mathcal{F}^{-1} = a^{-1} \circ j_U^{-1}$ et
$j_{\mathcal{F}, *} = j_{U, *} \circ a$. La description de
$j_{\mathcal{F}!}$ résulte de ce qui précède. Celle de
$j_\mathcal{F}^{-1}$ résulte du
Lemme \ref{lemma-compute-j-shriek-restrict}.
\end{proof}

\begin{lemma}
\label{lemma-localize-pushforward}
Dans la situation du Lemme \ref{lemma-localize-topos}, le foncteur
$j_{\mathcal{F}, *}$ associe à
$\varphi : \mathcal{G} \to \mathcal{F}$ le faisceau
$$
U
\longmapsto
\{\alpha : \mathcal{F}|_U \to \mathcal{G}|_U
\text{ telle que } \alpha \text{ est un inverse à droite de }\varphi|_U \}.
$$
\end{lemma}

\begin{proof}
Pour tout $\varphi : \mathcal{G} \to \mathcal{F}$, convenons de noter
$\mathcal{G}/\mathcal{F}$ l’objet correspondant de
$\Sh(\mathcal{C})/\mathcal{F}$. On a
$$
(j_{\mathcal{F}, *}(\mathcal{G}/\mathcal{F}))(U) =
\Mor_{\Sh(\mathcal{C})}(h_U^\#, j_{\mathcal{F}, *}(\mathcal{G}/\mathcal{F})) =
\Mor_{\Sh(\mathcal{C})/\mathcal{F}}(j_{\mathcal{F}}^{-1}h_U^{\#},
(\mathcal{G}/\mathcal{F})).
$$
D’après le Lemme \ref{lemma-localize-topos}, cet ensemble est la fibre
au-dessus de l’élément $h_U^\# \times \mathcal{F} \to \mathcal{F}$ pour
l’application d’ensembles
$$
\Mor_{\Sh(\mathcal{C})}(h_U^\# \times \mathcal{F}, \mathcal{G})
\xrightarrow{\varphi \circ}
\Mor_{\Sh(\mathcal{C})}(h_U^\# \times \mathcal{F}, \mathcal{F}).
$$
Par l’adjonction du Lemme \ref{lemma-internal-hom-sheaf}, on a
\begin{align*}
\Mor_{\Sh(\mathcal{C})}(h_U^{\#}\times\mathcal{F}, \mathcal{G})
& =
\Mor_{\Sh(\mathcal{C})}(h_U^{\#},\SheafHom(\mathcal{F}, \mathcal{G})) \\
& =
\Mor_{\Sh(\mathcal{C}/U)}(\mathcal{F}|_{\mathcal{C}/U},
\mathcal{G}|_{\mathcal{C}/U}), \\
\Mor_{\Sh(\mathcal{C})}(h_U^{\#} \times \mathcal{F}, \mathcal{F})
& =
\Mor_{\Sh(\mathcal{C})}(h_U^{\#},\SheafHom(\mathcal{F},\mathcal{F})) \\
& =
\Mor_{\Sh(\mathcal{C}/U)}(\mathcal{F}|_{\mathcal{C}/U},
\mathcal{F}|_{\mathcal{C}/U}),
\end{align*}
et, sous cette adjonction, l’application $\varphi\circ$ devient
$$
\Mor_{\Sh(\mathcal{C}/U)}(\mathcal{F}|_{\mathcal{C}/U},
\mathcal{G}|_{\mathcal{C}/U})
\longrightarrow
\Mor_{\Sh(\mathcal{C}/U)}(\mathcal{F}|_{\mathcal{C}/U},
\mathcal{F}|_{\mathcal{C}/U}),\quad
\psi \longmapsto \varphi|_{\mathcal{C}/U} \circ \psi,
$$
tandis que l’élément $h_U^\# \times \mathcal{F} \to \mathcal{F}$ devient
$\text{id}_{\mathcal{F}|_{\mathcal{C}/U}}$.
Par conséquent, $(j_{\mathcal{F}, *}\mathcal{G}/\mathcal{F})(U)$ est
isomorphe à la fibre au-dessus de $\text{id}_{\mathcal{F}|_{\mathcal{C}/U}}$
pour l’application
$$
\Mor_{\Sh(\mathcal{C}/U)}(\mathcal{F}|_{\mathcal{C}/U},
\mathcal{G}|_{\mathcal{C}/U})
\xrightarrow{\varphi|_{\mathcal{C}/U}\circ}
\Mor_{\Sh(\mathcal{C}/U)}(\mathcal{F}|_{\mathcal{C}/U},
\mathcal{F}|_{\mathcal{C}/U}),
$$
laquelle est $\{\alpha : \mathcal{F}|_U \to \mathcal{G}|_U
\text{ telle que } \alpha \text{ est un inverse à droite de }\varphi|_U \}$,
comme voulu.
\end{proof}

\begin{lemma}
\label{lemma-localize-topos-site}
Soit $\mathcal{C}$ un site. Soit $\mathcal{F}$ un faisceau sur $\mathcal{C}$.
Soit $\mathcal{C}/\mathcal{F}$ la catégorie des couples $(U, s)$, où
$U \in \Ob(\mathcal{C})$ et $s \in \mathcal{F}(U)$. Déclarons qu’un
recouvrement de $\mathcal{C}/\mathcal{F}$ est une famille
$\{(U_i, s_i) \to (U, s)\}$ telle que $\{U_i \to U\}$ soit un recouvrement de
$\mathcal{C}$. Alors $j : \mathcal{C}/\mathcal{F} \to \mathcal{C}$ est un
foncteur de sites continu et cocontinu, qui induit un morphisme de topos
$j : \Sh(\mathcal{C}/\mathcal{F}) \to \Sh(\mathcal{C})$. En fait, il existe
une équivalence $\Sh(\mathcal{C}/\mathcal{F}) =
\Sh(\mathcal{C})/\mathcal{F}$ qui transforme $j$ en $j_\mathcal{F}$.
\end{lemma}

\begin{proof}
Nous omettons de vérifier que $\mathcal{C}/\mathcal{F}$ est un site et que
$j$ est continu et cocontinu. D’après le Lemme \ref{lemma-when-shriek}, il
existe un morphisme de topos $j$ comme indiqué, avec
$j^{-1}\mathcal{G}(U, s) = \mathcal{G}(U)$, et il existe un adjoint à gauche
$j_!$ de $j^{-1}$. Un morphisme $\varphi : * \to j^{-1}\mathcal{G}$ sur
$\mathcal{C}/\mathcal{F}$ revient à se donner une règle qui associe à chaque
couple $(U, s)$ une section $\varphi(s) \in \mathcal{G}(U)$ compatible avec les
applications de restriction. Cela revient donc à se donner un morphisme
$\varphi : \mathcal{F} \to \mathcal{G}$ sur $\mathcal{C}$.
On en conclut que $j_!* = \mathcal{F}$. En particulier, pour tout
$\mathcal{H} \in \Sh(\mathcal{C}/\mathcal{F})$, il existe un morphisme
canonique
$$
j_!\mathcal{H} \to j_!* = \mathcal{F}
$$
c’est-à-dire que l’on obtient un foncteur
$j'_! : \Sh(\mathcal{C}/\mathcal{F}) \to
\Sh(\mathcal{C})/\mathcal{F}$.
Un inverse de ce foncteur est la règle qui associe à un objet
$\varphi : \mathcal{G} \to \mathcal{F}$ de
$\Sh(\mathcal{C})/\mathcal{F}$ le faisceau
$$
a(\mathcal{G}/\mathcal{F}) :
(U, s) \longmapsto \{t \in \mathcal{G}(U) \mid \varphi(t) = s\}
$$
Nous omettons de vérifier que $a(\mathcal{G}/\mathcal{F})$ est un faisceau et
que $a$ est inverse de $j'_!$.
\end{proof}

\begin{definition}
\label{definition-localize-topos}
Soit $\mathcal{C}$ un site.
Soit $\mathcal{F}$ un faisceau sur $\mathcal{C}$.
\begin{enumerate}
\item Le topos $\Sh(\mathcal{C})/\mathcal{F}$ est appelé la
{\it localisation du topos $\Sh(\mathcal{C})$ en $\mathcal{F}$}.
\item Le morphisme de topos
$j_\mathcal{F} :
\Sh(\mathcal{C})/\mathcal{F}
\to
\Sh(\mathcal{C})$ du
Lemme \ref{lemma-localize-topos}
est appelé le {\it morphisme de localisation}.
\end{enumerate}
\end{definition}

\noindent
Nous allons montrer que, chaque fois que le faisceau $\mathcal{F}$ est égal à
$h_U^\#$ pour un certain objet $U$ du site, la localisation du topos est égale
à la catégorie des faisceaux sur la localisation du site en $U$. Nous allons
en outre vérifier que toutes les fonctorialités sont compatibles avec cette
identification.

\begin{lemma}
\label{lemma-localize-compare}
Soit $\mathcal{C}$ un site. Soit $\mathcal{F} = h_U^\#$ pour un certain objet
$U$ de $\mathcal{C}$. Alors le morphisme
$j_\mathcal{F} : \Sh(\mathcal{C})/\mathcal{F}
\to \Sh(\mathcal{C})$ construit dans le
Lemme \ref{lemma-localize-topos}
coïncide avec le morphisme de topos
$j_U : \Sh(\mathcal{C}/U) \to \Sh(\mathcal{C})$
construit dans la section \ref{section-localize} par l’identification
$\Sh(\mathcal{C}/U) = \Sh(\mathcal{C})/h_U^\#$ du
Lemme \ref{lemma-essential-image-j-shriek}.
\end{lemma}

\begin{proof}
Nous avons vu dans le Lemme \ref{lemma-essential-image-j-shriek} que le
composé
$\Sh(\mathcal{C}/U) \to \Sh(\mathcal{C})/h_U^\#
\to \Sh(\mathcal{C})$
est $j_{U!}$. Le foncteur
$\Sh(\mathcal{C})/h_U^\# \to \Sh(\mathcal{C})$
est $j_{\mathcal{F}!}$ d’après le Lemme \ref{lemma-localize-topos}.
Ainsi, $j_{\mathcal{F}!} = j_{U!}$ par l’identification.
Donc $j_\mathcal{F}^{-1} = j_U^{-1}$ (par adjonction), puis
$j_{\mathcal{F}, *} = j_{U, *}$ (encore par adjonction).
\end{proof}

\begin{lemma}
\label{lemma-relocalize-topos}
Soit $\mathcal{C}$ un site.
Si $s : \mathcal{G} \to \mathcal{F}$ est un morphisme de faisceaux sur
$\mathcal{C}$, alors il existe un diagramme naturel commutatif de morphismes de
topos
$$
\xymatrix{
\Sh(\mathcal{C})/\mathcal{G} \ar[rd]_{j_\mathcal{G}} \ar[rr]_j & &
\Sh(\mathcal{C})/\mathcal{F} \ar[ld]^{j_\mathcal{F}} \\
& \Sh(\mathcal{C}) &
}
$$
où $j = j_{\mathcal{G}/\mathcal{F}}$ est la localisation du topos
$\Sh(\mathcal{C})/\mathcal{F}$ en l’objet $\mathcal{G}/\mathcal{F}$. En
particulier, on a
$$
j^{-1}(\mathcal{H} \to \mathcal{F}) =
(\mathcal{H} \times_\mathcal{F} \mathcal{G} \to \mathcal{G})
$$
et
$$
j_!(\mathcal{E} \xrightarrow{e} \mathcal{G}) =
(\mathcal{E} \xrightarrow{s \circ e} \mathcal{F}).
$$
\end{lemma}

\begin{proof}
La description de $j^{-1}$ et de $j_!$ résulte de la description de ces
foncteurs dans le Lemme \ref{lemma-localize-topos}.
L’égalité de foncteurs $j_{\mathcal{G}!} = j_{\mathcal{F}!} \circ j_!$ résulte
clairement de leur description comme foncteurs d’oubli. Par adjonction, on
obtient également les égalités
$j_\mathcal{G}^{-1} = j^{-1} \circ j_\mathcal{F}^{-1}$ et
$j_{\mathcal{G}, *} = j_{\mathcal{F}, *} \circ j_*$.
\end{proof}

\begin{lemma}
\label{lemma-relocalize-compare}
Supposons que $\mathcal{C}$ et $s : \mathcal{G} \to \mathcal{F}$ soient comme
dans le Lemme \ref{lemma-relocalize-topos}.
Si $\mathcal{G} = h_V^\#$, $\mathcal{F} = h_U^\#$ et si
$s : \mathcal{G} \to \mathcal{F}$ provient d’un morphisme $V \to U$ de
$\mathcal{C}$, alors le diagramme du Lemme \ref{lemma-relocalize-topos}
s’identifie au diagramme (\ref{equation-relocalize}) par les identifications
$\Sh(\mathcal{C}/V) = \Sh(\mathcal{C})/h_V^\#$ et
$\Sh(\mathcal{C}/U) = \Sh(\mathcal{C})/h_U^\#$ du
Lemme \ref{lemma-essential-image-j-shriek}.
\end{lemma}

\begin{proof}
C’est vrai parce que les descriptions de $j^{-1}$ coïncident. Voir le
Lemme \ref{lemma-relocalize-explicit} et le
Lemme \ref{lemma-relocalize-topos}.
\end{proof}






\section{Localisation et morphismes de topos}
\label{section-localize-morphisms-topoi}

\noindent
Cette section est l’analogue de la section \ref{section-localize-morphisms}
pour les morphismes de topos.

\begin{lemma}
\label{lemma-localize-morphism-topoi}
Soit $f : \Sh(\mathcal{C}) \to \Sh(\mathcal{D})$ un morphisme de topos. Soit
$\mathcal{G}$ un faisceau sur $\mathcal{D}$. Posons
$\mathcal{F} = f^{-1}\mathcal{G}$. Il existe alors un diagramme commutatif de
topos
$$
\xymatrix{
\Sh(\mathcal{C})/\mathcal{F} \ar[r]_{j_\mathcal{F}} \ar[d]_{f'} &
\Sh(\mathcal{C}) \ar[d]^f \\
\Sh(\mathcal{D})/\mathcal{G} \ar[r]^{j_\mathcal{G}} &
\Sh(\mathcal{D}).
}
$$
Le morphisme $f'$ est caractérisé par la propriété
$$
(f')^{-1}(\mathcal{H} \xrightarrow{\varphi} \mathcal{G})
=
(f^{-1}\mathcal{H} \xrightarrow{f^{-1}\varphi} \mathcal{F})
$$
et l’on a $f'_*j_\mathcal{F}^{-1} = j_\mathcal{G}^{-1}f_*$.
\end{lemma}

\begin{proof}
Comme l’énoncé porte sur des topos et ne fait pas référence aux sites
sous-jacents, nous pouvons changer ceux-ci à volonté. D’après la discussion de
la Remarque \ref{remark-morphism-topoi-comes-from-morphism-sites}, nous pouvons
donc supposer que $f$ est donné par un foncteur continu
$u : \mathcal{D} \to \mathcal{C}$ satisfaisant aux hypothèses de la
Proposition \ref{proposition-get-morphism}, entre des sites possédant toutes
les limites finies et des topologies sous-canoniques, et tel que
$\mathcal{G} = h_V$ pour un certain objet $V$ de $\mathcal{D}$. Alors
$\mathcal{F} = f^{-1}h_V = h_{u(V)}$ d’après le
Lemme \ref{lemma-pullback-representable-sheaf}.
D’après le Lemme \ref{lemma-localize-morphism}, on obtient un diagramme
commutatif de morphismes de topos
$$
\xymatrix{
\Sh(\mathcal{C}/U) \ar[r]_{j_U} \ar[d]_{f'} &
\Sh(\mathcal{C}) \ar[d]^f \\
\Sh(\mathcal{D}/V) \ar[r]^{j_V} &
\Sh(\mathcal{D}),
}
$$
et l’on a $f'_*j_U^{-1} = j_V^{-1}f_*$. D’après le
Lemme \ref{lemma-localize-compare}, on peut identifier $j_\mathcal{F}$ à
$j_U$, et $j_\mathcal{G}$ à $j_V$. La description de $(f')^{-1}$ est donnée
dans le Lemme \ref{lemma-localize-morphism}.
\end{proof}

\begin{lemma}
\label{lemma-localize-morphism-compare}
Soit $f : \mathcal{C} \to \mathcal{D}$ un morphisme de sites donné par le
foncteur continu $u : \mathcal{D} \to \mathcal{C}$.
Soit $V$ un objet de $\mathcal{D}$. Posons $U = u(V)$.
Posons $\mathcal{G} = h_V^\#$ et
$\mathcal{F} = h_U^\# = f^{-1}h_V^\#$ (voir le
Lemme \ref{lemma-pullback-representable-sheaf}).
Alors le diagramme de morphismes de topos du
Lemme \ref{lemma-localize-morphism-topoi} coïncide avec le diagramme de
morphismes de topos du Lemme \ref{lemma-localize-morphism}, par les
identifications $j_\mathcal{F}= j_U$ et $j_\mathcal{G} = j_V$ du
Lemme \ref{lemma-localize-compare}.
\end{lemma}

\begin{proof}
Ce n’est pas tout à fait trivial, car le morphisme de sites donnant naissance
à $f$ choisi dans la démonstration du
Lemme \ref{lemma-localize-morphism-topoi} peut différer du morphisme de sites
qui nous est donné dans le lemme. Mais, dans les deux cas, le foncteur
$(f')^{-1}$ est décrit par la même règle. Les deux foncteurs coïncident donc,
ainsi que les morphismes de topos associés. Nous omettons certains détails.
\end{proof}

\begin{lemma}
\label{lemma-relocalize-morphism-topoi}
Soit $f : \Sh(\mathcal{C}) \to \Sh(\mathcal{D})$ un morphisme de topos.
Soient $\mathcal{G} \in \Sh(\mathcal{D})$,
$\mathcal{F} \in \Sh(\mathcal{C})$ et
$s : \mathcal{F} \to f^{-1}\mathcal{G}$ un morphisme de faisceaux.
Il existe un diagramme commutatif de topos
$$
\xymatrix{
\Sh(\mathcal{C})/\mathcal{F} \ar[r]_{j_\mathcal{F}} \ar[d]_{f_s} &
\Sh(\mathcal{C}) \ar[d]^f \\
\Sh(\mathcal{D})/\mathcal{G} \ar[r]^{j_\mathcal{G}} &
\Sh(\mathcal{D}).
}
$$
On a $f_s = f' \circ j_{\mathcal{F}/f^{-1}\mathcal{G}}$, où
$f' :
\Sh(\mathcal{C})/f^{-1}\mathcal{G}
\to
\Sh(\mathcal{D})/\mathcal{F}$
est comme dans le Lemme \ref{lemma-localize-morphism-topoi}, et
$j_{\mathcal{F}/f^{-1}\mathcal{G}} :
\Sh(\mathcal{C})/\mathcal{F}
\to
\Sh(\mathcal{C})/f^{-1}\mathcal{G}$
est comme dans le Lemme \ref{lemma-relocalize-topos}.
Le foncteur $(f_s)^{-1}$ est décrit par la règle
$$
(f_s)^{-1}(\mathcal{H} \xrightarrow{\varphi} \mathcal{G})
=
(f^{-1}\mathcal{H} \times_{f^{-1}\varphi, f^{-1}\mathcal{G}, s} \mathcal{F}
\rightarrow \mathcal{F}).
$$
Enfin, étant donnés des morphismes quelconques
$b : \mathcal{G}' \to \mathcal{G}$,
$a : \mathcal{F}' \to \mathcal{F}$ et
$s' : \mathcal{F}' \to f^{-1}\mathcal{G}'$ tels que
$$
\xymatrix{
\mathcal{F}' \ar[r]_-{s'} \ar[d]_a & f^{-1}\mathcal{G}' \ar[d]^{f^{-1}b} \\
\mathcal{F} \ar[r]^-s & f^{-1}\mathcal{G}
}
$$
soit commutatif, le diagramme
$$
\xymatrix{
\Sh(\mathcal{C})/\mathcal{F}'
\ar[r]_{j_{\mathcal{F}'/\mathcal{F}}} \ar[d]_{f_{s'}} &
\Sh(\mathcal{C})/\mathcal{F} \ar[d]^{f_s} \\
\Sh(\mathcal{D})/\mathcal{G}' \ar[r]^{j_{\mathcal{G}'/\mathcal{G}}} &
\Sh(\mathcal{D})/\mathcal{G}.
}
$$
est commutatif.
\end{lemma}

\begin{proof}
La commutativité du premier carré résulte de celle du diagramme du
Lemme \ref{lemma-relocalize-topos} et de celle du diagramme du
Lemme \ref{lemma-localize-morphism-topoi}. La description de $f_s^{-1}$
résulte de la combinaison de la description de $(f')^{-1}$ dans le
Lemme \ref{lemma-localize-morphism-topoi} et de la description de
$(j_{\mathcal{F}/f^{-1}\mathcal{G}})^{-1}$ dans le
Lemme \ref{lemma-relocalize-topos}. La commutativité du dernier carré résulte
alors de l’égalité
$$
f^{-1}\mathcal{H} \times_{f^{-1}\mathcal{G}, s} \mathcal{F}
\times_\mathcal{F} \mathcal{F}'
=
f^{-1}(\mathcal{H} \times_\mathcal{G} \mathcal{G}')
\times_{f^{-1}\mathcal{G}', s'} \mathcal{F}'
$$
qui est formelle.
\end{proof}

\begin{lemma}
\label{lemma-relocalize-morphism-compare}
Soit $f : \mathcal{C} \to \mathcal{D}$ un morphisme de sites donné par le
foncteur continu $u : \mathcal{D} \to \mathcal{C}$.
Soit $V$ un objet de $\mathcal{D}$. Soit $c : U \to u(V)$ un morphisme.
Posons $\mathcal{G} = h_V^\#$ et
$\mathcal{F} = h_U^\# = f^{-1}h_V^\#$.
Soit $s : \mathcal{F} \to f^{-1}\mathcal{G}$ l’application induite par $c$.
Alors le diagramme de morphismes de topos du
Lemme \ref{lemma-relocalize-morphism} coïncide avec le diagramme de morphismes
de topos du Lemme \ref{lemma-relocalize-morphism-topoi} par les identifications
$j_\mathcal{F} = j_U$ et $j_\mathcal{G} = j_V$ du
Lemme \ref{lemma-localize-compare}.
\end{lemma}

\begin{proof}
Cela résulte de la combinaison des Lemmes \ref{lemma-relocalize-compare} et
\ref{lemma-localize-morphism-compare}.
\end{proof}















\section{Points}
\label{section-points}

\begin{definition}
\label{definition-point-topos}
Soit $\mathcal{C}$ un site.
Un {\it point du topos $\Sh(\mathcal{C})$} est un morphisme de topos $p$ de
$\Sh(pt)$ vers $\Sh(\mathcal{C})$.
\end{definition}

\noindent
Nous définirons un point d’un site au moyen d’un foncteur
$u : \mathcal{C} \to \textit{Ensembles}$.
Nous verrons plus loin que $u$ définit un morphisme de sites qui donne
naissance à un point du topos associé à $\mathcal{C}$ ; voir le
Lemme \ref{lemma-site-point-morphism}.

\medskip\noindent
Soit $\mathcal{C}$ un site. Soit $p = u$ un foncteur
$u : \mathcal{C} \to \textit{Ensembles}$.
Nous introduisons cette curieuse terminologie parce qu’il paraît étrange de
parler des voisinages de foncteurs ; nous considérons donc un « point » $p$
comme un objet géométrique donné par une donnée catégorique, à savoir le
foncteur $u$. Le fait que $p$ soit effectivement égal à $u$ importe peu.
Un {\it voisinage} de $p$ est un couple $(U, x)$ avec
$U \in \Ob(\mathcal{C})$ et $x \in u(U)$.
Un {\it morphisme de voisinages} $(V, y) \to (U, x)$ est donné par un
morphisme $\alpha :V \to U$ de $\mathcal{C}$ tel que
$u(\alpha)(y) = x$. Remarquons que la catégorie des voisinages n’est pas une
catégorie « grande ».

\medskip\noindent
Nous définissons la {\it fibre} d’un préfaisceau $\mathcal{F}$ en $p$ par
\begin{equation}
\label{equation-stalk}
\mathcal{F}_p = \colim_{\{(U, x)\}^{opp}} \mathcal{F}(U).
\end{equation}
La limite inductive est prise sur l’opposée de la catégorie des
voisinages de $p$. Autrement dit, un élément de $\mathcal{F}_p$ est donné par
un triplet $(U, x, s)$, où $(U, x)$ est un voisinage de $p$ et
$s \in \mathcal{F}(U)$. L’égalité des triplets est la relation d’équivalence
engendrée par $(U, x, s) \sim (V, y, \alpha^*s)$ lorsque $\alpha$ est comme
ci-dessus.

\medskip\noindent
Remarquons que, si $\varphi : \mathcal{F} \to \mathcal{G}$ est un morphisme
de préfaisceaux d’ensembles, on obtient une application canonique entre les
fibres $\varphi_p : \mathcal{F}_p \to \mathcal{G}_p$. On obtient ainsi un
{\it foncteur fibre}
$$
\textit{PSh}(\mathcal{C}) \longrightarrow \textit{Ensembles}, \quad
\mathcal{F} \longmapsto \mathcal{F}_p.
$$
Nous avons défini le foncteur fibre à partir d’un foncteur quelconque
$p = u : \mathcal{C} \to \textit{Ensembles}$. Aucune condition n’est nécessaire
pour que la définition ait un sens\footnote{Il convient d’éviter le cas où
$u(U) = \emptyset$ pour tout $U$.}. En revanche, il vaut probablement mieux
ne pas employer cette notion à moins que $p$ ne soit effectivement un point
(voir la définition ci-dessous), car le foncteur fibre ne possède pas, en
général, de bonnes propriétés.

\begin{definition}
\label{definition-point}
Soit $\mathcal{C}$ un site. Un {\it point $p$ du site $\mathcal{C}$} est donné
par un foncteur $u : \mathcal{C} \to \textit{Ensembles}$ tel que
\begin{enumerate}
\item Pour tout recouvrement $\{U_i \to U\}$ de $\mathcal{C}$, l’application
$\coprod u(U_i) \to u(U)$ est surjective.
\item Pour tout recouvrement $\{U_i \to U\}$ de $\mathcal{C}$ et tout
morphisme $V \to U$, les applications
$u(U_i \times_U V) \to u(U_i) \times_{u(U)} u(V)$ sont bijectives.
\item Le foncteur fibre $\Sh(\mathcal{C}) \to \textit{Ensembles}$,
$\mathcal{F} \mapsto \mathcal{F}_p$, est exact à gauche.
\end{enumerate}
\end{definition}

\noindent
Ces conditions devraient être familières, car elles sont calquées sur celles
des Définitions \ref{definition-continuous} et
\ref{definition-morphism-sites}.
Remarquons que (3) implique que $*_p = \{*\}$ ; voir
l’Exemple \ref{example-singleton-sheaf}. Ainsi,
$u(U) \not = \emptyset$ pour au moins un $U$ (car une limite inductive vide
donne l’ensemble vide). Nous montrerons plus bas
(Lemme \ref{lemma-point-site-topos}) que cela donne bien naissance à un point
du topos $\Sh(\mathcal{C})$.
Avant cela, démontrons quelques lemmes pour des foncteurs $u$ généraux.

\begin{lemma}
\label{lemma-points-recover}
Soit $\mathcal{C}$ un site.
Soit $p = u : \mathcal{C} \to \textit{Ensembles}$ un foncteur.
Il existe des isomorphismes fonctoriels $(h_U)_p = u(U)$ pour
$U \in \Ob(\mathcal{C})$.
\end{lemma}

\begin{proof}
Un élément de $(h_U)_p$ est donné par un triplet $(V, y, f)$, où
$V \in \Ob(\mathcal{C})$, $y\in u(V)$ et
$f \in h_U(V) = \Mor_\mathcal{C}(V, U)$.
Deux triplets de ce type $(V, y, f)$, $(V', y', f')$ déterminent le même
élément s’il existe un morphisme $\phi : V \to V'$ tel que
$u(\phi)(y) = y'$ et $f' \circ \phi = f$ ; en général, il faut prendre des
chaînes d’identifications de ce type pour obtenir la bonne relation
d’équivalence. Quoi qu’il en soit, tout $(V, y, f)$ est équivalent à l’élément
$(U, u(f)(y), \text{id}_U)$. Si $\phi$ existe comme ci-dessus, alors les
triplets $(V, y, f)$, $(V', y', f')$ déterminent le même triplet
$(U, u(f)(y), \text{id}_U) = (U, u(f')(y'), \text{id}_U)$.
Cela montre que l’application
$u(U) \to (h_U)_p$, $x \mapsto \text{classe de }(U, x, \text{id}_U)$
est bijective.
\end{proof}

\noindent
Soit $\mathcal{C}$ un site. Soit
$p = u : \mathcal{C} \to \textit{Ensembles}$ un foncteur. Par analogie avec les
constructions de la section \ref{section-functoriality-PSh}, étant donné un
ensemble $E$, nous définissons un préfaisceau $u^pE$ par la règle
\begin{equation}
\label{equation-skyscraper}
U
\longmapsto
u^pE(U) = \Mor_{\textit{Ensembles}}(u(U), E) = \text{Map}(u(U), E).
\end{equation}
Cela définit un foncteur
$u^p : \textit{Ensembles} \to \textit{PSh}(\mathcal{C})$, $E \mapsto u^pE$.

\begin{lemma}
\label{lemma-adjoint-point-push-stalk}
Pour tout foncteur $u : \mathcal{C} \to \textit{Ensembles}$, le foncteur $u^p$ est
adjoint à droite du foncteur fibre sur les préfaisceaux.
\end{lemma}

\begin{proof}
Soit $\mathcal{F}$ un préfaisceau sur $\mathcal{C}$.
Soit $E$ un ensemble. Un morphisme $\mathcal{F} \to u^pE$ est donné par un
système compatible d’applications
$\mathcal{F}(U) \to \text{Map}(u(U), E)$, c’est-à-dire par un système
compatible d’applications $\mathcal{F}(U) \times u(U) \to E$.
Et, par définition de $\mathcal{F}_p$, une application
$\mathcal{F}_p \to E$ est donnée par une règle qui associe à chaque triplet
$(U, x, \sigma)$ un élément de $E$, de telle sorte que des triplets équivalents
aient la même image ; voir la discussion entourant
l’Équation (\ref{equation-stalk}).
C’est encore la donnée d’un système compatible d’applications
$\mathcal{F}(U) \times u(U) \to E$.
\end{proof}

\noindent
Par analogie avec la section \ref{section-continuous-functors}, on a le lemme
suivant.

\begin{lemma}
\label{lemma-point-pushforward-sheaf}
Soit $\mathcal{C}$ un site. Soit
$p = u : \mathcal{C} \to \textit{Ensembles}$ un foncteur. Supposons que, pour tout
recouvrement $\{U_i \to U\}$ de $\mathcal{C}$,
\begin{enumerate}
\item l’application $\coprod u(U_i) \to u(U)$ soit surjective, et
\item les applications
$u(U_i \times_U U_j) \to u(U_i) \times_{u(U)} u(U_j)$ soient surjectives.
\end{enumerate}
Alors
\begin{enumerate}
\item le préfaisceau $u^pE$ est un faisceau pour tout ensemble $E$ ;
notons-le $u^sE$,
\item le foncteur fibre $\Sh(\mathcal{C}) \to \textit{Ensembles}$ et le foncteur
$u^s: \textit{Ensembles} \to \Sh(\mathcal{C})$ sont adjoints, et
\item on a $\mathcal{F}_p = \mathcal{F}^\#_p$ pour tout préfaisceau
d’ensembles $\mathcal{F}$.
\end{enumerate}
\end{lemma}

\begin{proof}
La première assertion résulte immédiatement de la définition d’un faisceau,
des hypothèses (1) et (2), et de la définition de $u^pE$. La deuxième est une
reformulation de l’adjonction entre $u^p$ et le foncteur fibre
(Lemme \ref{lemma-adjoint-point-push-stalk}) restreinte aux faisceaux.
La troisième assertion résulte de ce que, pour tout ensemble $E$, on a
$$
\text{Map}(\mathcal{F}_p, E) =
\Mor_{\textit{PSh}(\mathcal{C})}(\mathcal{F}, u^pE) =
\Mor_{\Sh(\mathcal{C})}(\mathcal{F}^\#, u^sE) =
\text{Map}(\mathcal{F}^\#_p, E)
$$
par la propriété d’adjonction du passage au faisceau associé.
\end{proof}

\noindent
En particulier, le Lemme \ref{lemma-point-pushforward-sheaf} s’applique
lorsque $p = u$ est un point. Dans ce cas, nous considérons le faisceau $u^sE$
comme le faisceau « gratte-ciel » de valeur $E$ en $p$.

\begin{definition}
\label{definition-pushforward-point}
Soit $p$ un point du site $\mathcal{C}$ donné par le foncteur $u$.
Pour un ensemble $E$, nous définissons $p_*E = u^sE$ comme le faisceau décrit
dans le Lemme \ref{lemma-point-pushforward-sheaf} ci-dessus.
Nous l’appelons parfois un {\it faisceau gratte-ciel}.
\end{definition}

\noindent
On a en particulier la propriété d’adjonction suivante pour les faisceaux
gratte-ciel et les fibres :
$$
\Mor_{\Sh(\mathcal{C})}(\mathcal{F}, p_*E)
=
\text{Map}(\mathcal{F}_p, E)
$$
Cela motive la notation $p^{-1}\mathcal{F} = \mathcal{F}_p$ que nous
emploierons parfois.

\begin{lemma}
\label{lemma-point-site-topos}
Soit $\mathcal{C}$ un site.
\begin{enumerate}
\item Soit $p$ un point du site $\mathcal{C}$.
Alors le couple de foncteurs $(p_*, p^{-1})$ introduit ci-dessus définit un
morphisme de topos $\Sh(pt) \to \Sh(\mathcal{C})$.
\item Soit $p = (p_*, p^{-1})$ un point du topos $\Sh(\mathcal{C})$.
Alors le foncteur $u : U \mapsto p^{-1}(h_U^\#)$ donne naissance à un point
$p'$ du site $\mathcal{C}$ dont le morphisme de topos associé
$(p'_*, (p')^{-1})$ est égal à $p$.
\end{enumerate}
\end{lemma}

\begin{proof}
Démonstration de (1). D’après ce qui précède, les foncteurs $p_*$ et $p^{-1}$
sont adjoints. Le foncteur $p^{-1}$ est exact par la
Définition \ref{definition-point}.
Toutes les conditions imposées dans la Définition \ref{definition-topos} sont
donc satisfaites, ce qui démontre (1).

\medskip\noindent
Démonstration de (2). Soit $\{U_i \to U\}$ un recouvrement de $\mathcal{C}$.
Alors $\coprod (h_{U_i})^\# \to h_U^\#$ est surjective ; voir le
Lemme \ref{lemma-covering-surjective-after-sheafification}.
Comme $p^{-1}$ est exact (par définition d’un morphisme de topos), on en
conclut que $\coprod u(U_i) \to u(U)$ est surjective.
Cela démontre la partie (1) de la Définition \ref{definition-point}.
Le passage au faisceau associé est exact ; voir le
Lemme \ref{lemma-sheafification-exact}.
Ainsi, si $U \times_V W$ existe dans $\mathcal{C}$, alors
$$
h_{U \times_V W}^\# =  h_U^\# \times_{h_V^\#} h_W^\#
$$
et l’on voit que $u(U \times_V W) = u(U) \times_{u(V)} u(W)$, puisque
$p^{-1}$ est exact. Cela démontre la partie (2) de la
Définition \ref{definition-point}.
Posons $p' = u$ et soit $\mathcal{F}_{p'}$ le foncteur fibre défini par
l’Équation (\ref{equation-stalk}) au moyen de $u$. Il existe une application
canonique de comparaison
$c : \mathcal{F}_{p'} \to \mathcal{F}_p = p^{-1}\mathcal{F}$.
En effet, étant donné un triplet $(U, x, \sigma)$ représentant un élément
$\xi$ de $\mathcal{F}_{p'}$, nous considérons $\sigma$ comme une application
$\sigma : h_U^\# \to \mathcal{F}$ et nous pouvons poser
$c(\xi) = p^{-1}(\sigma)(x)$, puisque
$x \in u(U) = p^{-1}(h_U^\#)$. D’après le
Lemme \ref{lemma-points-recover}, on a $(h_U)_{p'} = u(U)$. Puisque les
conditions (1) et (2) de la Définition \ref{definition-point} sont satisfaites
pour $p'$, on a également $(h_U^\#)_{p'} = (h_U)_{p'}$ d’après le
Lemme \ref{lemma-point-pushforward-sheaf}. Ainsi,
$$
(h_U^\#)_{p'} = (h_U)_{p'} = u(U) = p^{-1}(h_U^\#)
$$
Nous affirmons que cette bijection est égale à l’application de comparaison
$c : (h_U^\#)_{p'} \to p^{-1}(h_U^\#)$ (vérification omise).
Tout faisceau sur $\mathcal{C}$ est un coégalisateur d’applications entre des
coproduits de faisceaux de la forme $h_U^\#$ ; voir le
Lemme \ref{lemma-sheaf-coequalizer-representable}.
Le foncteur fibre $\mathcal{F} \mapsto \mathcal{F}_{p'}$ et le foncteur
$p^{-1}$ commutent aux limites inductives arbitraires (car ils sont tous deux
adjoints à gauche).
On en conclut que $c$ est un isomorphisme pour tout faisceau $\mathcal{F}$.
Ainsi, le foncteur fibre $\mathcal{F} \mapsto \mathcal{F}_{p'}$ est isomorphe
à $p^{-1}$ et, en particulier, il est exact. Cela démontre que la condition
(3) de la Définition \ref{definition-point} est satisfaite et que $p'$ est un
point. La dernière assertion a déjà été démontrée ci-dessus, puisque nous
avons vu que $p^{-1} = (p')^{-1}$.
\end{proof}

\noindent
En fait, un point correspond toujours à un morphisme de sites, comme le montre
le lemme suivant.

\begin{lemma}
\label{lemma-site-point-morphism}
Soit $\mathcal{C}$ un site. Soit $p$ un point de $\mathcal{C}$ donné par
$u : \mathcal{C} \to \textit{Ensembles}$. Soit $S_0$ un ensemble infini tel que
$u(U) \subset S_0$ pour tout $U \in \Ob(\mathcal{C})$. Soit $\mathcal{S}$ le
site construit à partir de l’ensemble des parties
$S = \mathcal{P}(S_0)$ dans la Remarque \ref{remark-pt-topos}.
Alors
\begin{enumerate}
\item il existe une équivalence $i : \Sh(pt) \to \Sh(\mathcal{S})$,
\item le foncteur $u : \mathcal{C} \to \mathcal{S}$ induit un morphisme de
sites $f : \mathcal{S} \to \mathcal{C}$, et
\item le composé
$$
\Sh(pt) \to
\Sh(\mathcal{S}) \to
\Sh(\mathcal{C})
$$
est le morphisme de topos $(p_*, p^{-1})$ du
Lemme \ref{lemma-point-site-topos}.
\end{enumerate}
\end{lemma}

\begin{proof}
La partie (1) a été établie dans la Remarque \ref{remark-pt-topos}.
Rappelons en outre que l’équivalence associe à l’ensemble $E$ le faisceau
$i_*E$ sur $\mathcal{S}$ défini par la règle
$V \mapsto \Mor_{\textit{Ensembles}}(V, E)$.
La partie (2) résulte clairement de la définition d’un point de $\mathcal{C}$
(Définition \ref{definition-point}) et de la définition d’un morphisme de
sites (Définition \ref{definition-morphism-sites}).
Enfin, considérons $f_*i_*E$. Par construction, on a
$$
f_*i_*E(U) = i_*E(u(U)) = \Mor_{\textit{Ensembles}}(u(U), E)
$$
ce qui est égal à $p_*E(U)$ ; voir
l’Équation (\ref{equation-skyscraper}).
Cela démontre (3).
\end{proof}

\noindent
Contrairement à ce qui se passe dans le cas topologique, la fibre du faisceau
gratte-ciel de valeur $E$ n’est pas toujours égale à $E$. Voici ce qui est
toujours vrai.

\begin{lemma}
\label{lemma-stalk-skyscraper}
Soit $\mathcal{C}$ un site. Soit
$p : \Sh(pt) \to \Sh(\mathcal{C})$ un point du topos associé à
$\mathcal{C}$. Pour tout ensemble $E$, il existe des applications canoniques
$$
E \longrightarrow (p_*E)_p \longrightarrow E
$$
dont le composé est $\text{id}_E$.
\end{lemma}

\begin{proof}
Il existe toujours un morphisme d’adjonction
$(p_*E)_p = p^{-1}p_*E \to E$.
Ce morphisme est un isomorphisme lorsque $E = \{*\}$, car $p_*$ et $p^{-1}$
sont tous deux exacts à gauche et transforment donc l’objet final en l’objet
final. Étant donné $e \in E$, on peut donc considérer l’application
$i_e : \{*\} \to E$, qui donne
$$
\xymatrix{
p^{-1}p_*\{*\} \ar[rr]_{p^{-1}p_*i_e} \ar[d]_{\cong} & & p^{-1}p_*E \ar[d] \\
\{*\} \ar[rr]^{i_e} & & E
}
$$
d’où l’application $E \to (p_*E)_p = p^{-1}p_*E$ voulue.
\end{proof}

\begin{lemma}
\label{lemma-skyscraper-functor-exact}
Soit $\mathcal{C}$ un site. Soit
$p : \Sh(pt) \to \Sh(\mathcal{C})$ un point du topos associé à
$\mathcal{C}$. Le foncteur $p_* : \textit{Ensembles} \to \Sh(\mathcal{C})$
possède les propriétés suivantes : il commute aux limites projectives
arbitraires, il est exact à gauche, il est fidèle, il transforme les
surjections en surjections, il commute aux coégalisateurs, il reflète les
injections, il reflète les surjections et il reflète les isomorphismes.
\end{lemma}

\begin{proof}
Comme $p_*$ est adjoint à droite, il commute aux limites projectives
arbitraires et est exact à gauche. Le fait que
$p^{-1}p_*E \to E$ soit une surjection canoniquement scindée implique que
$p_*$ est fidèle, reflète les injections, reflète les surjections et reflète
les isomorphismes. D’après le Lemme \ref{lemma-point-site-topos}, on peut
supposer que $p$ provient d’un point
$u : \mathcal{C} \to \textit{Ensembles}$ du site sous-jacent $\mathcal{C}$.
Dans ce cas, le faisceau $p_*E$ est donné par
$$
p_*E(U) = \Mor_{\textit{Ensembles}}(u(U), E)
$$
voir l’Équation (\ref{equation-skyscraper}) et la
Définition \ref{definition-pushforward-point}.
Il résulte immédiatement de cette formule que $p_*$ transforme les
surjections en surjections et les coégalisateurs en coégalisateurs.
\end{proof}



\section{Construction de points}
\label{section-construct-points}

\noindent
Dans cette section, nous donnons des critères pour qu’un foncteur
d’un site vers la catégorie des ensembles définisse un point de ce site.

\begin{lemma}
\label{lemma-neighbourhoods-cofiltered}
Soit $\mathcal{C}$ un site. Soit $p = u : \mathcal{C} \to \textit{Ensembles}$
un foncteur. Si la catégorie des voisinages de $p$ est
cofiltrante, alors le foncteur fibre (\ref{equation-stalk})
est exact à gauche.
\end{lemma}

\begin{proof}
Soit $\mathcal{I} \to \Sh(\mathcal{C})$,
$i \mapsto \mathcal{F}_i$ un diagramme fini de faisceaux.
Nous devons montrer que la fibre de la limite de ce
système coïncide avec la limite des fibres.
Soit $\mathcal{F}$ la limite du système en tant que {\it préfaisceau}.
D’après le Lemme \ref{lemma-limit-sheaf}, c’est un faisceau et
c’est la limite dans la catégorie des faisceaux.
Il nous faut donc montrer que
$\mathcal{F}_p = \lim_\mathcal{I} \mathcal{F}_{i, p}$.
Rappelons aussi que $\mathcal{F}$ admet une description simple ; voir la
section \ref{section-limits-colimits-PSh}. Il nous faut ainsi montrer que
$$
\lim_i \colim_{\{(U, x)\}^{opp}} \mathcal{F}_i(U)
=
\colim_{\{(U, x)\}^{opp}} \lim_i \mathcal{F}_i(U).
$$
Cela résulte de Catégories, Lemme \ref{categories-lemma-directed-commutes},
car l’opposée de la catégorie des voisinages est filtrante
par hypothèse.
\end{proof}

\begin{lemma}
\label{lemma-neighbourhoods-directed}
Soit $\mathcal{C}$ un site. Supposons que $\mathcal{C}$ possède
un objet final $X$ et des produits fibrés.
Soit $p = u : \mathcal{C} \to \textit{Ensembles}$ un foncteur tel que
\begin{enumerate}
\item $u(X)$ soit un singleton, et
\item pour toute paire de morphismes $U \to W$ et $V \to W$ de
même but, l’application
$u(U \times_W V) \to u(U) \times_{u(W)} u(V)$ soit bijective.
\end{enumerate}
Alors la catégorie des voisinages de $p$ est cofiltrante
et, par conséquent, le foncteur fibre $\Sh(\mathcal{C}) \to \textit{Ensembles}$,
$\mathcal{F} \to \mathcal{F}_p$ commute aux limites finies.
\end{lemma}

\begin{proof}
Remarquons l’analogie avec le Lemme \ref{lemma-directed}.
Les hypothèses sur $\mathcal{C}$ impliquent que $\mathcal{C}$ admet les limites finies.
Voir Catégories, Lemme \ref{categories-lemma-finite-limits-exist}.
L’hypothèse (1) implique que la catégorie des voisinages
n’est pas vide. Soient $(U, x)$ et $(V, y)$ deux voisinages.
Alors
$u(U \times V) = u(U \times_X V) =
u(U) \times_{u(X)} u(V) = u(U) \times u(V)$ d’après (2).
Il existe donc un voisinage $(U \times_X V, z)$ se projetant
sur $(U, x)$ et sur $(V, y)$.
Soient $a, b : (V, y) \to (U, x)$ deux morphismes
de la catégorie des voisinages. Soit $W$ l’égalisateur de
$a, b : V \to U$. Comme dans la démonstration de
Catégories, Lemme \ref{categories-lemma-finite-limits-exist},
nous pouvons écrire $W$ au moyen de produits fibrés :
$$
W = (V \times_{a, U, b} V) \times_{(pr_1, pr_2), V \times V, \Delta} V
$$
La bijectivité de (2) garantit l’existence d’un élément $z \in u(W)$
dont l’image est $((y, y), y)$.
Alors $(W, z) \to (V, y)$ égalise $a, b$, comme voulu.
La conclusion « par conséquent » est le Lemme \ref{lemma-neighbourhoods-cofiltered}.
\end{proof}

\begin{proposition}
\label{proposition-point-limits}
Soit $\mathcal{C}$ un site. Supposons que les limites finies existent
dans $\mathcal{C}$. (Autrement dit, $\mathcal{C}$ possède des produits fibrés et un
objet final.) Un point $p$ d’un tel site $\mathcal{C}$
est donné par un foncteur $u : \mathcal{C} \to \textit{Ensembles}$ tel que
\begin{enumerate}
\item $u$ commute aux limites finies, et
\item si $\{U_i \to U\}$ est un recouvrement, alors
$\coprod_i u(U_i) \to u(U)$ est surjective.
\end{enumerate}
\end{proposition}

\begin{proof}
Supposons d’abord que $p$ soit un point (Définition \ref{definition-point})
donné par un foncteur $u$. La condition (2) résulte directement de
la définition d’un point. D’après le Lemme \ref{lemma-points-recover},
nous avons $(h_U)_p = u(U)$. D’après le Lemme \ref{lemma-point-pushforward-sheaf},
nous avons $(h_U^\#)_p = (h_U)_p$. Ainsi, $u$
est égal au composé des foncteurs
$$
\mathcal{C} \xrightarrow{h}
\textit{PSh}(\mathcal{C}) \xrightarrow{{}^\#}
\Sh(\mathcal{C}) \xrightarrow{()_p}
\textit{Ensembles}
$$
Chacun de ces foncteurs est exact à gauche ; ainsi, $u$ satisfait (1).

\medskip\noindent
Réciproquement, supposons que $u$ satisfasse (1) et (2).
Nous voyons aussitôt que $u$ satisfait les deux premières
conditions de la Définition \ref{definition-point}. De plus, son
foncteur fibre est exact, car c’est un adjoint à gauche d’après le
Lemme \ref{lemma-point-pushforward-sheaf}, et il commute
aux limites finies d’après le Lemme \ref{lemma-neighbourhoods-directed}.
\end{proof}

\begin{remark}
\label{remark-improve-proposition-points-limits}
En fait, soit $\mathcal{C}$ un site. Supposons que $\mathcal{C}$ possède un objet final
$X$ et des produits fibrés. Soit $p = u: \mathcal{C} \to \textit{Ensembles}$ un
foncteur tel que
\begin{enumerate}
\item $u(X) = \{*\}$ soit un singleton, et
\item pour toute paire de morphismes $U \to W$ et $V \to W$ de
même but, l’application
$u(U \times_W V) \to u(U) \times_{u(W)} u(V)$ soit surjective.
\item pour tout recouvrement $\{U_i \to U\}$, l’application
$\coprod u(U_i) \to u(U)$ soit surjective.
\end{enumerate}
Alors, en général, $p$ n’est {\bf pas} un point de $\mathcal{C}$.
Un exemple est la catégorie $\mathcal{C}$ à deux objets $\{U, X\}$
et une seule flèche non identique, à savoir $U \to X$. Munissons $\mathcal{C}$
de la topologie triviale, c’est-à-dire que les seuls recouvrements sont $\{U \to U\}$ et
$\{X \to X\}$. Un faisceau $\mathcal{F}$ est alors la même chose qu’un préfaisceau et
consiste en un triplet $(A, B, A \to B)$ : à savoir $A = \mathcal{F}(X)$,
$B = \mathcal{F}(U)$, où $A \to B$ est l’application de restriction correspondant
à $U \to X$. Remarquons que $U \times_X U = U$ ; les produits fibrés existent donc.
Considérons le foncteur $u = p$ défini par $u(X) = \{*\}$ et $u(U) = \{*_1, *_2\}$.
Il satisfait (1), (2) et (3), mais le foncteur fibre correspondant
(\ref{equation-stalk}) est le foncteur
$$
(A, B, A \to B) \longmapsto B \amalg_A B
$$
qui n’est pas exact. En effet, considérons
$(\emptyset, \{1\}, \emptyset \to \{1\}) \to (\{1\}, \{1\}, \{1\} \to \{1\})$,
qui est une application injective de faisceaux mais qui est transformée en l’application
non injective d’ensembles
$$
\{1\} \amalg \{1\} \longrightarrow \{1\} \amalg_{\{1\}} \{1\}
$$
par le foncteur fibre.
\end{remark}

\begin{example}
\label{example-point-topological}
Soit $X$ un espace topologique. Soit $X_{Zar}$ le site de
l’Exemple \ref{example-site-topological}.
Soit $x \in X$ un point. Considérons le foncteur
$$
u : X_{Zar} \longrightarrow \textit{Ensembles}, \quad
U \mapsto
\left\{
\begin{matrix}
\emptyset & \text{si} & x \not \in U \\
\{*\} & \text{si} & x \in U
\end{matrix}
\right.
$$
Ce foncteur commute aux produits et aux produits fibrés,
et transforme les recouvrements en familles surjectives d’applications.
Nous obtenons donc un point $p$ du site $X_{Zar}$.
On vérifie immédiatement que le foncteur fibre
coïncide avec la fibre en $x$ définie dans
Faisceaux, section
\ref{sheaves-section-stalks}.
\end{example}

\begin{example}
\label{example-point-topology}
Soit $X$ un espace topologique. Quels sont les points du topos
$\Sh(X)$ ? Pour le déterminer, soit $X_{Zar}$ le site de
l’Exemple \ref{example-site-topological}.
D’après le
Lemme \ref{lemma-point-site-topos},
un point de $\Sh(X)$ correspond à un point de ce site.
Soit $p$ un point du site $X_{Zar}$ donné par le foncteur
$u : X_{Zar} \to \textit{Ensembles}$. Nous allons utiliser la
caractérisation de tels $u$ donnée dans la
Proposition \ref{proposition-point-limits}.
Elle implique immédiatement que $u(\emptyset) = \emptyset$ et
$u(U \cap V) = u(U) \times u(V)$. En particulier,
$u(U) = u(U) \times u(U)$ par l’application diagonale, ce qui implique que $u(U)$
est soit un singleton, soit vide. En outre, si $U = \bigcup U_i$ est un
recouvrement ouvert, alors
$$
u(U) = \emptyset \Rightarrow \forall i, \ u(U_i) = \emptyset
\quad\text{et}\quad
u(U) \not = \emptyset \Rightarrow \exists i, \ u(U_i) \not = \emptyset.
$$
Nous en concluons qu’il existe un unique plus grand ouvert $W \subset X$ tel que
$u(W) = \emptyset$, à savoir la réunion de tous les ouverts $U$ tels que
$u(U) = \emptyset$. Posons $Z = X \setminus W$. Si $Z = Z_1 \cup Z_2$ avec
$Z_i \subset Z$ fermés, alors $W = (X \setminus Z_1) \cap (X \setminus Z_2)$ ;
ainsi $\emptyset = u(W) = u(X \setminus Z_1) \times u(X \setminus Z_2)$,
et nous concluons que $u(X \setminus Z_1) = \emptyset$ ou que
$u(X \setminus Z_2) = \emptyset$. Cela signifie que $X \setminus Z_1 = W$
ou que $X \setminus Z_2 = W$. Autrement dit, $Z$ est irréductible.
Nous voyons maintenant que $u$ est décrit par la règle
$$
u : X_{Zar} \longrightarrow \textit{Ensembles}, \quad
U \mapsto
\left\{
\begin{matrix}
\emptyset & \text{si} & Z \cap U = \emptyset \\
\{*\} & \text{si} & Z \cap U \not = \emptyset
\end{matrix}
\right.
$$
Remarquons que, pour tout fermé irréductible $Z \subset X$, ce
foncteur satisfait les hypothèses (1), (2) de la
Proposition \ref{proposition-point-limits}
et définit donc un point. Autrement dit, les points du
site $X_{Zar}$ sont en correspondance bijective avec les
sous-ensembles fermés irréductibles de $X$. En particulier, si $X$ est
un espace topologique sobre, alors les points de $X_{Zar}$ et les
points de $X$ sont en correspondance bijective ; voir
l’Exemple \ref{example-point-topological}.
\end{example}

\begin{example}
\label{example-point-G-sets}
Considérons le site $\mathcal{T}_G$ décrit dans
l’Exemple \ref{example-site-on-group} et la
section \ref{section-example-sheaf-G-sets}.
Le foncteur d’oubli $u : \mathcal{T}_G \to \textit{Ensembles}$
commute aux produits et aux produits fibrés, et transforme
les recouvrements en familles surjectives. Il définit donc un point
de $\mathcal{T}_G$. Identifions $\Sh(\mathcal{T}_G)$
et $G\textit{-Ensembles}$. Le foncteur fibre
$$
p^{-1} :
\Sh(\mathcal{T}_G) = G\textit{-Ensembles}
\longrightarrow
\textit{Ensembles}
$$
est le foncteur d’oubli. L’image directe $p_*$ est le
foncteur
$$
\textit{Ensembles}
\longrightarrow
\Sh(\mathcal{T}_G) = G\textit{-Ensembles}
$$
qui envoie un ensemble $S$ sur le $G$-ensemble $\text{Map}(G, S)$ muni de
l’action $g \cdot \psi = \psi \circ R_g$, où $R_g$ est la multiplication
à droite. En particulier,
$p^{-1}p_*S = \text{Map}(G, S)$ en tant qu’ensemble, et les applications
$S \to \text{Map}(G, S) \to S$ du
Lemme \ref{lemma-stalk-skyscraper}
sont les applications évidentes.
\end{example}

\begin{example}
\label{example-indiscrete-points}
Soit $\mathcal{C}$ une catégorie munie de la topologie grossière
(Exemple \ref{example-indiscrete}). Pour tout objet $U_0$ de $\mathcal{C}$,
le foncteur $u : U \mapsto \Mor_\mathcal{C}(U_0, U)$ définit un point
$p$ de $\mathcal{C}$. En effet, les conditions (1) et (2) de la
Définition \ref{definition-point} sont immédiates puisque les seuls recouvrements
sont donnés par les applications identiques. La condition (3) est satisfaite car
$\mathcal{F}_p = \mathcal{F}(U_0)$ et, la topologie étant grossière,
prendre les sections sur $U_0$ est un foncteur exact.
\end{example}





\section{Points et morphismes de topos}
\label{section-functorial-points}

\noindent
Dans cette section, nous formulons quelques remarques sur les points et les morphismes
de topos.

\begin{lemma}
\label{lemma-point-functor}
Soit $u : \mathcal{C} \to \mathcal{D}$ un foncteur. Soit
$v : \mathcal{D} \to \textit{Ensembles}$ un foncteur, et posons
$w = v \circ u$. Notons $q$, respectivement $p$, le foncteur fibre
(\ref{equation-stalk}) associé à $v$, respectivement à $w$.
Alors $(u_p\mathcal{F})_q = \mathcal{F}_p$, fonctoriellement par rapport au
préfaisceau $\mathcal{F}$ sur $\mathcal{C}$.
\end{lemma}

\begin{proof}
C’est un fait catégorique simple. Nous avons
\begin{align*}
(u_p\mathcal{F})_q
& =
\colim_{(V, y)} \colim_{U, \phi : V \to u(U)} \mathcal{F}(U) \\
& = \colim_{(V, y, U, \phi : V \to u(U))} \mathcal{F}(U) \\
& = \colim_{(U, x)} \mathcal{F}(U) \\
& = \mathcal{F}_p
\end{align*}
La première égalité résulte de la définition de $u_p$ et de la
définition du foncteur fibre. Observons que $y \in v(V)$.
Dans la deuxième égalité, nous réunissons simplement les limites inductives.
Pour établir la troisième égalité, appliquons
Catégories, Lemme \ref{categories-lemma-colimit-constant-connected-fibers},
au foncteur $F$ des catégories de diagrammes défini par la règle
$$
F((V, y, U, \phi : V \to u(U))) = (U, v(\phi)(y)).
$$
Cela a un sens car $w(U) = v(u(U))$. Vérifions les hypothèses de
Catégories, Lemme \ref{categories-lemma-colimit-constant-connected-fibers}.
Observons que $F$ possède un inverse à droite, à savoir
$(U, x) \mapsto (u(U), x, U, \text{id} : u(U) \to u(U))$.
Cela a encore un sens car $x \in w(U) = v(u(U))$. D’autre part,
il existe toujours un morphisme
$$
(V, y, U, \phi : V \to u(U))
\longrightarrow
(u(U), v(\phi)(y), U, \text{id} : u(U) \to u(U))
$$
dans la catégorie fibre au-dessus de $(U, x)$, ce qui montre que les catégories fibres
sont connexes. La quatrième et dernière égalité est claire.
\end{proof}

\begin{lemma}
\label{lemma-point-morphism-sites}
\begin{slogan}
Un morphisme de sites définit un morphisme sur les points, et l’image inverse respecte les fibres.
\end{slogan}
Soit $f : \mathcal{D} \to \mathcal{C}$ un morphisme de sites
donné par un foncteur continu $u : \mathcal{C} \to \mathcal{D}$.
Soit $q$ un point de $\mathcal{D}$ donné par le foncteur
$v : \mathcal{D} \to \textit{Ensembles}$ ; voir la
Définition \ref{definition-point}.
Alors le foncteur $v \circ u : \mathcal{C} \to \textit{Ensembles}$
définit un point $p$ de $\mathcal{C}$ et il existe en outre
une identification canonique
$$
(f^{-1}\mathcal{F})_q = \mathcal{F}_p
$$
pour tout faisceau $\mathcal{F}$ sur $\mathcal{C}$.
\end{lemma}

\begin{proof}[Première démonstration du Lemme \ref{lemma-point-morphism-sites}]
Puisque $u$ est continu et que $v$ définit un point,
il est immédiat que $v \circ u$ satisfait les conditions (1) et (2) de la
Définition \ref{definition-point}. Démontrons l’égalité affichée.
Soit $\mathcal{F}$ un faisceau sur $\mathcal{C}$. Alors
$$
(f^{-1}\mathcal{F})_q = (u_s\mathcal{F})_q =
(u_p \mathcal{F})_q = \mathcal{F}_p
$$
La première égalité tient à $f^{-1} = u_s$, la deuxième au
Lemme \ref{lemma-point-pushforward-sheaf}, et la troisième
au Lemme \ref{lemma-point-functor}.
Nous voyons alors que $p$ satisfait aussi la condition (3) de la
Définition \ref{definition-point},
car son foncteur fibre est un composé de foncteurs exacts. Cela achève la démonstration.
\end{proof}

\begin{proof}[Seconde démonstration du Lemme \ref{lemma-point-morphism-sites}]
D’après le
Lemme \ref{lemma-site-point-morphism},
nous pouvons factoriser $(q_*, q^{-1})$ sous la forme
$$
\Sh(pt) \xrightarrow{i} \Sh(\mathcal{S})
\xrightarrow{h} \Sh(\mathcal{D})
$$
où le second morphisme de topos provient d’un morphisme de sites
$h : \mathcal{S} \to \mathcal{D}$ induit par le foncteur
$v : \mathcal{D} \to \mathcal{S}$ (ce qui a un sens puisque
$\mathcal{S} \subset \textit{Ensembles}$ est une sous-catégorie pleine contenant
tout objet de l’image de $v$). D’après le
Lemme \ref{lemma-composition-morphisms-sites},
le composé $v \circ u : \mathcal{C} \to \mathcal{S}$
définit un morphisme de sites $g : \mathcal{S} \to \mathcal{C}$.
En particulier, le foncteur
$v \circ u : \mathcal{C} \to \mathcal{S}$
est continu, ce qui, par la définition des recouvrements
de $\mathcal{S}$ (voir la
Remarque \ref{remark-pt-topos}),
signifie que $v \circ u$ satisfait les conditions (1) et (2) de la
Définition \ref{definition-point}.
D’autre part, nous voyons que
$$
g_*i_*E(U) = i_*E(v(u(U)) = \Mor_{\textit{Ensembles}}(v(u(U)), E)
$$
par la construction de $i$ dans la
Remarque \ref{remark-pt-topos}.
Remarquons qu’il s’agit précisément de la formule pour
$(v \circ u)^pE$ ; voir l’Équation (\ref{equation-skyscraper}).
D’après le
Lemme \ref{lemma-point-pushforward-sheaf},
le foncteur $g_*i_* = (v \circ u)^p = (v \circ u)^s$
est adjoint à droite du foncteur fibre
$\mathcal{F} \mapsto \mathcal{F}_q$.
Nous voyons donc que le foncteur fibre $p^{-1}$ est canoniquement
isomorphe à $i^{-1} \circ g^{-1}$. Il est donc exact, et
nous concluons que $p$ est un point. Enfin, puisque
$g = f \circ h$ par construction, nous voyons que
$p^{-1} = i^{-1} \circ h^{-1} \circ f^{-1} = q^{-1} \circ f^{-1}$,
c’est-à-dire que nous avons la formule affichée du lemme.
\end{proof}

\begin{lemma}
\label{lemma-point-morphism-topoi}
Soit $f : \Sh(\mathcal{D}) \to \Sh(\mathcal{C})$
un morphisme de topos. Soit $q : \Sh(pt) \to \Sh(\mathcal{D})$
un point. Alors $p = f \circ q$ est un point du topos
$\Sh(\mathcal{C})$ et nous avons
une identification canonique
$$
(f^{-1}\mathcal{F})_q = \mathcal{F}_p
$$
pour tout faisceau $\mathcal{F}$ sur $\mathcal{C}$.
\end{lemma}

\begin{proof}
Cela résulte immédiatement des définitions et du fait que nous pouvons
composer les morphismes de topos.
\end{proof}





\section{Localisation et points}
\label{section-localize-points}

\noindent
Dans cette section, nous montrons que les points d’une localisation $\mathcal{C}/U$
se construisent simplement à partir des points de $\mathcal{C}$.

\begin{lemma}
\label{lemma-point-localize}
Soit $\mathcal{C}$ un site. Soit $p$ un point de $\mathcal{C}$ donné par
$u : \mathcal{C} \to \textit{Ensembles}$. Soit $U$ un objet de $\mathcal{C}$
et soit $x \in u(U)$. Le foncteur
$$
v : \mathcal{C}/U \longrightarrow \textit{Ensembles}, \quad
(\varphi : V \to U) \longmapsto \{y \in u(V) \mid u(\varphi)(y) = x\}
$$
définit un point $q$ du site $\mathcal{C}/U$ tel que le diagramme
$$
\xymatrix{
& \Sh(pt) \ar[d]^p \ar[ld]_q \\
\Sh(\mathcal{C}/U) \ar[r]^{j_U} &
\Sh(\mathcal{C})
}
$$
soit commutatif. Autrement dit,
$\mathcal{F}_p = (j_U^{-1}\mathcal{F})_q$ pour tout
faisceau sur $\mathcal{C}$.
\end{lemma}

\begin{proof}
Choisissons $S$ et $\mathcal{S}$ comme dans le
Lemme \ref{lemma-site-point-morphism}.
Nous pouvons identifier $\Sh(pt) = \Sh(\mathcal{S})$
comme dans ce lemme, et écrire
$p = f : \Sh(\mathcal{S}) \to \Sh(\mathcal{C})$
pour le morphisme de topos induit par $u$.
D’après le
Lemme \ref{lemma-localize-morphism},
nous obtenons un diagramme commutatif de topos
$$
\xymatrix{
\Sh(\mathcal{S}/u(U)) \ar[r]_-{j_{u(U)}} \ar[d]_{p'} &
\Sh(\mathcal{S}) \ar[d]^p \\
\Sh(\mathcal{C}/U) \ar[r]^{j_U} &
\Sh(\mathcal{C}),
}
$$
où $p'$ est donné par le foncteur $u' : \mathcal{C}/U \to \mathcal{S}/u(U)$,
$V/U \mapsto u(V)/u(U)$.
Considérons le foncteur $j_x : \mathcal{S} \cong \mathcal{S}/x$ obtenu
en associant à un ensemble $E$ l’ensemble $E$ muni de l’application constante
$E \to u(U)$ de valeur $x$.
Alors $j_x$ est un foncteur cocontinu pleinement fidèle qui possède un
adjoint à droite continu
$v_x : (\psi : E \to u(U)) \mapsto \psi^{-1}(\{x\})$.
Remarquons que $j_{u(U)} \circ j_x = \text{id}_\mathcal{S}$ et que
$v_x \circ u' = v$.
Ces observations impliquent le diagramme commutatif
de topos suivant :
$$
\xymatrix{
\Sh(\mathcal{S}) \ar[rd]^a \ar[rdd]_q
\ar `r[rrr] `d[dd]^p [rrdd] & & & \\
& \Sh(\mathcal{S}/u(U)) \ar[r]_-{j_{u(U)}} \ar[d]^{p'} &
\Sh(\mathcal{S}) \ar[d]^p & \\
& \Sh(\mathcal{C}/U) \ar[r]^{j_U} &
\Sh(\mathcal{C}) &
}
$$
En effet :
\begin{enumerate}
\item Le morphisme
$a : \Sh(\mathcal{S}) \to \Sh(\mathcal{S}/u(U))$
est le morphisme de topos associé au foncteur cocontinu
$j_x$, lequel est égal au morphisme associé au foncteur continu
$v_x$ ; voir le
Lemme \ref{lemma-cocontinuous-morphism-topoi}
et la
section \ref{section-cocontinuous-adjoint}.
\item Le composé $p \circ j_{u(U)} \circ a = p$ puisque
$j_{u(U)} \circ j_x = \text{id}_\mathcal{S}$.
\item Le composé $p' \circ a$ donne un morphisme de topos.
De plus, c’est le morphisme de topos associé au foncteur continu
$v_x \circ u' = v$. Ainsi, $v$ définit bien un point $q$ de
$\mathcal{C}/U$, qui s’insère par construction dans le diagramme ci-dessus.
\end{enumerate}
Cela achève la démonstration du lemme.
\end{proof}

\begin{lemma}
\label{lemma-points-above-point}
Soient $\mathcal{C}$, $p$, $u$, $U$ comme dans le
Lemme \ref{lemma-point-localize}.
La construction du
Lemme \ref{lemma-point-localize}
donne une correspondance bijective entre les points $q$
de $\mathcal{C}/U$ situés au-dessus de $p$ et les éléments $x$ de $u(U)$.
\end{lemma}

\begin{proof}
Soit $q$ un point de $\mathcal{C}/U$ donné par le foncteur
$v : \mathcal{C}/U \to \textit{Ensembles}$ tel que $j_U \circ q = p$
comme morphismes de topos. Rappelons que $u(V) = p^{-1}(h_V^\#)$ pour tout
objet $V$ de $\mathcal{C}$ ; voir le
Lemme \ref{lemma-point-site-topos}.
De même, $v(V/U) = q^{-1}(h_{V/U}^\#)$ pour tout objet
$V/U$ de $\mathcal{C}/U$.
Considérons les deux diagrammes suivants :
$$
\vcenter{
\xymatrix{
\Mor_{\mathcal{C}/U}(W/U, V/U) \ar[r] \ar[d] &
\Mor_\mathcal{C}(W, V) \ar[d] \\
\Mor_{\mathcal{C}/U}(W/U, U/U) \ar[r] &
\Mor_\mathcal{C}(W, U)
}
}
\quad
\vcenter{
\xymatrix{
h_{V/U}^\# \ar[r] \ar[d] & j_U^{-1}(h_V^\#) \ar[d] \\
h_{U/U}^\# \ar[r] & j_U^{-1}(h_U^\#)
}
}
$$
Le diagramme de droite est le faisceau associé au diagramme des
préfaisceaux sur $\mathcal{C}/U$ qui envoie $W/U$ sur le diagramme de gauche
d’ensembles. (Il y a ici un petit point technique :
nous avons $(j_U^{-1}h_V)^\# = j_U^{-1}(h_V^\#)$, et de même pour $h_U$ ; voir le
Lemme \ref{lemma-technical-pu}.)
Remarquons que le diagramme de gauche d’ensembles est cartésien.
Comme le passage au faisceau associé est exact
(Lemme \ref{lemma-sheafification-exact}),
nous en concluons que le diagramme de droite est cartésien.

\medskip\noindent
Appliquons le foncteur exact $q^{-1}$ au diagramme de droite
pour obtenir un diagramme cartésien
$$
\xymatrix{
v(V/U) \ar[r] \ar[d] & u(V) \ar[d] \\
v(U/U) \ar[r] & u(U)
}
$$
d’ensembles. Nous avons utilisé ici
$q^{-1} \circ j^{-1} = p^{-1}$. Puisque $U/U$ est un objet final de
$\mathcal{C}/U$, nous voyons que $v(U/U)$ est un singleton. L’image de
$v(U/U)$ dans $u(U)$ est donc un élément $x$, et l’application horizontale
supérieure donne une bijection
$v(V/U) \to  \{y \in u(V) \mid y \mapsto x\text{ dans }u(U)\}$,
comme voulu.
\end{proof}

\begin{lemma}
\label{lemma-stalk-j-shriek}
Soit $\mathcal{C}$ un site. Soit $p$ un point de $\mathcal{C}$ donné par
$u : \mathcal{C} \to \textit{Ensembles}$. Soit $U$ un objet de $\mathcal{C}$.
Pour tout faisceau $\mathcal{G}$ sur $\mathcal{C}/U$, nous avons
$$
(j_{U!}\mathcal{G})_p =
\coprod\nolimits_q \mathcal{G}_q
$$
où la somme disjointe porte sur les points $q$ de $\mathcal{C}/U$
associés aux éléments $x \in u(U)$ comme dans le
Lemme \ref{lemma-point-localize}.
\end{lemma}

\begin{proof}
Nous utilisons la description de $j_{U!}\mathcal{G}$ comme le faisceau associé
au préfaisceau
$V \mapsto
\coprod\nolimits_{\varphi \in \Mor_\mathcal{C}(V, U)}
\mathcal{G}(V/_\varphi U)$
du
Lemme \ref{lemma-describe-j-shriek}.
De plus, la fibre de $j_{U!}\mathcal{G}$ en $p$ est égale à la
fibre de ce préfaisceau ; voir le
Lemme \ref{lemma-point-pushforward-sheaf}.
Nous voyons donc que
$$
(j_{U!}\mathcal{G})_p =
\colim_{(V, y)} \coprod\nolimits_{\varphi : V \to U} \mathcal{G}(V/_\varphi U)
$$
À chaque élément $(V, y, \varphi, s)$ de cette limite inductive, nous pouvons associer
$x = u(\varphi)(y) \in u(U)$. Nous obtenons donc
$$
(j_{U!}\mathcal{G})_p =
\coprod\nolimits_{x \in u(U)}
\colim_{(\varphi : V \to U, y), \ u(\varphi)(y) = x} \mathcal{G}(V/_\varphi U).
$$
Cette expression est égale à celle du lemme par notre
construction des points $q$.
\end{proof}

\begin{remark}
\label{remark-not-pushforward}
Attention : le résultat du
Lemme \ref{lemma-stalk-j-shriek}
n’a pas d’analogue pour $j_{U, *}$.
\end{remark}




\section{2-flèches de topos}
\label{section-2-category}

\noindent
Voici une brève section consacrée à la notion de $2$-flèche
de topos.

\begin{definition}
\label{definition-2-morphism-topoi}
Soient $f, g : \Sh(\mathcal{C}) \to \Sh(\mathcal{D})$
deux morphismes de topos. Une {\it 2-flèche de $f$ vers $g$}
est donné par une transformation de foncteurs $t : f_* \to g_*$.
\end{definition}

\noindent
Nous représentons parfois $t$ par le diagramme suivant :
$$
\xymatrix{
\Sh(\mathcal{C})
\rrtwocell^f_g{t}
&
&
\Sh(\mathcal{D})
}
$$
Puisque $f^{-1}$ est adjoint à gauche de $f_*$ et que
$g^{-1}$ est adjoint à gauche de $g_*$, nous voyons que $t$ induit aussi
une transformation de foncteurs
$t : g^{-1} \to f^{-1}$ (notée généralement du même symbole),
caractérisée de manière unique par la commutativité du diagramme
$$
\xymatrix{
\Mor_{\Sh(\mathcal{D})}(\mathcal{G}, f_*\mathcal{F})
\ar[d]_{t \circ -} \ar@{=}[r] &
\Mor_{\Sh(\mathcal{C})}(f^{-1}\mathcal{G}, \mathcal{F})
\ar[d]^{- \circ t}
\\
\Mor_{\Sh(\mathcal{D})}(\mathcal{G}, g_*\mathcal{F})
\ar@{=}[r] &
\Mor_{\Sh(\mathcal{C})}(g^{-1}\mathcal{G}, \mathcal{F})
}
$$
En raison de difficultés ensemblistes (voir la
Remarque \ref{remark-morphism-topoi-big}),
nous n’obtenons pas une 2-catégorie de topos. Nous pouvons néanmoins définir
les compositions horizontale et verticale, et montrer que les axiomes d’une
2-catégorie stricte énumérés dans
Catégories, section \ref{categories-section-2-categories},
sont satisfaits.
En effet, la composition verticale des 2-flèches est claire (il suffit de composer
les transformations de foncteurs), la composition des 1-morphismes a été
définie dans la
Définition \ref{definition-topos},
et la composition horizontale de
$$
\xymatrix{
\Sh(\mathcal{C})
\rtwocell^f_g{t}
&
\Sh(\mathcal{D})
\rtwocell^{f'}_{g'}{s}
&
\Sh(\mathcal{E})
}
$$
est définie par la transformation de foncteurs $s \star t$
introduite dans
Catégories, Définition \ref{categories-definition-horizontal-composition}.
Explicitement, $s \star t$ est donnée par
$$
\xymatrix{
f'_*f_*\mathcal{F}
\ar[r]^{f'_*t} &
f'_*g_*\mathcal{F}
\ar[r]^s &
g'_*g_*\mathcal{F}
}
\quad\text{ou}\quad
\xymatrix{
f'_*f_*\mathcal{F}
\ar[r]^s &
g'_*f_*\mathcal{F}
\ar[r]^{g'_*t} &
g'_*g_*\mathcal{F}
}
$$
(ces applications sont égales). Comme ces définitions coïncident avec celles de
Catégories, section \ref{categories-section-formal-cat-cat},
il résulte de
Catégories, Lemme \ref{categories-lemma-properties-2-cat-cats},
que les axiomes d’une 2-catégorie stricte sont satisfaits avec ces définitions.




\section{Morphismes entre points}
\label{section-morphisms-points}

\begin{lemma}
\label{lemma-maps-u-points}
Soit $\mathcal{C}$ un site.
Soient $u, u' : \mathcal{C} \to \textit{Ensembles}$ deux
foncteurs, et soit $t : u' \to u$ une transformation de foncteurs.
Nous obtenons alors une transformation canonique de foncteurs fibres
$t_{stalk} : \mathcal{F}_{p'} \to \mathcal{F}_p$
qui coïncide avec $t$ par les identifications du
Lemme \ref{lemma-points-recover}.
\end{lemma}

\begin{proof}
Omis.
\end{proof}

\begin{definition}
\label{definition-morphism-points}
Soit $\mathcal{C}$ un site. Soient $p, p'$ des points de $\mathcal{C}$
donnés par des foncteurs $u, u' : \mathcal{C} \to \textit{Ensembles}$.
Un {\it morphisme $f : p \to p'$} est donné par une transformation de
foncteurs
$$
f_u : u' \to u.
$$
\end{definition}

\noindent
Remarquons que la transformation de foncteurs va dans le sens opposé.
Cela s’explique, comme nous le verrons plus loin, en considérant
le morphisme $f$ comme une sorte de $2$-flèche, représentée par le diagramme
suivant :
$$
\xymatrix{
\textit{Ensembles}
=
\Sh(pt)
\rrtwocell^p_{p'}{f}
&
&
\Sh(\mathcal{C})
}
$$
En effet, nous verrons plus loin que $f_u$ induit une transformation
canonique de foncteurs $p_* \to p'_*$ entre les
constructions de faisceaux gratte-ciel.

\medskip\noindent
Cette notion est assez importante et mérite ici un traitement plus complet.
Liste des compléments souhaitables :
\begin{enumerate}
\item Décrire les automorphismes du point de $\mathcal{T}_G$
décrit dans l’Exemple \ref{example-point-G-sets}.
\item Décrire $\Mor(p, p')$ au moyen de $\Mor(p_*, p'_*)$.
\item Spécialisation des points dans les espaces topologiques.
Montrer que, si $x' \in \overline{\{x\}}$ dans l’espace topologique
$X$, alors il existe un morphisme $p \to p'$, où $p$ (respectivement $p'$)
est le point de $X_{Zar}$ associé à $x$ (respectivement $x'$).
\end{enumerate}




\section{Sites ayant assez de points}
\label{section-sites-enough-points}

\begin{definition}
\label{definition-enough-points}
Soit $\mathcal{C}$ un site.
\begin{enumerate}
\item Une famille de points $\{p_i\}_{i\in I}$ est dite {\it conservative}
si toute application de faisceaux $\phi : \mathcal{F} \to \mathcal{G}$
qui est un isomorphisme sur toutes les fibres $\mathcal{F}_{p_i}
\to \mathcal{G}_{p_i}$ est un isomorphisme.
\item On dit que $\mathcal{C}$ {\it a assez de points}
s’il existe une famille conservative de points.
\end{enumerate}
\end{definition}

\noindent
On peut alors vérifier l’« exactitude » sur les fibres.

\begin{lemma}
\label{lemma-exactness-stalks}
Soit $\mathcal{C}$ un site, et soit $\{p_i\}_{i\in I}$ une famille conservative
de points. Alors :
\begin{enumerate}
\item Pour toute application de faisceaux $\varphi : \mathcal{F} \to \mathcal{G}$,
l’assertion « $\forall i, \varphi_{p_i}$ est injective » implique que $\varphi$ est injective.
\item Pour toute application de faisceaux $\varphi : \mathcal{F} \to \mathcal{G}$,
l’assertion « $\forall i, \varphi_{p_i}$ est surjective » implique que $\varphi$ est surjective.
\item Pour toute paire d’applications de faisceaux
$\varphi_1, \varphi_2 : \mathcal{F} \to \mathcal{G}$,
l’égalité $\forall i, \varphi_{1, p_i} = \varphi_{2, p_i}$
implique $\varphi_1 = \varphi_2$.
\item Soient un diagramme fini $\mathcal{G} : \mathcal{J}
\to \Sh(\mathcal{C})$, un faisceau $\mathcal{F}$ et des morphismes
$q_j : \mathcal{F} \to \mathcal{G}_j$. Alors $(\mathcal{F}, q_j)$
est une limite du diagramme si et seulement si, pour tout $i$, la fibre
$(\mathcal{F}_{p_i}, (q_j)_{p_i})$ en est une.
\item Soient un diagramme fini $\mathcal{F} : \mathcal{J}
\to \Sh(\mathcal{C})$, un faisceau $\mathcal{G}$ et des morphismes
$e_j : \mathcal{F}_j \to \mathcal{G}$. Alors $(\mathcal{G}, e_j)$
est une limite inductive du diagramme si et seulement si, pour tout $i$, la fibre
$(\mathcal{G}_{p_i}, (e_j)_{p_i})$ en est une.
\end{enumerate}
\end{lemma}

\begin{proof}
Nous utiliserons à plusieurs reprises le fait que tous les foncteurs fibres commutent
aux limites et limites inductives finies, donc aux produits, produits fibrés,
etc. Nous utiliserons aussi le fait que les applications injectives sont les
monomorphismes, et que les applications surjectives sont les épimorphismes.
Une application de faisceaux $\varphi : \mathcal{F} \to \mathcal{G}$
est injective si et seulement si
$\mathcal{F} \to \mathcal{F} \times_\mathcal{G} \mathcal{F}$
est un isomorphisme. D’où (1).
De même, $\varphi : \mathcal{F} \to \mathcal{G}$
est surjective si et seulement si
$\mathcal{G} \amalg_\mathcal{F} \mathcal{G} \to \mathcal{G}$
est un isomorphisme. D’où (2).
Les applications $a, b : \mathcal{F} \to \mathcal{G}$
sont égales si et seulement si
$\text{Égalisateur}(a, b : \mathcal{F} \to \mathcal{G}) \to \mathcal{F}$
est un isomorphisme\footnote{L’égalisateur s’exprime
au moyen de produits et de produits fibrés ; voir la démonstration de
Catégories, Lemme \ref{categories-lemma-finite-limits-exist}.}. D’où (3).
Les assertions (4) et (5) résultent immédiatement des définitions
et des remarques faites au début de cette démonstration.
\end{proof}

\begin{lemma}
\label{lemma-enough}
Soit $\mathcal{C}$ un site, et soit $\{(p_i, u_i)\}_{i\in I}$ une
famille de points. Cette famille est conservative si et seulement si, pour tout
faisceau $\mathcal{F}$, tout $U\in \Ob(\mathcal{C})$ et toute
paire de sections distinctes $s, s' \in \mathcal{F}(U)$, $s \not = s'$, il
existe $i$ et $x\in u_i(U)$ tels que les triplets
$(U, x, s)$ et $(U, x, s')$ définissent des éléments distincts de
$\mathcal{F}_{p_i}$.
\end{lemma}

\begin{proof}
Supposons la famille conservative, et soient $\mathcal{F}$, $U$ et
$s, s'$ comme dans le lemme. Les sections $s$, $s'$ définissent des applications
$a, a' : (h_U)^\# \to \mathcal{F}$ qui sont distinctes. D’après le Lemme
\ref{lemma-exactness-stalks}, il existe donc $i$ tel que $a_{p_i}
\not = a'_{p_i}$. Rappelons que $(h_U)^\#_{p_i} = u_i(U)$, d’après les
Lemmes \ref{lemma-points-recover} et
\ref{lemma-point-pushforward-sheaf}.
Il existe donc $x \in u_i(U)$ tel que $a_{p_i}(x)
\not = a'_{p_i}(x)$ dans $\mathcal{F}_{p_i}$.
En déroulant les définitions, on voit que $(U, x, s)$
et $(U, x, s')$ sont comme dans l’énoncé du lemme.

\medskip\noindent
Pour démontrer la réciproque, supposons satisfaite la condition d’existence de
points donnée dans le lemme. Soit $\phi : \mathcal{F} \to \mathcal{G}$
une application de faisceaux qui est un isomorphisme sur toutes les fibres.
Nous devons montrer que $\phi$ est à la fois
injective et surjective ; voir le Lemme \ref{lemma-mono-epi-sheaves}.
L’injectivité résulte immédiatement de l’hypothèse.
Pour montrer la surjectivité, il nous faut montrer que
$\mathcal{G} \amalg_\mathcal{F} \mathcal{G} \to \mathcal{G}$
est un isomorphisme
(Catégories, Lemme \ref{categories-lemma-characterize-mono-epi}).
Comme cette application est manifestement surjective, il suffit d’en vérifier l’injectivité,
qui résulte du fait que $\mathcal{G} \amalg_\mathcal{F} \mathcal{G} \to \mathcal{G}$
est injective sur toutes les fibres par hypothèse.
\end{proof}

\noindent
Dans le lemme suivant, les points $q_{i, x}$ sont exactement tous les points
de $\mathcal{C}/U$ situés au-dessus du point $p_i$, d’après le
Lemme \ref{lemma-points-above-point}.

\begin{lemma}
\label{lemma-localize-enough}
Soit $\mathcal{C}$ un site. Soit $U$ un objet de $\mathcal{C}$.
Soit $\{(p_i, u_i)\}_{i\in I}$ une famille de points de $\mathcal{C}$.
Pour $x \in u_i(U)$, soit $q_{i, x}$ le point de $\mathcal{C}/U$
construit dans le
Lemme \ref{lemma-point-localize}.
Si $\{p_i\}$ est une famille conservative de points, alors
$\{q_{i, x}\}_{i \in I, x \in u_i(U)}$ est une famille conservative
de points de $\mathcal{C}/U$.
En particulier, si $\mathcal{C}$ a assez de points, il en va de même
de toute localisation $\mathcal{C}/U$.
\end{lemma}

\begin{proof}
Nous savons que $j_{U!}$ induit une équivalence
$j_{U!} : \Sh(\mathcal{C}/U) \to \Sh(\mathcal{C})/h_U^\#$ ; voir le
Lemme \ref{lemma-essential-image-j-shriek}.
De plus, nous savons que
$(j_{U!}\mathcal{G})_{p_i} = \coprod_x \mathcal{G}_{q_{i, x}}$ ; voir le
Lemme \ref{lemma-stalk-j-shriek}.
Le résultat en découle donc formellement.
\end{proof}

\noindent
Le lemme suivant affirme que l’existence de points peut être vérifiée
localement sur le site.

\begin{lemma}
\label{lemma-enough-points-local}
Soit $\mathcal{C}$ un site. Soit $\{U_i\}_{i \in I}$ une famille
d’objets de $\mathcal{C}$. Supposons que
\begin{enumerate}
\item $\coprod h_{U_i}^\# \to *$ soit une application surjective de faisceaux, et
\item chaque localisation $\mathcal{C}/U_i$ ait assez de points.
\end{enumerate}
Alors $\mathcal{C}$ a assez de points.
\end{lemma}

\begin{proof}
Pour tout $i \in I$, soit $\{p_j\}_{j \in J_i}$ une famille conservative
de points de $\mathcal{C}/U_i$. Pour $j \in J_i$, notons
$q_j : \Sh(pt) \to \Sh(\mathcal{C})$ le composé
de $p_j$ avec le morphisme de localisation
$\Sh(\mathcal{C}/U_i) \to \Sh(\mathcal{C})$.
Alors $q_j$ est un point ; voir le
Lemme \ref{lemma-point-morphism-topoi}.
Nous affirmons que la famille de points $\{q_j\}_{j \in \coprod J_i}$
est conservative.
Soit en effet $\mathcal{F} \to \mathcal{G}$ une application de faisceaux
sur $\mathcal{C}$ telle que $\mathcal{F}_{q_j} \to \mathcal{G}_{q_j}$
soit un isomorphisme pour tout $j \in \coprod J_i$.
Soit $W$ un objet de $\mathcal{C}$.
D’après l’hypothèse (1), il existe un recouvrement $\{W_a \to W\}$ et
des morphismes $W_a \to U_{i(a)}$.
Comme $(\mathcal{F}|_{\mathcal{C}/U_{i(a)}})_{p_j} = \mathcal{F}_{q_j}$
et $(\mathcal{G}|_{\mathcal{C}/U_{i(a)}})_{p_j} = \mathcal{G}_{q_j}$ d’après le
Lemme \ref{lemma-point-morphism-topoi},
nous voyons que
$\mathcal{F}|_{U_{i(a)}} \to \mathcal{G}|_{U_{i(a)}}$ est un isomorphisme,
puisque la famille de points $\{p_j\}_{j \in J_{i(a)}}$ est conservative.
Ainsi, $\mathcal{F}(W_a) \to \mathcal{G}(W_a)$ est bijective pour tout $a$.
De même, $\mathcal{F}(W_a \times_W W_b) \to \mathcal{G}(W_a \times_W W_b)$
est bijective pour tous $a, b$.
La condition de faisceau montre alors que
$\mathcal{F}(W) \to \mathcal{G}(W)$ est bijective, c’est-à-dire que
$\mathcal{F} \to \mathcal{G}$ est un isomorphisme.
\end{proof}

\begin{lemma}
\label{lemma-check-morphism-sites}
Soit $u : \mathcal{C} \to \mathcal{D}$ un foncteur continu de sites.
Soit $\{(q_i, v_i)\}_{i\in I}$ une famille conservative de points de
$\mathcal{D}$. Si chaque foncteur $u_i = v_i \circ u$ définit
un point de $\mathcal{C}$,
alors $u$ définit un morphisme de sites $f : \mathcal{D} \to \mathcal{C}$.
\end{lemma}

\begin{proof}
Notons $p_i$ le foncteur fibre (\ref{equation-stalk}) sur
$\textit{PSh}(\mathcal{C})$ correspondant au foncteur $u_i$. Nous avons
$$
(f^{-1}\mathcal{F})_{q_i} =
(u_s\mathcal{F})_{q_i} =
(u_p\mathcal{F})_{q_i} =
\mathcal{F}_{p_i}
$$
La première égalité tient à $f^{-1} = u_s$, la deuxième au
Lemme \ref{lemma-point-pushforward-sheaf}, et la troisième
au Lemme \ref{lemma-point-functor}.
Ainsi, si $p_i$ est un point, alors prendre l’image inverse par $f$
puis les fibres en $q_i$ est un foncteur exact.
Comme la famille de points $\{q_i\}$ est conservative, cela
implique que $f^{-1}$ est un foncteur exact ; nous voyons donc
que $f$ est un morphisme de sites par la Définition \ref{definition-morphism-sites}.
\end{proof}





\section{Critère d’existence de points}
\label{section-criterion-points}

\noindent
Cette section correspond à l’appendice de Deligne à \cite[Expos\'e VI]{SGA4}.
Elle lui est en fait presque littéralement identique.

\medskip\noindent
Soit $\mathcal{C}$ un site.
Supposons que $(I, \geq)$ soit un ensemble ordonné filtrant,
et que $(U_i, f_{ii'})$ soit un système projectif sur $I$ ; voir
Catégories, Définition \ref{categories-definition-system-over-poset}.
Aux données $(I, \geq, U_i, f_{ii'})$, nous associons le foncteur
$$
u : \mathcal{C} \longrightarrow \textit{Ensembles}, \quad
u(V) = \colim_i \Mor_\mathcal{C}(U_i , V)
$$
Soit $\mathcal{F} \mapsto \mathcal{F}_p$ le foncteur fibre
associé à $u$ comme dans la section \ref{section-points}.
Il résulte directement de la définition que
$$
\mathcal{F}_p = \colim_i \mathcal{F}(U_i)
$$
dans ce cas particulier.
Remarquons que $u$ commute à toutes les limites finies (c’est-à-dire à celles qui
sont représentables dans $\mathcal{C}$), car chacun des foncteurs
$V \mapsto \Mor_\mathcal{C}(U_i , V)$ le fait ; voir
Catégories, Lemme \ref{categories-lemma-directed-commutes}.

\medskip\noindent
On dit qu’un système $(I, \geq, U_i, f_{ii'})$
est un {\it raffinement} de $(J, \geq, V_j, g_{jj'})$ si
$J \subset I$, si l’ordre sur $J$ est celui induit par l’ordre de $I$,
et si $V_j = U_j$, $g_{jj'} = f_{jj'}$ (autrement dit, si le système projectif
sur $J$ est induit par celui sur $I$). Soit $u$ le foncteur
associé à $(I, \geq, U_i, f_{ii'})$, et soit $u'$ le
foncteur associé à $(J, \geq, V_j, g_{jj'})$.
Cela induit une transformation de foncteurs
$$
u' \longrightarrow u
$$
simplement parce que les limites inductives définissant $u'$ portent sur un sous-système
des systèmes intervenant dans les limites inductives définissant $u$.
En particulier, nous obtenons une transformation associée des
foncteurs fibres $\mathcal{F}_{p'} \to \mathcal{F}_p$ ;
voir le Lemme \ref{lemma-maps-u-points}.

\begin{lemma}
\label{lemma-refine}
Soit $\mathcal{C}$ un site.
Soit $(J, \geq, V_j, g_{jj'})$ un système comme ci-dessus, muni de la
paire de foncteurs associée $(u', p')$.
Soit $\mathcal{F}$ un faisceau sur $\mathcal{C}$.
Soient $s, s' \in \mathcal{F}_{p'}$ des éléments distincts.
Soit $\{W_k \to W\}$ un recouvrement fini de $\mathcal{C}$.
Soit $f \in u'(W)$.
Il existe un raffinement $(I, \geq, U_i, f_{ii'})$
de $(J, \geq, V_j, g_{jj'})$ tel que les images de $s, s'$ dans
$\mathcal{F}_p$ soient distinctes et que
l’image de $f$ dans $u(W)$ appartienne à l’image de l’un des
$u(W_k)$.
\end{lemma}

\begin{proof}
Il existe $j_0 \in J$ tel que $f$ soit défini par $f' : V_{j_0} \to W$.
Pour $j \geq j_0$, posons $V_{j, k} = V_j \times_{f'\circ f_{j j_0}, W} W_k$.
Alors $\{V_{j, k} \to V_j\}$ est un recouvrement fini du site
$\mathcal{C}$. Par conséquent,
$\mathcal{F}(V_j) \subset \prod_k \mathcal{F}(V_{j, k})$.
En appliquant encore Catégories, Lemme \ref{categories-lemma-directed-commutes},
nous voyons que
$$
\mathcal{F}_{p'} =
\colim_j \mathcal{F}(V_j)
\longrightarrow
\prod\nolimits_k \colim_j \mathcal{F}(V_{j, k})
$$
est injective. Il existe donc $k$ tel que $s$ et $s'$ aient des images
distinctes dans $\colim_j \mathcal{F}(V_{j, k})$.
Posons $J_0 = \{j \in J, j \geq j_0\}$ et $I = J \amalg J_0$.
Ordonnons $I$ de sorte qu’aucun élément du second terme
ne soit inférieur à un élément du premier terme, et utilisons par ailleurs
l’ordre sur $J$. Si $j \in I$ appartient au premier
terme, nous prenons $V_j$ ; si $j \in I$ appartient au second terme,
nous prenons $V_{j, k}$. Nous omettons la définition
des morphismes de transition du système projectif. Il résulte de ce qui précède
que les images de $s, s'$ dans $\mathcal{F}_p$ sont distinctes.
En outre, la restriction de $f'$ à $V_{j, k}$ se factorise
par $W_k$ par construction.
\end{proof}

\begin{lemma}
\label{lemma-refine-all-at-once}
Soit $\mathcal{C}$ un site.
Soit $(J, \geq, V_j, g_{jj'})$ un système comme ci-dessus, muni de la
paire de foncteurs associée $(u', p')$.
Soit $\mathcal{F}$ un faisceau sur $\mathcal{C}$.
Soient $s, s' \in \mathcal{F}_{p'}$ des éléments distincts.
Il existe un raffinement $(I, \geq, U_i, f_{ii'})$
de $(J, \geq, V_j, g_{jj'})$ tel que les images de $s, s'$ dans
$\mathcal{F}_p$ soient distinctes et que,
pour tout recouvrement fini $\{W_k \to W\}$ du site
$\mathcal{C}$ et tout $f \in u'(W)$, l’image de $f$ dans $u(W)$
appartienne à l’image de l’un des $u(W_k)$.
\end{lemma}

\begin{proof}
Soit $E$ l’ensemble des couples $(\{W_k \to W\}, f\in u'(W))$.
Considérons les couples $(E' \subset E, (I, \geq, U_i, f_{ii'}))$
tels que
\begin{enumerate}
\item $(I, \geq, U_i, g_{ii'})$ soit un raffinement de $(J, \geq, V_j, g_{jj'})$,
\item les images de $s, s'$ dans $\mathcal{F}_p$ soient distinctes, et
\item pour tout couple $(\{W_k \to W\}, f\in u'(W)) \in E'$, l’image de
$f$ dans $u(W)$ appartienne à l’image de l’un des $u(W_k)$.
\end{enumerate}
Ordonnons ces couples par inclusion sur le premier facteur et
par raffinement sur le second. Notons $\mathcal{S}$ la classe
de tous les couples $(E' \subset E, (I, \geq, U_i, f_{ii'}))$ ci-dessus.
Nous affirmons que l’hypothèse du lemme de Zorn est satisfaite pour $\mathcal{S}$.
Supposons en effet que $(E'_a, (I_a, \geq, U_i, f_{ii'}))_{a \in A}$
soit une partie totalement ordonnée de $\mathcal{S}$. Nous pouvons définir
$E' = \bigcup_{a \in A} E'_a$ et poser $I = \bigcup_{a \in A} I_a$.
Nous affirmons que le couple correspondant
$(E' , (I, \geq, U_i, f_{ii'}))$ est un élément de $\mathcal{S}$.
Les conditions (1) et (3) sont claires. Pour la condition (2), remarquons
que
$$
u = \colim_{a \in A} u_a
\text{ et, par conséquent, }
\mathcal{F}_p = \colim_{a \in A} \mathcal{F}_{p_a}
$$
Le fait que les images de $s, s'$ dans cette fibre soient distinctes résulte
de la description d’une limite inductive filtrante d’ensembles ; voir
Catégories, section \ref{categories-section-directed-colimits}.
Nous écrirons simplement
$(E', (I, \ldots)) = \bigcup_{a \in A}(E'_a, (I_a, \ldots))$
dans cette situation.

\medskip\noindent
Le lemme de Zorn s’appliquerait donc si $\mathcal{S}$ était un ensemble,
et, combiné au Lemme \ref{lemma-refine} ci-dessus, il démontrerait aisément
le lemme. Mais $\mathcal{S}$ est une classe. Pour contourner cette difficulté,
choisissons un bon ordre sur $E$.
Pour $e \in E$, posons $E'_e = \{e' \in E \mid e' \leq e\}$.
Par récurrence transfinie, construisons des couples
$(E'_e, (I_e, \ldots)) \in \mathcal{S}$ tels que
$e_1 \leq e_2 \Rightarrow (E'_{e_1}, (I_{e_1}, \ldots))
\leq (E'_{e_2}, (I_{e_2}, \ldots))$.
Soit $e \in E$, disons $e = (\{W_k \to W\}, f\in u'(W))$.
Si $e$ possède un prédécesseur $e - 1$, prenons
$(I_e, \ldots)$ comme un raffinement de $(I_{e - 1}, \ldots)$
comme dans le Lemme \ref{lemma-refine}, relativement au système
$e = (\{W_k \to W\}, f\in u'(W))$.
Si $e$ ne possède pas de prédécesseur, prenons
$(I_e, \ldots)$ comme un raffinement de $\bigcup_{e' < e} (I_{e'}, \ldots)$
relativement au système
$e = (\{W_k \to W\}, f\in u'(W))$.
Enfin, la réunion $\bigcup_{e \in E} I_e$ est une solution au
problème posé dans le lemme.
\end{proof}

\begin{proposition}
\label{proposition-criterion-points}
\begin{reference}
\cite[Expos\'e VI, Appendice de Deligne, Proposition 9.0]{SGA4}
\end{reference}
Soit $\mathcal{C}$ un site. Supposons que
\begin{enumerate}
\item les limites finies existent dans $\mathcal{C}$, et que
\item tout recouvrement $\{U_i \to U\}_{i \in I}$
admette un raffinement par un recouvrement fini de $\mathcal{C}$.
\end{enumerate}
Alors $\mathcal{C}$ a assez de points.
\end{proposition}

\begin{proof}
Nous devons montrer que, pour tout faisceau
$\mathcal{F}$ sur $\mathcal{C}$, tout $U \in \Ob(\mathcal{C})$
et toutes sections distinctes $s, s' \in \mathcal{F}(U)$, il existe
un point $p$ tel que les images de $s, s'$ dans
$\mathcal{F}_p$ soient distinctes. Voir le Lemme \ref{lemma-enough}.
Considérons le système $(J, \geq, V_j, g_{jj'})$
défini par $J = \{1\}$, $V_1 = U$, $g_{11} = \text{id}_U$.
Appliquons le Lemme \ref{lemma-refine-all-at-once}.
Ce lemme fournit un système
$(I, \geq, U_i, f_{ii'})$ raffinant le nôtre de sorte que
$s_p \not = s'_p$ et que, de plus, pour tout
recouvrement fini $\{W_k \to W\}$ du site $\mathcal{C}$, l’application
$\coprod_k u(W_k) \to u(W)$ soit surjective.
Comme tout recouvrement de $\mathcal{C}$ admet un raffinement par
un recouvrement fini, nous en concluons que
$\coprod_k u(W_k) \to u(W)$ est surjective pour {\it tout}
recouvrement $\{W_k \to W\}$ du site $\mathcal{C}$.
Cela implique que $u = p$ est un point ; voir la
Proposition \ref{proposition-point-limits} (ainsi que la discussion
au début de cette section, qui garantit que $u$
commute aux limites finies).
\end{proof}

\begin{lemma}
\label{lemma-criterion-points}
Soit $\mathcal{C}$ un site. Soit $I$ un ensemble et, pour
$i \in I$, soit $U_i$ un objet de $\mathcal{C}$ tel que
\begin{enumerate}
\item $\coprod h_{U_i}$ se projette surjectivement sur
l’objet final de $\Sh(\mathcal{C})$, et
\item $\mathcal{C}/U_i$ satisfasse les hypothèses de la
Proposition \ref{proposition-criterion-points}.
\end{enumerate}
Alors $\mathcal{C}$ a assez de points.
\end{lemma}

\begin{proof}
D’après l’hypothèse (2) et la proposition, $\mathcal{C}/U_i$ a assez de points.
Les points de $\mathcal{C}/U_i$ donnent des points de $\mathcal{C}$
par le procédé du Lemme \ref{lemma-point-morphism-sites}.
Il suffit donc de montrer ceci : si $\phi : \mathcal{F} \to \mathcal{G}$
est une application de faisceaux sur $\mathcal{C}$ telle que $\phi|_{\mathcal{C}/U_i}$
soit un isomorphisme pour tout $i$, alors $\phi$ est un isomorphisme.
D’après l’hypothèse (1), pour tout objet $W$ de $\mathcal{C}$,
il existe un recouvrement $\{W_j \to W\}_{j \in J}$
tel que, pour $j \in J$, il existe $i \in I$ et un morphisme
$f_j : W_j \to U_i$. Les applications
$\mathcal{F}(W_j) \to \mathcal{G}(W_j)$
sont alors bijectives, de même que
$\mathcal{F}(W_j \times_W W_{j'}) \to \mathcal{G}(W_j \times_W W_{j'})$.
La condition de faisceau nous dit que $\mathcal{F}(W) \to \mathcal{G}(W)$
est bijective, comme voulu.
\end{proof}




\section{Objets faiblement contractiles}
\label{section-w-contractible}

\noindent
Un objet d’un site est dit {\it faiblement contractile} s’il satisfait
aux conditions équivalentes du lemme suivant.

\begin{lemma}
\label{lemma-w-contractible}
Soit $\mathcal{C}$ un site. Soit $U$ un objet de $\mathcal{C}$.
Les conditions suivantes sont équivalentes :
\begin{enumerate}
\item Pour toute famille couvrante $\{U_i \to U\}$, il existe un morphisme de
faisceaux $h_U^\# \to \coprod h_{U_i}^\#$ inverse à droite du morphisme de
faisceaux associé à $\coprod h_{U_i} \to h_U$.
\item Pour toute surjection de faisceaux d’ensembles $\mathcal{F} \to \mathcal{G}$,
l’application $\mathcal{F}(U) \to \mathcal{G}(U)$ est surjective.
\end{enumerate}
\end{lemma}

\begin{proof}
Supposons (1) et soit $\mathcal{F} \to \mathcal{G}$ un morphisme surjectif de
faisceaux d’ensembles. Pour $s \in \mathcal{G}(U)$, il existe une famille couvrante
$\{U_i \to U\}$ et des $t_i \in \mathcal{F}(U_i)$ dont l’image est
$s|_{U_i}$, voir la définition \ref{definition-sheaves-injective-surjective}.
Considérons $t_i$ comme un morphisme $t_i : h_{U_i}^\# \to \mathcal{F}$ grâce à
(\ref{equation-map-representable-into-sheaf}).
En précomposant alors $\coprod t_i : \coprod h_{U_i}^\# \to \mathcal{F}$
par le morphisme $h_U^\# \to \coprod h_{U_i}^\#$ fourni par (1),
nous obtenons une section $t \in \mathcal{F}(U)$ dont l’image est $s$.
Ainsi (2) est vérifiée.

\medskip\noindent
Supposons (2). Soit $\{U_i \to U\}$ une famille couvrante.
Alors $\coprod h_{U_i}^\# \to h_U^\#$ est surjectif
(lemme \ref{lemma-covering-surjective-after-sheafification}).
Il existe donc par (2) une section $s$ de $\coprod h_{U_i}^\#$ dont l’image
est la section $\text{id}_U$ de $h_U^\#$. Cette section correspond à un morphisme
$h_U^\# \to \coprod h_{U_i}^\#$ qui est inverse à droite du morphisme de
faisceaux associé à $\coprod h_{U_i} \to h_U$, ce qui prouve (1).
\end{proof}

\begin{definition}
\label{definition-w-contractible}
Soit $\mathcal{C}$ un site.
\begin{enumerate}
\item On dit qu’un objet $U$ de $\mathcal{C}$ est {\it faiblement contractile}
si les conditions équivalentes du lemme \ref{lemma-w-contractible} sont vérifiées.
\item On dit qu’un site possède {\it suffisamment d’objets faiblement contractiles}
si tout objet $U$ de $\mathcal{C}$ admet une famille couvrante $\{U_i \to U\}$
telle que $U_i$ soit faiblement contractile pour tout $i$.
\item Plus généralement, si $P$ est une propriété des objets de $\mathcal{C}$,
on dit que $\mathcal{C}$ possède {\it suffisamment d’objets vérifiant $P$} si tout objet $U$ de
$\mathcal{C}$ admet une famille couvrante $\{U_i \to U\}$ telle que $U_i$ vérifie $P$
pour tout $i$.
\end{enumerate}
\end{definition}

\noindent
Le petit site \'etale de $\mathbf{A}^1_\mathbf{C}$ ne possède aucun
objet faiblement contractile. En revanche, le petit site
pro-\'etale de tout schéma possède suffisamment d’objets contractiles.







\section{Propriétés d’exactitude de l’image directe}
\label{section-pushforward}

\noindent
Soit $f$ un morphisme de topos. En général, le foncteur $f_*$ n’est
qu’exact à gauche. On peut imposer à ce foncteur de nombreuses conditions
supplémentaires afin de distinguer certaines classes de morphismes de topos.
Nous les rassemblons ici et relevons quelques-unes de leurs implications logiques.
Certaines parties du lemme suivant sont de nature purement catégorique (c’est-à-dire
qu’elles ne dépendent pas de l’existence d’un morphisme de topos, l’existence d’une paire
de foncteurs adjoints suffit).

\begin{lemma}
\label{lemma-exactness-properties}
Soit $f : \Sh(\mathcal{C}) \to \Sh(\mathcal{D})$ un
morphisme de topos. Considérons les propriétés suivantes (pour les faisceaux
d’ensembles) :
\begin{enumerate}
\item $f_*$ est fidèle,
\item $f_*$ est pleinement fidèle,
\item $f^{-1}f_*\mathcal{F} \to \mathcal{F}$ est surjectif pour
tout $\mathcal{F}$ de $\Sh(\mathcal{C})$,
\item $f_*$ transforme les surjections en surjections,
\item $f_*$ commute aux coégalisateurs,
\item $f_*$ commute aux sommes amalgamées,
\item $f^{-1}f_*\mathcal{F} \to \mathcal{F}$ est un isomorphisme pour
tout $\mathcal{F}$ de $\Sh(\mathcal{C})$,
\item $f_*$ reflète les injections,
\item $f_*$ reflète les surjections,
\item $f_*$ reflète les bijections, et
\item pour toute surjection $\mathcal{F} \to f^{-1}\mathcal{G}$, il
existe une surjection $\mathcal{G}' \to \mathcal{G}$ telle que
$f^{-1}\mathcal{G}' \to f^{-1}\mathcal{G}$ se factorise par
$\mathcal{F} \to f^{-1}\mathcal{G}$.
\end{enumerate}
On a alors les implications suivantes :
\begin{enumerate}
\item[(a)] (2) $\Rightarrow$ (1),
\item[(b)] (3) $\Rightarrow$ (1),
\item[(c)] (7) $\Rightarrow$ (1), (2), (3), (8), (9), (10).
\item[(d)] (3) $\Leftrightarrow$ (9),
\item[(e)] (6) $\Rightarrow$ (4) et (5) $\Rightarrow$ (4),
\item[(f)] (4) $\Leftrightarrow$ (11),
\item[(g)] (9) $\Rightarrow$ (8), (10), et
\item[(h)] (2) $\Leftrightarrow$ (7).
\end{enumerate}
Le diagramme est le suivant :
$$
\xymatrix{
(6) \ar@{=>}[rd] & & & & & (9) \ar@{=>}[r] \ar@{=>}[rd] & (8) \\
& (4) \ar@{<=>}[r] & (11) &
(2) \ar@{<=>}[r] &
(7) \ar@{=>}[ru] \ar@{=>}[rd] & & (10) \\
(5) \ar@{=>}[ur] & & & & & (3) \ar@{=>}[r] & (1)
}
$$
\end{lemma}

\begin{proof}
Démonstration de (a) : cela résulte immédiatement des définitions.

\medskip\noindent
Démonstration de (b). Supposons que $a, b : \mathcal{F} \to \mathcal{F}'$ soient des
morphismes de faisceaux sur $\mathcal{C}$. Si $f_*a = f_*b$, alors
$f^{-1}f_*a = f^{-1}f_*b$. Considérons le diagramme commutatif
$$
\xymatrix{
\mathcal{F} \ar@<-1ex>[r] \ar@<1ex>[r] & \mathcal{F}' \\
f^{-1}f_*\mathcal{F} \ar@<-1ex>[r] \ar@<1ex>[r] \ar[u] &
f^{-1}f_*\mathcal{F}' \ar[u]
}
$$
Si les deux flèches inférieures sont égales et les flèches verticales surjectives,
alors les deux flèches supérieures sont égales. D’où (b).

\medskip\noindent
Démonstration de (c). Soit $a : \mathcal{F} \to \mathcal{F}'$ un
morphisme de faisceaux sur $\mathcal{C}$. Considérons le diagramme commutatif
$$
\xymatrix{
\mathcal{F} \ar[r] & \mathcal{F}' \\
f^{-1}f_*\mathcal{F} \ar[r] \ar[u] &
f^{-1}f_*\mathcal{F}' \ar[u]
}
$$
Si (7) est vérifiée, les flèches verticales sont des isomorphismes.
Ainsi, si $f_*a$ est injectif (resp.\ surjectif, resp.\ bijectif),
alors la flèche inférieure est injective (resp.\ surjective, resp.\ bijective), et
il en va donc de même pour la flèche supérieure.
On voit ainsi que (7) implique (8), (9), (10). Il est clair que (7) implique (3).
Les implications (7) $\Rightarrow$ (2), (1) résultent de (a) et (h), que
nous verrons ci-dessous.

\medskip\noindent
Démonstration de (d). Supposons (3). Soit $a : \mathcal{F} \to \mathcal{F}'$
un morphisme de faisceaux sur $\mathcal{C}$ tel que $f_*a$ soit surjectif.
Comme $f^{-1}$ est exact, il en résulte que
$f^{-1}f_*a : f^{-1}f_*\mathcal{F} \to f^{-1}f_*\mathcal{F}'$
est surjectif. Avec (3), cela implique que $a$ est surjectif.
Ainsi (9) est vérifiée.
Supposons (9). Soit $\mathcal{F}$ un faisceau sur $\mathcal{C}$.
Il faut montrer que le morphisme $f^{-1}f_*\mathcal{F} \to \mathcal{F}$ est
surjectif. Il suffit de montrer que
$f_*f^{-1}f_*\mathcal{F} \to f_*\mathcal{F}$ est surjectif.
Or cela est vrai puisqu’il existe un morphisme canonique
$f_*\mathcal{F} \to f_*f^{-1}f_*\mathcal{F}$ qui est un inverse d’un côté de ce morphisme.

\medskip\noindent
Démonstration de (e). Nous utilisons
Catégories, lemme \ref{categories-lemma-characterize-mono-epi}
sans autre mention.
Si $\mathcal{F} \to \mathcal{F}'$ est surjectif, alors
$\mathcal{F}' \amalg_\mathcal{F} \mathcal{F}' \to \mathcal{F}'$
est un isomorphisme. Par conséquent, (6) implique que
$$
f_*\mathcal{F}' \amalg_{f_*\mathcal{F}} f_*\mathcal{F}' =
f_*(\mathcal{F}' \amalg_\mathcal{F} \mathcal{F}')
\longrightarrow
f_*\mathcal{F}'
$$
est également un isomorphisme. Cela implique à son tour que
$f_*\mathcal{F} \to f_*\mathcal{F}'$ est surjectif.
On voit donc que (6) implique (4). Si $\mathcal{F} \to \mathcal{F}'$
est surjectif, alors $\mathcal{F}'$ est le coégalisateur des deux
projections $\mathcal{F} \times_{\mathcal{F}'} \mathcal{F} \to \mathcal{F}$
d’après le lemme \ref{lemma-coequalizer-surjection}.
Par conséquent, si (5) est vérifiée, alors $f_*\mathcal{F}'$ est le coégalisateur
des deux projections
$$
f_*(\mathcal{F} \times_{\mathcal{F}'} \mathcal{F}) =
f_*\mathcal{F} \times_{f_*\mathcal{F}'} f_*\mathcal{F}
\longrightarrow
f_*\mathcal{F}
$$
ce qui signifie clairement que $f_*\mathcal{F} \to f_*\mathcal{F}'$ est
surjectif. Ainsi (5) implique également (4).

\medskip\noindent
Démonstration de (f). Supposons (4). Soit $\mathcal{F} \to f^{-1}\mathcal{G}$
un morphisme surjectif de faisceaux sur $\mathcal{C}$. Par (4), on voit que
$f_*\mathcal{F} \to f_*f^{-1}\mathcal{G}$ est surjectif. Soit
$\mathcal{G}'$ le produit fibré
$$
\xymatrix{
f_*\mathcal{F} \ar[r] & f_*f^{-1}\mathcal{G} \\
\mathcal{G}' \ar[u] \ar[r] & \mathcal{G} \ar[u]
}
$$
de sorte que $\mathcal{G}' \to \mathcal{G}$ est lui aussi surjectif. Considérons
le diagramme commutatif
$$
\xymatrix{
\mathcal{F} \ar[r] & f^{-1}\mathcal{G} \\
f^{-1}f_*\mathcal{F} \ar[r] \ar[u] & f^{-1}f_*f^{-1}\mathcal{G} \ar[u] \\
f^{-1}\mathcal{G}' \ar[u] \ar[r] & f^{-1}\mathcal{G} \ar[u]
}
$$
qui donne le résultat voulu. Réciproquement, supposons (11).
Soit $a : \mathcal{F} \to \mathcal{F}'$ un morphisme surjectif de faisceaux
sur $\mathcal{C}$. Considérons le diagramme cartésien
$$
\xymatrix{
\mathcal{F} \ar[r] & \mathcal{F}' \\
\mathcal{F}'' \ar[u] \ar[r] & f^{-1}f_*\mathcal{F}' \ar[u]
}
$$
Comme la flèche horizontale inférieure est surjective, (11) permet de
trouver une surjection $\gamma : \mathcal{G}' \to f_*\mathcal{F}'$ telle
que $f^{-1}\gamma$ se factorise par $\mathcal{F}'' \to f^{-1}f_*\mathcal{F}'$ :
$$
\xymatrix{
& \mathcal{F} \ar[r] & \mathcal{F}' \\
f^{-1}\mathcal{G}' \ar[r] & \mathcal{F}'' \ar[u] \ar[r] &
f^{-1}f_*\mathcal{F}' \ar[u]
}
$$
En appliquant $f_*$, on obtient un diagramme commutatif
$$
\xymatrix{
& f_*\mathcal{F} \ar[r] & f_*\mathcal{F}' \\
f_*f^{-1}\mathcal{G}' \ar[r] & f_*\mathcal{F}'' \ar[u] \ar[r] &
f_*f^{-1}f_*\mathcal{F}' \ar[u] \\
\mathcal{G}' \ar[u] \ar[rr] & & f_*\mathcal{F}' \ar[u]
}
$$
ce qui prouve (4).

\medskip\noindent
Démonstration de (g). Supposons (9). Nous utilisons
Catégories, lemme \ref{categories-lemma-characterize-mono-epi}
sans autre mention. Soit $a : \mathcal{F} \to \mathcal{F}'$ un morphisme
de faisceaux sur $\mathcal{C}$ tel que $f_*a$ soit injectif. Cela signifie que
$f_*\mathcal{F} \to f_*\mathcal{F} \times_{f_*\mathcal{F}'} f_*\mathcal{F} =
f_*(\mathcal{F} \times_{\mathcal{F}'} \mathcal{F})$ est un isomorphisme.
Ainsi, par (9), on voit que
$\mathcal{F} \to \mathcal{F} \times_{\mathcal{F}'} \mathcal{F}$ est
surjectif, c’est-à-dire un isomorphisme. Ainsi $a$ est injectif, c’est-à-dire que (8) est vérifiée.
Puisque (10) est trivialement équivalente à (8) $+$ (9), cela achève la démonstration de (g).

\medskip\noindent
Démonstration de (h). C’est
Catégories, lemme \ref{categories-lemma-adjoint-fully-faithful}.
\end{proof}

\noindent
Voici une condition portant sur un morphisme de sites qui garantit que
le foncteur $f_*$ transforme les morphismes surjectifs en morphismes surjectifs.

\begin{lemma}
\label{lemma-weaker}
Soit $f : \mathcal{D} \to \mathcal{C}$ un morphisme de sites associé au
foncteur continu $u : \mathcal{C} \to \mathcal{D}$.
Supposons que, pour tout objet $U$ de $\mathcal{C}$ et toute famille couvrante
$\{V_j \to u(U)\}$ de $\mathcal{D}$, il existe une famille couvrante $\{U_i \to U\}$
de $\mathcal{C}$ telle que le morphisme de faisceaux
$$
\coprod h_{u(U_i)}^\# \to h_{u(U)}^\#
$$
se factorise par le morphisme de faisceaux
$$
\coprod h_{V_j}^\# \to h_{u(U)}^\#.
$$
Alors $f_*$ transforme les morphismes surjectifs de faisceaux en morphismes surjectifs de
faisceaux.
\end{lemma}

\begin{proof}
Soit $a : \mathcal{F} \to \mathcal{G}$ un morphisme surjectif de faisceaux sur
$\mathcal{D}$. Soit $U$ un objet de $\mathcal{C}$ et soit
$s \in f_*\mathcal{G}(U) = \mathcal{G}(u(U))$. Par hypothèse, il existe
une famille couvrante $\{V_j \to u(U)\}$ et des sections $s_j \in \mathcal{F}(V_j)$
telles que $a(s_j) = s|_{V_j}$. On peut maintenant considérer les sections $s$,
$s_j$ et $a$ comme définissant un diagramme commutatif de morphismes de faisceaux
$$
\xymatrix{
\coprod h_{V_j}^\# \ar[r]_-{\coprod s_j} \ar[d] & \mathcal{F} \ar[d]^a \\
h_{u(U)}^\# \ar[r]^s & \mathcal{G}
}
$$
Par hypothèse, il existe une famille couvrante $\{U_i \to U\}$ qui permet
d’agrandir comme suit le diagramme commutatif ci-dessus :
$$
\xymatrix{
& \coprod h_{V_j}^\# \ar[r]_-{\coprod s_j} \ar[d] & \mathcal{F} \ar[d]^a \\
\coprod h_{u(U_i)}^\# \ar[r] \ar[ur] &
h_{u(U)}^\# \ar[r]^s & \mathcal{G}
}
$$
Comme $\mathcal{F}$ est un faisceau, le morphisme du coin inférieur gauche vers
le coin supérieur droit correspond à une famille de sections
$s_i \in \mathcal{F}(u(U_i))$, c’est-à-dire à des sections $s_i \in f_*\mathcal{F}(U_i)$.
La commutativité du diagramme implique que $a(s_i)$ est égal à la
restriction de $s$ à $U_i$. Autrement dit, nous avons montré que $f_*a$ est un
morphisme surjectif de faisceaux.
\end{proof}

\begin{example}
\label{example-weaker-not-coequalizer}
Supposons que $f : \mathcal{D} \to \mathcal{C}$ satisfasse aux hypothèses du
lemme \ref{lemma-weaker}.
En général, $f_*$ ne commute alors ni aux coégalisateurs
ni aux sommes amalgamées. En effet, supposons que $f$ soit le morphisme de sites associé
au morphisme d’espaces topologiques $X = \{1, 2\} \to Y = \{*\}$ (voir
l’exemple \ref{example-continuous-map}),
où $Y$ est un espace réduit à un point et $X = \{1, 2\}$ est un espace discret
à deux points. Un faisceau $\mathcal{F}$ sur $X$ est donné par une paire
$(A_1, A_2)$ d’ensembles. Alors $f_*\mathcal{F}$ correspond à
l’ensemble $A_1 \times A_2$. Par conséquent, si
$a = (a_1, a_2), b = (b_1, b_2) : (A_1, A_2) \to (B_1, B_2)$
sont des morphismes de faisceaux sur $X$, alors le
coégalisateur de $a, b$ est $(C_1, C_2)$, où $C_i$ est le coégalisateur de
$a_i, b_i$, tandis que le coégalisateur de $f_*a, f_*b$ est le coégalisateur de
$$
a_1 \times a_2, b_1 \times b_2 :
A_1 \times A_2 \longrightarrow B_1 \times B_2
$$
qui est en général différent de $C_1 \times C_2$. En effet, si
$A_2 = \emptyset$, alors $A_1 \times A_2 = \emptyset$, et par conséquent le
coégalisateur des flèches affichées est $B_1 \times B_2$, mais en général
$C_1 \not = B_1$. Un exemple analogue vaut pour les sommes amalgamées.
\end{example}

\noindent
Le lemme suivant donne un critère assurant que, pour un morphisme de sites,
le foncteur $f_*$ reflète les injections et les surjections.
Remarquons que cela implique aussi que $f_*$ est fidèle et que le
morphisme $f^{-1}f_*\mathcal{F} \to \mathcal{F}$ est toujours surjectif.

\begin{lemma}
\label{lemma-cover-from-below}
Soit $f : \mathcal{D} \to \mathcal{C}$ un morphisme de sites défini
par le foncteur $u : \mathcal{C} \to \mathcal{D}$.
Supposons que, pour tout objet $V$ de $\mathcal{D}$, il existe des objets
$U_i$ de $\mathcal{C}$ et des morphismes $u(U_i) \to V$ tels que
$\{u(U_i) \to V\}$ soit une famille couvrante de $\mathcal{D}$. Dans ce cas, le foncteur
$f_* : \Sh(\mathcal{D}) \to \Sh(\mathcal{C})$ reflète les
injections et les surjections.
\end{lemma}

\begin{proof}
Soit $\alpha : \mathcal{F} \to \mathcal{G}$ un morphisme de faisceaux
sur $\mathcal{D}$. Par hypothèse, pour tout objet $V$ de $\mathcal{D}$,
on a $\mathcal{F}(V) \subset \prod \mathcal{F}(u(U_i)) =
\prod f_*\mathcal{F}(U_i)$ en vertu de la condition de faisceau,
pour certains $U_i \in \Ob(\mathcal{C})$, et de même pour $\mathcal{G}$.
Il est donc clair que, si $f_*\alpha$ est injectif, alors
$\alpha$ est injectif. Autrement dit, $f_*$ reflète les injections.

\medskip\noindent
Supposons que $f_*\alpha$ soit surjectif. Pour $V, U_i, u(U_i) \to V$
comme ci-dessus et une section $s \in \mathcal{G}(V)$, il existe alors des familles couvrantes
$\{U_{ij} \to U_i\}$ telles que $s|_{u(U_{ij})}$ appartienne à l’image
de $\mathcal{F}(u(U_{ij}))$. Comme $\{u(U_{ij}) \to V\}$ est une famille couvrante
(puisque $u$ est continu et d’après les axiomes d’un site), on conclut que
$s$ appartient localement à l’image. Ainsi $\alpha$ est surjectif. Autrement
dit, $f_*$ reflète les surjections.
\end{proof}

\begin{example}
\label{example-cover-from-below}
Construisons un morphisme $f : \mathcal{D} \to \mathcal{C}$ qui satisfait
aux hypothèses du lemme \ref{lemma-cover-from-below}. Soit donc
$\varphi : G \to H$ un morphisme de groupes finis. Considérons les sites
$\mathcal{D} = \mathcal{T}_G$ et $\mathcal{C} = \mathcal{T}_H$ des $G$-ensembles et
$H$-ensembles dénombrables, dont les familles couvrantes sont les familles dénombrables de
morphismes conjointement surjectifs (exemple \ref{example-site-on-group}).
Soit $u : \mathcal{T}_H \to \mathcal{T}_G$ le foncteur décrit
dans la section \ref{section-G-sets-morphisms},
et $f : \mathcal{T}_G \to \mathcal{T}_H$ le
morphisme de sites correspondant. Si $\varphi$ est injectif, alors
tout $G$-ensemble dénombrable est, en tant que $G$-ensemble, le quotient d’un
$H$-ensemble dénombrable (cela est faux si $\varphi$ n’est pas injectif). Ainsi $f$ satisfait
à l’hypothèse du lemme \ref{lemma-cover-from-below}.
Si le faisceau $\mathcal{F}$ sur $\mathcal{T}_G$ correspond au $G$-ensemble
$S$, alors le morphisme canonique
$$
f^{-1}f_*\mathcal{F} \longrightarrow \mathcal{F}
$$
correspond à l’application
$$
\text{Map}_G(H, S) \longrightarrow S,\quad
a \longmapsto a(1_H)
$$
Si $\varphi$ est injectif mais non surjectif, cette application est surjective
(comme il se doit d’après le lemme \ref{lemma-cover-from-below}),
mais elle n’est pas injective en général
(par exemple, prendre $G = \{1\}$, $H = \{1, \sigma\}$ et $S = \{1, 2\}$).
De plus, le foncteur $f_*$ ne commute ni aux coégalisateurs ni aux
sommes amalgamées (pour $G = \{1\}$ et $H = \{1, \sigma\}$).
\end{example}





\section{Foncteurs presque cocontinus}
\label{section-almost-cocontinuous}

\noindent
Soit $\mathcal{C}$ un site. La catégorie $\textit{PSh}(\mathcal{C})$ possède
un objet initial, à savoir le préfaisceau qui associe l’ensemble vide
à chaque objet de $\mathcal{C}$. Notons ce préfaisceau
$\emptyset$. Il résulte des propriétés du passage au faisceau associé que le
faisceau $\emptyset^\#$ associé à $\emptyset$ est un objet initial
de la catégorie $\Sh(\mathcal{C})$ des faisceaux sur $\mathcal{C}$.

\begin{definition}
\label{definition-empty}
Soit $\mathcal{C}$ un site. On dit qu’un objet $U$ de $\mathcal{C}$
est {\it vide au sens des faisceaux} si $\emptyset^\# \to h_U^\#$
est un isomorphisme de faisceaux.
\end{definition}

\noindent
Le lemme suivant explicite cette notion.

\begin{lemma}
\label{lemma-characterize-empty}
Soit $\mathcal{C}$ un site. Soit $U$ un objet de $\mathcal{C}$.
Les conditions suivantes sont équivalentes :
\begin{enumerate}
\item $U$ est vide au sens des faisceaux,
\item $\mathcal{F}(U)$ est un ensemble à un élément pour tout faisceau $\mathcal{F}$,
\item $\emptyset^\#(U)$ est un ensemble à un élément,
\item $\emptyset^\#(U)$ est non vide, et
\item la famille vide est une famille couvrante de $U$ dans $\mathcal{C}$.
\end{enumerate}
De plus, si $U$ est vide au sens des faisceaux, alors, pour tout morphisme
$U' \to U$ de $\mathcal{C}$, l’objet $U'$ est vide au sens des faisceaux.
\end{lemma}

\begin{proof}
Pour tout faisceau $\mathcal{F}$, on a
$\mathcal{F}(U) = \Mor_{\Sh(\mathcal{C})}(h_U^\#, \mathcal{F})$.
On voit donc que (1) et (2) sont équivalentes.
Il est clair que (2) implique (3), qui implique (4).
Si toute famille couvrante de $U$ est une famille non vide,
alors $\emptyset^+(U)$ est vide par définition de la construction plus.
Remarquons que $\emptyset^+ = \emptyset^\#$, puisque $\emptyset$ est un
préfaisceau séparé, voir le
théorème \ref{theorem-plus}.
On voit ainsi que (4) implique (5). Si (5) est vérifiée, alors
$\mathcal{F}(U)$ est un ensemble à un élément pour tout faisceau $\mathcal{F}$
en vertu de la condition de faisceau pour $\mathcal{F}$, voir la
remarque \ref{remark-sheaf-condition-empty-covering}.
Ainsi (5) implique (2), et (1) -- (5) sont équivalentes. La dernière
assertion du lemme résulte de l’axiome (3) de la
définition \ref{definition-site}
appliqué à la famille couvrante vide de $U$.
\end{proof}

\begin{definition}
\label{definition-almost-cocontinuous}
Soient $\mathcal{C}$ et $\mathcal{D}$ des sites.
Soit $u : \mathcal{C} \to \mathcal{D}$ un foncteur.
On dit que $u$ est {\it presque cocontinu} si, pour tout
objet $U$ de $\mathcal{C}$ et toute famille couvrante
$\{V_j \to u(U)\}_{j \in J}$, il existe une famille couvrante
$\{U_i \to U\}_{i \in I}$ dans $\mathcal{C}$ telle que,
pour tout $i$ dans $I$, l’une au moins des deux conditions suivantes soit vérifiée :
\begin{enumerate}
\item $u(U_i)$ est vide au sens des faisceaux, ou
\item le morphisme $u(U_i) \to u(U)$ se factorise par $V_j$ pour un certain $j \in J$.
\end{enumerate}
\end{definition}

\noindent
La motivation de cette définition provient d’une immersion fermée
$i : Z \to X$ d’espaces topologiques. Comme on l’a vu dans
l’exemple \ref{example-closed-map-cocontinuous-false},
le foncteur continu
$X_{Zar} \to Z_{Zar}$, $U \mapsto Z \cap U$, n’est
pas cocontinu. Mais il est presque cocontinu au sens défini ci-dessus.
Nous savons que $i_*$, bien qu’il ne soit pas exact sur les faisceaux d’ensembles, est
exact sur les faisceaux de groupes abéliens, voir
Faisceaux, remarque \ref{sheaves-remark-i-star-not-exact}.
Et cela vaut en général pour les foncteurs continus et presque cocontinus.

\begin{lemma}
\label{lemma-almost-cocontinuous-sheafification}
Soient $\mathcal{C}$ et $\mathcal{D}$ des sites.
Soit $u : \mathcal{C} \to \mathcal{D}$ un foncteur.
Supposons que $u$ soit continu et presque cocontinu.
Soit $\mathcal{G}$ un préfaisceau sur $\mathcal{D}$ tel que $\mathcal{G}(V)$ soit
un ensemble à un élément dès que $V$ est vide au sens des faisceaux.
Alors $(u^p\mathcal{G})^\# = u^p(\mathcal{G}^\#)$.
\end{lemma}

\begin{proof}
Soit $U \in \Ob(\mathcal{C})$. Il faut montrer que
$(u^p\mathcal{G})^\#(U) = u^p(\mathcal{G}^\#)(U)$.
Il suffit de montrer que
$(u^p\mathcal{G})^+(U) = u^p(\mathcal{G}^+)(U)$,
puisque $\mathcal{G}^+$ est un autre préfaisceau qui satisfait à
l’hypothèse du lemme. On a
$$
u^p(\mathcal{G}^+)(U) =
\mathcal{G}^+(u(U)) =
\colim_\mathcal{V} \check H^0(\mathcal{V}, \mathcal{G})
$$
où la limite inductive porte sur les familles couvrantes $\mathcal{V}$ de $u(U)$ dans
$\mathcal{D}$. D’autre part, on a
$$
u^p(\mathcal{G})^+(U) =
\colim_\mathcal{U} \check H^0(u(\mathcal{U}), \mathcal{G})
$$
où la limite inductive porte sur la catégorie des familles couvrantes
$\mathcal{U} = \{U_i \to U\}_{i \in I}$ de $U$ dans $\mathcal{C}$, et
$u(\mathcal{U}) = \{u(U_i) \to u(U)\}_{i \in I}$. La continuité
de $u$ signifie que chaque $u(\mathcal{U})$ est une famille couvrante.
Écrivons $I = I_1 \amalg I_2$, où
$$
I_2 = \{i \in I \mid u(U_i)\text{ est vide au sens des faisceaux}\}
$$
Alors $u(\mathcal{U})' = \{u(U_i) \to u(U)\}_{i \in I_1}$ est encore une famille couvrante de
car chacune des autres composantes peut être recouverte par la famille vide
et peut donc être omise en vertu de l’axiome (2) de la
définition \ref{definition-site}.
En outre,
$\check H^0(u(\mathcal{U}), \mathcal{G}) =
\check H^0(u(\mathcal{U})', \mathcal{G})$
par notre hypothèse sur $\mathcal{G}$. Enfin, la presque cocontinuité de $u$
implique que, pour toute famille couvrante $\mathcal{V}$
de $u(U)$, il existe une famille couvrante $\mathcal{U}$ de $U$ telle que
$u(\mathcal{U})'$ raffine $\mathcal{V}$. Il s’ensuit que les deux limites inductives
affichées ci-dessus ont la même valeur, comme souhaité.
\end{proof}

\begin{lemma}
\label{lemma-continuous-almost-cocontinuous}
Soient $\mathcal{C}$ et $\mathcal{D}$ des sites.
Soit $u : \mathcal{C} \to \mathcal{D}$ un foncteur.
Supposons que $u$ soit continu et presque cocontinu.
Alors $u^s = u^p : \Sh(\mathcal{D}) \to \Sh(\mathcal{C})$
commute aux sommes amalgamées et aux coégalisateurs (et, plus généralement,
aux limites inductives connexes finies).
\end{lemma}

\begin{proof}
Soit $\mathcal{I}$ une catégorie d’indices finie et connexe.
Soit $\mathcal{I} \to \Sh(\mathcal{D})$,
$i \mapsto \mathcal{G}_i$ un diagramme. On sait que la limite inductive de
ce diagramme est le faisceau associé à la limite inductive dans la catégorie des
préfaisceaux, voir le
lemme \ref{lemma-colimit-sheaves}.
Notons $\colim^{Psh}$ la limite inductive dans la catégorie
des préfaisceaux. Puisque $\mathcal{I}$ est finie et connexe,
on voit que $\colim^{Psh}_i \mathcal{G}_i$
est un préfaisceau satisfaisant aux hypothèses du
lemme \ref{lemma-almost-cocontinuous-sheafification}
(car une limite inductive finie connexe d’ensembles à un élément est un
ensemble à un élément). Ce lemme donne donc
\begin{align*}
u^s(\colim_i \mathcal{G}_i) & =
u^s((\colim^{Psh}_i \mathcal{G}_i)^\#) \\
& = (u^p(\colim^{Psh}_i \mathcal{G}_i))^\# \\
& = (\colim^{PSh}_i u^p(\mathcal{G}_i))^\# \\
& = \colim_i u^s(\mathcal{G}_i)
\end{align*}
comme voulu.
\end{proof}

\begin{lemma}
\label{lemma-morphism-of-sites-almost-cocontinuous}
Soit $f : \mathcal{D} \to \mathcal{C}$ un morphisme de sites
associé au foncteur continu $u : \mathcal{C} \to \mathcal{D}$.
Si $u$ est presque cocontinu, alors $f_*$ commute aux
sommes amalgamées et aux coégalisateurs (et, plus généralement, aux limites inductives connexes finies).
\end{lemma}

\begin{proof}
C’est un cas particulier du lemme \ref{lemma-continuous-almost-cocontinuous}.
\end{proof}








\section{Sous-topos}
\label{section-subtopoi}

\noindent
Voici la définition.

\begin{definition}
\label{definition-embedding}
Soient $\mathcal{C}$ et $\mathcal{D}$ des sites.
Un morphisme de topos $f : \Sh(\mathcal{D}) \to \Sh(\mathcal{C})$
est appelé un {\it plongement} si $f_*$ est pleinement fidèle.
\end{definition}

\noindent
D'après le lemme \ref{lemma-exactness-properties},
cela équivaut à demander que le
morphisme d'adjonction $f^{-1}f_*\mathcal{F} \to \mathcal{F}$
soit un isomorphisme pour tout faisceau $\mathcal{F}$ sur $\mathcal{D}$.

\begin{definition}
\label{definition-subtopos}
Soit $\mathcal{C}$ un site. Une sous-catégorie strictement pleine
$E \subset \Sh(\mathcal{C})$ est un {\it sous-topos} s'il
existe un plongement de topos $f : \Sh(\mathcal{D}) \to \Sh(\mathcal{C})$
tel que $E$ soit égale à l'image essentielle du foncteur $f_*$.
\end{definition}

\noindent
Les sous-topos construits dans le lemme suivant seront qualifiés
d'« ouverts » dans la définition donnée plus loin.

\begin{lemma}
\label{lemma-open-subtopos}
Soit $\mathcal{C}$ un site. Soit $\mathcal{F}$ un faisceau sur
$\mathcal{C}$. Les assertions suivantes sont équivalentes :
\begin{enumerate}
\item $\mathcal{F}$ est un sous-objet de l'objet final de
$\Sh(\mathcal{C})$, et
\item le topos $\Sh(\mathcal{C})/\mathcal{F}$ est un sous-topos de
$\Sh(\mathcal{C})$.
\end{enumerate}
\end{lemma}

\begin{proof}
Nous avons vu dans le lemme \ref{lemma-localize-topos} que
$\Sh(\mathcal{C})/\mathcal{F}$ est un topos. Rappelons en fait la
démonstration. Nous appliquons d'abord le lemme \ref{lemma-topos-good-site}
pour voir que nous pouvons supposer que $\mathcal{C}$ est un site à topologie
sous-canonique, possédant les produits fibrés, un objet final $X$ et un objet $U$ tel que
$\mathcal{F} = h_U$. La démonstration du
lemme \ref{lemma-localize-topos}
montre que le morphisme de topos
$j_\mathcal{F} : \Sh(\mathcal{C})/\mathcal{F} \to \Sh(\mathcal{C})$
est égal (à certaines identifications près) au morphisme de localisation
$j_U : \Sh(\mathcal{C}/U) \to \Sh(\mathcal{C})$.

\medskip\noindent
Supposons (2). Cela signifie que $j_U^{-1}j_{U, *}\mathcal{G} \to \mathcal{G}$
est un isomorphisme pour tout faisceau $\mathcal{G}$ sur $\mathcal{C}/U$.
Pour tout objet $Z/U$ de $\mathcal{C}/U$, nous avons
$$
(j_{U, *}h_{Z/U})(U) = \Mor_{\mathcal{C}/U}(U \times_X U/U, Z/U)
$$
d'après le lemme \ref{lemma-localize-given-products}.
En posant $\mathcal{G} = h_{Z/U}$ dans l'égalité ci-dessus, nous obtenons
$$
\Mor_{\mathcal{C}/U}(U \times_X U/U, Z/U) = \Mor_{\mathcal{C}/U}(U, Z/U)
$$
pour tout $Z/U$. D'après le lemme de Yoneda
(Catégories, Lemme \ref{categories-lemma-yoneda}),
cela implique $U \times_X U = U$. D'après
Catégories, Lemme \ref{categories-lemma-characterize-mono-epi},
$U \to X$ est un monomorphisme, autrement dit (1) est satisfaite.

\medskip\noindent
Supposons (1). Alors $j_U^{-1} j_{U, *} = \text{id}$ d'après le
lemme \ref{lemma-restrict-back}.
\end{proof}

\begin{definition}
\label{definition-open-subtopos}
Soit $\mathcal{C}$ un site. Une sous-catégorie strictement pleine
$E \subset \Sh(\mathcal{C})$ est un {\it sous-topos ouvert}
s'il existe un sous-faisceau $\mathcal{F}$ de l'objet final
de $\Sh(\mathcal{C})$ tel que $E$ soit le sous-topos
$\Sh(\mathcal{C})/\mathcal{F}$ décrit dans le lemme \ref{lemma-open-subtopos}.
\end{definition}

\noindent
Cela signifie qu'il existe une bijection entre la collection des sous-topos ouverts
de $\Sh(\mathcal{C})$ et l'ensemble des sous-objets de l'objet final de
$\Sh(\mathcal{C})$. À tout sous-topos ouvert correspond un complémentaire « fermé ».

\begin{lemma}
\label{lemma-closed-subtopos}
Soit $\mathcal{C}$ un site. Soit $\mathcal{F}$ un sous-faisceau de l'objet final
$*$ de $\Sh(\mathcal{C})$. La sous-catégorie pleine des faisceaux
$\mathcal{G}$ tels que $\mathcal{F} \times \mathcal{G} \to \mathcal{F}$
soit un isomorphisme est un sous-topos de $\Sh(\mathcal{C})$.
\end{lemma}

\begin{proof}
Nous appliquons le lemme \ref{lemma-topos-good-site} pour voir que nous pouvons supposer
que $\mathcal{C}$ est un site possédant les propriétés énumérées dans ce lemme.
En particulier, $\mathcal{C}$ possède un objet final $X$ (de sorte que
$* = h_X$) et un objet $U$ tel que $\mathcal{F} = h_U$.

\medskip\noindent
Soit $\mathcal{D} = \mathcal{C}$ en tant que catégorie, mais déclarons qu'un recouvrement
est une famille $\{V_j \to V\}$ de morphismes telle que
$\{V_i \to V\} \cup \{U \times_X V \to V\}$ soit un recouvrement.
Par notre choix de $\mathcal{C}$, cela signifie exactement que
$$
h_{U \times_X V} \amalg \coprod h_{V_i} \longrightarrow h_V
$$
est surjectif. Nous affirmons que $\mathcal{D}$ est un site, c'est-à-dire que les recouvrements
satisfont les conditions (1), (2), (3) de la définition \ref{definition-site}.
La condition (1) est satisfaite. Pour la condition (2), supposons que
$\{V_i \to V\}$ et $\{V_{ij} \to V_i\}$ soient des recouvrements de $\mathcal{D}$.
Alors la composée
$$
\coprod \left(
h_{U \times_X V_i} \amalg \coprod h_{V_{ij}}
\right) \longrightarrow
h_{U \times_X V} \amalg \coprod h_{V_i} \longrightarrow h_V
$$
est surjective. Comme chacun des morphismes $U \times_X V_i \to V$
se factorise par $U \times_X V$, nous voyons que
$$
h_{U \times_X V} \amalg \coprod h_{V_{ij}} \longrightarrow h_V
$$
est surjectif, c'est-à-dire que $\{V_{ij} \to V\}$ est un recouvrement de $V$ dans
$\mathcal{D}$. La condition (3) s'en déduit de même, puisque tout changement de base
d'un morphisme surjectif de faisceaux est surjectif.

\medskip\noindent
Notons que le foncteur (identité) $u : \mathcal{C} \to \mathcal{D}$ est
continu et commute aux produits fibrés et aux objets finaux. Nous obtenons donc
un morphisme $f : \mathcal{D} \to \mathcal{C}$ de sites
(proposition \ref{proposition-get-morphism}).
Remarquons que $f_*$ est le foncteur identité sur les
préfaisceaux sous-jacents, donc qu'il est pleinement fidèle. Pour achever la démonstration, il nous faut
montrer que l'image essentielle de $f_*$ est la sous-catégorie pleine
$E \subset \Sh(\mathcal{C})$ distinguée dans le lemme. Pour cela, notons
que $\mathcal{G} \in \Ob(\Sh(\mathcal{C}))$ appartient à $E$ si et seulement si
$\mathcal{G}(U \times_X V)$ est un singleton pour tout objet
$V$ de $\mathcal{C}$. Un tel faisceau satisfait donc la
propriété de faisceau pour tous les recouvrements de $\mathcal{D}$ (argument omis).
Réciproquement, si $\mathcal{G}$ satisfait la propriété de faisceau
pour tous les recouvrements de $\mathcal{D}$, alors $\mathcal{G}(U \times_X V)$
est un singleton, puisque, dans $\mathcal{D}$, l'objet $U \times_X V$ est
recouvert par le recouvrement vide.
\end{proof}

\begin{definition}
\label{definition-closed-subtopos}
Soit $\mathcal{C}$ un site. Une sous-catégorie strictement pleine
$E \subset \Sh(\mathcal{C})$ est un {\it sous-topos fermé}
s'il existe un sous-faisceau $\mathcal{F}$ de l'objet final
de $\Sh(\mathcal{C})$ tel que $E$ soit le sous-topos
décrit dans le lemme \ref{lemma-closed-subtopos}.
\end{definition}

\noindent
Nous pouvons maintenant définir ce que sont une immersion fermée
et une immersion ouverte de topos.

\begin{definition}
\label{definition-immersion-topoi}
Soit $f : \Sh(\mathcal{D}) \to \Sh(\mathcal{C})$ un morphisme de topos.
\begin{enumerate}
\item Nous disons que $f$ est une {\it immersion ouverte} si $f$ est un plongement
et si l'image essentielle de $f_*$ est un sous-topos ouvert.
\item Nous disons que $f$ est une {\it immersion fermée} si $f$ est un plongement
et si l'image essentielle de $f_*$ est un sous-topos fermé.
\end{enumerate}
\end{definition}

\begin{lemma}
\label{lemma-closed-immersion}
Soit $i : \Sh(\mathcal{D}) \to \Sh(\mathcal{C})$ une immersion fermée de
topos. Alors $i_*$ est pleinement fidèle, transforme les surjections en surjections,
commute aux coégalisateurs, commute aux sommes amalgamées, reflète les injections,
reflète les surjections et reflète les bijections.
\end{lemma}

\begin{proof}
Soit $\mathcal{F}$ un sous-faisceau de l'objet final $*$ de $\Sh(\mathcal{C})$
et soit $E \subset \Sh(\mathcal{C})$ la sous-catégorie pleine formée
des $\mathcal{G}$ tels que
$\mathcal{F} \times \mathcal{G} \to \mathcal{F}$ soit un isomorphisme.
D'après le lemme \ref{lemma-closed-subtopos},
le foncteur $i_*$ est isomorphe au
foncteur d'inclusion $\iota : E \to \Sh(\mathcal{C})$.

\medskip\noindent
Soit $j_{\mathcal{F}} : \Sh(\mathcal{C})/\mathcal{F} \to \Sh(\mathcal{C})$
le foncteur de localisation (lemme \ref{lemma-localize-topos}).
Notons que $E$ peut également être décrite comme
la collection des faisceaux $\mathcal{G}$ tels que
$j_\mathcal{F}^{-1}\mathcal{G} = *$.

\medskip\noindent
Soient $a, b : \mathcal{G}_1 \to \mathcal{G}_2$ deux morphismes de $E$.
Pour prouver que $\iota$ commute aux coégalisateurs, il suffit de montrer que
le coégalisateur de $a$, $b$ dans $\Sh(\mathcal{C})$ appartient à $E$.
C'est clair, car
le coégalisateur de deux morphismes $* \to *$ est $*$ et parce que
$j_\mathcal{F}^{-1}$ est exact. Il en va de même pour les sommes amalgamées.

\medskip\noindent
Ainsi $i_*$ satisfait les propriétés (5), (6) et (7) du
lemme \ref{lemma-exactness-properties} et, par suite,
le morphisme $i$ satisfait toutes les propriétés mentionnées dans ce
lemme, en particulier celles qui sont mentionnées dans le présent lemme.
\end{proof}








\section{Faisceaux de structures algébriques}
\label{section-sheaves-algebraic-structures}

\noindent
Dans Faisceaux, section \ref{sheaves-section-algebraic-structures},
nous avons défini un type de structure algébrique comme une paire
$(\mathcal{A}, s)$, où $\mathcal{A}$ est une catégorie
et $s : \mathcal{A} \to \textit{Ensembles}$ est un foncteur tel
que
\begin{enumerate}
\item $s$ soit fidèle,
\item $\mathcal{A}$ possède les limites et $s$ commute aux limites,
\item $\mathcal{A}$ possède les colimites filtrantes et $s$ commute à celles-ci, et
\item $s$ reflète les isomorphismes.
\end{enumerate}
Pour un tel type de structure algébrique, nous avons vu qu'un préfaisceau
$\mathcal{F}$ à valeurs dans $\mathcal{A}$ sur un espace $X$ est un faisceau si et
seulement si le préfaisceau d'ensembles associé est un faisceau. De plus,
nous avons élaboré la notion de fibre et, étant donnée une application continue
$f : X \to Y$, nous avons défini les foncteurs adjoints image directe et image inverse
sur les faisceaux de structures algébriques, lesquels coïncident avec l'image directe
et l'image inverse sur les faisceaux d'ensembles sous-jacents. En outre, le prolongement
d'un faisceau de structures algébriques d'une base à tous les ouverts
d'un espace se comporte comme prévu.

\medskip\noindent
Une partie de ces résultats reste valable dans le cadre des sites et des faisceaux.
Soit $(\mathcal{A}, s)$ un type de structure algébrique.
Soit $\mathcal{C}$ un site. Notons
$\textit{PSh}(\mathcal{C}, \mathcal{A})$,
resp.\ $\Sh(\mathcal{C}, \mathcal{A})$ la catégorie
des préfaisceaux, resp.\ des faisceaux à valeurs dans $\mathcal{A}$ sur $\mathcal{C}$.
\begin{itemize}
\item[] ($\alpha$) Un préfaisceau à valeurs dans $\mathcal{A}$ est
un faisceau si et seulement si son préfaisceau d'ensembles sous-jacent est un faisceau.
Voir la démonstration de Faisceaux, Lemme \ref{sheaves-lemma-sheaves-structure}.
\item[] ($\beta$) Étant donné un préfaisceau $\mathcal{F}$ à valeurs dans
$\mathcal{A}$, le préfaisceau ${\mathcal{F}}^\# = (\mathcal{F}^+)^+$
est un faisceau. Cela est vrai parce que les colimites intervenant dans la faisceautisation
sont filtrantes, et même des colimites sur des ensembles ordonnés filtrants (voir
la section \ref{section-sheafification}, notamment la démonstration du
lemme \ref{lemma-sheafification-exact}),
et parce que $s$ commute aux colimites filtrantes.
\item[] ($\gamma$) Nous obtenons le diagramme commutatif suivant
$$
\xymatrix{
\Sh(\mathcal{C}, \mathcal{A}) \ar@<1ex>[r] \ar[d]_s &
\textit{PSh}(\mathcal{C}, \mathcal{A}) \ar@<1ex>[l]^{{}^\#} \ar[d]^s\\
\Sh(\mathcal{C}) \ar@<1ex>[r] &
\textit{PSh}(\mathcal{C}) \ar@<1ex>[l]
}
$$
\item[] ($\delta$) Nous avons $\mathcal{F} = \mathcal{F}^\#$ si et seulement si
$\mathcal{F}$ est un faisceau de structures algébriques.
\item[] ($\epsilon$) Le foncteur ${}^\#$ est adjoint au foncteur d'inclusion :
$$
\Mor_{\textit{PSh}(\mathcal{C}, \mathcal{A})}(\mathcal{G}, \mathcal{F})
=
\Mor_{\Sh(\mathcal{C}, \mathcal{A})}(\mathcal{G}^\#, \mathcal{F})
$$
La démonstration est la même que celle de la
proposition \ref{proposition-sheafification-adjoint}.
\item[] ($\zeta$) Le foncteur
$\mathcal{F} \mapsto \mathcal{F}^\#$ est exact à gauche.
La démonstration est la même que celle du lemme \ref{lemma-sheafification-exact}.
\end{itemize}

\begin{definition}
\label{definition-pushforward-algebraic-structures}
Soit $f : \mathcal{D} \to \mathcal{C}$ un morphisme de sites
donné par un foncteur $u : \mathcal{C} \to \mathcal{D}$.
Nous définissons le foncteur {\it image directe} pour les préfaisceaux de structures algébriques
par la règle $u^p\mathcal{F}(U) = \mathcal{F}(uU)$,
et pour les faisceaux de structures algébriques par la même règle, à savoir
$f_*\mathcal{F}(U) = \mathcal{F}(uU)$.
\end{definition}

\noindent
Le problème survient lorsqu'on cherche à définir l'image inverse.
La raison en est que les colimites définissant
le foncteur $u_p$ dans la section \ref{section-functoriality-PSh}
ne sont pas nécessairement filtrantes. Les axiomes ci-dessus ne suffisent donc pas
en général pour définir l'image inverse d'un (pré)faisceau de structures
algébriques. Néanmoins, dans presque tous les cas, le lemme
suivant suffit pour définir l'image directe et l'image inverse
des (pré)faisceaux de structures algébriques.

\begin{lemma}
\label{lemma-push-pull-good-case}
Supposons que le foncteur $u : \mathcal{C} \to \mathcal{D}$ satisfasse
aux hypothèses de la proposition \ref{proposition-get-morphism}
et donne donc naissance à un morphisme de sites
$f : \mathcal{D} \to \mathcal{C}$. Dans ce cas,
le foncteur image inverse $f^{-1}$ (resp.\ $u_p$) et le foncteur image directe
$f_*$ (resp. $u^p$) s'étendent en une paire de foncteurs adjoints sur
les catégories de faisceaux (resp.\ de préfaisceaux) de structures algébriques.
De plus, ces foncteurs commutent au passage au
faisceau (resp.\ au préfaisceau) d'ensembles sous-jacent.
\end{lemma}

\begin{proof}
Nous avons défini $f_* = u^p$ ci-dessus.
Au cours de la démonstration de la proposition \ref{proposition-get-morphism},
nous avons vu que toutes les colimites utilisées pour définir $u_p$ sont
filtrantes sous les hypothèses de la proposition.
Nous déduisons donc de la définition d'un type de
structure algébrique que nous pouvons définir $u_p$, en tant que foncteur sur
les préfaisceaux de structures algébriques, au moyen exactement des mêmes colimites.
L'adjonction de $u_p$ et $u^p$ se démontre exactement de la
même manière que le lemme \ref{lemma-adjoints-u}.
La discussion ci-dessus de la faisceautisation des préfaisceaux de
structures algébriques implique alors que nous pouvons définir
$f^{-1}(\mathcal{F}) = (u_p\mathcal{F})^\#$.
\end{proof}

\noindent
Nous discutons brièvement une méthode pour traiter l'image inverse et
l'image directe pour un morphisme général de sites et, plus généralement,
pour tout morphisme de topos.

\medskip\noindent
Soit $\mathcal{C}$ un site.
Dans le cas $\mathcal{A} = \textit{Ab}$,
nous pouvons considérer un (pré)faisceau abélien sur $\mathcal{C}$
comme un quadruplet $(\mathcal{F}, +, 0, i)$.
Les données sont les suivantes :
\begin{enumerate}
\item[(D1)] $\mathcal{F}$ est un faisceau d'ensembles,
\item[(D2)] $+ : \mathcal{F} \times \mathcal{F} \to \mathcal{F}$ est
un morphisme de faisceaux d'ensembles,
\item[(D3)] $0 : * \to \mathcal{F}$ est un morphisme depuis le
faisceau singleton (voir l'exemple \ref{example-singleton-sheaf})
vers $\mathcal{F}$, et
\item[(D4)] $i : \mathcal{F} \to \mathcal{F}$ est un morphisme de faisceaux
d'ensembles.
\end{enumerate}
Ces données doivent satisfaire les axiomes suivants :
\begin{enumerate}
\item[(A1)] $+$ est associative et commutative,
\item[(A2)] $0$ est un élément neutre pour $+$, et
\item[(A3)] $+ \circ (1, i) = 0 \circ (\mathcal{F} \to *)$.
\end{enumerate}
Comparer avec Faisceaux, Lemme \ref{sheaves-lemma-abelian-presheaves}.
Soit $f : \mathcal{D} \to \mathcal{C}$ un morphisme de sites.
Notons que, puisque $f^{-1}$ est exact, nous avons
$f^{-1}* = *$ et
$f^{-1}(\mathcal{F} \times \mathcal{F}) =
f^{-1}\mathcal{F} \times f^{-1}\mathcal{F}$.
Nous pouvons donc définir $f^{-1}\mathcal{F}$ simplement comme le quadruplet
$(f^{-1}\mathcal{F}, f^{-1}+, f^{-1}0, f^{-1}i)$. Les axiomes
seront préservés parce que $f^{-1}$ est un foncteur
qui commute aux limites finies. Enfin, il n'est pas difficile
de vérifier que $f_*$ et $f^{-1}$ sont adjoints comme à l'ordinaire.

\medskip\noindent
Dans \cite{SGA4}, cette méthode est utilisée. On y introduit ce
qui est appelé une « {\it esp\`ece de structure alg\'ebrique $\ll$d\'efinie
par limites projectives finies$\gg$} ». Pour une telle esp\`ece, on
peut utiliser la méthode décrite ci-dessus pour définir une paire de foncteurs adjoints
$f^{-1}$ et $f_*$ comme ci-dessus. Cela fonctionne manifestement pour la plupart
des structures algébriques que l'on rencontre en pratique.
Au lieu de formaliser cette construction, nous énumérons simplement les
structures algébriques pour lesquelles cette méthode fonctionne (ce qui doit être
vérifié cas par cas). En fait, cette méthode fonctionne pour tout
morphisme de topos.

\begin{proposition}
\label{proposition-functoriality-algebraic-structures-topoi}
\begin{slogan}
Les morphismes de topos préservent les structures algébriques.
\end{slogan}
Soient $\mathcal{C}$ et $\mathcal{D}$ des sites.
Soit $f = (f^{-1}, f_*)$ un morphisme de topos
de $\Sh(\mathcal{D}) \to \Sh(\mathcal{C})$.
La méthode introduite ci-dessus donne naissance à une
paire de foncteurs adjoints $(f^{-1}, f_*)$ sur les faisceaux de structures algébriques,
compatible avec le passage aux faisceaux d'ensembles sous-jacents,
pour les types suivants de structures algébriques :
\begin{enumerate}
\item ensembles pointés,
\item groupes abéliens,
\item groupes,
\item monoïdes,
\item anneaux,
\item modules sur un anneau fixé, et
\item algèbres de Lie sur un corps fixé.
\end{enumerate}
De plus, dans chacun de ces cas, les résultats ci-dessus étiquetés ($\alpha$),
($\beta$), ($\gamma$), ($\delta$), ($\epsilon$) et ($\zeta$) sont valables.
\end{proposition}

\begin{proof}
La dernière assertion de la proposition est satisfaite simplement parce que chacune des
catégories énumérées, munie du foncteur d'oubli évident, est bien un type de
structure algébrique au sens expliqué au début de cette section.
Voir Faisceaux, Lemme \ref{sheaves-lemma-list-algebraic-structures}.

\medskip\noindent
Démonstration de (2). Nous considérons un faisceau de groupes abéliens comme
un quadruplet $(\mathcal{F}, +, 0, i)$, ainsi qu'il a été expliqué dans la discussion précédant
la proposition.
Si $(\mathcal{F}, +, 0, i)$ est défini sur $\mathcal{C}$, alors son image inverse
est définie comme $(f^{-1}\mathcal{F}, f^{-1}+, f^{-1}0, f^{-1}i)$.
Si $(\mathcal{G}, +, 0, i)$ est défini sur $\mathcal{D}$, alors son image directe
est définie comme $(f_*\mathcal{G}, f_*+, f_*0, f_*i)$. Cette définition convient puisque
$f_*(\mathcal{G} \times \mathcal{G}) = f_*\mathcal{G} \times f_*\mathcal{G}$.
L'adjonction résulte de l'adjonction des foncteurs sur les ensembles,
puisque
$$
\Mor_{\textit{Ab}(\mathcal{C})}
((\mathcal{F}_1, +, 0, i), (\mathcal{F}_2, +, 0, i))
=
\left\{
\begin{matrix}
\varphi \in \Mor_{\Sh(\mathcal{C})}
(\mathcal{F}_1, \mathcal{F}_2) \\
\varphi \text{ est compatible avec }+, 0, i
\end{matrix}
\right\}
$$
Les détails sont laissés au lecteur.

\medskip\noindent
Cette méthode fonctionne également pour les faisceaux d'anneaux lorsque l'on considère
un faisceau d'anneaux (unitaires) comme un sextuplet
$(\mathcal{O}, + , 0, i, \cdot, 1)$ satisfaisant une liste
d'axiomes que l'on peut trouver dans tout ouvrage élémentaire
d'algèbre.

\medskip\noindent
Un faisceau d'ensembles pointés est une paire $(\mathcal{F}, p)$, où
$\mathcal{F}$ est un faisceau d'ensembles et $p : * \to \mathcal{F}$
est un morphisme de faisceaux d'ensembles.

\medskip\noindent
Un faisceau de groupes est donné par un quadruplet $(\mathcal{F}, \cdot, 1, i)$
muni des axiomes appropriés.

\medskip\noindent
Un faisceau de monoïdes est donné par une paire $(\mathcal{F}, \cdot)$
munie de l'axiome approprié.

\medskip\noindent
Soit $R$ un anneau. Un faisceau de $R$-modules est donné par
un quintuplet $(\mathcal{F}, +, 0, i, \{\lambda_r\}_{r \in R})$,
où le quadruplet $(\mathcal{F}, +, 0, i)$ est un faisceau de
groupes abéliens comme ci-dessus et $\lambda_r : \mathcal{F} \to \mathcal{F}$
est une famille de morphismes de faisceaux d'ensembles
telle que
$\lambda_r \circ 0 = 0$,
$\lambda_r \circ + = + \circ (\lambda_r, \lambda_r)$,
$\lambda_{r + r'} =
+ \circ \lambda_r \times \lambda_{r'} \circ (\text{id}, \text{id})$,
$\lambda_{rr'} = \lambda_r \circ \lambda_{r'}$,
$\lambda_1 = \text{id}$, $\lambda_0 = 0 \circ (\mathcal{F} \to *)$.
\end{proof}

\noindent
Nous étudierons la catégorie des faisceaux de modules sur un faisceau d'anneaux dans
Modules sur les sites, section \ref{sites-modules-section-sheaves-modules}.

\begin{remark}
\label{remark-no-pullback-presheaves}
Soient $\mathcal{C}$ et $\mathcal{D}$ des sites.
Soit $u : \mathcal{D} \to \mathcal{C}$ un foncteur continu
qui donne naissance à un morphisme de sites $\mathcal{C} \to \mathcal{D}$.
Notons que, même dans le cas des groupes abéliens, nous n'avons pas défini
de foncteur image inverse pour les préfaisceaux de groupes abéliens.
Puisque toutes les colimites sont représentables dans
la catégorie des groupes abéliens, nous pouvons assurément définir
un foncteur $u_p^{ab}$ sur les préfaisceaux abéliens au moyen des mêmes colimites
que celles que nous avons utilisées pour définir $u_p$ sur les préfaisceaux d'ensembles.
Il sera également vrai que $u_p^{ab}$ est adjoint à
$u^p$ sur les catégories de préfaisceaux abéliens.
Cependant, $u_p^{ab}$ ne coïncidera pas toujours avec
$u_p$ sur les préfaisceaux d'ensembles sous-jacents.
\end{remark}












\section{Applications d'image inverse}
\label{section-pullback}

\noindent
Il arrive parfois qu'un site $\mathcal{C}$ ne possède
pas d'objet final. Dans ce cas, nous définissons le
foncteur des sections globales comme suit.

\begin{definition}
\label{definition-global-sections}
L'{\it ensemble des sections globales} d'un préfaisceau d'ensembles $\mathcal{F}$ sur un
site $\mathcal{C}$ est l'ensemble
$$
\Gamma(\mathcal{C}, \mathcal{F}) =
\Mor_{\textit{PSh}(\mathcal{C})}(*, \mathcal{F})
$$
où $*$ est l'objet final de la catégorie des préfaisceaux sur
$\mathcal{C}$, c'est-à-dire le préfaisceau qui associe à tout objet
un singleton.
\end{definition}

\noindent
Bien entendu, la même définition s'applique aussi aux faisceaux.
Voici une manière de calculer les sections globales.

\begin{lemma}
\label{lemma-compute-global-sections}
Soit $\mathcal{C}$ un site. Soient $a, b : V \to U$ des objets de $\mathcal{C}$
tels que
$$
\xymatrix{
h_V^\# \ar@<1ex>[r] \ar@<-1ex>[r] & h_U^\# \ar[r] & {*}
}
$$
soit un coégalisateur dans $\Sh(\mathcal{C})$. Alors
$\Gamma(\mathcal{C}, \mathcal{F})$ est l'égalisateur de
$a^*, b^* : \mathcal{F}(U) \to \mathcal{F}(V)$.
\end{lemma}

\begin{proof}
Puisque $\Mor_{\Sh(\mathcal{C})}(h_U^\#, \mathcal{F}) = \mathcal{F}(U)$,
cela résulte immédiatement des définitions.
\end{proof}

\noindent
Soit maintenant $f : \Sh(\mathcal{D}) \to \Sh(\mathcal{C})$ un morphisme de topos.
Alors, pour tout faisceau $\mathcal{F}$ sur $\mathcal{C}$, il existe une application d'image inverse
$$
f^{-1} :
\Gamma(\mathcal{C}, \mathcal{F})
\longrightarrow
\Gamma(\mathcal{D}, f^{-1}\mathcal{F})
$$
En effet, comme $f^{-1}$ est exact, il transforme $*$ en $*$. Ainsi, une section
globale $s$ de $\mathcal{F}$ sur $\mathcal{C}$, qui est un morphisme
de faisceaux $s : * \to \mathcal{F}$, a pour image inverse le morphisme
$f^{-1}s : * = f^{-1}* \to f^{-1}\mathcal{F}$.

\medskip\noindent
Nous pouvons généraliser légèrement cette construction en considérant une paire de faisceaux
$\mathcal{F}$, $\mathcal{G}$ sur $\mathcal{C}$, $\mathcal{D}$,
munie d'un morphisme $f^{-1}\mathcal{F} \to \mathcal{G}$. Nous composons alors
la construction ci-dessus avec le morphisme évident
$\Gamma(\mathcal{D}, f^{-1}\mathcal{F}) \to \Gamma(\mathcal{D}, \mathcal{G})$
pour obtenir un morphisme
$$
\Gamma(\mathcal{C}, \mathcal{F})
\longrightarrow
\Gamma(\mathcal{D}, \mathcal{G})
$$
Ce morphisme est parfois lui aussi appelé application d'image inverse.

\medskip\noindent
Une construction un peu plus générale, qui intervient fréquemment de façon naturelle,
est la suivante. Supposons donné un diagramme commutatif de
morphismes de topos
$$
\xymatrix{
\Sh(\mathcal{D}) \ar[rd]_h \ar[rr]_f & &
\Sh(\mathcal{C}) \ar[ld]^g \\
& \Sh(\mathcal{B})
}
$$
Supposons ensuite donné un faisceau $\mathcal{F}$ sur $\mathcal{C}$.
Il existe alors une {\it application d'image inverse}
$$
f^{-1} : g_*\mathcal{F} \longrightarrow h_*f^{-1}\mathcal{F}
$$
En effet, il s'agit simplement du morphisme provenant de l'identification
$g_*f_*f^{-1}\mathcal{F} = h_*f^{-1}\mathcal{F}$ et de
l'application de $g_*$ au morphisme canonique $\mathcal{F} \to f_*f^{-1}\mathcal{F}$.
Si $g$ est l'identité, alors ce morphisme sur les sections globales coïncide
avec l'application d'image inverse ci-dessus.

\medskip\noindent
Dans la situation du paragraphe précédent, supposons que
nous ayons une paire de faisceaux $\mathcal{F}$, $\mathcal{G}$
sur $\mathcal{C}$, $\mathcal{D}$, munie d'un morphisme
$f^{-1}\mathcal{F} \to \mathcal{G}$ ; nous composons alors
l'application d'image inverse ci-dessus avec l'application de $h_*$ au morphisme
$f^{-1}\mathcal{F} \to \mathcal{G}$
pour obtenir un morphisme
$$
g_*\mathcal{F} \longrightarrow h_*\mathcal{G}
$$
En se restreignant aux sections sur un objet de $\mathcal{B}$, on retrouve
l'« application d'image inverse » sur les sections globales discutée ci-dessus (pour des
choix convenables de sites).

\medskip\noindent
Une situation encore plus générale est celle où nous avons un diagramme commutatif de topos
$$
\xymatrix{
\Sh(\mathcal{D}) \ar[d]_h \ar[r]_f &
\Sh(\mathcal{C}) \ar[d]^g \\
\Sh(\mathcal{B}) \ar[r]^e & \Sh(\mathcal{A})
}
$$
et un faisceau $\mathcal{G}$ sur $\mathcal{C}$. Il existe alors un
{\it morphisme de changement de base}
$$
e^{-1}g_*\mathcal{G} \longrightarrow h_*f^{-1}\mathcal{G}.
$$
En effet, ce morphisme est adjoint à un morphisme
$g_*\mathcal{G} \to e_*h_*f^{-1}\mathcal{G} = (e \circ h)_*f^{-1}\mathcal{G}$,
qui est l'application d'image inverse que nous venons de décrire.

\begin{remark}
\label{remark-compose-base-change}
Considérons un diagramme commutatif
$$
\xymatrix{
\Sh(\mathcal{B}') \ar[r]_k \ar[d]_{f'} &
\Sh(\mathcal{B}) \ar[d]^f \\
\Sh(\mathcal{C}') \ar[r]^l \ar[d]_{g'} &
\Sh(\mathcal{C}) \ar[d]^g \\
\Sh(\mathcal{D}') \ar[r]^m &
\Sh(\mathcal{D})
}
$$
de topos. Alors les morphismes de changement de base des deux carrés se composent pour donner le
morphisme de changement de base du rectangle extérieur. Plus précisément, la composée
\begin{align*}
m^{-1} \circ (g \circ f)_*
& =
m^{-1} \circ g_* \circ f_* \\
& \to g'_* \circ l^{-1} \circ f_* \\
& \to g'_* \circ f'_* \circ k^{-1} \\
& = (g' \circ f')_* \circ k^{-1}
\end{align*}
est le morphisme de changement de base du rectangle. Nous omettons la vérification.
\end{remark}

\begin{remark}
\label{remark-compose-base-change-horizontal}
Considérons un diagramme commutatif
$$
\xymatrix{
\Sh(\mathcal{C}'') \ar[r]_{g'} \ar[d]_{f''} &
\Sh(\mathcal{C}') \ar[r]_g \ar[d]_{f'} &
\Sh(\mathcal{C}) \ar[d]^f \\
\Sh(\mathcal{D}'') \ar[r]^{h'} &
\Sh(\mathcal{D}') \ar[r]^h &
\Sh(\mathcal{D})
}
$$
de topos annelés. Alors les morphismes de changement de base
des deux carrés se composent pour donner le morphisme de
changement de base du rectangle extérieur. Plus précisément,
la composée
\begin{align*}
(h \circ h')^{-1} \circ f_*
& =
(h')^{-1} \circ h^{-1} \circ f_* \\
& \to (h')^{-1} \circ f'_* \circ g^{-1} \\
& \to f''_* \circ (g')^{-1} \circ g^{-1} \\
& = f''_* \circ (g \circ g')^{-1}
\end{align*}
est le morphisme de changement de base du rectangle. Nous omettons la vérification.
\end{remark}






\section{Comparaison avec SGA4}
\label{section-notation-comparison}

\noindent
Nos notations pour les foncteurs
$u^p$ et $u_p$ de la section \ref{section-functoriality-PSh} ainsi que
$u^s$ et $u_s$ de la section \ref{section-continuous-functors}
sont empruntées à \cite[pages 14 et 42]{ArtinTopologies}. Ces choix étant faits,
la notation pour le foncteur
${}_pu$ de la section \ref{section-more-functoriality-PSh} et
${}_su$ de la section \ref{section-cocontinuous-functors}
paraît raisonnable.
Dans cette section, nous comparons nos notations à celles de SGA4.

\medskip\noindent
{\bf Préfaisceaux :} Soit $u : \mathcal{C} \to \mathcal{D}$ un foncteur
entre catégories.
Le foncteur $u^p$ est noté $u^*$ dans \cite[Exposé I, section 5]{SGA4}.
Le foncteur $u_p$ est noté $u_!$ dans \cite[Exposé I, proposition 5.1]{SGA4}.
Le foncteur ${}_pu$ est noté $u_*$ dans
\cite[Exposé I, proposition 5.1]{SGA4}.
Autrement dit, nous avons
$$
u_p, u^p, {}_pu\quad(SP)
\quad\text{contre}\quad
u_!, u^*, u_*\quad(SGA4)
$$
Nous attirons l'attention du lecteur sur le fait que des notations différentes
sont employées pour ces foncteurs dans différentes parties de SGA4.

\medskip\noindent
{\bf Faisceaux et foncteurs continus :}
Supposons que $\mathcal{C}$ et $\mathcal{D}$ soient des sites et que
$u : \mathcal{C} \to \mathcal{D}$ soit
un foncteur continu (définition \ref{definition-continuous}).
Le foncteur $u^s$ est noté $u_s$ dans \cite[Exposé III, 1.11]{SGA4}.
Le foncteur $u_s$ est noté $u^s$ dans
\cite[Exposé III, proposition 1.2]{SGA4}.
Autrement dit, nous avons
$$
u_s, u^s\quad(SP)
\quad\text{contre}\quad
u^s, u_s\quad(SGA4)
$$
Lorsque $u$ définit un morphisme de sites
$f : \mathcal{D} \to \mathcal{C}$
(définition \ref{definition-morphism-sites}),
nous voyons que le morphisme de topos associé
(lemme \ref{lemma-morphism-sites-topoi}) est le même que celui de
\cite[Exposé IV, (4.9.1.1)]{SGA4}.

\medskip\noindent
{\bf Faisceaux et foncteurs cocontinus :}
Supposons que $\mathcal{C}$ et $\mathcal{D}$ soient des sites et que
$u : \mathcal{C} \to \mathcal{D}$ soit
un foncteur cocontinu (définition \ref{definition-cocontinuous}).
Le foncteur ${}_su$ (lemme \ref{lemma-pu-sheaf}) est noté $u_*$ dans
\cite[Exposé III, proposition 2.3]{SGA4}.
Le foncteur $(u^p\ )^\#$ est noté $u^*$ dans
\cite[Exposé III, proposition 2.3]{SGA4}.
Autrement dit, nous avons
$$
(u^p\ )^\#, {}_su\quad(SP)
\quad\text{contre}\quad
u^*, u_*\quad(SGA4)
$$
Ainsi, le morphisme de topos associé à $u$ dans le
lemme \ref{lemma-cocontinuous-morphism-topoi}
est le même que celui de \cite[Exposé IV, 4.7]{SGA4}.

\medskip\noindent
{\bf Morphismes de topos :} Si $f$ est un morphisme de topos donné
par les foncteurs $(f^{-1}, f_*)$, alors le foncteur $f^{-1}$
est noté $f^*$ dans \cite[Exposé IV, définition 3.1]{SGA4}.
Nous emploierons $f^{-1}$ pour désigner l'image inverse des faisceaux d'ensembles
ou, plus généralement, des faisceaux de structures algébriques
(section \ref{section-sheaves-algebraic-structures}).
Nous emploierons $f^*$ pour désigner l'image inverse des faisceaux de modules
pour un morphisme de topos annelés
(Modules sur les sites, définition \ref{sites-modules-definition-pushforward}).



\section{Topologies}
\label{section-topologies}

\noindent
Dans cette section, nous définissons ce qu'est une topologie sur une catégorie,
au sens de \cite{SGA4}. On peut développer dans ce langage toute la théorie des
faisceaux et des topos. On trouvera un exposé moderne de cette matière
dans \cite{KS}. Toutefois, le cas qui intéresse le plus la géométrie algébrique
est celui de la topologie définie par un site sur sa catégorie sous-jacente.
Nous conseillons donc vivement au lecteur qui aborde ce texte pour la première fois de
{\bf sauter cette section, ainsi que toutes les autres sections de ce chapitre} !

\begin{definition}
\label{definition-sieve}
Soit $\mathcal{C}$ une catégorie. Soit $U \in \Ob(\mathcal{C})$.
Un {\it crible $S$ sur $U$} est un sous-préfaisceau $S \subset h_U$.
\end{definition}

\noindent
Autrement dit, un crible sur $U$ détermine, pour chaque objet
$T \in \Ob(\mathcal{C})$, un sous-ensemble $S(T)$ de l'ensemble
de tous les morphismes $T \to U$. En fait, l'unique condition
imposée à la famille de sous-ensembles
$S(T) \subset h_U(T) = \Mor_\mathcal{C}(T, U)$
est la règle suivante :
\begin{equation}
\label{equation-property-sieve}
\left.
\begin{matrix}
(\alpha : T \to U) \in S(T) \\
g : T' \to T
\end{matrix}
\right\} \Rightarrow
(\alpha \circ g : T' \to U) \in S(T')
\end{equation}
Une bonne représentation mentale consiste à considérer
l'application $S \to h_U$ comme un ``morphisme de $S$ vers $U$''.

\begin{lemma}
\label{lemma-sieves-set}
Soit $\mathcal{C}$ une catégorie. Soit $U \in \Ob(\mathcal{C})$.
\begin{enumerate}
\item La collection des cribles sur $U$ est un ensemble.
\item L'inclusion définit un ordre partiel sur cet ensemble.
\item Les réunions et les intersections de cribles sont des cribles.
\item
\label{item-sieve-generated}
Étant donnée une famille de morphismes $\{U_i \to U\}_{i\in I}$
de $\mathcal{C}$ de but $U$,
il existe un unique plus petit crible $S$ sur $U$ tel que
chaque $U_i \to U$ appartienne à $S(U_i)$.
\item Le crible $S = h_U$ est le crible maximal.
\item Le sous-préfaisceau vide est le crible minimal.
\end{enumerate}
\end{lemma}

\begin{proof}
D'après notre définition d'un sous-préfaisceau, la collection de
tous les sous-préfaisceaux d'un préfaisceau $\mathcal{F}$ est une partie de
$\prod_{U \in \Ob(\mathcal{C})} \mathcal{P}(\mathcal{F}(U))$.
Or ceci est un ensemble. (Ici $\mathcal{P}(A)$ désigne
l'ensemble des parties de $A$.) Par conséquent, la collection des cribles sur $U$
est un ensemble.

\medskip\noindent
L'ordre partiel est défini par : $S \leq S'$ si et seulement si
$S(T) \subset S'(T)$ pour tout $T \to U$. Notation : $S \subset S'$.

\medskip\noindent
Étant donnée une famille de cribles $S_i$, $i \in I$, sur $U$, nous pouvons
définir $\bigcup S_i$ comme le crible prenant les valeurs
$(\bigcup S_i)(T) = \bigcup S_i(T)$ pour tout
$T \in \Ob(\mathcal{C})$.
Nous définissons l'intersection $\bigcap S_i$ de la même manière.

\medskip\noindent
Étant donnée $\{U_i \to U\}_{i\in I}$ comme dans l'énoncé, considérons
les morphismes de préfaisceaux $h_{U_i} \to h_U$. Nous définissons simplement
$S$ comme la réunion des images (définition \ref{definition-image})
de ces applications de préfaisceaux.

\medskip\noindent
Les deux dernières assertions du lemme sont évidentes.
\end{proof}

\begin{definition}
\label{definition-sieve-generated}
Soit $\mathcal{C}$ une catégorie.
Étant donnée une famille de morphismes $\{f_i : U_i \to U\}_{i\in I}$
de $\mathcal{C}$ de but $U$, nous disons que le crible
$S$ sur $U$ décrit dans le lemme \ref{lemma-sieves-set},
partie (\ref{item-sieve-generated}), est le {\it crible sur $U$
engendré par les morphismes $f_i$}.
\end{definition}

\begin{definition}
\label{definition-pullback-sieve}
Soit $\mathcal{C}$ une catégorie.
Soit $f : V \to U$ un morphisme de $\mathcal{C}$.
Soit $S \subset h_U$ un crible. Nous définissons
{\it l'image réciproque de $S$ par $f$} comme le crible
$S \times_U V$ de $V$ défini par la règle
$$
(\alpha : T \to V) \in (S \times_U V)(T)
\Leftrightarrow
(f \circ \alpha : T \to U) \in S(T)
$$
\end{definition}

\noindent
Nous laissons au lecteur le soin de vérifier qu'il s'agit bien d'un crible
(indication : utiliser l'équation \ref{equation-property-sieve}).
Nous appelons aussi parfois $S \times_U V$ le {\it changement de base}
de $S$ par $f : V \to U$.

\begin{lemma}
\label{lemma-pullback-sieve-section}
Soit $\mathcal{C}$ une catégorie.
Soit $U \in \Ob(\mathcal{C})$.
Soit $S$ un crible sur $U$.
Si $f : V \to U$ appartient à $S$, alors
$S \times_U V = h_V$ est maximal.
\end{lemma}

\begin{proof}
Cela résulte immédiatement des définitions.
\end{proof}

\begin{definition}
\label{definition-topology}
Soit $\mathcal{C}$ une catégorie. Une {\it topologie sur $\mathcal{C}$} consiste
en une règle qui associe à tout $U \in \Ob(\mathcal{C})$
une partie $J(U)$ de l'ensemble de tous les cribles sur $U$, satisfaisant
aux conditions suivantes :
\begin{enumerate}
\item Pour tout morphisme $f : V \to U$ de $\mathcal{C}$ et
tout élément $S \in J(U)$, l'image réciproque $S \times_U V$
est un élément de $J(V)$.
\item Si $S$ et $S'$ sont des cribles sur $U \in \Ob(\mathcal{C})$,
si $S \in J(U)$, et si, pour tout $f \in S(V)$, l'image réciproque
$S' \times_U V$ appartient à $J(V)$, alors $S'$ appartient à $J(U)$.
\item Pour tout $U \in \Ob(\mathcal{C})$, le
crible maximal $S = h_U$ appartient à $J(U)$.
\end{enumerate}
\end{definition}

\noindent
Dans ce cas, les cribles appartenant à $J(U)$ sont appelés
les {\it cribles couvrants}.

\begin{lemma}
\label{lemma-topology-basic}
Soit $\mathcal{C}$ une catégorie.
Soit $J$ une topologie sur $\mathcal{C}$.
Soit $U \in \Ob(\mathcal{C})$.
\begin{enumerate}
\item Les intersections finies d'éléments de $J(U)$ appartiennent à $J(U)$.
\item Si $S \in J(U)$ et $S' \supset S$, alors $S' \in J(U)$.
\end{enumerate}
\end{lemma}

\begin{proof}
Démonstration de (1). Le cas de l'intersection de la famille vide résulte de la condition
(3) de la définition \ref{definition-topology}.
Soient $S, S' \in J(U)$. Considérons $S'' = S \cap S'$. Pour tout
$V \to U$ appartenant à $S(V)$, nous avons
$$
S' \times_U V = S'' \times_U V
$$
simplement parce que $V \to U$ appartient déjà à $S$. Le deuxième
axiome de la définition montre donc que $S'' \in J(U)$.

\medskip\noindent
Démonstration de (2). Soient $S \in J(U)$ et $S' \supset S$. Pour tout
$V \to U$ appartenant à $S(V)$, nous avons $S' \times_U V = h_V$ d'après le
lemme \ref{lemma-pullback-sieve-section}. Ainsi,
$S' \times_U V \in J(V)$ par le troisième axiome. Par conséquent,
$S' \in J(U)$ par le deuxième axiome.
\end{proof}

\begin{definition}
\label{definition-finer}
Soit $\mathcal{C}$ une catégorie. Soient $J$, $J'$ deux
topologies sur $\mathcal{C}$. Nous disons que $J$ est
{\it plus fine} ou {\it plus forte} que $J'$ si et seulement si, pour tout objet
$U$ de $\mathcal{C}$, nous avons $J'(U) \subset J(U)$.
Dans ce cas, nous disons aussi que $J'$ est
{\it moins fine} ou {\it plus faible} que $J$.
\end{definition}

\noindent
Autrement dit, tout crible couvrant pour $J'$ est un
crible couvrant pour $J$. Il existe une topologie la plus fine
sur $\mathcal{C}$, à savoir celle pour laquelle tout crible
est un crible couvrant. On l'appelle la
{\it topologie discrète} de $\mathcal{C}$.
Il existe également une topologie la moins fine.
Il s'agit de la topologie pour laquelle $J(U) = \{h_U\}$
pour tout objet $U$. On l'appelle la
{\it topologie chaotique} ou {\it topologie indiscrète}.

\begin{lemma}
\label{lemma-play-with-topologies}
Soit $\mathcal{C}$ une catégorie.
Soit $\{J_i\}_{i\in I}$ un ensemble de topologies.
\begin{enumerate}
\item La règle $J(U) = \bigcap J_i(U)$ définit
une topologie sur $\mathcal{C}$.
\item Il existe une topologie la moins fine parmi celles qui sont plus fines que
toutes les topologies $J_i$.
\end{enumerate}
\end{lemma}

\begin{proof}
La première assertion résulte directement des définitions.
La seconde s'obtient en prenant l'intersection
de toutes les topologies plus fines que chacun des $J_i$.
\end{proof}

\noindent
Nous pouvons à présent définir,
sans aucune motivation, ce qu'est un faisceau.

\begin{definition}
\label{definition-sheaf-sets-topology}
Soit $\mathcal{C}$ une catégorie munie d'une
topologie $J$. Soit $\mathcal{F}$ un préfaisceau d'ensembles
sur $\mathcal{C}$.
Nous disons que $\mathcal{F}$ est un
{\it faisceau} sur $\mathcal{C}$
si, pour tout $U \in \Ob(\mathcal{C})$ et pour
tout crible couvrant $S$ sur $U$, l'application canonique
$$
\Mor_{\textit{PSh}(\mathcal{C})}(h_U, \mathcal{F})
\longrightarrow
\Mor_{\textit{PSh}(\mathcal{C})}(S, \mathcal{F})
$$
est bijective.
\end{definition}

\noindent
Rappelons que le membre de gauche de la formule
affichée est égal à $\mathcal{F}(U)$. Autrement dit, $\mathcal{F}$
est un faisceau si et seulement si une section de $\mathcal{F}$
sur $U$ revient au même qu'une famille compatible de sections
$s_{T, \alpha} \in \mathcal{F}(T)$ paramétrée par
$(\alpha : T \to U) \in S(T)$, et ce pour tout crible couvrant $S$
sur $U$.

\begin{lemma}
\label{lemma-topology-presheaves-sheaves}
Soit $\mathcal{C}$ une catégorie. Soit $\{ \mathcal{F}_i \}_{i\in I}$ une
famille de préfaisceaux d'ensembles sur $\mathcal{C}$. Pour chaque
$U \in \Ob(\mathcal{C})$, notons
$J(U)$ l'ensemble des cribles $S$ sur $U$ possédant la propriété suivante :
Pour tout morphisme $V \to U$, les applications
$$
\Mor_{\textit{PSh}(\mathcal{C})}(h_V, \mathcal{F}_i)
\longrightarrow
\Mor_{\textit{PSh}(\mathcal{C})}(S \times_U V, \mathcal{F}_i)
$$
sont bijectives pour tout $i \in I$. Alors $J$ définit une
topologie sur $\mathcal{C}$. Cette topologie est la plus fine
parmi celles pour lesquelles tous les $\mathcal{F}_i$ sont des faisceaux.
\end{lemma}

\begin{proof}
Remarquons d'abord que, si $J$ est une topologie, la dernière assertion du
lemme en résulte immédiatement.

\medskip\noindent
Pour tout $i \in I$ et tout $U \in \Ob(\mathcal{C})$, soit $J_i(U)$
l'ensemble des cribles $S$ sur $U$ possédant la propriété suivante :
Pour tout morphisme $V \to U$, l'application
$$
\Mor_{\textit{PSh}(\mathcal{C})}(h_V, \mathcal{F}_i)
\longrightarrow
\Mor_{\textit{PSh}(\mathcal{C})}(S \times_U V, \mathcal{F}_i)
$$
est bijective. Remarquons que $J(U) = \bigcap J_i(U)$. Ainsi, si nous montrons
que $J_i$ définit une topologie, nous concluons par le
lemme \ref{lemma-play-with-topologies}.
Nous sommes donc ramenés au cas étudié dans le paragraphe suivant.

\medskip\noindent
Montrons que $J$ est une topologie lorsque notre famille est réduite à un seul
préfaisceau $\mathcal{F}$. Les premier et troisième axiomes d'une
topologie sont immédiatement vérifiés. Supposons donc que nous ayons
un objet $U$ et des cribles $S, S'$ sur $U$
tels que $S \in J(U)$ et que, pour tout $V \to U$ appartenant à $S(V)$,
nous ayons $S' \times_U V \in J(V)$. Nous devons montrer que
$S' \in J(U)$. Autrement dit, nous devons montrer que, pour
tout $f : W \to U$, l'application
$$
\mathcal{F}(W) =
\Mor_{\textit{PSh}(\mathcal{C})}(h_W, \mathcal{F})
\longrightarrow
\Mor_{\textit{PSh}(\mathcal{C})}(S' \times_U W, \mathcal{F})
$$
est bijective. Choisissons un élément $\varphi \in
\Mor_{\textit{PSh}(\mathcal{C})}(S' \times_U W, \mathcal{F})$.
Nous allons construire une section $s \in \mathcal{F}(W)$
qui s'envoie sur $\varphi$.

\medskip\noindent
Supposons que $\alpha : V \to W$ soit un élément de $S \times_U W$.
D'après la définition des images réciproques, nous voyons que
la composée $f \circ \alpha : V \to W  \to U$ appartient à $S$. Ainsi,
$S' \times_U V$ appartient à $J(V)$ par l'hypothèse faite sur le couple
de cribles $S, S'$. Nous avons alors un diagramme commutatif
de préfaisceaux
$$
\xymatrix{
S' \times_U V \ar[r] \ar[d] & h_V \ar[d] \\
S' \times_U W \ar[r] & h_W
}
$$
La restriction de $\varphi$ à $S' \times_U V$
correspond à un élément $s_{V, \alpha} \in \mathcal{F}(V)$.
Cela résulte de la définition de $J$
et du fait que $S' \times_U V$ appartient à $J(V)$.
Nous laissons au lecteur le soin de vérifier
que la règle $(V, \alpha) \mapsto s_{V, \alpha}$ définit
un élément
$\psi \in
\Mor_{\textit{PSh}(\mathcal{C})}(S \times_U W, \mathcal{F})$.
Puisque $S \in J(U)$, la définition de $J$ montre immédiatement
que $\psi$ correspond à un élément $s$ de $\mathcal{F}(W)$.

\medskip\noindent
Nous laissons au lecteur le soin de vérifier que la construction
$\varphi \mapsto s$ est inverse de l'application naturelle affichée ci-dessus.
\end{proof}

\begin{definition}
\label{definition-canonical-topology}
Soit $\mathcal{C}$ une catégorie.
La topologie la plus fine sur $\mathcal{C}$ pour laquelle
tous les préfaisceaux représentables sont des faisceaux, voir le
lemme \ref{lemma-topology-presheaves-sheaves},
est appelée la {\it topologie canonique} de $\mathcal{C}$.
\end{definition}














\section{La topologie définie par un site}
\label{section-topology-site}

\noindent
Supposons que $\mathcal{C}$ soit une catégorie et que
$\text{Cov}_1(\mathcal{C})$ et $\text{Cov}_2(\mathcal{C})$
soient des ensembles de recouvrements définissant chacun une structure de site
sur $\mathcal{C}$. Dans cette situation, il peut arriver que
les catégories de faisceaux (d'ensembles) pour $\text{Cov}_1(\mathcal{C})$
et $\text{Cov}_2(\mathcal{C})$ soient les mêmes ; voir, par exemple, le
lemme \ref{lemma-refine-same-topology}.

\medskip\noindent
Il se trouve que la catégorie des faisceaux sur $\mathcal{C}$
relativement à une topologie $J$
détermine la topologie $J$ et est déterminée par elle.
C'est un énoncé non trivial que nous traiterons plus loin ;
voir le théorème \ref{theorem-topology-and-topos}.

\medskip\noindent
En l'admettant pour le moment, il est naturel d'étudier la
topologie déterminée par un site.

\begin{lemma}
\label{lemma-site-gives-topology}
Soit $\mathcal{C}$ un site dont les recouvrements sont $\text{Cov}(\mathcal{C})$.
Pour tout objet $U$ de $\mathcal{C}$, notons $J(U)$
l'ensemble des cribles $S$ sur $U$ possédant la propriété suivante :
il existe un recouvrement
$\{f_i : U_i \to U\}_{i\in I} \in \text{Cov}(\mathcal{C})$
tel que le crible $S'$ engendré par les $f_i$ (voir la définition
\ref{definition-sieve-generated}) soit contenu dans $S$.
\begin{enumerate}
\item Ce $J$ est une topologie sur $\mathcal{C}$.
\item Un préfaisceau $\mathcal{F}$ est un faisceau pour cette topologie
(voir la définition \ref{definition-sheaf-sets-topology})
si et seulement si c'est un faisceau sur le site (voir la
définition \ref{definition-sheaf-sets}).
\end{enumerate}
\end{lemma}

\begin{proof}
Pour démontrer la première assertion, il suffit de remarquer que les axiomes
(1), (2) et (3) de la définition d'un site
(définition \ref{definition-site})
entraînent directement les axiomes
(3), (2) et (1) de la définition d'une topologie
(définition \ref{definition-topology}). À titre d'exemple, nous
montrons que $J$ possède la propriété (2). Soit donc $U$ un objet
de $\mathcal{C}$, et soient $S, S'$ des cribles sur $U$ tels que
$S \in J(U)$ et tels que, pour tout $V \to U$ appartenant à $S(V)$,
nous ayons $S' \times_U V \in J(V)$. Par définition de $J(U)$,
nous pouvons trouver un recouvrement $\{f_i : U_i \to U\}$ du site
tel que $S$ l'image de $h_{U_i} \to h_U$ soit contenue
dans $S$. Comme chaque $S'\times_U U_i$ appartient à $J(U_i)$, nous
voyons qu'il existe des recouvrements $\{U_{ij} \to U_i\}$ du
site tels que $h_{U_{ij}} \to h_{U_i}$ soit contenu
dans $S' \times_U U_i$. Par définition du changement de base,
cela signifie que $h_{U_{ij}} \to h_U$ est contenu
dans le sous-préfaisceau $S' \subset h_U$. Par l'axiome (2) pour les
sites, nous voyons que $\{U_{ij} \to U\}$ est un recouvrement de
$U$, et nous concluons que $S' \in J(U)$ par définition de $J$.

\medskip\noindent
Soit $\mathcal{F}$ un préfaisceau. Supposons que $\mathcal{F}$
soit un faisceau pour la topologie $J$. Nous allons montrer que $\mathcal{F}$
est aussi un faisceau sur le site. Soit $\{f_i : U_i \to U\}_{i\in I}$
un recouvrement du site. Soit $s_i \in \mathcal{F}(U_i)$ une
famille de sections telle que
$s_i|_{U_i \times_U U_j} = s_j|_{U_i \times_U U_j}$ pour tous
$i, j$. Nous devons montrer qu'il existe une unique section
$s \in \mathcal{F}(U)$ qui se restreint en $s_i$ sur $U_i$.
Soit $S \subset h_U$ le crible engendré par les $f_i$.
Remarquons que $S \in J(U)$ par définition. Au lieu de construire
$s$, il suffit, d'après la condition de faisceau pour la topologie,
de construire un élément
$$
\varphi \in \Mor_{\textit{PSh}(\mathcal{C})}(S, \mathcal{F}).
$$
Prenons $\alpha \in S(T)$ pour un objet $T \in \mathcal{U}$.
Cela signifie exactement que $\alpha : T \to U$ est un morphisme
qui se factorise par $f_i$ pour un certain $i\in I$ (et peut-être plus de $1$).
Choisissons un tel indice $i$ et une factorisation $\alpha = f_i \circ \alpha_i$.
Posons $\varphi(\alpha) = \alpha_i^* s_i$. Si $i'$ et $\alpha
= f_i \circ \alpha_{i'}'$ donnent un second choix, alors
$\alpha_i^* s_i = (\alpha_{i'}')^* s_{i'}$ précisément en vertu de notre
condition $s_i|_{U_i \times_U U_j} = s_j|_{U_i \times_U U_j}$ pour tous
$i, j$. Ainsi, $\varphi(\alpha)$ est bien défini. Nous laissons au lecteur
le soin de vérifier que $\varphi$, qui à son tour détermine $s$, convient,
au sens où $s$ se restreint en $s_i$.

\medskip\noindent
Soit $\mathcal{F}$ un préfaisceau. Supposons que $\mathcal{F}$
soit un faisceau sur le site $(\mathcal{C}, \text{Cov}(\mathcal{C}))$.
Nous allons montrer que $\mathcal{F}$ est aussi un faisceau pour la topologie $J$.
Soit $U$ un objet de $\mathcal{C}$. Soit
$S$ un crible couvrant sur $U$ relativement à la topologie $J$.
Soit
$$
\varphi \in \Mor_{\textit{PSh}(\mathcal{C})}(S, \mathcal{F}).
$$
Nous devons montrer qu'il existe un unique élément de
$\mathcal{F}(U) = \Mor_{\textit{PSh}(\mathcal{C})}(h_U, \mathcal{F})$
dont la restriction est $\varphi$. Par définition, il existe
un recouvrement $\{f_i : U_i \to U\}_{i\in I} \in \text{Cov}(\mathcal{C})$
tel que $f_i : U_i \to U$ appartienne à $S(U_i)$. Ainsi,
nous pouvons poser $s_i = \varphi(f_i) \in \mathcal{F}(U_i)$.
C'est alors un exercice agréable de vérifier que
$s_i|_{U_i \times_U U_j} = s_j|_{U_i \times_U U_j}$ pour tous
$i, j$. Nous obtenons donc la section voulue $s$ par la condition de
faisceau pour $\mathcal{F}$ sur le site
$(\mathcal{C}, \text{Cov}(\mathcal{C}))$.
Les détails sont laissés au lecteur.
\end{proof}

\begin{definition}
\label{definition-topology-associated-site}
Soit $\mathcal{C}$ un site dont les recouvrements sont $\text{Cov}(\mathcal{C})$.
La {\it topologie associée à $\mathcal{C}$} est la topologie
$J$ construite dans le lemme \ref{lemma-site-gives-topology} ci-dessus.
\end{definition}

\noindent
Soit $\mathcal{C}$ une catégorie.
Soient $\text{Cov}_1(\mathcal{C})$ et $\text{Cov}_2(\mathcal{C})$
deux recouvrements définissant chacun une structure de site sur $\mathcal{C}$.
Il peut fort bien arriver que les topologies ainsi définies
soient les mêmes. Lorsque cela se produit,
nous disons que $\text{Cov}_1(\mathcal{C})$ et
$\text{Cov}_2(\mathcal{C})$ {\it définissent la même topologie} sur
$\mathcal{C}$. Dans ce cas, les catégories de faisceaux
sont les mêmes, par le lemme \ref{lemma-site-gives-topology}.

\medskip\noindent
En général, nous ne nous intéressons qu'à la topologie définie
par une collection de recouvrements, et nous considérons la possibilité de choisir
différents ensembles de recouvrements comme un outil pour étudier la topologie.

\begin{remark}
\label{remark-enlarge-coverings}
Élargissement de la classe des recouvrements.
Il est clair que, si $\text{Cov}(\mathcal{C})$
définit une structure de site sur $\mathcal{C}$, nous pouvons
ajouter à $\mathcal{C}$ n'importe quel ensemble de familles de morphismes de but fixé
tautologiquement équivalentes
(voir la définition \ref{definition-combinatorial-tautological})
à des éléments de $\text{Cov}(\mathcal{C})$ sans changer la topologie.
\end{remark}

\begin{remark}
\label{remark-shrink-coverings}
Réduction de la classe des recouvrements. Soit $\mathcal{C}$ un site.
Considérons l'ensemble
$$
\mathcal{S} = P(\text{Arrows}(\mathcal{C})) \times \Ob(\mathcal{C})
$$
où $P(\text{Arrows}(\mathcal{C}))$ est l'ensemble des parties de l'ensemble des morphismes,
c'est-à-dire l'ensemble de tous les ensembles de morphismes.
Soit $\mathcal{S}_\tau \subset \mathcal{S}$
la partie formée des $(T, U) \in \mathcal{S}$ tels que
(a) tout $\varphi \in T$ ait pour but $U$,
(b) la famille $\{\varphi\}_{\varphi \in T}$ soit tautologiquement
équivalente (voir la définition \ref{definition-combinatorial-tautological})
à un recouvrement de $\text{Cov}(\mathcal{C})$.
Il est clair qu'en considérant les éléments de $\mathcal{S}_\tau$ comme
les recouvrements, nous n'obtenons pas exactement la notion de site
telle qu'elle est définie dans la définition \ref{definition-site}.
La structure $(\mathcal{C}, \mathcal{S}_\tau)$
ainsi obtenue satisfait des conditions légèrement modifiées. Ces
conditions modifiées sont les suivantes :
\begin{enumerate}
\item[(0')] $\text{Cov}(\mathcal{C}) \subset
P(\text{Arrows}(\mathcal{C})) \times \Ob(\mathcal{C})$,
\item[(1')] Si $V \to U$ est un isomorphisme, alors
$(\{V \to U\}, U) \in \text{Cov}(\mathcal{C})$.
\item[(2')] Si $(T, U) \in \text{Cov}(\mathcal{C})$
et si, pour $f : U' \to U$ dans $T$, on s'est donné
$(T_f, U') \in \text{Cov}(\mathcal{C})$,
alors, en posant $T' = \{f \circ f' \mid f \in T,\ f' \in T_f\}$,
nous obtenons $(T', U) \in \text{Cov}(\mathcal{C})$.
\item[(3')] Si $(T, U) \in \text{Cov}(\mathcal{C})$ et si $g : V \to U$
est un morphisme de $\mathcal{C}$, alors
\begin{enumerate}
\item $U' \times_{f, U, g} V$ existe pour $f : U' \to U$ dans $T$, et
\item en posant
$T' = \{\text{pr}_2 : U' \times_{f, U, g} V \to V \mid f : U' \to U \in T\}$
pour un certain choix de produits fibrés, nous obtenons $(T', V) \in \text{Cov}(\mathcal{C})$.
\end{enumerate}
\end{enumerate}
Il est alors facile de vérifier que, partant d'une structure satisfaisant
aux conditions (0') -- (3') ci-dessus, puis en élargissant convenablement
$\text{Cov}(\mathcal{C})$ (comparer avec Ensembles,
section \ref{sets-section-coverings-site}), nous obtenons un site.
Manifestement, il y a peu de différence entre cette notion et la
notion usuelle de site, du moins du point de vue de la
topologie. Elle présente deux avantages :
en raison de la condition (0') ci-dessus, les recouvrements forment automatiquement
un ensemble, et, toujours en raison de (0'), la totalité des structures
de ce type forme elle aussi un ensemble.
Le prix à payer est qu'il faut écrire partout
``tautologiquement équivalente''.
\end{remark}

















\section{Faisceautisation pour une topologie}
\label{section-sheafification-topology}

\noindent
Dans cette section, nous expliquons l'analogue de la
construction de faisceautisation pour une topologie.

\medskip\noindent
Soit $\mathcal{C}$ une catégorie.
Soit $J$ une topologie sur $\mathcal{C}$.
Soit $\mathcal{F}$ un préfaisceau d'ensembles.
Pour tout $U \in \Ob(\mathcal{C})$, nous
définissons
$$
L\mathcal{F}(U)
=
\colim_{S \in J(U)^{opp}}
\Mor_{\textit{PSh}(\mathcal{C})}(S, \mathcal{F})
$$
comme une colimite. Nous considérons ici $J(U)$ comme un ensemble
partiellement ordonné par inclusion, voir le lemme \ref{lemma-sieves-set}. Les
applications de transition du système sont définies comme suit.
Si $S \subset S'$ appartiennent à $J(U)$, alors $S \to S'$ est
un morphisme de préfaisceaux. Il existe donc une application naturelle
de restriction
$$
\Mor_{\textit{PSh}(\mathcal{C})}(S', \mathcal{F})
\longrightarrow
\Mor_{\textit{PSh}(\mathcal{C})}(S, \mathcal{F}).
$$
Ainsi, nous voyons que
$S \mapsto \Mor_{\textit{PSh}(\mathcal{C})}(S, \mathcal{F})$
est un système inductif au sens de Catégories,
définition \ref{categories-definition-system-over-poset},
à condition de renverser l'ordre
sur $J(U)$ (ce qu'indique l'exposant ${}^{opp}$).
En particulier, puisque $h_U \in J(U)$,
il existe une application canonique
$$
\ell : \mathcal{F}(U) \longrightarrow L\mathcal{F}(U)
$$
provenant de l'identification
$\mathcal{F}(U) = \Mor_{\textit{PSh}(\mathcal{C})}(h_U, \mathcal{F})$.
En outre, la colimite définissant $L\mathcal{F}(U)$ est une colimite filtrante,
car, pour toute paire de cribles couvrants $S, S'$ sur $U$, le
crible $S \cap S'$ est encore couvrant, voir le lemme \ref{lemma-sieves-set}.

\medskip\noindent
Soit $f : V \to U$ un morphisme de $\mathcal{C}$.
Soit $S \in J(U)$. Il existe un diagramme commutatif
$$
\xymatrix{
S \times_U V \ar[r] \ar[d] & h_V \ar[d] \\
S \ar[r] & h_U
}
$$
Nous pouvons utiliser la flèche verticale de gauche pour obtenir des applications canoniques de restriction
$$
\Mor_{\textit{PSh}(\mathcal{C})}(S, \mathcal{F})
\to \Mor_{\textit{PSh}(\mathcal{C})}(S \times_U V, \mathcal{F}).
$$
Le changement de base $S \mapsto S \times_U V$ induit une application
croissante $J(U) \to J(V)$. Les applications de restriction
définissent une transformation de foncteurs comme dans
Catégories, lemme \ref{categories-lemma-functorial-colimit}.
Nous obtenons donc une application naturelle de restriction
$$
L\mathcal{F}(U) \longrightarrow L\mathcal{F}(V).
$$

\begin{lemma}
\label{lemma-L-presheaf}
Dans la situation ci-dessus :
\begin{enumerate}
\item La correspondance $U \mapsto L\mathcal{F}(U)$, munie des
applications de restriction définies ci-dessus, est un préfaisceau.
\item Les applications $\ell$ se recollent en un morphisme de préfaisceaux
$\ell : \mathcal{F} \to L\mathcal{F}$.
\item La règle $\mathcal{F} \mapsto (\mathcal{F} \xrightarrow{\ell}
L\mathcal{F})$ est un foncteur.
\item Si $\mathcal{F}$ est un sous-préfaisceau de $\mathcal{G}$, alors
$L\mathcal{F}$ est un sous-préfaisceau de $L\mathcal{G}$.
\item L'application $\ell : \mathcal{F} \to L\mathcal{F}$ possède la
propriété suivante : pour toute section $s \in L\mathcal{F}(U)$,
il existe un crible couvrant $S$ sur $U$ et un élément
$\varphi \in \Mor_{\textit{PSh}(\mathcal{C})}(S, \mathcal{F})$
tels que $\ell(\varphi)$ soit égal à la restriction de
$s$ à $S$.
\end{enumerate}
\end{lemma}

\begin{proof}
La démonstration est omise.
\end{proof}

\begin{definition}
\label{definition-presheaf-separated-topology}
Soit $\mathcal{C}$ une catégorie.
Soit $J$ une topologie sur $\mathcal{C}$.
Nous disons qu'un préfaisceau d'ensembles $\mathcal{F}$
est {\it séparé} si, pour tout objet $U$ et
tout crible couvrant $S$ sur $U$, l'application canonique
$\mathcal{F}(U) \to \Mor_{\textit{PSh}(\mathcal{C})}(S, \mathcal{F})$
est injective.
\end{definition}

\begin{theorem}
\label{theorem-L-topology}
Soit $\mathcal{C}$ une catégorie.
Soit $J$ une topologie sur $\mathcal{C}$.
Soit $\mathcal{F}$ un préfaisceau d'ensembles.
\begin{enumerate}
\item Le préfaisceau $L\mathcal{F}$ est séparé.
\item Si $\mathcal{F}$ est séparé, alors $L\mathcal{F}$ est un faisceau
et le morphisme de préfaisceaux $\mathcal{F} \to L\mathcal{F}$ est injectif.
\item Si $\mathcal{F}$ est un faisceau, alors $\mathcal{F} \to L\mathcal{F}$
est un isomorphisme.
\item Le préfaisceau $LL\mathcal{F}$ est toujours un faisceau.
\end{enumerate}
\end{theorem}

\begin{proof}
L'assertion (3) résulte immédiatement de la définition de $L$ et
de celle d'un faisceau (définition \ref{definition-sheaf-sets-topology}).
L'assertion (4) résulte formellement des autres.

\medskip\noindent
Esquissons la démonstration de (1). Soit $S$ un crible couvrant
sur l'objet $U$. Supposons que $\varphi_i \in L\mathcal{F}(U)$,
$i = 1, 2$, aient la même image dans
$\Mor_{\textit{PSh}(\mathcal{C})}(S, L\mathcal{F})$.
Nous pouvons trouver un même crible couvrant $S'$ sur $U$
tel que les deux $\varphi_i$ soient représentés par des éléments
$\varphi_i \in \Mor_{\textit{PSh}(\mathcal{C})}(S', \mathcal{F})$.
Nous pouvons supposer que $S' = S$ en remplaçant $S$ et $S'$ par
$S' \cap S$, qui est encore un crible couvrant, voir le lemme \ref{lemma-sieves-set}.
Soient $V\in \Ob(\mathcal{C})$ et
$\alpha : V \to U$ dans $S(V)$.
Alors $S \times_U V = h_V$,
voir le lemme \ref{lemma-pullback-sieve-section}. Ainsi, les restrictions
des $\varphi_i$ par $V \to U$ correspondent à des sections $s_{i, V, \alpha}$
de $\mathcal{F}$ sur $V$. L'hypothèse affirme qu'il existe
un crible couvrant $S_{V, \alpha}$ sur $V$ tel que
les $s_{i, V, \alpha}$ se restreignent au même élément de
$\Mor_{\textit{PSh}(\mathcal{C})}(S_{V, \alpha}, \mathcal{F})$.
Considérons le crible $S''$ sur $U$ défini par la règle
\begin{eqnarray}
\label{equation-S-prime-prime}
(f : T \to U) \in S''(T)
& \Leftrightarrow &
\exists\ V , \ \alpha : V \to U, \ \alpha \in S(V), \nonumber \\
& &
\exists\ g : T \to V, \ g \in S_{V, \alpha}(T), \\
& &
f = \alpha \circ g \nonumber
\end{eqnarray}
L'axiome (2) d'une topologie montre que $S''$ est un crible
couvrant sur $U$. Par construction, $\varphi_1$
et $\varphi_2$ se restreignent au même élément de
$\Mor_{\textit{PSh}(\mathcal{C})}(S'', L\mathcal{F})$,
comme souhaité.

\medskip\noindent
Esquissons la démonstration de (2). Supposons que $\mathcal{F}$ soit un
préfaisceau d'ensembles séparé sur $\mathcal{C}$ relativement à
la topologie $J$.
Soit $S$ un crible couvrant sur l'objet $U$ de $\mathcal{C}$.
Supposons que $\varphi \in \Mor_\mathcal{C}(S, L\mathcal{F})$.
Nous devons trouver un élément $s \in L\mathcal{F}(U)$ dont la restriction
soit $\varphi$. Soient $V\in \Ob(\mathcal{C})$ et
$\alpha : V \to U$ dans $S(V)$. La valeur $\varphi(\alpha)
\in L\mathcal{F}(V)$ est donnée par un crible couvrant
$S_{V, \alpha}$ sur $V$ et un morphisme de préfaisceaux
$\varphi_{V, \alpha} : S_{V, \alpha} \to \mathcal{F}$.
Comme dans la démonstration précédente, définissons un crible couvrant $S''$ sur $U$ par
l'équation (\ref{equation-S-prime-prime}). Nous définissons
$$
\varphi'' : S'' \longrightarrow \mathcal{F}
$$
par la règle simple suivante : pour tout $f : T \to U$,
$f \in S''(T)$, choisissons $V, \alpha, g$ comme dans
l'équation (\ref{equation-S-prime-prime}). Posons alors
$$
\varphi''(f) = \varphi_{V, \alpha}(g).
$$
Nous affirmons que ceci est indépendant du
choix de $V, \alpha, g$.
Considérons un second choix$ V', \alpha', g'$.
Les restrictions de $\varphi_{V, \alpha}$ et de
$\varphi_{V', \alpha'}$ à l'intersection
des cribles couvrants suivants sur $T$
$$
(S_{V, \alpha} \times_{V, g} T) \cap (S_{V', \alpha'} \times_{V', g'} T)
$$
coïncident. En effet, ces deux restrictions correspondent à la
restriction de $\varphi$ à $T$ (par $f$), et l'égalité voulue
résulte de ce que $\mathcal{F}$ est séparé.
Notons $\psi$ la restriction commune.
L'indépendance du choix résulte de
$\varphi_{V, \alpha}(g) = \psi(\text{id}_T) =
\varphi_{V', \alpha'}(g')$. Ainsi, $\varphi''$
définit un élément $s \in L\mathcal{F}(U)$. Nous laissons au
lecteur le soin de vérifier que la restriction de $s$ est $\varphi$.
\end{proof}

\begin{definition}
\label{definition-associated-sheaf-topology}
Soit $\mathcal{C}$ une catégorie munie d'une topologie $J$.
Soit $\mathcal{F}$ un préfaisceau d'ensembles sur $\mathcal{C}$.
Le faisceau $\mathcal{F}^\# := LL\mathcal{F}$,
muni du morphisme canonique $\mathcal{F} \to \mathcal{F}^\#$,
est appelé le {\it faisceau associé à $\mathcal{F}$}.
\end{definition}

\begin{proposition}
\label{proposition-sheafification-adjoint-topology}
Soit $\mathcal{C}$ une catégorie munie d'une topologie.
Soit $\mathcal{F}$ un préfaisceau d'ensembles sur $\mathcal{C}$.
Le morphisme canonique $\mathcal{F} \to \mathcal{F}^\#$ possède la
propriété universelle suivante : pour tout morphisme
$\mathcal{F} \to \mathcal{G}$,
où $\mathcal{G}$ est un faisceau d'ensembles, il existe un unique morphisme
$\mathcal{F}^\# \to \mathcal{G}$ tel que $\mathcal{F} \to \mathcal{F}^\#
\to \mathcal{G}$ soit le morphisme donné.
\end{proposition}

\begin{proof}
La démonstration est la même que celle de la proposition \ref{proposition-sheafification-adjoint}.
\end{proof}














\section{Topologies et faisceaux}
\label{section-topology-and-sheaves}

\begin{lemma}
\label{lemma-sieve-sheafification}
Soit $\mathcal{C}$ une catégorie munie d'une topologie $J$.
Soit $U$ un objet de $\mathcal{C}$.
Soit $S$ un crible sur $U$. Les propriétés suivantes sont équivalentes :
\begin{enumerate}
\item Le crible $S$ est couvrant.
\item La faisceautisation $S^\# \to h_U^\#$
du morphisme $S \to h_U$ est un isomorphisme.
\end{enumerate}
\end{lemma}

\begin{proof}
Commençons par quelques remarques générales.
Nous utiliserons les égalités $S^\# = LLS$ et $h_U^\# = LLh_U$.
En particulier, le lemme \ref{lemma-L-presheaf} montre que
$S^\# \to h_U^\#$ est injective. Remarquons que
$\text{id}_U \in h_U(U)$. Cet élément donne donc naissance à
des sections de $Lh_U$ et de $h_U^\# = LLh_U$ sur $U$ que
nous noterons encore $\text{id}_U$.

\medskip\noindent
Supposons que $S$ soit un crible couvrant. Il suffit manifestement de
trouver un morphisme $h_U \to S^\#$ tel que le composé
$h_U \to h_U^\#$ soit le morphisme canonique. Pour trouver un tel morphisme,
il suffit de trouver une section $s \in S^\#(U)$ dont la restriction
soit $\text{id}_U$. Or, puisque $S$ est un
crible couvrant, l'élément
$\text{id}_S \in \Mor_{\textit{PSh}(\mathcal{C})}(S, S)$
donne naissance à une section de $LS$ sur $U$ dont la restriction
est $\text{id}_U$ dans $Lh_U$. Ceci conclut.

\medskip\noindent
Supposons que $S^\# \to h_U^\#$ soit un isomorphisme.
Soit $1 \in S^\#(U)$ l'élément correspondant à
$\text{id}_U$ dans $h_U^\#(U)$. Puisque $S^\# = LLS$,
il existe un crible couvrant $S'$ sur $U$ tel que
$1$ provienne d'un élément
$$
\varphi \in \Mor_{\textit{PSh}(\mathcal{C})}(S', LS).
$$
Cela signifie à son tour que, pour tout $\alpha : V \to U$,
$\alpha\in S'(V)$, il existe un crible couvrant $S_{V, \alpha}$
sur $V$ tel que $\varphi(\alpha)$ corresponde à
un morphisme de préfaisceaux $S_{V, \alpha} \to S$. Autrement dit,
$S_{V, \alpha}$ est contenu dans $S \times_U V$. Le second
axiome d'une topologie montre que $S$ est un crible couvrant.
\end{proof}

\begin{theorem}
\label{theorem-topology-and-topos}
Soit $\mathcal{C}$ une catégorie.
Soient $J$ et $J'$ des topologies sur $\mathcal{C}$.
Les propriétés suivantes sont équivalentes :
\begin{enumerate}
\item $J = J'$,
\item les faisceaux pour la topologie $J$ sont les mêmes que
les faisceaux pour la topologie $J'$.
\end{enumerate}
\end{theorem}

\begin{proof}
Il est tautologique que, si $J = J'$, les notions de faisceau
coïncident. Réciproquement, le lemme \ref{lemma-sieve-sheafification}
caractérise les cribles couvrants au moyen du foncteur de faisceautisation.
Or le foncteur de faisceautisation
$\textit{PSh}(\mathcal{C}) \to \Sh(\mathcal{C}, J)$
est l'adjoint à gauche du foncteur d'inclusion
$\Sh(\mathcal{C}, J) \to \textit{PSh}(\mathcal{C})$.
Par conséquent, si les sous-catégories
$\Sh(\mathcal{C}, J)$ et
$\Sh(\mathcal{C}, J')$ sont les mêmes, alors les foncteurs de faisceautisation
sont les mêmes, et les familles de cribles couvrants sont donc
les mêmes.
\end{proof}

\begin{lemma}
\label{lemma-finer-topology}
Hypothèses et notations comme dans le théorème \ref{theorem-topology-and-topos}.
Alors $J \subset J'$ si et seulement si tout faisceau pour la
topologie $J'$ est un faisceau pour la topologie $J$.
\end{lemma}

\begin{proof}
Une implication est claire. Pour la réciproque, supposons que
$\Sh(\mathcal{C}, J') \subset \Sh(\mathcal{C}, J)$.
Par les propriétés formelles, cela implique
que, si $\mathcal{F}$ est un préfaisceau d'ensembles,
et si $\mathcal{F} \to \mathcal{F}^\#$,
resp.\ $\mathcal{F} \to \mathcal{F}^{\#, \prime}$,
est la faisceautisation relativement à $J$, resp.\ $J'$, alors il existe
un morphisme canonique $\mathcal{F}^\# \to \mathcal{F}^{\#, \prime}$
tel que
$\mathcal{F} \to \mathcal{F}^\# \to \mathcal{F}^{\#, \prime}$
soit le morphisme canonique $\mathcal{F} \to \mathcal{F}^{\#, \prime}$.
Bien entendu, $\mathcal{F}^\# \to \mathcal{F}^{\#, \prime}$
identifie le second faisceau à la faisceautisation du premier
relativement à la topologie $J'$.
Appliquons ceci au morphisme $S \to h_U$ du
lemme \ref{lemma-sieve-sheafification}. Nous obtenons un diagramme
commutatif
$$
\xymatrix{
S \ar[r] \ar[d] &
S^\# \ar[r] \ar[d] &
S^{\#, \prime} \ar[d] \\
h_U \ar[r] &
h_U^\# \ar[r] &
h_U^{\#, \prime}
}
$$
De plus, si $S$ est un crible couvrant pour la topologie $J$,
alors la flèche verticale médiane est un isomorphisme (par le lemme),
et nous en déduisons que la flèche verticale de droite est un isomorphisme,
car elle est la faisceautisation de celle du milieu relativement à $J'$.
Le lemme montre de nouveau que $S$ est aussi un crible couvrant
pour $J'$.
\end{proof}




\section{Topologies et foncteurs continus}
\label{section-topologies-continuous-functors}

\noindent
Expliquer comment un foncteur continu donne une paire adjointe
de foncteurs sur les faisceaux.





\section{Points et topologies}
\label{section-points-topologies}

\noindent
Rappelons, d'après la section \ref{section-points}, qu'étant donné un foncteur
$p = u : \mathcal{C} \to \textit{Ensembles}$, nous pouvons définir
un foncteur fibre
$$
\textit{PSh}(\mathcal{C}) \longrightarrow \textit{Ensembles},
\mathcal{F} \longmapsto \mathcal{F}_p.
$$

\begin{definition}
\label{definition-point-topology}
Soit $\mathcal{C}$ une catégorie.
Soit $J$ une topologie sur $\mathcal{C}$.
Un {\it point $p$} de la topologie est la donnée d'un foncteur
$u : \mathcal{C} \to \textit{Ensembles}$ tel que
\begin{enumerate}
\item Pour tout crible couvrant $S$ sur $U$, l'application
$S_p \to (h_U)_p$ est surjective.
\item Le foncteur fibre $\Sh(\mathcal{C}) \to \textit{Ensembles}$,
$\mathcal{F} \to \mathcal{F}_p$, est exact.
\end{enumerate}
\end{definition}












\begin{multicols}{2}[\section{Autres chapitres}]
\noindent
Préliminaires
\begin{enumerate}
\item \hyperref[introduction-section-phantom]{Introduction}
\item \hyperref[conventions-section-phantom]{Conventions}
\item \hyperref[sets-section-phantom]{Théorie des ensembles}
\item \hyperref[categories-section-phantom]{Catégories}
\item \hyperref[topology-section-phantom]{Topologie}
\item \hyperref[sheaves-section-phantom]{Faisceaux sur les espaces}
\item \hyperref[sites-section-phantom]{Sites et faisceaux}
\item \hyperref[stacks-section-phantom]{Champs}
\item \hyperref[fields-section-phantom]{Corps}
\item \hyperref[algebra-section-phantom]{Algèbre commutative}
\item \hyperref[brauer-section-phantom]{Groupes de Brauer}
\item \hyperref[homology-section-phantom]{Algèbre homologique}
\item \hyperref[derived-section-phantom]{Catégories dérivées}
\item \hyperref[simplicial-section-phantom]{Méthodes simpliciales}
\item \hyperref[more-algebra-section-phantom]{Compléments d'algèbre}
\item \hyperref[smoothing-section-phantom]{Lissage des morphismes d'anneaux}
\item \hyperref[modules-section-phantom]{Faisceaux de modules}
\item \hyperref[sites-modules-section-phantom]{Modules sur les sites}
\item \hyperref[injectives-section-phantom]{Objets injectifs}
\item \hyperref[cohomology-section-phantom]{Cohomologie des faisceaux}
\item \hyperref[sites-cohomology-section-phantom]{Cohomologie sur les sites}
\item \hyperref[dga-section-phantom]{Algèbre différentielle graduée}
\item \hyperref[dpa-section-phantom]{Algèbre à puissances divisées}
\item \hyperref[sdga-section-phantom]{Faisceaux différentiels gradués}
\item \hyperref[hypercovering-section-phantom]{Hyperrecouvrements}
\end{enumerate}
Schémas
\begin{enumerate}
\setcounter{enumi}{25}
\item \hyperref[schemes-section-phantom]{Schémas}
\item \hyperref[constructions-section-phantom]{Constructions de schémas}
\item \hyperref[properties-section-phantom]{Propriétés des schémas}
\item \hyperref[morphisms-section-phantom]{Morphismes de schémas}
\item \hyperref[coherent-section-phantom]{Cohomologie des schémas}
\item \hyperref[divisors-section-phantom]{Diviseurs}
\item \hyperref[limits-section-phantom]{Limites de schémas}
\item \hyperref[varieties-section-phantom]{Variétés}
\item \hyperref[topologies-section-phantom]{Topologies sur les schémas}
\item \hyperref[descent-section-phantom]{Descente}
\item \hyperref[perfect-section-phantom]{Catégories dérivées de schémas}
\item \hyperref[more-morphisms-section-phantom]{Compléments sur les morphismes}
\item \hyperref[flat-section-phantom]{Compléments sur la platitude}
\item \hyperref[groupoids-section-phantom]{Schémas en groupoïdes}
\item \hyperref[more-groupoids-section-phantom]{Compléments sur les schémas en groupoïdes}
\item \hyperref[etale-section-phantom]{Morphismes étales de schémas}
\end{enumerate}
Thèmes de théorie des schémas
\begin{enumerate}
\setcounter{enumi}{41}
\item \hyperref[chow-section-phantom]{Homologie de Chow}
\item \hyperref[intersection-section-phantom]{Théorie de l'intersection}
\item \hyperref[pic-section-phantom]{Schémas de Picard des courbes}
\item \hyperref[weil-section-phantom]{Théories cohomologiques de Weil}
\item \hyperref[adequate-section-phantom]{Modules adéquats}
\item \hyperref[dualizing-section-phantom]{Complexes dualisants}
\item \hyperref[duality-section-phantom]{Dualité pour les schémas}
\item \hyperref[discriminant-section-phantom]{Discriminants et différentes}
\item \hyperref[derham-section-phantom]{Cohomologie de de Rham}
\item \hyperref[local-cohomology-section-phantom]{Cohomologie locale}
\item \hyperref[algebraization-section-phantom]{Géométrie algébrique et formelle}
\item \hyperref[curves-section-phantom]{Courbes algébriques}
\item \hyperref[resolve-section-phantom]{Résolution des surfaces}
\item \hyperref[models-section-phantom]{Réduction semi-stable}
\item \hyperref[functors-section-phantom]{Foncteurs et morphismes}
\item \hyperref[equiv-section-phantom]{Catégories dérivées de variétés}
\item \hyperref[pione-section-phantom]{Groupes fondamentaux de schémas}
\item \hyperref[etale-cohomology-section-phantom]{Cohomologie étale}
\item \hyperref[crystalline-section-phantom]{Cohomologie cristalline}
\item \hyperref[proetale-section-phantom]{Cohomologie pro-étale}
\item \hyperref[relative-cycles-section-phantom]{Cycles relatifs}
\item \hyperref[more-etale-section-phantom]{Compléments de cohomologie étale}
\item \hyperref[trace-section-phantom]{La formule des traces}
\end{enumerate}
Espaces algébriques
\begin{enumerate}
\setcounter{enumi}{64}
\item \hyperref[spaces-section-phantom]{Espaces algébriques}
\item \hyperref[spaces-properties-section-phantom]{Propriétés des espaces algébriques}
\item \hyperref[spaces-morphisms-section-phantom]{Morphismes d'espaces algébriques}
\item \hyperref[decent-spaces-section-phantom]{Espaces algébriques décents}
\item \hyperref[spaces-cohomology-section-phantom]{Cohomologie des espaces algébriques}
\item \hyperref[spaces-limits-section-phantom]{Limites d'espaces algébriques}
\item \hyperref[spaces-divisors-section-phantom]{Diviseurs sur les espaces algébriques}
\item \hyperref[spaces-over-fields-section-phantom]{Espaces algébriques sur des corps}
\item \hyperref[spaces-topologies-section-phantom]{Topologies sur les espaces algébriques}
\item \hyperref[spaces-descent-section-phantom]{Descente et espaces algébriques}
\item \hyperref[spaces-perfect-section-phantom]{Catégories dérivées d'espaces}
\item \hyperref[spaces-more-morphisms-section-phantom]{Compléments sur les morphismes d'espaces}
\item \hyperref[spaces-flat-section-phantom]{Platitude sur les espaces algébriques}
\item \hyperref[spaces-groupoids-section-phantom]{Groupoïdes dans les espaces algébriques}
\item \hyperref[spaces-more-groupoids-section-phantom]{Compléments sur les groupoïdes dans les espaces}
\item \hyperref[bootstrap-section-phantom]{Amorçage}
\item \hyperref[spaces-pushouts-section-phantom]{Sommes amalgamées d'espaces algébriques}
\end{enumerate}
Thèmes de géométrie
\begin{enumerate}
\setcounter{enumi}{81}
\item \hyperref[spaces-chow-section-phantom]{Groupes de Chow des espaces}
\item \hyperref[groupoids-quotients-section-phantom]{Quotients de groupoïdes}
\item \hyperref[spaces-more-cohomology-section-phantom]{Compléments sur la cohomologie des espaces}
\item \hyperref[spaces-simplicial-section-phantom]{Espaces simpliciaux}
\item \hyperref[spaces-duality-section-phantom]{Dualité pour les espaces}
\item \hyperref[formal-spaces-section-phantom]{Espaces algébriques formels}
\item \hyperref[restricted-section-phantom]{Algébrisation des espaces formels}
\item \hyperref[spaces-resolve-section-phantom]{Résolution des surfaces revisitée}
\end{enumerate}
Théorie des déformations
\begin{enumerate}
\setcounter{enumi}{89}
\item \hyperref[formal-defos-section-phantom]{Théorie formelle des déformations}
\item \hyperref[defos-section-phantom]{Théorie des déformations}
\item \hyperref[cotangent-section-phantom]{Le complexe cotangent}
\item \hyperref[examples-defos-section-phantom]{Problèmes de déformation}
\end{enumerate}
Champs algébriques
\begin{enumerate}
\setcounter{enumi}{93}
\item \hyperref[algebraic-section-phantom]{Champs algébriques}
\item \hyperref[examples-stacks-section-phantom]{Exemples de champs}
\item \hyperref[stacks-sheaves-section-phantom]{Faisceaux sur les champs algébriques}
\item \hyperref[criteria-section-phantom]{Critères de représentabilité}
\item \hyperref[artin-section-phantom]{Axiomes d'Artin}
\item \hyperref[quot-section-phantom]{Espaces de Quot et de Hilbert}
\item \hyperref[stacks-properties-section-phantom]{Propriétés des champs algébriques}
\item \hyperref[stacks-morphisms-section-phantom]{Morphismes de champs algébriques}
\item \hyperref[stacks-limits-section-phantom]{Limites de champs algébriques}
\item \hyperref[stacks-cohomology-section-phantom]{Cohomologie des champs algébriques}
\item \hyperref[stacks-perfect-section-phantom]{Catégories dérivées de champs}
\item \hyperref[stacks-introduction-section-phantom]{Introduction aux champs algébriques}
\item \hyperref[stacks-more-morphisms-section-phantom]{Compléments sur les morphismes de champs}
\item \hyperref[stacks-geometry-section-phantom]{La géométrie des champs}
\end{enumerate}
Thèmes de théorie des espaces de modules
\begin{enumerate}
\setcounter{enumi}{107}
\item \hyperref[moduli-section-phantom]{Champs de modules}
\item \hyperref[moduli-curves-section-phantom]{Espaces de modules de courbes}
\end{enumerate}
Divers
\begin{enumerate}
\setcounter{enumi}{109}
\item \hyperref[examples-section-phantom]{Exemples}
\item \hyperref[exercises-section-phantom]{Exercices}
\item \hyperref[guide-section-phantom]{Guide bibliographique}
\item \hyperref[desirables-section-phantom]{Desiderata}
\item \hyperref[coding-section-phantom]{Style de programmation}
\item \hyperref[obsolete-section-phantom]{Obsolète}
\item \hyperref[fdl-section-phantom]{Licence de documentation libre GNU}
\item \hyperref[index-section-phantom]{Index généré automatiquement}
\end{enumerate}
\end{multicols}


\bibliography{my}
\bibliographystyle{amsalpha}

\end{document}
